\documentclass[11pt, a4paper]{ctexart}
\usepackage{../common/ncd-style}

\hypersetup{
  pdftitle={新概念丹麦语 3 · Hverdag og Samfund},
  pdfauthor={Jaehaerys},
  pdfsubject={Ny Koncept Dansk 3}
}

\begin{document}

% =====================================================================
% Book Cover Page
% =====================================================================
\begin{titlepage}
  \centering
  \vspace*{2.5cm}
  {\color{NCDRed}\rule{\linewidth}{2.5pt}}\\[1.2cm]
  {\fontsize{34}{40}\selectfont\bfseries\color{NCDNavy} 新概念丹麦语 3}\\[0.5cm]
  {\huge\itshape\color{NCDTeal} Ny Koncept Dansk 3 · Hverdag og Samfund}\\[0.8cm]
  {\large\bfseries\color{NCDGray} 社会篇（全 20 课完整版） \quad|\quad 对标 NCE 2 后半--NCE 3 前半 \& DU3 Modul 3 (B1)}\\[1.2cm]
  {\color{NCDRed}\rule{\linewidth}{2.5pt}}\\[2cm]

  \begin{tcolorbox}[
    colback=NCDBgBlue,
    colframe=NCDNavy,
    width=0.90\linewidth,
    arc=4pt,
    boxrule=0.8pt,
    left=14pt, right=14pt, top=12pt, bottom=12pt
  ]
    \textbf{\color{NCDNavy}本册学习目标：}实现从“日常叙事”到“社会观察与制度理解”的 B1 跃升（每课 200--300 词）。每课采用“B1 社会制度/文化深读课文 + DU3 Modul 3 真题场景应用”双轨编排，重点攻克 \textbf{\textit{-s} 被动 vs \textit{blive} 被动}、\textbf{介词 + 形式宾语 \textit{det} + \textit{at/hv-} 从句}、\textbf{分词作定语}、\textbf{省略 \textit{hvis} 条件倒装}与\textbf{复合名词构词法}，全面胜任 Modul 3 投诉/推荐与 90+ 词邮件写作、聊天/段落补全阅读与双人互动口试。
  \end{tcolorbox}

  \vfill
  {\large\color{NCDGray} 编写体例：导读 $\rightarrow$ 社会深读与真题文体 $\rightarrow$ 屈折词汇表 $\rightarrow$ 句法精讲 $\rightarrow$ 句型操练 $\rightarrow$ Modul 3 实战与详解}
\end{titlepage}

% Table of Contents
\tableofcontents
\clearpage

% =====================================================================
% Unit 1
% =====================================================================
\ncdunit{1}{Arbejdsliv og Uddannelse}{丹麦职场与教育体系}{Lektion 1--5：被动语态（-s vs blive）· 整句关系从句 · 分词定语 · 无 hvis 条件倒装 · 职场介词搭配}

\ncdlesson{1}{Det flade hierarki på danske arbejdspladser}{丹麦职场的扁平化层级}
{DU3 Modul 3（B1）｜对标 NCE 2 进阶叙事 \& Modul 3 \textit{Arbejdsliv \& Skrivning Opgave 2}}
{丹麦语被动语态两大核心区分：\textbf{\textit{-s} 被动}（客观规则/常态/习惯） vs \textbf{\textit{blive} 被动}（具体事件/动作过程）· \textbf{状态被动（\textit{er} + 过去分词）} · \textit{der} 引出的无人称被动}
{欢迎进入《新概念丹麦语 3》！从本册起，重点从“日常叙事”跃升为“职场与社会制度观察”。在丹麦语讨论规章制度或汇报工作进展时，说话人常刻意淡化“谁做的”，而突出“事情怎么被处理”。彻底分清 \textbf{\textit{-s} 被动}与 \textbf{\textit{blive} 被动}，是 Modul 3 阅读理解与高分写作的第一块基石。}

% ---------------------------------------------------------------------
\begin{foerduleaeser}
带着下面两个问题阅读 Tekst A 和 Tekst B：
\begin{enumerate}
  \item \textit{Hvordan adskiller Nadias danske arbejdsplads sig fra arbejdspladsen i hendes hjemland, og hvad skete der med hendes forslag i sidste uge?}（与祖国相比，娜迪亚所在的丹麦职场有什么不同？上周她的建议经历了什么？）
  \item \textit{Hvor holdes fredagsmøderne i denne uge ifølge Tekst B, og hvorfor blev kundemødet fra Odense flyttet?}（根据 Tekst B，本周的周五例会在哪里举行？为什么与欧登塞客户的会议被改期了？）
\end{enumerate}
\end{foerduleaeser}

% ---------------------------------------------------------------------
\ncdsubsec{二、Tekst A \& Tekst B（社会深读与 Modul 3 真题场景）}

\subsection*{Tekst A：\textit{Det flade hierarki på danske arbejdspladser}（社会文化短文 · 约 245 词）}

\begin{tipbox}{语篇说明}
本篇聚焦丹麦职场最经典的文化现象——“扁平化层级（\textit{det flade hierarki}）”。阅读时请紧盯两种色彩标记：\textbf{\color{NCDNavy}深蓝色加粗的 \textit{-s} 被动}（普遍规则与日常惯例）与\textbf{\color{NCDRed}红色加粗的 \textit{blive} 被动}（具体发生的事件与推进过程）。
\end{tipbox}

Når man begynder at arbejde i en dansk virksomhed, lægger man hurtigt mærke til det flade hierarki. Da den polske ingeniør Nadia for seks måneder siden \dabold{blev ansat} i et it-firma i Aarhus, var hun først lidt usikker på de uskrevne regler. I hendes hjemland \dabold{blev} alle vigtige beslutninger altid \dabold{truffet} af direktøren, og chefen \dabold{blev} konsekvent \dabold{tiltalt} formelt med titel og efternavn.

På hendes nye danske arbejdsplads fungerer det helt anderledes. Her \da{tiltales} både teamledere og direktører ved fornavn, og dørene til kontorerne står som regel åbne. Ugentlige møder \da{holdes} i en uformel stemning, hvor nye ideer \da{diskuteres} åbent over en kop kaffe. Arbejdsopgaverne \da{fordeles} i fællesskab, og det \da{forventes}, at hver medarbejder tager selvstændigt ansvar og tør sige sin ærlige mening --- også selvom man er uenig med sin leder.

I sidste uge oplevede Nadia denne kultur i praksis. Et vigtigt kundesystem skulle \da{opdateres}, så et ekstra møde \dabold{blev indkaldt} tirsdag formiddag. Under debatten \dabold{blev} to tidligere planer \dabold{afvist} af gruppen, mens Nadias eget forslag om et enklere brugerdesign \dabold{blev taget} alvorligt og hurtigt \dabold{blev valgt} af hele afdelingen. Projektets deadline \dabold{er} nu \dabold{blevet flyttet} fem dage frem, så de nye løsninger kan \da{testes} grundigt, før systemet \da{åbnes} for kunderne.

Nu \dabold{er} hele tidsplanen endelig \dabold{godkendt}, og arbejdet er i fuld gang. Nadia kan tydeligt mærke, at folk arbejder med større arbejdsglæde, når tillid og ansvar \da{deles} af hele teamet.

\begin{transbox}
\small
\textbf{\color{NCDTeal}[Tekst A 中文精译]}\\
当一个人在丹麦企业开始工作时，很快就会注意到那种扁平化的层级文化。六个月前，波兰工程师娜迪亚被奥胡斯一家 IT 公司录用时，一开始对这里的潜规则还有些拿不准。在她的祖国，所有重要决定总是都由总经理拍板做出，而且大家称呼老板时一贯要正式地加上头衔和姓氏。

在她新的丹麦工作单位，运转方式却完全不同。在这里，无论是团队主管还是公司总裁，一律直呼其名，办公室的门也通常敞开着。每周例会在轻松非正式的氛围中举行，大家一边喝着咖啡，一边坦诚讨论新点子。工作任务由团队共同分配，而且公司期待每位员工都能独立承担责任，并敢于说出自己的真实想法——哪怕观点与自己的领导不同。

上周，娜迪亚切身体验了这种职场文化。由于一套重要的客户系统需要更新，周二上午公司临时召集了一场特别会议。讨论中，团队否决了两项旧方案，而娜迪亚自己提出的更简洁用户设计建议则受到了认真对待，并很快被整个部门选中。项目的截止日期现已顺延五天，以便新方案在系统向客户开放前得到充分测试。

如今整份时间表终于获批，工作正全面推进。娜迪亚深切体会到，当信任与责任由整个团队共同分担时，大家工作起来会更有幸福感与动力。
\end{transbox}

\subsection*{Tekst B：\textit{Intern mailudveksling på kontoret}（对标 DU3 Modul 3 职场协作与问答）}

\begin{emailbox}{来信｜Fra: Nadia \quad Til: Mads (Teamleder) \quad Emne: Spørgsmål om ugemøde og ferie}
\textbf{Hej Mads}\\
Tak for et rigtig godt møde i går! Jeg vil lige høre, hvordan vores feriedage \da{registreres} i det nye system, og om fredagsmøderne stadig \da{holdes} i kantinen kl. 9.00. Desuden lagde jeg mærke til, at mødet med kunden fra Odense \dabold{blev aflyst} i morges. Ved du, hvornår der \dabold{bliver fundet} en ny dato?\\
\textbf{Venlig hilsen, Nadia}
\end{emailbox}

\begin{emailbox}{回复｜Fra: Mads \quad Til: Nadia \quad Emne: Sv: Spørgsmål om ugemøde og ferie}
\textbf{Hej Nadia}\\
Alle feriedage \da{registreres} direkte i medarbejderportalen, som \da{opdateres} hver mandag. Fredagsmøderne \da{holdes} nu i mødelokale B, fordi kantinen \dabold{bliver malet} i denne uge. Mødet med kunden i Odense \dabold{blev} desværre \dabold{flyttet} på grund af sygdom, men en ny tid \dabold{bliver sendt} ud i eftermiddag. Sig endelig til, hvis noget skal \da{forklares} nærmere!\\
\textbf{Bedste hilsner, Mads}
\end{emailbox}

\textbf{Tekst B 阅读核对（Rigtigt / Forkert / Kort svar，答案见第六部分）：}
\begin{enumerate}
  \item \textit{Medarbejderportalen opdateres hver fredag.} \underline{\hspace{2.2cm}}
  \item \textit{Hvorfor kan fredagsmødet ikke holdes i kantinen i denne uge?} \underline{\hspace{5.5cm}}
  \item \textit{Hvornår får Nadia besked om den nye tid til kundemødet?} \underline{\hspace{4.5cm}}
\end{enumerate}

\begin{tipbox}{本课发音与吞音要点（Udtale \& Reduktioner）}
\begin{itemize}
  \item \textbf{\textit{-s} 被动词尾的真实读音}：现在时 \textit{-s} 被动中，词尾 \textit{-es} 全部弱读为轻柔的 \ipa{əs}，且词干里的弱辅音往往直接消失！如：\textit{tiltales} \ipa{ˈtˢelˌtˢæːləs}、\textit{holdes} \ipa{ˈhʌləs}（\textbf{d 不发音}！）、\textit{bruges} \ipa{ˈbʁuːəs}（\textbf{g 不发音}！）、\textit{træffes} \ipa{ˈtˢʁafəs}、\textit{forklares} \ipa{fʌˈkʰlɑːˀɐs}。
  \item \textbf{助动词 \textit{blive} 全家桶读音}：\textit{blive} \ipa{ˈbliːə}（v 在口语常脱落）、\textit{bliver} \ipa{ˈbliːˀɐ}（带喉塞音 \textit{stød}）、\textit{blev} \ipa{ˈbleʊ̯ˀ}、\textit{blevet} \ipa{ˈbleːˀəð}（结尾为软 d \ipa{ð}）。
  \item \textbf{职场高频核心词}：\textit{et hierarki} \ipa{hiɐ̯ɑˈkʰiːˀ}（重音在最后音节！）、\textit{en arbejdsplads} \ipa{ˈɑːbɑi̯ˀdsˌpʰlas}（d 在 \textit{-ds} 中常弱化）、\textit{en beslutning} \ipa{beˈsludneŋ}、\textit{en afdeling} \ipa{ˈɑʊ̯ˌdeːˀleŋ}（\textit{af} 读如 \ipa{ɑʊ̯}）。
\end{itemize}
\end{tipbox}

\begin{transbox}
\small
\textbf{\color{NCDTeal}[Tekst B 中文对照]}\\
\textbf{娜迪亚的邮件}：你好，马斯。谢谢昨天那场很棒的会议！我想顺便问一下，我们的年假天数在新系统里是如何登记的，以及周五例会是否仍在九点于食堂举行。另外我注意到，原定今早与欧登塞客户的会议被取消了。你知道什么时候能确定一个新日期吗？诚挚的问候，娜迪亚。

\textbf{马斯的回复}：你好，娜迪亚。所有假期天数都直接在员工门户网站登记，系统每周一更新。周五例会现在改在 B 会议室举行，因为食堂这周正在粉刷装修。与欧登塞客户的会议因为生病原因遗憾改期了，不过新的时间安排今天下午就会发出来。如果还有什么需要进一步解释的，请随时告诉我！谨致问候，马斯。
\end{transbox}

% ---------------------------------------------------------------------
\ncdsubsec{三、Gloser（核心生词表）}

\begin{tipbox}{词表重点}
B1 阶段务必随名词一起背诵单复数定冠词，并记牢不规则动词过去分词（用于构成 \textit{blive} 被动！）。注意 \textit{at træffe} 搭配 \textit{en beslutning} 是丹麦语“做决定”的地道强搭配；而 \textit{at blive} 自身作系动词或被动助动词时，完成时必须用 \dabold{[er]}。
\end{tipbox}

\begin{glosertable}
名词 & et hierarki & \textit{(hierarkiet, hierarkier, hierarkierne)} & 等级制度，层级 \\
名词 & en arbejdsplads & \textit{(arbejdspladsen, -pladser, -erne)} & 工作场所，单位 \\
名词 & en beslutning & \textit{(beslutningen, beslutninger, -erne)} & 决定（\textit{træffe en beslutning}） \\
名词 & et ansvar & \textit{(ansvaret, kun ental)} & 责任（\textit{tage ansvar for}） \\
名词 & en medarbejder & \textit{(medarbejderen, -e, -ne)} & 员工，职员 \\
名词 & en afdeling & \textit{(afdelingen, afdelinger, -erne)} & 部门，科室 \\
名词 & et forslag & \textit{(forslaget, forslag, forslagene)} & 建议，提案（单复同形） \\
名词 & en stemning & \textit{(stemningen, stemninger, -erne)} & 气氛，氛围 \\
\midrule
动词 [har] & at tiltale & \textit{(tiltaler, tiltalte, har tiltalt)} & 称呼，呼唤 \\
动词 [har] & at træffe & \textit{(træffer, traf, har truffet)} & 做出（决定）；遇见（不规则） \\
动词 [har] & at fordele & \textit{(fordeler, fordelte, har fordelt)} & 分配，分发 \\
动词 [har] & at forvente & \textit{(forventer, forventede, har forventet)} & 期待，预期 \\
动词 [har] & at indkalde (til) & \textit{(indkalder, indkaldte, har indkaldt)} & 召集（会议） \\
动词 [har] & at afvise & \textit{(afviser, afviste, har afvist)} & 拒绝，驳回 \\
动词 [har] & at godkende & \textit{(godkender, godkendte, har godkendt)} & 批准，认可 \\
动词 \dabold{[er]} & at blive & \textit{(bliver, blev, \dabold{er blevet})} & 成为；被（构成动作被动） \\
\midrule
形容词 & flad / selvstændig & \textit{flad (-t, -e)} / \textit{selvstændig (-t, -e)} & 扁平的 / 独立的，自主的 \\
\end{glosertable}

% ---------------------------------------------------------------------
\ncdsubsec{四、Grammatik \& Faste vendinger（核心语法与地道表达精讲）}

\begin{grammatikbox}{语法点 1：被动语态两大主角——\textit{-s} 被动 vs \textit{blive} 被动}
丹麦语表达“某事被做”，绝不是只靠一种形式，而是根据\textbf{“客观常规 vs 具体事件”}分成两条并行轨道：

\vspace{0.3em}
\renewcommand{\arraystretch}{1.25}
{\small
\begin{tabularx}{\linewidth}{@{} >{\bfseries\raggedright\arraybackslash}p{2.3cm} >{\raggedright\arraybackslash}p{4.1cm} >{\raggedright\arraybackslash}p{4.3cm} >{\raggedright\arraybackslash}X @{}}
\toprule
\textbf{被动类型} & \textbf{构成公式} & \textbf{核心适用场景} & \textbf{典型例句} \\
\midrule
\da{1. \textit{-s} 被动}\newline （\textit{s-passiv}） &
\ding{172} 现在时：原形去 \textit{-r} 加 \da{-s}（\textit{taler $\rightarrow$ tales}）\newline
\ding{173} 情态动词后：直接加 \da{-s}（\textit{skal gøres}） &
\ding{172} \textbf{普遍规律、常态、制度说明}\newline
\ding{173} \textbf{高频搭配 \textit{skal / kan / må / bør}}（必须/可以/应当被……） &
\textit{Chefen \da{tiltales} ved fornavn.}\newline
\textit{Ferie \da{registreres} i portalen.}\newline
\textit{Løsningen skal \da{testes} først.} \\
\midrule
\dabold{2. \textit{blive} 被动}\newline （\textit{blive-passiv}） &
\dabold{bliver / blev / er blevet / vil blive} + \textbf{过去分词} &
\ding{172} \textbf{具体的单个事件、正在进行或已完成的动作过程}\newline
\ding{173} \textbf{过去时、完成时与将来事件}（日常写作文最主力！） &
\textit{Kantinen \dabold{bliver malet} i denne uge.}\newline
\textit{Mødet \dabold{blev aflyst} i morges.}\newline
\textit{Deadlinen \dabold{er blevet flyttet}.} \\
\bottomrule
\end{tabularx}
}
\vspace{0.3em}

\textbf{无人称被动（\textit{Upersonlig passiv}）加分句型}：当动作没有具体名词作主语（如“这里说丹麦语”“会议上讨论了很多问题”）时，丹麦语在句首借用形式主语 \da{der}：\textit{Der \da{tales} dansk på kontoret.} / \textit{Der \dabold{blev} ikke \dabold{truffet} nogen beslutning i går.}（注意否定词 \textit{ikke} 放在 \textit{blev} 后面、分词前面！）。
\end{grammatikbox}

\begin{grammatikbox}{语法点 2：中/英对比排雷——为什么不能照搬英语 \textit{is / was done}？}
学过英语的同学写丹麦语被动时，最容易踩中下面三个致命陷阱：
\begin{enumerate}
  \item \textbf{陷阱一｜动作过程 vs 结果状态（\textit{bliver lukket} vs \textit{er lukket}）}：英语说 \textit{The door is closed}，既可以指“门正被人关上”，也可以指“门现在关着”。丹麦语把这三层意思撕得一清二楚：
    \begin{itemize}
      \item 每天规律（\textit{-s} 被动）：\textit{Kontoret \da{lukkes} kl. 17.00.}（办公室每天 17 点关门——规章制度。）
      \item 正在发生的动作（\textit{blive} 被动）：\textit{Døren \dabold{bliver lukket} nu.}（门现在正在被关上——有人在执行动作。）
      \item 已经形成的静态结果（\textit{er} + 分词，状态被动）：\textit{Døren \textbf{er lukket}.}（门是关着的——只描述状态，切勿用来翻译“门昨天被关了”！）
    \end{itemize}
  \item \textbf{陷阱二｜英语 \textit{was done} 别用 \textit{-tes}，请用 \textit{blev}！}：虽然丹麦语理论上存在过去时 \textit{-tes}（如 \textit{mødet holdtes}），但现代日常语和 B1 写作中，\textbf{具体一次过去事件几乎一律用 \dabold{blev + 过去分词}}（\textit{Mødet blev holdt i går}），不要把过去时 \textit{-s} 当万能模板！
  \item \textbf{陷阱三｜被动句中的“被/由（施动者）”用 \dabold{af}，绝不用 \textit{by / ved / fra}}：英语用 \textit{by the manager}，丹麦语固定用介词 \dabold{af}：\textit{Forslaget blev godkendt \dabold{af} chefen.}（注意：如果动作的发出者不重要或人所共知，地道丹麦语直接省略 \textit{af...}）。
\end{enumerate}
\end{grammatikbox}

\begin{grammatikbox}{地道职场表达（Faste vendinger i arbejdslivet）}
\begin{itemize}
  \item \da{at lægge mærke til noget}（注意到某事）：\textit{Man lægger hurtigt mærke til den åbne kultur.}
  \item \da{at tage noget alvorligt}（认真对待某事）：\textit{Hendes forslag blev taget alvorligt.}
  \item \da{at træffe en beslutning}（做出决定）：比直译 \textit{*lave en beslutning} 地道十倍！
  \item \da{at tiltale nogen ved fornavn}（直呼某人的名字）：丹麦扁平职场文化的标志表达。
\end{itemize}
\end{grammatikbox}

% ---------------------------------------------------------------------
\ncdsubsec{五、Mønsterøvelser（句型操练）}

\textbf{练习 A：语态手术——按括号提示将主动句改写为指定被动句（有施动者时用 \textit{af}，无人称用 \textit{Der}）}
\begin{enumerate}
  \item \textit{Man taler både dansk og engelsk i vores afdeling.}（改 \textit{-s} 被动，以 \textit{Der} 开头）$\rightarrow$
  \item \textit{Chefen aflyste mødet i går morges.}（改 \textit{blive} 过去时被动，保留施动者）$\rightarrow$
  \item \textit{Vi skal aflevere rapporten inden fredag.}（改情态动词 + \textit{-s} 被动，省略施动者）$\rightarrow$
  \item \textit{Afdelingen har nu godkendt det nye forslag.}（改 \textit{blive} 现在完成时被动）$\rightarrow$
  \item \textit{Medarbejderne registrerer altid deres ferie i systemet.}（改 \textit{-s} 被动，省略施动者）$\rightarrow$
  \item \textit{De maler kantinen i denne uge.}（改 \textit{blive} 现在进行过程被动）$\rightarrow$
\end{enumerate}

\textbf{练习 B：精准选填（从 \textit{-s} 被动、\textit{blive} 被动与状态被动中选择唯一准确形式）}
\begin{enumerate}
  \item På denne arbejdsplads \underline{\hspace{2.4cm}} (\textit{tiltales / bliver tiltalt}) alle ledere ved fornavn.
  \item Da Nadia startede for et halvt år siden, \underline{\hspace{2.4cm}} (\textit{indkaldtes / blev hun kaldt}) ind til et velkomstmøde den første dag.
  \item Du kan ikke gå ind i kantinen lige nu, for gulvet \underline{\hspace{2.4cm}} (\textit{vaskes / bliver vasket}) netop nu af rengøringen.
  \item Nu \underline{\hspace{2.4cm}} (\textit{bliver / er}) hele bygningen endelig færdig og malet, så vi kan flytte ind i morgen.
  \item Alle spørgsmål om løn og arbejdstid skal \underline{\hspace{2.4cm}} (\textit{diskutere / diskuteres}) på mødet på torsdag.
  \item I sidste uge \underline{\hspace{2.4cm}} (\textit{flyttes / blev}) vores projektdeadline flyttet frem til den 20. maj.
  \item Der \underline{\hspace{2.4cm}} (\textit{blev ikke truffet / ikke blev truffet}) nogen endelig beslutning i går eftermiddags.
  \item I Danmark \underline{\hspace{2.4cm}} (\textit{forventes / er forventet}) det ofte, at medarbejdere selv tager ansvar.
\end{enumerate}

\textbf{练习 C：中译丹实战（活用两种被动与职场固定表达）}
\begin{enumerate}
  \item 在丹麦企业里，老板通常被直呼其名。
  \item 这项任务必须在周一下午四点之前提交。
  \item 昨天上午召开了一场关于新项目的会议。
  \item 她的建议在会议上被团队认真对待了。
  \item 在我们部门，所有重要决定都是大家共同讨论的。
  \item 幸好办公室的门现在终于修好了（\textit{reparere}）。
\end{enumerate}

% ---------------------------------------------------------------------
\ncdsubsec{六、DU3 Modul 3 实战微任务与参考答案}

\begin{emailbox}{DU3 Modul 3 微型写作任务（对标 Skrivning Opgave 2 · 私人回复邮件）}
\textbf{Situation}: Din ven Kasper er lige blevet ansat i en stor dansk virksomhed i København og har skrevet til dig, fordi han er lidt nervøs før sin første uge.\\
\textbf{Opgave}: Skriv et opmuntrende svar på \textbf{100--130 ord}. Fortæl om det flade hierarki på danske arbejdspladser (hvordan man tiltaler chefen, hvordan møder holdes), giv et konkret eksempel på noget, der blev besluttet på din arbejdsplads for nylig, og brug både \textbf{\textit{-s} passiv} og \textbf{\textit{blive}-passiv}.
\end{emailbox}

\begin{answerbox}
\textbf{Før du læser 参考回答}
\begin{enumerate}
  \item \textit{I Danmark tiltales både ledere og direktører ved fornavn, møder holdes i en uformel stemning, og det forventes, at alle siger deres ærlige mening. I sidste uge blev Nadias forslag taget alvorligt og valgt af hele afdelingen.}
  \item \textit{Fredagsmøderne holdes i mødelokale B i denne uge, fordi kantinen bliver malet. Kundemødet fra Odense blev flyttet på grund af sygdom.}
\end{enumerate}

\textbf{Tekst B 阅读核对答案}
\begin{enumerate}
  \item \da{Forkert}（门户网站每周一更新：\textit{opdateres hver mandag}）
  \item \da{Fordi kantinen bliver malet i denne uge.}
  \item \da{I eftermiddag.}（\textit{en ny tid bliver sendt ud i eftermiddag}）
\end{enumerate}

\textbf{练习 A 答案（语态手术）}
\begin{enumerate}
  \item \textit{Der \da{tales} både dansk og engelsk i vores afdeling.}（无人称常态用 \textit{Der + -s} 被动）
  \item \textit{Mødet \dabold{blev aflyst} af chefen i går morges.}（具体过去事件用 \textit{blev + 分词}，施动者前用 \textit{af}）
  \item \textit{Rapporten skal \da{afleveres} inden fredag.}（情态动词 \textit{skal} 后接 \textit{-s} 被动原形）
  \item \textit{Det nye forslag \dabold{er} nu \dabold{blevet godkendt} af afdelingen.}（现在完成时被动：\textit{er blevet + 分词}）
  \item \textit{Ferie \da{registreres} altid i systemet.}（普遍规章制度用 \textit{-s} 被动）
  \item \textit{Kantinen \dabold{bliver malet} i denne uge.}（正在进行的事件过程用 \textit{bliver + 分词}）
\end{enumerate}

\textbf{练习 B 答案}
\begin{enumerate}
  \item \da{tiltales}（职场普遍惯例）
  \item \da{blev hun kaldt}（半年前第一天发生的具体事件，且时间从句在句首触发主句 V2 倒装）
  \item \da{bliver vasket}（此时此刻正在发生的动作过程）
  \item \da{er}（状态被动：大楼现已完工并粉刷好了）
  \item \da{diskuteres}（情态动词 \textit{skal} 后必须用原形加 \textit{-s}）
  \item \da{blev}（具体过去时间 \textit{I sidste uge} 触发倒装且需用 \textit{blev}）
  \item \da{blev ikke truffet}（主句中否定词 \textit{ikke} 位于变位助动词 \textit{blev} 之后）
  \item \da{forventes}（丹麦社会的普遍期待）
\end{enumerate}

\textbf{练习 C 答案}
\begin{enumerate}
  \item \textit{I danske virksomheder tiltales chefen normalt ved fornavn.}
  \item \textit{Denne opgave skal afleveres inden mandag klokken seksten.}
  \item \textit{I går formiddags blev der indkaldt til et møde om det nye projekt.}
  \item \textit{Hendes forslag blev taget alvorligt af teamet på mødet.}
  \item \textit{I vores afdeling diskuteres alle vigtige beslutninger i fællesskab.}
  \item \textit{Heldigvis er døren til kontoret nu endelig repareret.}
\end{enumerate}

\textbf{DU3 Modul 3 高分范文（115 词，不含署名）}\\
\textit{Hej Kasper}

\textit{Stort tillykke med det nye job! Jeg kan godt forstå, at du er lidt spændt, men du behøver slet ikke være nervøs for det flade hierarki i Danmark.}

\textit{På danske arbejdspladser \da{tiltales} alle ved fornavn, og stemningen er som regel meget uformel. Møder \da{holdes} ofte med åbne diskussioner, hvor det \da{forventes}, at også nye medarbejdere deler deres ideer. Hvis en opgave skal \da{løses}, \da{fordeles} ansvaret normalt i fællesskab.}

\textit{I sidste uge oplevede jeg det selv: Vores gamle plan \dabold{blev afvist} på tirsdagsmødet, og i stedet \dabold{blev} mit forslag \dabold{valgt} af hele gruppen! Alt \dabold{blev diskuteret} roligt og respektfuldt. Vær bare nysgerrig og stil spørgsmål, så skal det nok gå rigtig godt.}

\textit{Bedste hilsner\\
\mbox{[Dit navn]}}

\small\textbf{名师点评}：115 词（不含署名）精准落地 B1 要求！前段用 \textit{tiltales, holdes, forventes, skal løses, fordeles} 等 \textit{-s} 被动勾勒普遍规则，后段用 \textit{blev afvist, blev valgt, blev diskuteret} 叙述上周亲历的具体事件，两条被动轨道分明，且在从句（\textit{hvor det forventes, at...}）与倒装（\textit{I sidste uge oplevede jeg...}）上毫无瑕疵。
\end{answerbox}

\ncdlesson{2}{Fagforening, A-kasse og "flexicurity"}{工会、失业保险与“灵活安全”模式}
{DU3 Modul 3（B1）｜对标 NCE 2 社会制度说明文 \& Modul 3 \textit{Læsning Opgave 2 (Chat) \& Opgave 3}}
{\textbf{指代整句的关系从句}（\textit{..., hvilket betyder, at...}）· \textbf{介词复合关系代词}（\textit{hvori, hvormed, hvoraf, hvorfor}）vs 口语“\textit{som + 后置介词}”结构}
{丹麦语的 \textit{fagforening}（工会）与 \textit{a-kasse}（失业保险基金）是 DU3 Modul 3 阅读与口试的高频社会议题。语法上，学会用 \textbf{\textit{, hvilket ...}} 将两个独立短句丝滑升格为带有逻辑评价的 B1 复合长句，能立刻让你的书面表达摆脱“流水账”质感！}

% ---------------------------------------------------------------------
\begin{foerduleaeser}
带着下面两个问题阅读 Tekst A 和 Tekst B：
\begin{enumerate}
  \item \textit{Hvad er forskellen på en fagforening og en a-kasse i Danmark, og hvori består ”flexicurity-modellen”?}（在丹麦，工会和失业保险基金有什么区别？“灵活安全模式”的核心是什么？）
  \item \textit{Hvad er Jonas i tvivl om med hensyn til sin nye jobkontrakt i Tekst B, og hvad råder Freja ham til?}（在 Tekst B 中，乔纳斯对自己的新劳动合同有什么疑惑？弗蕾娅给了他什么建议？）
\end{enumerate}
\end{foerduleaeser}

% ---------------------------------------------------------------------
\ncdsubsec{二、Tekst A \& Tekst B（社会深读与 Modul 3 真题场景）}

\subsection*{Tekst A：\textit{Fagforening, A-kasse og ”flexicurity”}（社会制度说明文 · 约 245 词）}

\begin{tipbox}{语篇说明}
本文剖析丹麦劳动力市场的三大支柱。阅读时请特别留意\textbf{\color{NCDNavy}深蓝色加粗的 \textit{, hvilket ...}}（回指前面一整件事，译为“这一点使得/意味着……”），以及\textbf{\color{NCDRed}红色加粗的复合关系词 \textit{hvormed / hvori / hvorfor}}。
\end{tipbox}

Det danske arbejdsmarked er kendt verden over for den såkaldte ”flexicurity-model”, som kombinerer fleksibilitet for virksomheder med tryghed for lønmodtagere. I Danmark er det forholdsvis nemt at ansætte og opsige medarbejdere, \da{hvilket gør}, at virksomheder hurtigt kan tilpasse sig ændringer i økonomien.

Til gengæld er de fleste medarbejdere beskyttet af to vigtige institutioner, som udlændinge ofte forveksler: en fagforening og en a-kasse. En fagforening forhandler løn, arbejdstid og pension med arbejdsgiverne gennem kollektive overenskomster. Danmark har nemlig ingen lovbestemt mindsteløn, \da{hvilket betyder, at} en stærk fagforening spiller en helt central rolle for gode arbejdsvilkår. Mange fagforeninger tjekker også nye kontrakter gratis, før man skriver under, \da{hvilket giver} stor tryghed for nyuddannede og internationale kolleger. Derudover tilbyder fagforeningen juridisk rådgivning, hvis man bliver uenig med sin chef eller får en uretfærdig opsigelse.

En a-kasse fungerer derimod som en forsikring, \dabold{hvormed} man kan sikre sig økonomisk støtte i form af dagpenge, hvis man mister sit arbejde. For at få ret til fulde dagpenge skal man dog normalt have været medlem i mindst et år og have arbejdet et vist antal timer.

Samlet set danner disse ordninger et system, \dabold{hvori} høj mobilitet forenes med et stærkt socialt sikkerhedsnet. Når arbejdsløse samtidig får aktiv vejledning og videreuddannelse via jobcentret, finder mange hurtigt en ny stilling, \da{hvilket mindsker} utrygheden ved at skifte job. Det er netop hovedårsagen, \dabold{hvorfor} de fleste nye medarbejdere vælger at melde sig ind i både en fagforening og en a-kasse, så snart de underskriver deres første danske kontrakt.

\begin{transbox}
\small
\textbf{\color{NCDTeal}[Tekst A 中文精译]}\\
丹麦劳动力市场因所谓的“灵活安全模式（flexicurity）”而闻名世界，该模式将企业的用工灵活性与雇员的保障感结合在了一起。在丹麦，招聘和解雇员工相对容易，这使得企业能够迅速适应经济变化。

作为平衡，大多数雇员都受到两个常被外国人混淆的重要机构的保护：工会（\textit{fagforening}）与失业保险基金（\textit{a-kasse}）。工会通过集体劳资协议与雇主谈判工资、工时和养老金。因为丹麦没有法定的全国最低工资标准，这意味着强大的工会在争取良好工作条件方面扮演着核心角色。许多工会还会在员工签字前免费审核新合同，这为刚毕业的新人和国际同事带来了极大的安心感。此外，如果员工与上司发生分歧或遭遇不公正解雇，工会还提供法律咨询。

相比之下，失业保险基金则像一份保险，借助它人们可以在失业时为自己确保以失业救济金（\textit{dagpenge}）形式发放的经济支持。不过，为了获得领取全额救济金的权利，通常必须已入会满一年并工作过一定时长。

总体而言，这些制度构成了一个体系，在其中高流动性与强大的社会安全网融为一体。当失业者同时通过就业中心获得积极指导和进修培训时，许多人能迅速找到新职位，这降低了跳槽带来的不安全感。这正是绝大多数新员工一签下第一份丹麦劳动合同，就选择同时加入工会和失业保险基金的主要原因。
\end{transbox}

\subsection*{Tekst B：\textit{Chat om kontrakt, fagforening og a-kasse}（对标 DU3 Modul 3 Læsning Opgave 2 聊天补全）}

\begin{tcolorbox}[enhanced, breakable, colback=NCDBgBlue, colframe=NCDNavy, coltitle=white, fonttitle=\bfseries, title={同事聊天气泡补全（请从下方 A--D 中选出 2 句填入 \textbf{(1)} 和 \textbf{(2)}，有 2 项多余）}, arc=2pt, boxrule=0.8pt, left=10pt, right=10pt, top=7pt, bottom=7pt]
\small
\textbf{Jonas}: Hej Freja! Jeg har lige fået mit første faste jobtilbud i Danmark! Men jeg læser kontrakten nu, \da{hvilket} gør mig lidt forvirret. Har man pligt til at melde sig ind i en fagforening?\\[3pt]
\textbf{Freja}: Tillykke, Jonas! \textbf{(Z 例)} \textit{Nej, det er helt frivilligt, men de fleste gør det alligevel.}\\[3pt]
\textbf{Jonas}: Okay, godt at vide. Men er fagforeningen og a-kassen ikke det samme firma?\\[3pt]
\textbf{Freja}: \textbf{(1)} \underline{\hspace{5.5cm}}\\[3pt]
\textbf{Jonas}: Ah, så giver det mening! Hvad med kontrakten --- hvordan ved jeg, om min løn og ferie passer?\\[3pt]
\textbf{Freja}: \textbf{(2)} \underline{\hspace{5.5cm}}\\[3pt]
\textbf{Jonas}: Perfekt! Det gør jeg allerede i morgen, inden jeg skriver under. Tak for hjælpen!

\vspace{0.4em}
\hrule
\vspace{0.35em}
\textbf{备选回复（A--D）}：\\
\textbf{A.} \textit{Send kontrakten til din fagforening først, så de kan tjekke alle punkterne, hvori dine rettigheder står.}\\
\textbf{B.} \textit{Nej, slet ikke. Fagforeningen hjælper med selve kontrakten og arbejdsvilkårene, mens a-kassen udbetaler dagpenge, hvis du bliver ledig.}\\
\textbf{C.} \textit{Jo, hvis du arbejder mere end 37 timer om ugen, skal du automatisk betale til begge dele fra næste måned.}\\
\textbf{D.} \textit{Jeg har faktisk glemt at forlænge mit medlemskab sidste år, hvilket desværre kostede mig et nyt kursus.}
\end{tcolorbox}

\begin{transbox}
\small
\textbf{\color{NCDTeal}[Tekst B 中文对照]}\\
\textbf{乔纳斯}：嗨，弗蕾娅！我刚拿到在丹麦的第一份正式全职工作录用通知！但我现在看合同，这让我有点糊涂。人有义务必须加入工会吗？\\
\textbf{弗蕾娅}：\textbf{（例 Z）}没有，那完全是自愿的，不过大多数人还是会加入。\\
\textbf{乔纳斯}：好，明白了。但工会和失业保险基金难道不是同一家机构吗？\\
\textbf{备选核对}：\textbf{A.} 先把合同发给你的工会，让他们核对写有你各项权利的条款。\quad \textbf{B.} 完全不是。工会负责帮你处理合同本身和工作条件，而当你失业时，失业保险基金会向你发放救济金。\quad \textbf{C.} 是的，如果你每周工作超过 37 小时，下个月起就必须自动向两边缴费。\quad \textbf{D.} 其实我去年忘了续期我的会员资格，这遗憾地让我错过了一门新课程。
\end{transbox}

\begin{tipbox}{本课发音与吞音要点（Udtale \& Reduktioner）}
\begin{itemize}
  \item \textbf{\textit{hv-} 关系词家族全员“h 不发音”}：\textit{hvilket} \ipa{ˈvelgəð}（词首 \textbf{h 静音}，词尾 \textit{-et} 弱读为 \ipa{əð}）；\textit{hvori} \ipa{vɒˈiːˀ}（重音在后面的 \textit{-i}！）；\textit{hvormed} \ipa{vɒˈmɛð}；\textit{hvorfor} \ipa{ˈvɒːfʌ}（重音在前面 \textit{hvor-}）。
  \item \textbf{软化 \textit{g} 变半元音 \ipa{ʊ̯}}：注意 \textit{fagforening} \ipa{ˈfɑʊ̯fʌˌeːˀneŋ} 和 \textit{dagpenge} \ipa{ˈdɑʊ̯ˌpʰɛŋə} 中的 \textit{fag} 与 \textit{dag}，字母 \textit{g} 都不发硬音 \ipa{ɡ}，而是滑向 \ipa{ɑʊ̯}！
  \item \textbf{职场社制长词重音}：\textit{overenskomst} \ipa{ʌʊ̯ɐˈenˀsˌkʰʌmsd}（主重音在 \textit{-ens-}，带喉塞音）；\textit{arbejdsmarked} \ipa{ˈɑːbɑi̯ˀdsˌmɑːɡəð}；\textit{tryghed} \ipa{ˈtˢʁøɡˌheːˀð}；\textit{sikkerhed} \ipa{ˈseɡɐˌheːˀð}。
\end{itemize}
\end{tipbox}

% ---------------------------------------------------------------------
\ncdsubsec{三、Gloser（核心生词表）}

\begin{tipbox}{词表重点}
\textit{dagpenge}（失业救济金）在丹麦语中\textbf{只有复数形式（kun flertal）}，定指形为 \textit{dagpengene}，作主语或搭配修饰语时也要用复数形态（如 \textit{fulde dagpenge}）。此外，“加入（协会/工会/俱乐部）”的黄金动词短语是 \textit{at melde sig ind i...}，而“退出”则是 \textit{at melde sig ud af...}！
\end{tipbox}

\begin{glosertable}
名词 & en fagforening & \textit{(fagforeningen, -foreninger, -erne)} & 工会 \\
名词 & en a-kasse & \textit{(a-kassen, a-kasser, a-kasserne)} & 失业保险基金（arbejdsløshedskasse） \\
名词 & et arbejdsmarked & \textit{(arbejdsmarkedet, -markeder, -erne)} & 劳动力市场 \\
名词 & en overenskomst & \textit{(overenskomsten, -er, -erne)} & 集体劳资协议 \\
名词 & dagpenge & \textit{(dagpengene, kun flertal)} & 失业救济金（仅复数） \\
名词 & en opsigelse & \textit{(opsigelsen, opsigelser, -erne)} & 解雇通知；辞职信 \\
名词 & en løn & \textit{(lønnen, lønninger, lønningerne)} & 工资，薪水 \\
名词 & en tryghed & \textit{(trygheden, kun ental)} & 安全感，保障感 \\
名词 & en kontrakt & \textit{(kontrakten, kontrakter, -erne)} & 合同，契约 \\
名词 & en lønmodtager & \textit{(lønmodtageren, -e, -ne)} & 领薪雇员，受薪劳动者 \\
\midrule
动词 [har] & at ansætte & \textit{(ansætter, ansatte, har ansat)} & 雇佣，录用 \\
动词 [har] & at opsige & \textit{(opsiger, opsagde, har opsagt)} & 解雇；终止（合同/会员） \\
动词 [har] & at forhandle & \textit{(forhandler, forhandlede, har forhandlet)} & 谈判，协商 \\
动词 [har] & at sikre (sig) & \textit{(sikrer, sikrede, har sikret)} & 确保，保障 \\
动词 [har] & at melde sig ind (i) & \textit{(melder, meldte, har meldt sig ind)} & 报名加入（工会/协会） \\
\midrule
形容词 & fleksibel / arbejdsløs & \textit{fleksibelt, fleksible} / \textit{arbejdsløst, -e} & 灵活的 / 失业的（同义：\textit{ledig}） \\
\end{glosertable}

% ---------------------------------------------------------------------
\ncdsubsec{四、Grammatik \& Faste vendinger（核心语法与地道表达精讲）}

\begin{grammatikbox}{语法点 1：整句回指神器 \textit{, hvilket ...}（“这意味/使得……”）}
在 A2 阶段，学生常写出两句断开的话：\textit{Danmark har ingen mindsteløn. Det betyder, at...}。进入 B1，丹麦人只需一个\textbf{逗号 + \textit{hvilket}} 就能把它拧成紧凑的高分长句：
\begin{itemize}
  \item \textbf{核心规则}：\da{, hvilket} 只能回指\textbf{前面一整句话/整件事}，不可用来代替一个具体名词（除非极少数中性名词的书面修饰，但日常应尽量只把 \textit{hvilket} 用于回指整句！）。
  \item \textbf{从句语序铁律}：\textit{hvilket} 引导的是\textbf{非限制性关系从句}，内部必须遵守\textbf{从句语序}（若有 \textit{ikke / desværre / ofte}，必须放在\textbf{动词前}）！
  \begin{itemize}
    \item \ding{51}\ \textit{Han har ikke været medlem længe nok, \da{hvilket} desværre \dabold{betyder}, at han ikke får dagpenge.}
  \end{itemize}
  \item \textbf{英丹对比致命雷区（英文 \textit{which / what} 陷阱）}：
    \begin{itemize}
      \item 英语用 \textit{which} 既指代物（\textit{the contract, which...}）又指代整句（\textit{..., which means...}）。但在丹麦语里，\textbf{回指具体名词（无论 \textit{en/et} 还是复数）请老老实实用 \textit{som / der}}；\textbf{只有回指前面整个分句时，才用 \dabold{hvilket}}！
      \item 英语学习者或德语学习者常误用 *\textit{..., hvad betyder, at...} —— \textbf{绝对错误}！丹麦语回指整句只能用 \da{hvilket}，不能用 \textit{hvad}（除非说极少见的固定感叹，标准写作一律用 \textit{hvilket}）。
    \end{itemize}
\end{itemize}
\end{grammatikbox}

\begin{grammatikbox}{语法点 2：介词 + 关系代词（\textit{hvori, hvormed, hvoraf, hvorfor}）vs 口语“介词后置”}
当关系从句中需要“在其中”“借助它”“从中”“因此原因”等介词含义时，丹麦语提供了一套\textbf{书面雅致体}与一套\textbf{日常口语体}，两者都能在 DU3 拿高分，但切忌混杂写错：

\vspace{0.3em}
\renewcommand{\arraystretch}{1.25}
{\small
\begin{tabularx}{\linewidth}{@{} >{\bfseries\raggedright\arraybackslash}p{2.2cm} >{\raggedright\arraybackslash}p{4.3cm} >{\raggedright\arraybackslash}p{4.5cm} >{\raggedright\arraybackslash}X @{}}
\toprule
\textbf{复合关系词} & \textbf{字面拆解与含义} & \textbf{书面高分例句（Modul 3 推荐）} & \textbf{日常自然等价句（\textit{som + 尾介词}）} \\
\midrule
\da{hvori} & \textit{i hvilken/hvilket}\newline （在其中，里面有……） & \textit{Et system, \da{hvori} tryghed og frihed forenes.} & \textit{Et system, \textbf{som} tryghed og frihed forenes \textbf{i}.} \\
\midrule
\da{hvormed} & \textit{med hvilken/hvilket}\newline （借助它，用它来……） & \textit{En ordning, \da{hvormed} man sikrer sin løn.} & \textit{En ordning, \textbf{som} man sikrer sin løn \textbf{med}.} \\
\midrule
\da{hvoraf} & \textit{af hvilken/hvilket}\newline （其中有一部分，由此） & \textit{Ti ansøgere, \da{hvoraf} tre har erfaring.} & \textit{Ti ansøgere, \textbf{som} tre \textbf{af} har erfaring.}（此义多用 \textit{hvoraf}） \\
\midrule
\da{hvorfor} & （……的原因；“为何”） & \textit{Hovedårsagen, \da{hvorfor} folk melder sig ind.} & \textit{Hovedårsagen \textbf{til, at} folk melder sig ind.}（更地道！） \\
\bottomrule
\end{tabularx}
}
\vspace{0.3em}

\textbf{两大易错提醒}：\ding{172} 书面用 \textit{hvori / hvormed} 时，\textbf{从句末尾不能再重复一次 \textit{i / med}}（\ding{55}\ *\textit{kontrakten, hvori min løn står i}）；\ding{173} 表达“……的原因”时，除了 \textit{årsagen, hvorfor...}，丹麦本地人最顺耳的搭配其实是 \da{grunden til, at...} 或 \da{årsagen til, at...}，写作文时可以交替使用！
\end{grammatikbox}

% ---------------------------------------------------------------------
\ncdsubsec{五、Mønsterøvelser（句型操练）}

\textbf{练习 A：B1 句法升格——按要求合并下列句子（注意从句中否定词与副词位置）}
\begin{enumerate}
  \item \textit{Virksomheder kan nemt ansætte nye folk. Det gør arbejdsmarkedet fleksibelt.}（用 \textit{, hvilket...} 合并）$\rightarrow$
  \item \textit{Man får ikke altid dagpenge fra dag ét. Det betyder ofte, at folk sparer penge op.}（用 \textit{, hvilket ofte...} 合并）$\rightarrow$
  \item \textit{Jeg har læst kontrakten igennem. Alle mine vilkår står i den.}（用 \textit{hvori} 合并）$\rightarrow$
  \item \textit{A-kassen er en forsikring. Man kan sikre sin økonomi med den.}（用 \textit{hvormed} 合并）$\rightarrow$
  \item \textit{Der var tyve ansøgere til stillingen. Fem af dem blev inviteret til samtale.}（用 \textit{hvoraf} 合并）$\rightarrow$
  \item \textit{Jeg har en gammel kollega. Jeg talte med ham i går.}（用口语自然式 \textit{som ... med} 合并）$\rightarrow$
\end{enumerate}

\textbf{练习 B：精准辨析选填（从 \textit{der, som, hvilket, hvori, hvorfor, hvad} 中选择最准确一项）}
\begin{enumerate}
  \item Han glemte at melde sig ind i en a-kasse i tide, \underline{\hspace{2.2cm}} desværre kostede ham dyrt.
  \item Kontrakten, \underline{\hspace{2.2cm}} han fik tilsendt i går, indeholder en god pensionsordning.
  \item Fagforeningen er en organisation, \underline{\hspace{2.2cm}} kæmper for medlemmernes rettigheder.
  \item Vi arbejder i et land, \underline{\hspace{2.2cm}} tillid mellem arbejdsgiver og ansat er meget vigtig.
  \item Det er præcis den vigtigste grund til, \underline{\hspace{2.2cm}} mange vælger at skifte job.
  \item En overenskomst er en aftale, \underline{\hspace{2.2cm}} løn og arbejdstid bliver fastlagt.
  \item Man ved aldrig på forhånd, \underline{\hspace{2.2cm}} der sker på arbejdsmarkedet næste år.
  \item Virksomheden mistede tre store kunder, \underline{\hspace{2.2cm}} førte til, at to medarbejdere blev opsagt.
\end{enumerate}

\textbf{练习 C：中译丹实战（练熟 \textit{, hvilket...} 与职场社会核心词）}
\begin{enumerate}
  \item 丹麦没有法定最低工资，这意味着工会扮演着重要的角色。
  \item 他去年加入了失业保险基金，这给了他很大的安全感。
  \item 我收到了一份新合同，里面写明了我的工资和工作时间。（用 \textit{hvori}）
  \item 如果一个人失业了，失业保险基金可以通过救济金提供支持。
  \item 在签订新合同之前，让工会帮你审核一下是明智的。
  \item 我们的团队有十二名成员，其中四名来自国外。（用 \textit{hvoraf}）
\end{enumerate}

% ---------------------------------------------------------------------
\ncdsubsec{六、DU3 Modul 3 实战微任务与参考答案}

\begin{tcolorbox}[enhanced, breakable, colback=NCDBgGray, colframe=NCDGray, coltitle=white, fonttitle=\bfseries, title={DU3 Modul 3 实战微任务（对标 Mundtlig Opgave 1 观点陈述 / 书面短论 · 100+ 词）}, arc=2pt, boxrule=0.8pt, borderline west={3pt}{0pt}{NCDRed}, left=10pt, right=10pt, top=7pt, bottom=7pt]
\small
\textbf{任务情境（Emne: Det danske arbejdsmarked）}：在 Modul 3 口语 Opgave 1 追问环节或社会短文写作中，考官常问：\\
\textit{”Hvorfor tror du, at så mange lønmodtagere i Danmark vælger at være medlem af både en fagforening og en a-kasse? Synes du selv, at det danske system er godt?”}\\[4pt]
请写一篇 \textbf{105--130 词}的 B1 口语陈述稿（或小短论），先说明两者的各自作用，再阐述你对“灵活安全模式”的看法。要求至少使用 \textbf{2 处 \textit{, hvilket ...}} 和 \textbf{1 处 \textit{hvori / hvormed}}！
\end{tcolorbox}

\begin{answerbox}
\textbf{Før du læser 参考回答}
\begin{enumerate}
  \item \textit{En fagforening forhandler løn og arbejdsvilkår gennem overenskomster og giver juridisk hjælp, mens en a-kasse fungerer som en forsikring, der udbetaler dagpenge ved arbejdsløshed. ”Flexicurity-modellen” består i, at det er nemt for virksomheder at ansætte og opsige folk, samtidig med at medarbejderne har et stærkt sikkerhedsnet.}
  \item \textit{Jonas er i tvivl om, hvorvidt man har pligt til at melde sig ind i en fagforening, og om fagforeningen og a-kassen er det samme. Freja råder ham til at sende sin kontrakt til fagforeningen, så de kan tjekke den, før han skriver under.}
\end{enumerate}

\textbf{Tekst B 聊天补全答案}
\begin{enumerate}
  \item \da{B}（回答两者职能区别）
  \item \da{A}（回应如何核对合同条款，并衔接下句“明天签字前照做”）
\end{enumerate}

\textbf{练习 A 答案（B1 句法升格）}
\begin{enumerate}
  \item \textit{Virksomheder kan nemt ansætte nye folk, \da{hvilket} gør arbejdsmarkedet fleksibelt.}
  \item \textit{Man får ikke altid dagpenge fra dag ét, \da{hvilket} ofte betyder, at folk sparer penge op.}
  \item \textit{Jeg har læst kontrakten igennem, \dabold{hvori} alle mine vilkår står.}
  \item \textit{A-kassen er en forsikring, \dabold{hvormed} man kan sikre sin økonomi.}
  \item \textit{Der var tyve ansøgere til stillingen, \dabold{hvoraf} fem blev inviteret til samtale.}
  \item \textit{Jeg har en gammel kollega, \da{som} jeg talte \da{med} i går.}
\end{enumerate}

\textbf{练习 B 答案}
\begin{enumerate}
  \item \da{hvilket}（回指整件事；\textit{desværre} 在动词前）
  \item \da{som}（指代名词 \textit{kontrakten} 作宾语）
  \item \da{der / som}（指代 \textit{organisation} 作主语）
  \item \da{hvori}（“在那个国家里面”，口语可用 \textit{hvor}）
  \item \da{at}（固定句型 \textit{grunden til, at...}；无 \textit{til} 才可接 \textit{hvorfor}）
  \item \da{hvori}（“在协议中规定……”）
  \item \da{hvad}（间接疑问宾语“会发生什么”）
  \item \da{hvilket}（回指“公司失去三个大客户”整件事）
\end{enumerate}

\textbf{练习 C 答案}
\begin{enumerate}
  \item \textit{Danmark har ingen lovbestemt mindsteløn, hvilket betyder, at fagforeningen spiller en vigtig rolle.}
  \item \textit{Han meldte sig ind i en a-kasse sidste år, hvilket gav ham en stor tryghed.}
  \item \textit{Jeg har fået en ny kontrakt, hvori min løn og arbejdstid står.}
  \item \textit{Hvis man bliver arbejdsløs, kan a-kassen støtte med dagpenge.}
  \item \textit{Før man skriver under på en ny kontrakt, er det klogt at lade fagforeningen tjekke den.}
  \item \textit{Vores team har tolv medlemmer, hvoraf fire kommer fra udlandet.}
\end{enumerate}

\textbf{DU3 Modul 3 高分陈述范文（111 词）}\\
\textit{Efter min mening er det danske arbejdsmarked meget velfungerende, fordi det skaber en god balance mellem frihed og sikkerhed. På den ene side kan en virksomhed nemt opsige medarbejdere i svære tider, \da{hvilket} gør økonomien mere dynamisk. På den anden side vælger mange lønmodtagere både en fagforening og en a-kasse for at beskytte sig selv.}

\textit{Fagforeningen sikrer fair overenskomster og hjælper, hvis der opstår problemer på jobbet. A-kassen er derimod en ordning, \dabold{hvormed} man kan modtage dagpenge under ledighed, \da{hvilket} giver ro til at søge det rigtige nye job. Jeg synes personligt, at ”flexicurity-modellen” er et stærkt system, \dabold{hvori} både arbejdsgivere og medarbejdere tager ansvar for, at alle kommer hurtigt videre.}

\small\textbf{名师点评}：111 词紧凑扎实，逻辑呈“\textit{På den ene side ... På den anden side ... Jeg synes personligt...}”经典 B1 论证结构，精准融入两处 \textit{, hvilket...} 与 \textit{hvormed / hvori} 书面关系结构，直接满足 Modul 3 口语深度问答与写作高分标准。
\end{answerbox}

\ncdlesson{3}{Fra børnehave til universitet: Livslang læring}{从幼儿园到大学：终身学习}
{DU3 Modul 3（B1）｜对标 NCE 2 文化教育说明文 \& Modul 3 \textit{Læsning Opgave 4 (Tre personer) \& Mundtlig Opgave 1}}
{\textbf{现在分词（\textit{-ende}）作定语（主动/不变形）} vs \textbf{过去分词（\textit{-et/-ede} 与 \textit{-t/-te}）作定语（被动/需屈折）}· 英语 \textit{-ing / -ed} 分词负迁移排雷}
{教育与终身学习（\textit{uddannelse og efteruddannelse}）是 DU3 Modul 3 口语 \textit{Opgave 1}（自选主题 \textit{Min uddannelse}）与阅读 \textit{Opgave 4} 的核心话题。熟练用 \textbf{\textit{-ende}} 与 \textbf{\textit{-et/-ede}} 分词打包信息，就能把干瘪的从句压缩成地道干练的 B1 定语短语！}

% ---------------------------------------------------------------------
\begin{foerduleaeser}
带着下面两个问题阅读 Tekst A 和 Tekst B：
\begin{enumerate}
  \item \textit{Hvilke to hovedveje kan unge i Danmark vælge efter folkeskolen, og hvorfor spiller efteruddannelse en vigtig rolle for voksne?}（丹麦年轻人在义务教育毕业后可以选择哪两条主要路径？为什么继续教育对成年人也很重要？）
  \item \textit{Hvilken uddannelse har Mikkel, Sofia og Omar valgt i Tekst B, og hvordan får hver af dem støtte til hverdagen eller økonomien?}（在 Tekst B 中，米克尔、索菲亚和奥马尔分别选择了什么教育？他们各自如何在生活或经济上获得支持？）
\end{enumerate}
\end{foerduleaeser}

% ---------------------------------------------------------------------
\ncdsubsec{二、Tekst A \& Tekst B（社会深读与 Modul 3 真题场景）}

\subsection*{Tekst A：\textit{Fra børnehave til universitet: Livslang læring}（社会教育短文 · 约 240 词）}

\begin{tipbox}{语篇说明}
本篇梳理丹麦从幼儿教育到成人终身学习的完整版图。阅读时请对比两种修饰语：\textbf{\color{NCDNavy}深蓝色加粗的现在分词（\textit{-ende}）}（词尾永远不变，表主动或“令人……”）与\textbf{\color{NCDRed}红色加粗的过去分词（\textit{-et/-ede/-t/-te}）}（表被动或“已被……”，复数与定指必须变形！）。
\end{tipbox}

I Danmark begynder vejen gennem uddannelsessystemet allerede i børnehaven, hvor \dabold{engagerede} pædagoger skaber et trygt og \da{stimulerende} miljø gennem leg og fællesskab. Når børnene fylder seks år, starter de i folkeskolen, som lægger stor vægt på \da{udfordrende} gruppearbejde og kritisk tænkning frem for blot at lære udenad.

Efter folkeskolen står de unge over for et \da{spændende} valg. Nogle vælger et treårigt gymnasium for at få adgang til universitetet, mens andre tager en praktisk erhvervsuddannelse, der kombinerer skoleundervisning med en \dabold{lønnet} læreplads i en virksomhed. Begge veje kan føre til \dabold{velbetalte} og meningsfulde stillinger, og fordi alle \dabold{optagne} studerende får statslig uddannelsesstøtte (SU), kan \dabold{motiverede} unge fra alle samfundslag gennemføre en videregående uddannelse uden \da{stigende} studieafgifter.

Uddannelse slutter dog langt fra, når man står med sit eksamensbevis i hånden. På et arbejdsmarked med \da{voksende} krav om grøn omstilling og digitale kompetencer spiller efteruddannelse en afgørende rolle for \dabold{uddannede} voksne i alle aldre. Uanset om man er en \dabold{nyuddannet} sygeplejerske eller en erfaren tømrer, er \dabold{opdaterede} færdigheder nødvendige for at følge med i udviklingen. Derudover vælger mange voksne at gå på aftenskole eller tage et ophold på en dansk folkehøjskole --- en \da{fascinerende} skoleform uden eksamener, hvor \da{lærende} i alle aldre bor og studerer sammen. Her handler undervisningen lige så meget om personlig dannelse og \dabold{delte} oplevelser som om nye faglige emner.

For danskerne er læring kort sagt ikke kun en forberedelse til det første job, men en \da{berigende} og livslang rejse for alle \dabold{interesserede} borgere.

\begin{transbox}
\small
\textbf{\color{NCDTeal}[Tekst A 中文精译]}\\
在丹麦，教育体系之路早在幼儿园便已开启，那里富有热忱的幼教老师通过游戏和集体生活营造出安全而富有启发性的环境。年满六岁时，孩子们进入国民基础学校（\textit{folkeskole}），学校非常注重富有挑战性的小组合作与批判性思考，而不仅仅是死记硬背。

基础学校毕业后，年轻人面临一个激动人心的选择。一部分人选择三年制普通高中以获得进入大学的资格，另一些人则选择实用型职业教育，将课堂教学与企业里的带薪学徒岗位结合起来。两条道路都能通往高薪且有意义的职位；而且由于所有被录取的大学生都能获得国家助学金（SU），来自各个社会阶层的有志青年都无需面对不断上涨的学费，便能完成高等教育。

然而，手握毕业证书远不意味着教育的终结。在一个对绿色转型和数字能力要求不断增长的劳动力市场上，继续教育对各个年龄段已经受过教育的成年人来说都起着决定性作用。无论你是刚毕业的护士，还是经验丰富的木匠，保持更新的技能对于跟上时代发展都是必不可少的。此外，许多成年人还选择上夜校，或在丹麦民众高等学院（\textit{folkehøjskole}）住读一段时间——这是一种不设考试的迷人学校形式，各年龄段的学习者同吃同住、共同学习。在这里，教学既关乎新知识，也同样关乎个人修养与共享的人生体验。

简而言之，对丹麦人而言，学习不只是为了第一份工作的准备，更是所有心怀兴趣的公民的一场充实而终身的旅程。
\end{transbox}

\subsection*{Tekst B：\textit{Tre personer om deres uddannelsesvalg}（对标 DU3 Modul 3 Læsning Opgave 4 三人对比）}

\begin{tcolorbox}[enhanced, breakable, colback=NCDBgBlue, colframe=NCDNavy, coltitle=white, fonttitle=\bfseries, title={人物独白对比（阅读以下三位学习者的自述，完成下方 A / B / C 配对）}, arc=2pt, boxrule=0.8pt, left=10pt, right=10pt, top=7pt, bottom=7pt]
\small
\textbf{A. Mikkel (22 år)}: Jeg valgte en erhvervsuddannelse som elektriker, fordi jeg trives bedst med praktisk arbejde. Allerede fra andet år fik jeg en \dabold{lønnet} læreplads i et lokalt firma, hvilket gav mig både erfaring og en god økonomi.\\[4pt]
\textbf{B. Sofia (29 år)}: Efter fem år som sælger trængte jeg til en \da{spændende} forandring. Derfor søgte jeg ind på et \dabold{forkortet} it-kursus på aftenskole, mens jeg stadig arbejdede på deltid. Det var \da{krævende}, men meget \da{motiverende}.\\[4pt]
\textbf{C. Omar (24 år)}: Jeg læser til læge på universitetet og modtager SU hver måned. Selvom pensum er stort, møder jeg mange \dabold{engagerede} undervisere og \da{hjælpsomme} medstuderende, og vi køber ofte \dabold{brugte} lærebøger for at spare penge.
\end{tcolorbox}

\textbf{Tekst B 阅读配对（为每题选出对应人物 A、B 或 C，答案见第六部分）：}
\begin{enumerate}
  \item \textit{Hvem kombinerer et deltidsjob med undervisning om aftenen?} \underline{\hspace{2.0cm}}
  \item \textit{Hvem har fået betaling for at arbejde i en virksomhed som en del af sin uddannelse?} \underline{\hspace{2.0cm}}
  \item \textit{Hvem sparer penge på studiematerialer og får økonomisk støtte fra staten?} \underline{\hspace{2.0cm}}
\end{enumerate}

\begin{transbox}
\small
\textbf{\color{NCDTeal}[Tekst B 中文对照]}\\
\textbf{A. 米克尔（22 岁）}：我选择了电工职业教育，因为我最适合动手实践的工作。从第二年起，我就在当地一家公司拿到了带薪学徒名额，这既给了我经验，也带来了不错的收入。\\
\textbf{B. 索菲亚（29 岁）}：做了五年销售后，我渴望来一次令人振奋的改变。于是我报名了一门夜校的浓缩 IT 课程，同时仍然兼职工作。这虽然很花精力，但非常有激励感。\\
\textbf{C. 奥马尔（24 岁）}：我在大学攻读医学专业，每月领取国家助学金 SU。虽然课业繁重，但我遇到了许多投入敬业的讲师和乐于助人的同学，而且我们常买二手教材来省钱。
\end{transbox}

\begin{tipbox}{本课发音与吞音要点（Udtale \& Reduktioner）}
\begin{itemize}
  \item \textbf{现在分词 \textit{-ende} 的必杀铁律——中间的 \textit{d} 永远彻底静音！}：绝不要读出 /d/ 音，而是顺滑连成 \ipa{ənə}（或 \ipa{nə}）。例如：\textit{spændende} \ipa{ˈsbɛnənə}、\textit{voksende} \ipa{ˈvʌɡsənə}、\textit{stigende} \ipa{ˈsdiːənə}（\textbf{g 与 d 都不发音}！）、\textit{studerende} \ipa{sduˈdeːˀɐnə}、\textit{motiverende} \ipa{motˢiˈveːˀɐnə}。
  \item \textbf{过去分词 \textit{-et} vs \textit{-ede} 对比}：\textit{uddannet} \ipa{ˈuðˌdanˀəð} $\rightarrow$ 复数/定指 \textit{uddannede} \ipa{ˈuðˌdanˀəðə}（柔滑软 d）；\textit{forkortet} \ipa{fʌˈkʰʌɐ̯dəð} $\rightarrow$ \textit{forkortede} \ipa{fʌˈkʰʌɐ̯dəðə}；\textit{brugt} \ipa{ˈbʁʌɡd} $\rightarrow$ \textit{brugte} \ipa{ˈbʁʌɡdə}。
  \item \textbf{教育机构核心词}：\textit{en børnehave} \ipa{ˈbɶɐ̯nəˌhæːvə}、\textit{en folkeskole} \ipa{ˈfʌlɡəˌsgoːlə}、\textit{et gymnasium} \ipa{ɡymˈnæːˀsium}（\textit{g} 读硬音 \ipa{ɡ}，重音在第二音节！）、\textit{en erhvervsuddannelse} \ipa{ɛɐ̯ˈvɛɐ̯ˀvsˌuðdanˀəlsə}（第一音节的 h 不发音！）。
\end{itemize}
\end{tipbox}

% ---------------------------------------------------------------------
\ncdsubsec{三、Gloser（核心生词表）}

\begin{tipbox}{词表重点}
注意 \textit{en studerende} 是由现在分词直接名词化的词，单数不定指为 \textit{en studerende}，单数定指为 \textit{den studerende}，复数不定指仍为 \textit{studerende}，复数定指为 \textit{de studerende}（切忌再加 \textit{-r}！）。另外，“读大学某专业”的地道说法是 \textit{at læse til læge / ingeniør}（或 \textit{at studere medicin}）。
\end{tipbox}

\begin{glosertable}
名词 & en børnehave & \textit{(børnehaven, børnehaver, -ne)} & 幼儿园 \\
名词 & en folkeskole & \textit{(folkeskolen, -skoler, -ne)} & 丹麦九年制公立基础学校 \\
名词 & et gymnasium & \textit{(gymnasiet, gymnasier, -erne)} & 普通高中（三年制） \\
名词 & en erhvervsuddannelse & \textit{(erhvervsuddannelsen, -r, -rne)} & 职业教育，技工教育 \\
名词 & en læreplads & \textit{(lærepladsen, lærepladser, -erne)} & 学徒岗位，实习职位 \\
名词 & en efteruddannelse & \textit{(efteruddannelsen, -r, -rne)} & 继续教育，在职进修 \\
名词 & en folkehøjskole & \textit{(folkehøjskolen, -skoler, -ne)} & 丹麦民众高等学院（无考试住读） \\
名词 & en adgang & \textit{(adgangen, kun ental)} & 入学资格；进入途径（\textit{adgang til}） \\
名词 & en studerende & \textit{(den studerende, studerende, de -)} & 大学生，在读学生（分词名词化） \\
名词 & en færdighed & \textit{(færdigheden, færdigheder, -erne)} & 技能，本领 \\
\midrule
动词 [har] & at uddanne (sig) & \textit{(uddanner, uddannede, har uddannet)} & 培养，教育；进修（反身） \\
动词 [har] & at undervise (i) & \textit{(underviser, underviste, har undervist)} & 教授，授课（\textit{undervise i dansk}） \\
动词 [har] & at gennemføre & \textit{(gennemfører, -førte, har gennemført)} & 完成，修完，贯彻执行 \\
动词 [har] & at opdatere & \textit{(opdaterer, opdaterede, har opdateret)} & 更新，升级 \\
\midrule
现在分词 & spændende / stigende & \textit{(现在分词作定语，词尾永不变)} & 令人兴奋的 / 不断上涨的 \\
过去分词 & uddannet / brugt & \textit{uddannet (-, uddannede)} / \textit{brugt (-, brugte)} & 受过教育的 / 二手的，用过的 \\
\end{glosertable}

% ---------------------------------------------------------------------
\ncdsubsec{四、Grammatik \& Faste vendinger（核心语法与地道表达精讲）}

\begin{grammatikbox}{语法点 1：现在分词（\textit{-ende}）vs 过去分词（\textit{-et / -ede / -t / -te}）作定语全景表}
想在 Modul 3 写出紧密凝练的名词短语，必须精准掌握两种分词作前置定语时的\textbf{语态分工}与\textbf{词尾屈折规则}：

\vspace{0.3em}
\renewcommand{\arraystretch}{1.25}
{\small
\begin{tabularx}{\linewidth}{@{} >{\bfseries\raggedright\arraybackslash}p{2.3cm} >{\raggedright\arraybackslash}p{4.0cm} >{\raggedright\arraybackslash}p{4.4cm} >{\raggedright\arraybackslash}X @{}}
\toprule
\textbf{分词类型} & \textbf{词尾变化铁律} & \textbf{语义核心（主动 vs 被动）} & \textbf{高频搭配例词} \\
\midrule
\da{现在分词}\newline （\textit{-ende}） &
动词词干 + \da{-ende}\newline \textbf{永远不变形}！不加 \textit{-t} 也不加 \textit{-e} &
\textbf{主动 / 进行 / “令人……的”}\newline 修饰“发出动作”或“带来该感受”的事物 &
\textit{et spændend\da{e} valg}（令人兴奋的选择）\newline
\textit{voksend\da{e} krav}（不断增长的要求）\newline
\textit{en smilend\da{e} lærer}（微笑着的老师） \\
\midrule
\dabold{过去分词 I}\newline （以 \textit{-et} 结尾） &
不定单数（\textit{en/et} 均用）：\dabold{-et}\newline 复数或定指（\textit{de/den/det/min}）：\dabold{-ede} &
\textbf{被动 / 完成 / “已被……”}\newline 修饰“承受动作”或“处于某状态”的人或物（多源自 1 类规则动词） &
\textit{en motiver\dabold{et} elev} / \textit{de motiver\dabold{ede} elever}\newline
\textit{et forkort\dabold{et} kursus} / \textit{opdater\dabold{ede} færdigheder} \\
\midrule
\dabold{过去分词 II}\newline （以 \textit{-t} 结尾） &
不定单数（\textit{en/et} 均用）：\dabold{-t}\newline 复数或定指（\textit{de/den/det/min}）：\dabold{-te} &
\textbf{被动 / 完成 / “已被……”}\newline 修饰“承受动作”或“处于某状态”的人或物（多源自 2 类规则动词） &
\textit{en brug\dabold{t} bog} / \textit{brug\dabold{te} bøger}\newline
\textit{et velbetal\dabold{t} job} / \textit{velbetal\dabold{te} stillinger} \\
\bottomrule
\end{tabularx}
}
\vspace{0.3em}

\textbf{额外识别提示（不规则强变化 \textit{-en / -et} 分词作定语）}：少数强变化动词作复数/定指前置定语时词尾变 \textit{-ne}（而非 \textit{-ede}），如 \textit{skrive $\rightarrow$ en skreven regel / de skrev\dabold{ne} regler}；\textit{stjæle $\rightarrow$ en stjålen cykel / de stjål\dabold{ne} cykler}；\textit{dele}（1 类弱变化）$\rightarrow$ \textit{del\dabold{te} oplevelser}（读作规则 \textit{-t/-te}）。
\end{grammatikbox}

\begin{grammatikbox}{语法点 2：中/英对比排雷——切勿把英语 \textit{-ing / -ed} 规则直接套进丹麦语！}
\begin{enumerate}
  \item \textbf{雷区一｜丹麦语绝对没有“\textit{er + -ende} 进行时”！}：英语说 \textit{I am studying Danish}，千万不能在丹麦语写成 *\textit{Jeg er studerende dansk}！丹麦语表达“正在做某事”只需用\textbf{一般现在时}（\textit{Jeg læser dansk}）或结构 \da{Jeg er i gang med at... / Jeg er ved at...}。在丹麦语里，\textit{-ende} 只作\textbf{定语、少数表语（\textit{Bogen er spændende}）或指人名词（\textit{en studerende}）}！
  \item \textbf{雷区二｜“让人感到……”（\textit{-ende}）vs “自己感到……”（\textit{-et/-t}）}：
    \begin{itemize}
      \item \textit{en \da{skuffende} karakter}（令人失望的分数——分数导致失望） vs \textit{en \dabold{skuffet} studerende}（感到失望的学生——人承受失望）
      \item \textit{en \da{inspirerende} lærer}（有启发性的老师） vs \textit{en \dabold{engageret} pædagog}（全身心投入的幼教）
    \end{itemize}
  \item \textbf{雷区三｜过去分词单复数忘变词尾}：英语的 \textit{used} 无论修饰单数还是复数都长一样（\textit{a used book / used books}）；丹麦语必须严格屈折：\textit{en brug\dabold{t} bog} $\leftrightarrow$ \textit{brug\dabold{te} bøger}；\textit{en uddann\dabold{et} læge} $\leftrightarrow$ \textit{de uddann\dabold{ede} læger}！
\end{enumerate}
\end{grammatikbox}

% ---------------------------------------------------------------------
\ncdsubsec{五、Mønsterøvelser（句型操练）}

\textbf{练习 A：词形转化——将括号内的动词变为正确的现在分词（\textit{-ende}）或过去分词定语（注意单复数与定指屈折！）}
\begin{enumerate}
  \item Hun er en meget \underline{\hspace{2.6cm}} (\textit{at engagere}，过去分词) underviser på universitetet.
  \item Vi købte tre \underline{\hspace{2.6cm}} (\textit{at bruge}，过去分词) lærebøger på nettet i går.
  \item Virksomheden søger en \underline{\hspace{2.6cm}} (\textit{at uddanne}，过去分词) elektriker med erfaring.
  \item Der er \underline{\hspace{2.6cm}} (\textit{at vokse}，现在分词) behov for digitale kompetencer i dag.
  \item Alle \underline{\hspace{2.6cm}} (\textit{at interessere}，过去分词) borgere kan melde sig til et gratis aftenskolekursus.
  \item Han drømmer om et \underline{\hspace{2.6cm}} (\textit{at spænde}，现在分词) og \underline{\hspace{2.6cm}} (\textit{at velbetale} $\rightarrow$ \textit{vel- + betale}，过去分词) job.
  \item De \underline{\hspace{2.6cm}} (\textit{at optage}，过去分词，注意不规则！) studerende modtager et velkomstbrev i juli.
  \item På grund af den \underline{\hspace{2.6cm}} (\textit{at stige}，现在分词) efterspørgsel åbner skolen et nyt hold.
\end{enumerate}

\textbf{练习 B：语义辨析与英式错误纠正（选出正确一项，或改正句子中的英语迁移错误）}
\begin{enumerate}
  \item 选择正确形式：Eleverne var meget \underline{\hspace{2.6cm}} (\textit{motiverende / motiverede}), fordi læreren holdt et \underline{\hspace{2.6cm}} (\textit{inspirerende / inspireret}) oplæg.
  \item 选择正确形式：Det var en \underline{\hspace{2.6cm}} (\textit{skuffende / skuffet}) eksamen for hele klassen.
  \item 选择正确形式：Alle medarbejdere skal have \underline{\hspace{2.6cm}} (\textit{opdateret / opdaterede}) oplysninger i systemet.
  \item 改正中/英式错误：\textit{*Jeg er læsende en spændende bog om dansk kultur lige nu.} $\rightarrow$
  \item 改正屈折错误：\textit{*Vi har ansat to nyuddannet ingeniører fra Aarhus.} $\rightarrow$
  \item 改正屈折错误：\textit{*Han har fundet et spændendet studie på universitetet.} $\rightarrow$
\end{enumerate}

\textbf{练习 C：中译丹实战（活用两种分词定语与教育主题词汇）}
\begin{enumerate}
  \item 在丹麦，所有被录取的大学生都可以领取国家助学金。
  \item 我们班有一位令人鼓舞、富有热忱的老师。
  \item 许多成年人上夜校是为了获得更新的技能。
  \item 为了省钱，大学生们常常购买二手的教材。
  \item 对实用型职业教育的需求正在不断增长。（用 \textit{stigende}）
  \item 他刚刚完成了一门令人兴奋的继续教育课程。
\end{enumerate}

% ---------------------------------------------------------------------
\ncdsubsec{六、DU3 Modul 3 实战微任务与参考答案}

\begin{emailbox}{DU3 Modul 3 微型写作任务（对标 Skrivning Opgave 2 · 私人回信 100+ 词）}
\textbf{Situation}: Din veninde Laura overvejer at tage en ny uddannelse eller et aftenskolekursus i Danmark og har spurgt dig om, hvordan det danske uddannelsessystem fungerer.\\
\textbf{Opgave}: Skriv et svar på \textbf{105--130 ord}. Beskriv mulighederne (fx erhvervsuddannelse, universitet med SU eller efteruddannelse), fortæl kort om din egen uddannelse eller et kursus, du har taget, og brug mindst \textbf{2 \textit{-ende}分词定语} og \textbf{2 过去分词定语（含复数/定指变化）}！
\end{emailbox}

\begin{answerbox}
\textbf{Før du læser 参考回答}
\begin{enumerate}
  \item \textit{Efter folkeskolen kan unge vælge mellem et almindeligt gymnasium (adgang til universitetet) eller en praktisk erhvervsuddannelse med lønnet læreplads. Efteruddannelse er vigtig for voksne, fordi arbejdsmarkedet har voksende krav om digitale og grønne kompetencer.}
  \item \textit{Mikkel har valgt en erhvervsuddannelse som elektriker og får løn på sin læreplads. Sofia tager et it-kursus på aftenskole, mens hun arbejder på deltid. Omar læser til læge på universitetet, får SU og køber brugte lærebøger.}
\end{enumerate}

\textbf{Tekst B 阅读配对答案}
\begin{enumerate}
  \item \da{B}（Sofia：\textit{på aftenskole, mens jeg stadig arbejdede på deltid}）
  \item \da{A}（Mikkel：\textit{en lønnet læreplads i et lokalt firma}）
  \item \da{C}（Omar：\textit{modtager SU hver måned ... køber ofte brugte lærebøger for at spare penge}）
\end{enumerate}

\textbf{练习 A 答案（分词词形转化）}
\begin{enumerate}
  \item \dabold{engageret}（修饰单数不定指 \textit{en ... underviser}，词尾 \textit{-et}）
  \item \dabold{brugte}（修饰复数 \textit{tre ... lærebøger}，由 \textit{brugt} 变复数 \textit{-te}）
  \item \dabold{uddannet}（修饰单数不定指 \textit{en ... elektriker}）
  \item \da{voksende}（现在分词词尾永远不变，哪怕后接中性名词 \textit{behov}）
  \item \dabold{interesserede}（修饰复数 \textit{alle ... borgere}，由 \textit{-et} 变 \textit{-ede}）
  \item \da{spændende}；\dabold{velbetalt}（现在分词不变；过去分词修饰中性 \textit{et ... job} 保持 \textit{-t}）
  \item \dabold{optagne}（不规则过去分词 \textit{optaget} 的复数定指形）
  \item \da{stigende}（现在分词定语“不断上涨的”，永不变形）
\end{enumerate}

\textbf{练习 B 答案}
\begin{enumerate}
  \item \dabold{motiverede}（学生们“感到被激励”，复数加 \textit{-ede}）；\da{inspirerende}（报告本身“富有启发性”）
  \item \da{skuffende}（一场“令人失望的”考试，用主动 \textit{-ende}）
  \item \dabold{opdaterede}（修饰复数 \textit{oplysninger}，变 \textit{-ede}）
  \item \textit{Jeg \da{læser} (eller \da{er ved at læse}) en spændende bog om dansk kultur lige nu.}（取消中英混杂进行时 *\textit{er læsende}）
  \item \textit{Vi har ansat to \dabold{nyuddannede} ingeniører fra Aarhus.}（复数加 \textit{-ede}）
  \item \textit{Han har fundet et \da{spændende} studie på universitetet.}（现在分词 \textit{-ende} 绝对不能加 \textit{-t}！）
\end{enumerate}

\textbf{练习 C 答案}
\begin{enumerate}
  \item \textit{I Danmark kan alle optagne studerende få statslig uddannelsesstøtte (SU).}
  \item \textit{I vores klasse har vi en inspirerende og engageret lærer.}
  \item \textit{Mange voksne går på aftenskole for at få opdaterede færdigheder.}
  \item \textit{For at spare penge køber studerende ofte brugte lærebøger.}
  \item \textit{Efterspørgslen efter en praktisk erhvervsuddannelse er stadig stigende.}
  \item \textit{Han har lige gennemført et spændende efteruddannelseskursus.}
\end{enumerate}

\textbf{DU3 Modul 3 高分范文（113 词，不含署名）}\\
\textit{Hej Laura}

\textit{Hvor er det dejligt at høre, at du overvejer at studere videre! I Danmark findes der heldigvis mange gode muligheder for livslang læring. Unge kan vælge mellem et gymnasium og en praktisk erhvervsuddannelse med en \dabold{lønnet} læreplads, og alle \dabold{optagne} studerende på universitetet kan få SU, hvilket gør det muligt at koncentrere sig om studiet.}

\textit{For voksne med \da{voksende} familieforpligtelser findes der også \dabold{forkortede} kurser og fleksibel efteruddannelse om aftenen. Sidste år gennemførte jeg selv et \da{spændende} sprogkursus med meget \dabold{engagerede} undervisere. Vi arbejdede med \dabold{opdaterede} tekster, og det var en virkelig \da{motiverende} oplevelse, som hjalp mig videre i mit job.}

\textit{Jeg kan varmt anbefale dig at søge ind!}

\textit{Kærlig hilsen\\
\mbox{[Dit navn]}}

\small\textbf{名师点评}：113 词（不含署名）流畅覆盖全部要点！自然嵌合现在分词（\textit{voksende, spændende, motiverende}）与过去分词定语（\textit{lønnet, optagne, forkortede, engagerede, opdaterede}），且精准回扣 L02 整句关系从句（\textit{, hvilket gør det muligt...}），体现出螺旋式递进的 B1 书面语实力。
\end{answerbox}

\ncdlesson{4}{Når arbejdsdagen slutter klokken fire}{当下午四点下班铃响起}
{DU3 Modul 3（B1）｜对标 NCE 2 社会习惯剖析 \& Modul 3 \textit{Læsning Opgave 3 \& Mundtlig Opgave 2}}
{\textbf{省略 \textit{hvis} 的条件倒装句（V1 条件从句）}（\textit{Har man børn, skal man hente dem inden kl. 17}）· \textbf{条件从句+主句“双倒装”}与否定词 \textit{ikke} 位置铁律}
{下午四点下班去托儿所接孩子、晚上参加社团活动、不盲目赞美加班——这是理解丹麦“工作与生活平衡（\textit{work-life balance}）”的关键窗口。语法上，把冗长的 \textit{Hvis man...} 压缩成“\textbf{动词抢占句首}”的条件倒装句，瞬间就能写出地道丹麦语的节奏感！}

% ---------------------------------------------------------------------
\begin{foerduleaeser}
带着下面两个问题阅读 Tekst A 和 Tekst B：
\begin{enumerate}
  \item \textit{Hvorfor går mange danske medarbejdere hjem allerede ved firetiden, og hvordan ser kolleger ofte på meget overarbejde?}（为什么许多丹麦员工四点左右就下班回家？同事们通常如何看待频繁加班？）
  \item \textit{Hvilke to fleksible ordninger tilbyder virksomheden NordTek sine ansatte ifølge Tekst B?}（根据 Tekst B，NordTek 公司向员工提供了哪两项灵活工作制度？）
\end{enumerate}
\end{foerduleaeser}

% ---------------------------------------------------------------------
\ncdsubsec{二、Tekst A \& Tekst B（社会深读与 Modul 3 真题场景）}

\subsection*{Tekst A：\textit{Når arbejdsdagen slutter klokken fire}（社会观察短文 · 约 245 词）}

\begin{tipbox}{语篇说明}
本篇深入丹麦家庭与职场的时间节拍。阅读时请特别关注\textbf{\color{NCDNavy}深蓝色加粗的“省略 \textit{hvis}”句首动词}与紧跟在逗号后面的\textbf{\color{NCDRed}红色加粗的主句倒装动词}——两半句竟然都是“动词在主语前”！
\end{tipbox}

Hvis man besøger et dansk kontor klokken halv fire om eftermiddagen, bliver man ofte overrasket. Mange medarbejdere pakker allerede deres computere sammen, tager jakken på og går hjem. I Danmark varer en normal arbejdsuge på fuld tid 37 timer, og de fleste starter tidligt om morgenen for at kunne gå tidligt.

\da{Har man} små børn, \dabold{skal man} nemlig nå at hente dem i vuggestuen eller børnehaven, inden institutionen lukker klokken fem. Mange familier spiser også aftensmad allerede omkring klokken seks, så børnene kan komme i seng i ordentlig tid. Samtidig spiller fritiden en kæmpe rolle i det danske samfund. \da{Går man} til sport i en forening eller \da{dyrker man} musik i et kor, \dabold{begynder} aktiviteterne ofte allerede klokken atten.

Derfor ser danskerne sjældent på overarbejde som et tegn på særlig flid. \da{Bliver en kollega} ofte siddende på kontoret til klokken otte om aftenen, \dabold{tænker andre} snarere, at arbejdsbyrden er for stor, eller at dagen ikke er planlagt effektivt nok.

På mange arbejdspladser har man i stedet indført fleksibel arbejdstid. \da{Har man} en dag brug for at hente børn tidligt eller gå til lægen, \dabold{kan man} arbejde lidt mindre om eftermiddagen og senere svare på et par mails hjemmefra, når børnene sover. Omvendt forventes det også, at man koncentrerer sig fuldt ud i arbejdstiden. \da{Planlægger man} \dabold{ikke} sin arbejdsdag godt nok, \dabold{bliver det} svært at nå alle opgaver før fyraften. Netop denne tillidsbaserede balance mellem arbejde og fritid gør, at mange medarbejdere oplever mindre stress i hverdagen.

\begin{transbox}
\small
\textbf{\color{NCDTeal}[Tekst A 中文精译]}\\
如果一个人在下午三点半走进一家丹麦办公室，常常会大吃一惊。许多员工已经在收拾电脑、穿上外套准备回家了。在丹麦，标准全职工作周为 37 小时，大多数人早晨很早开始工作，就是为了能早点下班。

因为如果家里有年幼的孩子，人就得赶在托儿所或幼儿园五点关门前把他们接走。许多家庭也会在六点左右就吃晚饭，好让孩子们按时上床睡觉。与此同时，业余时间在丹麦社会中占有极其重要的地位。如果你在俱乐部踢球练体育，或者在合唱团玩音乐，活动往往晚上六点就已经开始了。

因此，丹麦人很少把加班看作格外勤奋的象征。要是某位同事经常在办公室待到晚上八点，其他人反而会认为是不是工作负荷过重，或者当天的时间规划得不够高效。

相反，许多工作单位引入了弹性工时制度。如果某天你需要早点去接孩子或看医生，你可以下午少工作一会儿，等晚上孩子睡着后在家回复几封邮件。反之，公司也期待员工在工作时段内全神贯注。如果一个人没有把工作日规划妥当，就很难在下班前完成全部任务。正是这种基于信任的工作与生活平衡，让许多员工在日常生活中感受到更少的压力。
\end{transbox}

\subsection*{Tekst B：\textit{Fleksible regler hos NordTek}（对标 DU3 Modul 3 Læsning Opgave 3 段落补全）}

\begin{tcolorbox}[enhanced, breakable, colback=NCDBgBlue, colframe=NCDNavy, coltitle=white, fonttitle=\bfseries, title={企业员工手册节选（阅读两段短文，各自从 A--D 中选出 1 句最契合上下文的句子）}, arc=2pt, boxrule=0.8pt, left=10pt, right=10pt, top=7pt, bottom=7pt]
\small
\textbf{1. Fleksibel arbejdstid og hjemmearbejde}\\
Hos NordTek mener vi, at et godt arbejdsliv hænger tæt sammen med tid til familien. Derfor er det kun i særlige perioder, at der arbejdes over på kontoret. \textbf{(1)} \underline{\hspace{4.5cm}} Derudover har alle fastansatte ret til op til to hjemmearbejdsdage om ugen efter aftale med teamet.

\begin{itemize}
  \item[\textbf{A.}] \textit{Alligevel må ingen medarbejdere spise frokost i mødelokalerne.}
  \item[\textbf{B.}] \textit{Vil man hellere møde klokken syv og gå klokken tre, kan man frit aftale det.}
  \item[\textbf{C.}] \textit{Desværre betaler virksomheden ikke for offentlig transport i weekenden.}
  \item[\textbf{D.}] \textit{Derfor blev alle medarbejdere nødt til at arbejde indtil klokken ni hver dag.}
\end{itemize}

\vspace{0.3em}
\textbf{2. Når hverdagen med børn er uforudsigelig}\\
Vi ved også, at familielivet kræver fleksibilitet, især når man har små børn i daginstitution. \textbf{(2)} \underline{\hspace{4.5cm}} På den måde slipper forældre for unødig stress, og kollegerne ved, hvornår man er tilbage igen.

\begin{itemize}
  \item[\textbf{A.}] \textit{Bliver dit barn sygt om morgenen, kan du derfor tage ”barnets første sygedag” med fuld løn.}
  \item[\textbf{B.}] \textit{Således må man kun holde sommerferie i de tre sidste uger af august.}
  \item[\textbf{C.}] \textit{Selvom vejret er dårligt, foretrækker de fleste at cykle hele vejen hjem.}
  \item[\textbf{D.}] \textit{Fordi kontoret lukker sent, kan børnene sove i kantinen om eftermiddagen.}
\end{itemize}
\end{tcolorbox}

\begin{transbox}
\small
\textbf{\color{NCDTeal}[Tekst B 中文对照]}\\
\textbf{1. 弹性工作时间与居家办公}：在 NordTek，我们认为良好的职业生活与陪伴家人的时间密不可分。因此，办公室仅在特殊时期才需要加班。\textbf{（空 1）} 此外，所有正式员工与团队协商后，每周有权享受最多两天居家办公。\\
\textbf{2. 当育儿日常出现突发状况}：我们也理解，家庭生活需要灵活性，尤其是家里有孩子上托儿机构的时候。\textbf{（空 2）} 这样一来，父母能免于不必要的焦虑，同事们也清楚你何时能回归岗位。
\end{transbox}

\begin{tipbox}{本课发音与吞音要点（Udtale \& Reduktioner）}
\begin{itemize}
  \item \textbf{V1 条件倒装句开头的“轻快助动词连读”}：当 \textit{Har / Skal / Vil / Kan / Bliver} 抢占句首接 \textit{man} 时，助动词常快速轻读并粘连后面的 \textit{man}：\textit{Har man} \ipa{ˈhɑːˀ man}（或快读 \ipa{hɑ man}）、\textit{Skal man} \ipa{sga man}（\textbf{l 不发音}！）、\textit{Vil man} \ipa{ve man}（\textbf{l 不发音}！）、\textit{Bliver man} \ipa{bliɐ̯ man}。
  \item \textbf{家庭与时间高频词}：\textit{klokken fire} \ipa{ˈkʰlʌɡən ˈfiːɐ}；\textit{klokken atten} \ipa{ˈkʰlʌɡən ˈadən}；\textit{vuggestue} \ipa{ˈvɔɡəˌsduːə}；\textit{aftensmad} \ipa{ˈɑfdənsˌmað}（结尾软 d）；\textit{arbejdsbyrde} \ipa{ˈɑːbɑi̯ˀdsˌbyɐ̯də}。
  \item \textbf{注意静音辅音}：\textit{sjældent} \ipa{ˈɕɛlənˀd}（\textbf{d 不发音}，\textit{sj-} 读如轻柔 \ipa{ɕ}）；\textit{effektiv} \ipa{efɛɡˈtˢiːˀʊ̯}；\textit{hjemmefra} \ipa{ˈjɛməˌfʁɑːˀ}（词首 \textbf{h 静音}！）。
\end{itemize}
\end{tipbox}

% ---------------------------------------------------------------------
\ncdsubsec{三、Gloser（核心生词表）}

\begin{tipbox}{词表重点}
丹麦劳动制度中有一个极具特色的词：\textit{barnets første sygedag}（子女病假第一天），即孩子突然生病时，家长可以带薪请假一天在家照顾孩子并安排后续照料。此外，\textit{at nå} 既可接直接宾语表“赶上（时间/车）”，也可接 \textit{at +} 动词原形表“来得及做某事”（如 \textit{nå at hente børn}）。
\end{tipbox}

\begin{glosertable}
名词 & en arbejdsuge & \textit{(arbejdsugen, arbejdsuger, -ne)} & 工作周 \\
名词 & et overarbejde & \textit{(overarbejdet, kun ental)} & 加班，超时工作 \\
名词 & en arbejdsbyrde & \textit{(arbejdsbyrden, -byrder, -ne)} & 工作负荷，工作量 \\
名词 & en fritid & \textit{(fritiden, kun ental)} & 业余时间，空闲时间 \\
名词 & en vuggestue & \textit{(vuggestuen, vuggestuer, -ne)} & 托儿所（0--3 岁） \\
名词 & en forening & \textit{(foreningen, foreninger, -erne)} & 协会，体育/文化俱乐部 \\
名词 & en aftensmad & \textit{(aftensmaden, kun ental)} & 晚餐 \\
名词 & en balance & \textit{(balancen, balancer, -rne)} & 平衡（\textit{balance mellem A og B}） \\
名词 & en sygedag & \textit{(sygedagen, sygedage, -ne)} & 病假（\textit{barnets første sygedag}） \\
名词 & en tillid & \textit{(tilliden, kun ental)} & 信任（\textit{have tillid til}） \\
\midrule
动词 [har] & at hente & \textit{(henter, hentede, har hentet)} & 接（孩子）；取（包裹） \\
动词 [har] & at nå (at) & \textit{(når, nåede, har nået)} & 赶上，来得及（完成） \\
动词 [har] & at planlægge & \textit{(planlægger, planlagde, har planlagt)} & 计划，规划 \\
动词 [har] & at koncentrere sig & \textit{(koncentrerer, -ede, har koncentreret)} & 专心，集中注意力（\textit{+ om}） \\
动词 [har] & at indføre & \textit{(indfører, indførte, har indført)} & 引入，实行（制度） \\
动词 [har] & at vare & \textit{(varer, varede, har varet)} & 持续，历时 \\
\midrule
形容词/副词 & effektiv / hjemmefra & \textit{effektivt, effektive} / \textit{副词} & 高效的 / 从家里（远程地） \\
\end{glosertable}

% ---------------------------------------------------------------------
\ncdsubsec{四、Grammatik \& Faste vendinger（核心语法与地道表达精讲）}

\begin{grammatikbox}{语法点 1：省略 \textit{hvis} 的条件倒装句（V1 条件从句）完全公式}
在 A2 阶段，表达“如果……就……”只能依赖连词 \textit{hvis}。到了 B1，丹麦语母语者常常直接\textbf{砍掉句首的 \textit{hvis}}，让条件从句的变位动词直接冲到整句话第 1 位，形成紧凑利落的“\textbf{双倒装条件句}”：

\vspace{0.3em}
\renewcommand{\arraystretch}{1.26}
{\small
\begin{tabularx}{\linewidth}{@{} >{\bfseries\raggedright\arraybackslash}p{2.3cm} >{\raggedright\arraybackslash}p{5.2cm} >{\raggedright\arraybackslash}X @{}}
\toprule
\textbf{句型类型} & \textbf{标准 \textit{Hvis} 句式} & \textbf{省略 \textit{hvis} 的 V1 条件倒装句} \\
\midrule
\textbf{肯定条件}\newline （如果……，就……） &
\textit{\textbf{Hvis} man har børn, skal man hente dem tidligt.} &
\textit{\da{Har man} børn, \dabold{skal man} hente dem tidligt.}\newline （也可以在逗号后加连接词 \textit{så}：\textit{Har man børn, \textbf{så skal} man...}） \\
\midrule
\textbf{情态条件}\newline （要想/必须……） &
\textit{\textbf{Hvis} man vil gå kl. 15, kan man møde kl. 7.} &
\textit{\da{Vil man} gå kl. 15, \dabold{kan man} møde kl. 7.} \\
\midrule
\textbf{否定条件}\newline （\dabold{极易错！注意 \textit{ikke}}） &
\textit{\textbf{Hvis} man \dabold{ikke planlægger} dagen, bliver det svært.} &
\textit{\da{Planlægger man} \dabold{ikke} dagen, \dabold{bliver det} svært.}\newline （\textbf{没有 \textit{hvis}，从句已倒装，所以 \textit{ikke} 必须退回主语后面！}） \\
\bottomrule
\end{tabularx}
}
\vspace{0.3em}

\textbf{识别小秘诀}：如果一个句子的前半部分长得像“是非疑问句（动词开头：\textit{Har man..., Bliver dit barn...}）”，但逗号后面没有问号、反而又紧跟另一个倒装主句，那么它 100\% 就是“如果……”的条件句！
\end{grammatikbox}

\begin{grammatikbox}{语法点 2：两大铁律与英丹对比排雷（为什么英语做不到，丹麦语却天天用？）}
\begin{enumerate}
  \item \textbf{英丹对比｜英语极受限，丹麦语几乎“全动词通吃”！}：
    英语只有极少数正式书面虚拟词（\textit{Had I known... / Were she here... / Should you need...}）才能省略 \textit{if} 倒装，绝不能说 *\textit{Have you children, must you...}。但在丹麦语里，\textbf{几乎任何动词的现在时或过去时都能省略 \textit{hvis} 倒装}（\textit{Har man..., Bor man..., Arbejder man..., Får man..., Bliver barnet...}），既口语又地道，Modul 3 写作中使用 1--2 句极其加分！
  \item \textbf{位置铁律｜条件放在句子后半截时，绝对禁止省略 \textit{hvis}！}：
    \begin{itemize}
      \item \ding{51}\ 条件在前：\textit{\da{Har man} tid, kan man hjælpe.}
      \item \ding{51}\ 条件在后（必须有 \textit{hvis}）：\textit{Man kan hjælpe, \da{hvis} man har tid.}
      \item \ding{55}\ 致命错句：\textit{*Man kan hjælpe, har man tid.}（后置时去掉 \textit{hvis} 属于严重语病！）
    \end{itemize}
  \item \textbf{否定词 \textit{ikke} 归位对比}：有 \textit{hvis} 时是标准从句（\textit{Hvis man \dabold{ikke} har tid}）；删掉 \textit{hvis} 变成 V1 倒装后，\textit{ikke} 必须立刻跳到主语后（\textit{Har man \dabold{ikke} tid}）。考场上只要记住“\textbf{有连词 \textit{ikke} 在动词前，无连词倒装 \textit{ikke} 在主语后}”，就绝不会乱！
\end{enumerate}
\end{grammatikbox}

% ---------------------------------------------------------------------
\ncdsubsec{五、Mønsterøvelser（句型操练）}

\textbf{练习 A：条件倒装变形——把下列带 \textit{Hvis} 的句子改写为“省略 \textit{hvis} 的 V1 条件倒装句”（特别留意 \textit{ikke} 位置！）}
\begin{enumerate}
  \item \textit{Hvis man arbejder i Danmark, har man normalt fem ugers ferie om året.} $\rightarrow$
  \item \textit{Hvis du starter tidligt om morgenen, kan du gå hjem allerede klokken tre.} $\rightarrow$
  \item \textit{Hvis man ikke henter sit barn inden klokken fem, ringer pædagogen til forældrene.} $\rightarrow$
  \item \textit{Hvis vores afdeling får flere opgaver i næste måned, bliver vi nødt til at ansætte en ny kollega.} $\rightarrow$
  \item \textit{Hvis du bliver færdig med rapporten før frokost, må du gerne tage tidligt hjem.} $\rightarrow$
  \item \textit{Hvis man aldrig holder pauser i løbet af dagen, bliver man hurtigt træt og mindre effektiv.} $\rightarrow$
  \item \textit{Hvis dit barn bliver sygt om natten, kan du tage barnets første sygedag.} $\rightarrow$
\end{enumerate}

\textbf{练习 B：语序判断与还原（选出唯一合法语序，或按提示还原）}
\begin{enumerate}
  \item 选择正确句式：\underline{\hspace{4.8cm}} (\textit{Arbejder man ikke / Ikke arbejder man}) hjemmefra, skal man møde på kontoret kl. 8.30.
  \item 选择正确句式：Har du spørgsmål til den nye fleksaftale, \underline{\hspace{3.2cm}} (\textit{du kan / kan du}) altid skrive til chefen.
  \item 选择正确句式：Du skal give besked inden kl. 9, \underline{\hspace{3.6cm}} (\textit{bliver du / hvis du bliver}) syg.
  \item 改正错误（两处或一处）：\textit{*Hvis du har brug for hjælp, du skal bare sige til.} $\rightarrow$
  \item 改正错误（条件后置禁止倒装）：\textit{*Vi spiser aftensmad klokken seks, kommer alle hjem til tiden.} $\rightarrow$
  \item 还原为带 \textit{Hvis} 的标准从句（注意 \textit{ikke} 前移！）：\textit{Planlægger man ikke sin uge godt, får man for travlt.} $\rightarrow$
\end{enumerate}

\textbf{练习 C：中译丹实战（全部使用省略 \textit{hvis} 的条件倒装句）}
\begin{enumerate}
  \item 如果你有小孩，你必须在五点前去幼儿园接他们。
  \item 如果一个人经常加班，同事们就会认为工作量太大了。
  \item 如果你没有按时规划好你的工作日，就很难完成所有任务。
  \item 如果你在家里远程办公，就能省下很多通勤时间。
  \item 如果我们把生活和工作平衡好，日常生活中的压力就会小一些。
  \item 如果你有机会选择弹性工时，你愿意早上七点上班吗？
\end{enumerate}

% ---------------------------------------------------------------------
\ncdsubsec{六、DU3 Modul 3 实战微任务与参考答案}

\begin{tcolorbox}[enhanced, breakable, colback=NCDBgGray, colframe=NCDGray, coltitle=white, fonttitle=\bfseries, title={DU3 Modul 3 实战微任务（对标 Mundtlig Opgave 2 观点讨论 / 书面短论 · 100+ 词）}, arc=2pt, boxrule=0.8pt, borderline west={3pt}{0pt}{NCDRed}, left=10pt, right=10pt, top=7pt, bottom=7pt]
\small
\textbf{任务情境（Emne: Arbejdsliv og fritid）}：在 Modul 3 口语 \textit{Opgave 2} 讨论“弹性工时 vs 固定八小时”或考官个人追问时，你需要表明自己的偏好并说明原因。\\
请写一篇 \textbf{105--130 词}的 B1 观点短文，回答以下问题：\textit{Foretrækker du faste arbejdstider på kontoret, eller synes du bedst om fleksibel arbejdstid og mulighed for hjemmearbejde? Hvorfor?} 要求至少使用 \textbf{2 句省略 \textit{hvis} 的条件倒装句}（其中 1 句包含否定词 \textit{ikke}）！
\end{tcolorbox}

\begin{answerbox}
\textbf{Før du læser 参考回答}
\begin{enumerate}
  \item \textit{Mange går hjem ved firetiden, fordi en normal arbejdsuge er 37 timer, de starter tidligt, og mange skal hente små børn i institutionen før kl. 17 eller nå til fritidsaktiviteter. Meget overarbejde ses sjældent som flid, men snarere som tegn på en for stor arbejdsbyrde eller dårlig planlægning.}
  \item \textit{NordTek tilbyder både fleksibel arbejdstid (op til to hjemmearbejdsdage om ugen samt mulighed for at møde tidligt og gå tidligt) og ”barnets første sygedag” med fuld løn.}
\end{enumerate}

\textbf{Tekst B 阅读补全答案}
\begin{enumerate}
  \item \da{B}（前句说办公室很少加班，后句介绍员工可享受每周两天居家办公，B 项自然承接公司提供弹性上下班协议）
  \item \da{A}（紧接前句“家有幼儿时的突发状况”，A 项给出孩子生病时可休全薪“孩子病假第一天”的方案，并与后句“父母减少焦虑”精准呼应）
\end{enumerate}

\textbf{练习 A 答案（条件倒装变形）}
\begin{enumerate}
  \item \textit{\da{Arbejder man} i Danmark, \dabold{har man} normalt fem ugers ferie om året.}
  \item \textit{\da{Starter du} tidligt om morgenen, \dabold{kan du} gå hjem allerede klokken tre.}
  \item \textit{\da{Henter man} \dabold{ikke} sit barn inden klokken fem, \dabold{ringer} pædagogen til forældrene.}（注意 \textit{ikke} 位于主语 \textit{man} 之后！）
  \item \textit{\da{Får vores afdeling} flere opgaver i næste måned, \dabold{bliver vi} nødt til at ansætte en ny kollega.}
  \item \textit{\da{Bliver du} færdig med rapporten før frokost, \dabold{må du} gerne tage tidligt hjem.}
  \item \textit{\da{Holder man} \dabold{aldrig} pauser i løbet af dagen, \dabold{bliver man} hurtigt træt og mindre effektiv.}
  \item \textit{\da{Bliver dit barn} sygt om natten, \dabold{kan du} tage barnets første sygedag.}
\end{enumerate}

\textbf{练习 B 答案}
\begin{enumerate}
  \item \da{Arbejder man ikke}（无 \textit{hvis} 倒装，\textit{ikke} 放在主语 \textit{man} 后）
  \item \da{kan du}（条件从句在前，逗号后主句必须 V2 倒装）
  \item \da{hvis du bliver}（条件置于句尾时绝对不能省略 \textit{hvis}）
  \item \textit{Hvis du har brug for hjælp, \dabold{skal du} bare sige til.}（从句在前，主句动词 \textit{skal} 必须先于主语 \textit{du}！）
  \item \textit{Vi spiser aftensmad klokken seks, \da{hvis} alle kommer hjem til tiden.}（条件在后必须补回 \textit{hvis}）
  \item \textit{\da{Hvis} man \dabold{ikke} planlægger sin uge godt, får man for travlt.}（还原 \textit{Hvis} 后变为标准从句语序，\textit{ikke} 回到动词前！）
\end{enumerate}

\textbf{练习 C 答案}
\begin{enumerate}
  \item \textit{Har du små børn, skal du hente dem i børnehaven inden klokken fem.}
  \item \textit{Arbejder man ofte over, tænker kollegerne, at arbejdsbyrden er for stor.}
  \item \textit{Planlægger man ikke sin arbejdsdag godt nok, bliver det svært at nå alle opgaver.}
  \item \textit{Arbejder man hjemmefra, kan man spare meget transporttid.}
  \item \textit{Skaber vi en god balance mellem arbejde og fritid, får vi mindre stress i hverdagen.}
  \item \textit{Får du mulighed for at vælge fleksibel arbejdstid, vil du så gerne møde klokken syv om morgenen?}
\end{enumerate}

\textbf{DU3 Modul 3 高分范文（115 词）}\\
\textit{Personligt foretrækker jeg helt klart fleksibel arbejdstid med mulighed for at arbejde hjemmefra et par dage om ugen. Det passer rigtig godt til en travl hverdag, hvor man gerne vil skabe en sund balance mellem job og fritid.}

\textit{\da{Har man} for eksempel børn, eller bor man langt fra kontoret, \dabold{kan man} spare meget tid på transport om morgenen. Samtidig oplever jeg, at jeg koncentrerer mig bedre om svære opgaver i rolige omgivelser derhjemme. \da{Planlægger man} \dabold{ikke} sine faste møder godt sammen med teamet, \dabold{kan} hjemmearbejde selvfølgelig skabe lidt afstand mellem kollegerne. Men \da{mødes man} fysisk på kontoret tre dage om ugen, \dabold{fungerer} samarbejdet efter min mening fantastisk, \da{hvilket} giver både høj effektivitet og stor arbejdsglæde.}

\small\textbf{名师点评}：115 词结构平衡、立场鲜明！不仅自然嵌入三句省略 \textit{hvis} 的条件倒装句（含否定式 \textit{Planlægger man ikke..., kan...}），还用正反对比（\textit{Samtidig ... Men ...}）与 L02 的整句回指 \textit{, hvilket giver...} 收尾，是 Modul 3 口语与写作的双料高分样本。
\end{answerbox}

\ncdlesson{5}{En jobsamtale i en dansk virksomhed}{一家丹麦企业的求职面试}
{DU3 Modul 3（B1）｜对标 NCE 2 职场综合 \& Modul 3 \textit{Mundtlig Opgave 1 (Mit arbejde) \& Skrivning}}
{\textbf{职场核心动词 + 介词固定搭配}（\textit{have ansvar for, tage initiativ til, sætte pris på, bidrage til, lægge vægt på, passe ind i, se frem til}）· \textbf{介词 + \textit{at} + 动词原形}（严禁直译英语 \textit{V-ing}）}
{作为 Unit 1《丹麦职场与教育体系》的收官课，本课将前四课的被动语态、分词定语与职场文化融进一场真实的丹麦求职面试！记住：\textbf{动词后面的介词}在丹麦语中属于“半固定基因”，背单词时一定要连同介词一起记，才能彻底告别中式与英式硬伤！}

% ---------------------------------------------------------------------
\begin{foerduleaeser}
带着下面两个问题阅读 Tekst A 和 Tekst B：
\begin{enumerate}
  \item \textit{Hvad lægger danske arbejdsgivere stor vægt på til en jobsamtale ud over faglige kvalifikationer, og hvordan forberedte Martin sig?}（除了专业资质外，丹麦雇主在面试中还非常看重什么？马丁是如何做准备的？）
  \item \textit{Hvilke styrker fremhæver Martin i sin takkemail i Tekst B, og hvornår får han svar?}（在 Tekst B 的致谢邮件中，马丁重点强调了自己哪几项优势？他最晚什么时候能收到回复？）
\end{enumerate}
\end{foerduleaeser}

% ---------------------------------------------------------------------
\ncdsubsec{二、Tekst A \& Tekst B（社会深读与 Modul 3 真题场景）}

\subsection*{Tekst A：\textit{En jobsamtale i en dansk virksomhed}（职场叙事短文 · 约 250 词）}

\begin{tipbox}{语篇说明}
本篇追踪国际求职者马丁（Martin）在欧登塞参加丹麦企业面试的全过程。阅读时请用“组合视野”锁定\textbf{\color{NCDNavy}深蓝色的动词/短语核心}与紧随其后的\textbf{\color{NCDRed}红色固定介词}（如 \da{sætte pris} \dabold{på}、\da{bidrage} \dabold{til}）。
\end{tipbox}

For et par uger siden søgte Martin en stilling som projektkoordinator i en mellemstor dansk virksomhed, og kort efter blev han inviteret til en jobsamtale i Odense. Da det var hans første samtale i Danmark, \da{forberedte han sig} grundigt \dabold{på} at svare på spørgsmål om både sine faglige kvalifikationer og sin personlighed.

Da han trådte ind i mødelokalet, blev han straks mødt af en varm og uformel stemning. I stedet for stramme jakkesæt havde både HR-konsulenten og teamlederen Mette almindeligt tøj på, og de begyndte med at hente en kop kaffe og tale om cykelturen gennem byen. Alligevel mærkede Martin hurtigt, at samtalen var seriøs. Danske arbejdsgivere \da{lægger} nemlig stor vægt \dabold{på}, om en ny medarbejder \da{passer} godt ind \dabold{i} det eksisterende team og deler virksomhedens værdier.

Undervejs spurgte Mette, hvilke opgaver Martin helst arbejdede med i sit tidligere job. Han forklarede roligt, at han \da{er vant} \dabold{til} at arbejde selvstændigt og tidligere \da{har haft ansvar} \dabold{for} både projektstyring og kundekontakt. Derudover fortalte han, at han altid prøver at \da{tage initiativ} \dabold{til} at forbedre arbejdsgange, fordi han gerne vil \da{bidrage} \dabold{til} et stærkt og effektivt samarbejde. Til sidst nævnte Martin, at han virkelig \da{sætter pris} \dabold{på} at få ærlig feedback fra både kolleger og ledere.

Efter tre kvarter kørte Martin hjem med en god mavefornemmelse, og nu \da{ser han frem} \dabold{til} at høre virksomhedens svar i næste uge. Uanset resultatet har han lært, at åbenhed og god kemi tæller lige så højt som et flot CV.

\begin{transbox}
\small
\textbf{\color{NCDTeal}[Tekst A 中文精译]}\\
几周前，马丁申请了丹麦一家中型企业的项目协调员职位，不久后便受邀前往欧登塞参加求职面试。由于这是他在丹麦的第一次面试，他认真地做好准备，以应对关于专业资质和个人性格的双重提问。

当他走进会议室时，迎面而来的是热情而非正式的氛围。HR 顾问和团队主管梅特都没有穿紧绷的西装，而是穿着日常休闲服装，而且一上来先去倒了一杯咖啡，聊起骑车穿城的话题。尽管如此，马丁很快感觉到面试本身非常认真专业。因为丹麦雇主格外看重一位新员工是否能很好地融入现有团队，并认同企业的价值观。

过程中，梅特问马丁在上一份工作中最喜欢负责哪些任务。他沉着地解释说，自己习惯了独立工作，并且过去曾负责过项目管理和客户沟通。此外他还补充道，自己总是主动采取行动去改进工作流程，因为他希望能为强大而高效的团队协作做出贡献。最后马丁提到，自己非常看重并欣赏能从同事和领导那里获得坦诚的反馈。

四十五分钟后，马丁怀着不错的感觉骑车回家了，现在他正满怀期待地等着下周收到公司的答复。无论结果如何，他都已经体会到，坦诚开放的态度与良好的团队默契，和一份漂亮的简历同样重要。
\end{transbox}

\subsection*{Tekst B：\textit{Opfølgende takkemail efter jobsamtale}（对标 DU3 Modul 3 职场正式致谢邮件）}

\begin{emailbox}{求职者跟进致谢邮件｜Til: Mette (Teamleder) \quad Emne: Tak for en spændende jobsamtale}
\textbf{Kære Mette}\\
Tusind tak for en rigtig god og inspirerende samtale i dag. Jeg \da{satte} stor \da{pris} \dabold{på} at høre mere om jeres projekter og møde teamet. Som jeg nævnte under samtalen, har jeg solid erfaring med kundekontakt, og jeg er vant \dabold{til} at \da{have ansvar} \dabold{for} stramme deadlines. Jeg tror virkelig, at jeg kan \da{bidrage} \dabold{til} jeres afdeling ved at \da{tage initiativ} \dabold{til} bedre digitale løsninger. Jeg \da{ser} meget \da{frem} \dabold{til} at høre fra jer, når der bliver truffet en endelig beslutning i næste uge.\\
\textbf{Venlig hilsen, Martin}
\end{emailbox}

\begin{emailbox}{HR 简短回复｜Fra: Mette \quad Emne: Sv: Tak for en spændende jobsamtale}
\textbf{Hej Martin}\\
Selv tak for et rigtig godt besøg! Vi \da{sætter} også \da{pris} \dabold{på} din store interesse og vender tilbage med et svar senest onsdag eftermiddag.\\
\textbf{Bedste hilsner, Mette}
\end{emailbox}

\textbf{Tekst B 语言与信息核对（答案见第六部分）：}
\begin{enumerate}
  \item \textit{Hvordan mener Martin, at han kan bidrage til afdelingen?} \underline{\hspace{5.5cm}}
  \item \textit{Hvilket ord efterfølger præpositionen \textbf{ved} i ”ved at tage initiativ”? Hvad viser det om dansk grammatik?} \underline{\hspace{4.5cm}}
  \item \textit{Senest hvornår hører Martin fra virksomheden?} \underline{\hspace{3.5cm}}
\end{enumerate}

\begin{transbox}
\small
\textbf{\color{NCDTeal}[Tekst B 中文对照]}\\
\textbf{马丁的跟进致谢邮件}：亲爱的梅特，非常感谢今天一场精彩且富有启发的面试。能进一步了解你们的项目并与团队见面，我深感荣幸与欣喜。正如我在面试中提到的，我在客户沟通方面拥有扎实经验，也习惯于对紧凑的截止日期承担责任。我由衷相信，自己能够通过主动推进更优化的数字化方案为你们部门做出贡献。期待下周最终决定做出后收到你们的消息。诚挚的问候，马丁。

\textbf{梅特的回复}：你好，马丁。也感谢你非常棒的来访！我们同样十分欣赏你浓厚的兴趣，最迟将在周三下午给你回复。谨致问候，梅特。
\end{transbox}

\begin{tipbox}{本课发音与吞音要点（Udtale \& Reduktioner）}
\begin{itemize}
  \item \textbf{介词短语的节奏——主重音在名词/动词，介词可重可轻}：注意 \textit{bidrage} \ipa{ˈbiˌdʁɑʊ̯ə}（\textbf{\textit{g} 软化滑音}为 \ipa{ɑʊ̯}！）接弱读介词 \textit{til} \ipa{tˢe}（\textbf{l 经常脱落}！）；\textit{sætte pris på} \ipa{ˈsɛdə ˈpʰʁiːˀs pʰʌ}；\textit{have ansvar for} \ipa{hɑ ˈanˌsvɑːˀ fʌ}；\textit{lægge vægt på} \ipa{ˈlɛɡə ˈvɛɡd pʰʌ}。
  \item \textbf{外来专业词汇重音定位}：\textit{et initiativ} \ipa{initˢiaˈtˢiːˀʊ̯}（重音在最后的 \textit{-tiv}，词尾 \textit{v} 轻化为半元音）；\textit{en kvalifikation} \ipa{kʰvalifikaˈɕoːˀn}（重音在最后的 \textit{-tion}）；\textit{uformel} \ipa{ufʌˈmɛlˀ}（重音在后）。
  \item \textbf{静音与弱读提醒}：\textit{virksomhed} \ipa{ˈviɐ̯ɡsʌmˌheːˀð}（软 d）；\textit{selvstændig} \ipa{ˈsɛlˌsdɛnˀdi}（\textit{selv} 的 v 不发音，词尾 g 不发音）；\textit{vant til} \ipa{ˈvanˀd tˢe}。
\end{itemize}
\end{tipbox}

% ---------------------------------------------------------------------
\ncdsubsec{三、Gloser（核心生词表）}

\begin{tipbox}{词表重点}
本课动词请连同括号里的\textbf{固定介词}整体背诵！注意 \textit{en stilling}（职位）常配动词 \textit{at søge}（申请，注意：\textbf{丹麦语说 \textit{søge et job / en stilling}，中间不加 \textit{for}}！切勿照搬英语 \textit{*apply for} 而写成 \textit{*søge for}！）。
\end{tipbox}

\begin{glosertable}
名词 & en jobsamtale & \textit{(jobsamtalen, jobsamtaler, -ne)} & 求职面试 \\
名词 & en virksomhed & \textit{(virksomheden, -heder, -erne)} & 企业，公司 \\
名词 & en arbejdsgiver & \textit{(arbejdsgiveren, -e, -ne)} & 雇主，用人单位 \\
名词 & en ansøger & \textit{(ansøgeren, ansøgere, -ne)} & 申请者，求职者 \\
名词 & et initiativ & \textit{(initiativet, initiativer, -erne)} & 主动性，倡议（\textit{tage initiativ til}） \\
名词 & et samarbejde & \textit{(samarbejdet, samarbejder, -ne)} & 合作，协作 \\
名词 & en kvalifikation & \textit{(kvalifikationen, -er, -erne)} & 资质，资格条件 \\
名词 & en erfaring & \textit{(erfaringen, erfaringer, -erne)} & 经验（\textit{erfaring med}） \\
名词 & en stilling & \textit{(stillingen, stillinger, -erne)} & 职位，岗位 \\
名词 & en værdi & \textit{(værdien, værdier, værdierne)} & 价值观；价值 \\
\midrule
动词 [har] & at søge & \textit{(søger, søgte, har søgt)} & 申请（职位），寻找（直接加宾语） \\
动词 [har] & at bidrage (til) & \textit{(bidrager, bidrog, har bidraget)} & 做出贡献，出力（不规则过去时） \\
动词 [har] & at forberede sig (på/til) & \textit{(forbereder, -beredte, har forberedt)} & 做好准备 \\
动词 [har] & at sætte pris (på) & \textit{(sætter, satte, har sat)} & 欣赏，珍惜，看重 \\
动词 [har] & at lægge vægt (på) & \textit{(lægger, lagde, har lagt)} & 强调，重视，把重心放在 \\
动词 [har] & at passe ind (i) & \textit{(passer, passede, har passet)} & 融入，契合（团队/环境） \\
\midrule
形容词 & uformel / vant (til) & \textit{uformelt, uformelle} / \textit{vant, vante} & 非正式的 / 习惯于……的 \\
\end{glosertable}

% ---------------------------------------------------------------------
\ncdsubsec{四、Grammatik \& Faste vendinger（核心语法与地道表达精讲）}

\begin{grammatikbox}{语法点 1：丹麦职场八大“动词/短语 + 介词”黄金固定搭配（附英语陷阱对比）}
在 Modul 3 口试谈 \textit{Mit arbejde} 或写求职/职场信件时，下面这八个搭配使用率极高，且介词选择往往与英语截然不同：

\vspace{0.3em}
\renewcommand{\arraystretch}{1.24}
{\small
\begin{tabularx}{\linewidth}{@{} >{\bfseries\color{NCDNavy}\raggedright\arraybackslash}p{4.3cm} >{\raggedright\arraybackslash}p{3.0cm} >{\raggedright\arraybackslash}X @{}}
\toprule
\textbf{丹麦语固定搭配} & \textbf{中文含义} & \textbf{地道例句与英丹避坑对比} \\
\midrule
at have ansvar \dabold{for} & 对……负责 & \textit{Jeg har ansvar \dabold{for} projektet.}（同英语 \textit{responsible for}） \\
at tage initiativ \dabold{til} & 主动发起/采取行动去…… & \textit{Hun tog initiativ \dabold{til} et nyt system.}（英语多用 \textit{to}，丹麦语必用 \dabold{til}！） \\
at sætte pris \dabold{på} & 欣赏，感激，看重 & \textit{Jeg sætter stor pris \dabold{på} din hjælp.}（\textbf{英语 \textit{appreciate} 无介词，丹麦语必加 \dabold{på}！}） \\
at bidrage \dabold{til} & 为……做贡献 & \textit{Alle kan bidrage \dabold{til} et godt arbejdsmiljø.}（对应 \textit{contribute to}） \\
at lægge vægt \dabold{på} & 重视，看重，强调 & \textit{De lægger vægt \dabold{på} åbenhed.}（\textit{lægge} 不规则：\textit{lagde, lagt}） \\
at passe ind \dabold{i} & 融入，契合（团队/文化） & \textit{Han passer godt ind \dabold{i} vores afdeling.}（对应 \textit{fit into}） \\
at være vant \dabold{til} & 习惯于…… & \textit{Jeg er vant \dabold{til} at arbejde selvstændigt.}（区别：仅用 \textit{vant til}，绝不加 \textit{-ing}） \\
at se frem \dabold{til} & 期待，盼望 & \textit{Jeg ser frem \dabold{til} at møde jer.}（对应 \textit{look forward to}） \\
\bottomrule
\end{tabularx}
}
\vspace{0.3em}

\textbf{易混前置对比}：\textit{forberede sig \dabold{på}}（对某件事/挑战做心理或内容准备，如 \textit{forberede sig på spørgsmål}） vs \textit{forberede sig \dabold{til}}（为某场具体的活动/考试/会议做准备，如 \textit{forberede sig til jobsamtalen}）；而“有……经验”固定用 \textit{erfaring \dabold{med}}（如 \textit{erfaring med kundekontakt}，不要套用英语 \textit{*experience in/of}）！
\end{grammatikbox}

\begin{grammatikbox}{语法点 2：介词后面接“动作”怎么办？——铁律：\textbf{介词 + \textit{at} + 动词原形}}
这是所有学过英语的同学在丹麦语中最容易踩的大坑之一！请牢牢记住两条不可动摇的铁律：
\begin{enumerate}
  \item \textbf{严禁照搬英语的“介词 + V-ing”}：英语说 \textit{I look forward to \textbf{working} with you} 或 \textit{She is responsible for \textbf{planning} meetings}，但丹麦语\textbf{绝不能在介词后面接 \textit{-ende} 分词}！丹麦语必须用\textbf{\da{介词 + at + 动词原形}}：
    \begin{itemize}
      \item \ding{51}\ \textit{Jeg ser frem \dabold{til at arbejde} sammen med jer.}（\ding{55}\ *\textit{til arbejdende} 或 *\textit{til arbejde}！）
      \item \ding{51}\ \textit{Hun har ansvar \dabold{for at planlægge} møderne.}（\ding{55}\ *\textit{for planlæggende}！）
      \item \ding{51}\ \textit{Jeg sætter pris \dabold{på at få} ærlig feedback.}
      \item \ding{51}\ \textit{Vi kan forbedre samarbejdet \dabold{ved at tale} åbent sammen.}（\textit{ved at + 原形} =“通过做某事”，Modul 3 作文超高频！）
    \end{itemize}
  \item \textbf{Unit 1 收官小提示（预告 Unit 3 第 11 课）}：当介词后面跟的不是一个简单动作（\textit{at + 原形}），而是一整个带主语的句子（如“我很欣赏\textbf{你主动帮忙}”）时，正式丹麦语常在介词后补一个形式宾语 \textit{det} 再接从句（如 \textit{Jeg er glad for, \textbf{at} du hjælper / Jeg regner med, \textbf{(det)} at...}）。我们在 Unit 3 会专门拆解这套高阶从句！
\end{enumerate}
\end{grammatikbox}

% ---------------------------------------------------------------------
\ncdsubsec{五、Mønsterøvelser（句型操练）}

\textbf{练习 A：介词配对填空（从 \textit{for, til, på, i, med} 中选择最准确的介词，若不需介词填 $\times$）}
\begin{enumerate}
  \item Jeg har lige søgt \underline{\hspace{1.8cm}} en spændende stilling som projektleder.
  \item Til jobsamtalen lagde chefen stor vægt \underline{\hspace{1.8cm}} god kommunikation og samarbejde.
  \item Jeg har haft ansvaret \underline{\hspace{1.8cm}} nye kunders introduktion i tre år.
  \item Alle medarbejdere kan bidrage \underline{\hspace{1.8cm}} en bedre stemning på kontoret.
  \item Hun har allerede mange års erfaring \underline{\hspace{1.8cm}} at undervise voksne.
  \item Vi sætter virkelig pris \underline{\hspace{1.8cm}} din ærlige feedback på vores forslag.
  \item Jeg er sikker på, at Martin vil passe godt ind \underline{\hspace{1.8cm}} vores lille team.
  \item Vi ser meget frem \underline{\hspace{1.8cm}} at byde dig velkommen på mandag!
\end{enumerate}

\textbf{练习 B：句型整合与英式错误纠偏（特别注意“介词 + \textit{at} + 原形”与“直接及物”）}
\begin{enumerate}
  \item 用 \textit{ved at + 原形} 合并：\textit{Man kan lære meget om dansk arbejdskultur. Man stiller spørgsmål til sine kolleger.} $\rightarrow$
  \item 用 \textit{vant til at...} 改写：\textit{Jeg arbejder normalt selvstændigt med mine egne projekter.} $\rightarrow$ \textit{Jeg er...}
  \item 用 \textit{tage initiativ til at...} 改写：\textit{Hun foreslog ofte, at vi skulle forbedre vores digitale møder.} $\rightarrow$ \textit{Hun tog ofte...}
  \item 改正英式迁移错误：\textit{*Jeg ser frem til samarbejdende med jer i den nye afdeling.} $\rightarrow$
  \item 改正英式迁移错误：\textit{*Han vil gerne søge for det nye job i Aarhus.} $\rightarrow$
  \item 改正介词遗漏错误：\textit{*Jeg sætter stor pris din hurtige hjælp i går.} $\rightarrow$
\end{enumerate}

\textbf{练习 C：中译丹实战（活用本课职场固定搭配与 Unit 1 被动/分词句型）}
\begin{enumerate}
  \item 我对团队的日常规划负责。
  \item 丹麦公司很看重员工是否能融入团队。
  \item 我习惯于在压力下独立工作。
  \item 谢谢你为我们的合作做出积极贡献！
  \item 我非常期待下周再次见到您。
  \item 我们可以通过主动分享想法来改善工作氛围。（用 \textit{ved at...}）
\end{enumerate}

% ---------------------------------------------------------------------
\ncdsubsec{六、DU3 Modul 3 实战微任务与参考答案}

\begin{emailbox}{DU3 Modul 3 Unit 1 收官综合写作任务（对标 Skrivning Opgave 2 · 105--130 词）}
\textbf{Situation}: Din ven Jonas er blevet inviteret til sin første jobsamtale i en dansk virksomhed næste uge og har skrevet til dig for at få nogle gode råd.\\
\textbf{Opgave}: Skriv et hjælpsomt og motiverende svar på \textbf{105--130 ord}. Giv ham konkrete råd om den danske arbejdskultur (fx det flade hierarki, uformel stemning og hvad arbejdsgivere lægger vægt på). Brug mindst \textbf{4 af lektionens faste vendinger med præposition} samt mindst \textbf{1 passiv eller \textit{hvis}-løs betingelsessætning} fra Unit 1!
\end{emailbox}

\begin{answerbox}
\textbf{Før du læser 参考回答}
\begin{enumerate}
  \item \textit{Ud over faglige kvalifikationer lægger danske arbejdsgivere stor vægt på personlighed, og om man passer godt ind i teamet og deler virksomhedens værdier. Martin forberedte sig grundigt på at svare på spørgsmål om både sine kvalifikationer og sin personlighed.}
  \item \textit{Martin fremhæver sin solide erfaring med kundekontakt, at han er vant til at have ansvar for stramme deadlines, og at han kan bidrage med bedre digitale løsninger. Han får svar senest onsdag eftermiddag.}
\end{enumerate}

\textbf{Tekst B 语言与信息核对答案}
\begin{enumerate}
  \item \textit{Ved at tage initiativ til bedre digitale løsninger.}
  \item \da{at tage}（不定式小词 \textit{at} + 动词原形）。这体现了丹麦语“\textbf{介词 + \textit{at} + 原形}”铁律（绝不加 \textit{-ende}）。
  \item \da{Senest onsdag eftermiddag.}
\end{enumerate}

\textbf{练习 A 答案（介词配对）}
\begin{enumerate}
  \item \da{$\times$}（直接说 \textit{søge en stilling}，中间不加 \textit{for}！）
  \item \da{på}（\textit{lægge vægt på}）
  \item \da{for}（\textit{have ansvar for}）
  \item \da{til}（\textit{bidrage til}）
  \item \da{med}（\textit{erfaring med}）
  \item \da{på}（\textit{sætte pris på}）
  \item \da{i}（\textit{passe ind i}）
  \item \da{til}（\textit{se frem til}）
\end{enumerate}

\textbf{练习 B 答案}
\begin{enumerate}
  \item \textit{Man kan lære meget om dansk arbejdskultur \da{ved at stille} spørgsmål til sine kolleger.}
  \item \textit{Jeg er \da{vant til at arbejde} selvstændigt med mine egne projekter.}
  \item \textit{Hun tog ofte \da{initiativ til at forbedre} vores digitale møder.}
  \item \textit{Jeg ser frem \da{til at samarbejde} med jer i den nye afdeling.}（介词后改用 \textit{at + 原形}，去掉 \textit{-ende}）
  \item \textit{Han vil gerne \da{søge} det nye job i Aarhus.}（删除英式多余介词 \textit{for}）
  \item \textit{Jeg sætter stor \da{pris på} din hurtige hjælp i går.}（补全不可省略的 \textit{på}）
\end{enumerate}

\textbf{练习 C 答案}
\begin{enumerate}
  \item \textit{Jeg har ansvar for den daglige planlægning i teamet.}
  \item \textit{Danske virksomheder lægger stor vægt på, om en medarbejder passer ind i teamet.}
  \item \textit{Jeg er vant til at arbejde selvstændigt under pres.}
  \item \textit{Tak fordi du bidrager aktivt til vores samarbejde!}
  \item \textit{Jeg ser meget frem til at møde dig igen i næste uge.}
  \item \textit{Vi kan forbedre arbejdsmiljøet ved at tage initiativ til at dele vores ideer.}
\end{enumerate}

\textbf{DU3 Modul 3 高分范文（120 词，不含署名）}\\
\textit{Hej Jonas}

\textit{Hvor er det spændende, at du \dabold{er blevet inviteret} til din første jobsamtale! Jeg vil gerne dele et par hurtige råd med dig.}

\textit{Husk først, at stemningen ofte er ret uformel, og chefen \da{tiltales} ved fornavn. Danske arbejdsgivere \da{lægger} nemlig stor vægt \dabold{på}, om man \da{passer ind} \dabold{i} teamet. \da{Forbereder du dig} godt \dabold{på} at fortælle, hvordan du kan \da{bidrage} \dabold{til} samarbejdet, \dabold{giver det} et stærkt indtryk. Fortæl gerne om konkrete opgaver, som du har haft ansvar \dabold{for}, eller hvor du har taget initiativ \dabold{til} at løse et problem. Danske ledere \da{sætter} også stor pris \dabold{på} ærlighed og \da{engagerede} spørgsmål fra ansøgeren.}

\textit{Jeg krydser fingre for dig og \da{ser frem} \dabold{til} at høre, hvordan det går!}

\textit{Venlig hilsen\\
\mbox{[Dit navn]}}

\small\textbf{Unit 1 总结点评}：120 词（不含署名）完成 Unit 1 五大语法的全景式融合——L01 被动语态（\textit{er blevet inviteret, tiltales}）、L03 分词定语（\textit{engagerede}）、L04 无 \textit{hvis} 条件倒装（\textit{Forbereder du dig..., giver det...}）与 L05 六组职场介词搭配自然交织，标志着你已正式掌握 DU3 Modul 3 职场主题的核心写作力！
\end{answerbox}


% =====================================================================
% Unit 2
% =====================================================================
\ncdunit{2}{Kultur, Mentalitet og Tillid}{文化心理与社会默契}{Lektion 6--10：形容词名词化 · 因果进阶 · 情态推测与传闻 · 历史叙事时态 · 高级让步与对比}

\ncdlesson{6}{Hvad er egentlig "hygge"?}{究竟什么是“Hygge”？}
{DU3 Modul 3（B1）｜对标 NCE 2 后半 \& Modul 3 \textit{Kulturforståelse \& Skrivning Opgave 2}}
{\textbf{形容词名词化泛指结构（Substantiverede adjektiver）}：中性抽象泛指 \textit{det + 形容词-e}（\textit{det hyggelige, det vigtigste, det svære}）vs 复数群体泛指 \textit{de + 形容词-e}（\textit{de unge, de ældre, de ensomme}）}
{迈入 B1 阶段，谈论文化概念与社会群体时不能只会说 \textit{en vigtig ting} 或 \textit{unge mennesker}。熟练运用 \textbf{\textit{det vigtigste}}（最重要的事）与 \textbf{\textit{de unge / de ældre}}（年轻人 / 老年人）能让表达瞬间浓缩地道。记住铁律：名词化形容词一律加定指词尾 \textbf{\textit{-e}}！}

% ---------------------------------------------------------------------
\begin{foerduleaeser}
带着下面两个问题阅读 Tekst A 和 Tekst B：
\begin{enumerate}
  \item \textit{Hvad er det vigtigste ved dansk ”hygge” ifølge Tekst A, og hvordan oplever de unge og de ældre den?}（根据 Tekst A，丹麦“hygge”最核心的是什么？年轻人和老年人分别是如何体验它的？）
  \item \textit{Hvad synes Emma er det svære ved at oversætte ordet ”hygge”, og hvilket godt råd giver hun til Marco?}（在 Tekst B 中，艾玛认为翻译“hygge”一词最难的地方是什么？她给了马尔科什么好建议？）
\end{enumerate}
\end{foerduleaeser}

% ---------------------------------------------------------------------
\ncdsubsec{二、Tekst A \& Tekst B（文化深读与 Modul 3 真题场景）}

\subsection*{Tekst A：\textit{Hvad er egentlig ”hygge”?}（文化深读短文 · 约 250 词）}

\begin{tipbox}{语篇说明}
本篇深入解析丹麦最具代表性的文化概念 \textit{hygge}。阅读时请特别关注\textbf{\color{NCDNavy}深蓝色标记的抽象名词化结构（\textit{det + 形容词-e}）}与\textbf{\color{NCDRed}深红色标记的群体名词化结构（\textit{de + 形容词-e}）}。
\end{tipbox}

Mange udlændinge, der flytter til Danmark, møder hurtigt ordet ”hygge”. I udlandet tror man ofte, at hygge bare handler om stearinlys, varm te og uldsokker om vinteren. Men spørger man danskerne selv, er \da{det egentlige} ved hygge slet ikke de ydre ting. \da{Det vigtigste} er \da{det trygge} og \da{det afslappede}, som opstår, når man er sammen uden stress.

Hygge opleves dog forskelligt i forskellige generationer. \dabold{De unge} skaber ofte hygge i små lejligheder, hvor de laver mad sammen, spiller brætspil eller drikker kaffe efter eksamen. \dabold{De ældre} foretrækker måske en rolig søndag eftermiddag med hjemmebagt kage og lange samtaler i stuen. For både \dabold{de unge} og \dabold{de ældre} handler det om at lægge \da{det travle} til side for en stund og nyde \da{det enkle} i hverdagen.

Om vinteren samles man inde i varmen, mens man om sommeren tager tæpper med ud i parken eller ned til havnen. For \dabold{de ensomme} i samfundet kan en åben invitation fra naboer eller kolleger betyde alverden, fordi \da{det sociale} skaber et stærkt tilhørsforhold.

For nye tilflyttere kan kulturen omkring hygge alligevel virke lidt lukket i begyndelsen. \da{Det sværeste} er ofte at forstå, hvordan man bliver en del af \da{det fælles}, når danskere helst hygger sig med deres nære venner. Man lægger nemlig stor vægt på \da{det ligeværdige}: Ingen skal prale af \da{det dyre}, og \da{det gode} ved en hyggelig aften er netop, at alle kan være sig selv. Når man først oplever denne varme stemning, forstår man, hvorfor hygge betyder så meget i Danmark.

\begin{transbox}
\small
\textbf{\color{NCDTeal}[Tekst A 中文精译]}\\
许多搬到丹麦的外国人很快就会遇到“hygge”这个词。在国外，人们常常以为 hygge 不过是冬天里的蜡烛、热茶和羊毛袜。但如果问丹麦人自己，hygge 的真正本质根本不在于这些外在物件。最重要的是那种在没有压力的相处中自然生发出的安心感与松弛感。

不过，不同代际对 hygge 的体验方式各不相同。年轻人经常在小公寓里营造 hygge，他们一起做饭、玩桌游，或者在考完试后喝咖啡聊天。老年人则可能更喜欢在周日下午来点自制蛋糕，在客厅里安静地长谈。无论对年轻人还是老年人来说，hygge 的关键都在于暂时把忙碌放到一边，享受日常生活中的质朴与简单。

冬天大家聚在温暖的室内，夏天则带着毯子去公园或港口边。对于社会上的独居或孤独群体而言，来自邻居或同事的一份敞开怀抱的邀请意义非凡，因为这种社交往来能带来强烈的归属感。

尽管如此，对新搬来的人而言，围绕 hygge 的文化一开始可能显得有点封闭。最难的往往是弄清楚如何融入那种“共同感”，因为丹麦人通常更喜欢和亲密的老友一起享受 hygge。事实上，人们非常看重平等感：谁都不必炫耀昂贵的东西，一个惬意夜晚的美妙之处恰恰在于每个人都能做真实的自己。一旦亲身体验过这种温暖的氛围，你就会明白为什么 hygge 在丹麦如此重要。
\end{transbox}

\subsection*{Tekst B：\textit{Hygge i hverdagen --- En mail til en ny tilflytter}（对标 Modul 3 Skrivning Opgave 2）}

\textbf{情景说明：}你刚认识不久的朋友马尔科（Marco）搬到了奥胡斯，来信询问丹麦的 \textit{hygge} 到底是什么体验。

\begin{emailbox}{来信｜Fra: Marco \quad Emne: Spørgsmål om dansk hygge}
\textbf{Hej Emma}\\
Jeg hører hele tiden mine kolleger tale om ”hygge”. Hvad betyder det egentlig i praksis, og hvad er vigtigst for at skabe en hyggelig aften?\\
\textbf{Mange hilsner, Marco}
\end{emailbox}

\begin{emailbox}{跨文化讨论回信范本（116 词）｜Til: Marco \quad Emne: Sv: Spørgsmål om dansk hygge}
\textbf{Hej Marco}

Tak for din mail! Jeg forstår godt, at du spørger, for \da{det svære} ved ordet ”hygge” er, at man ikke kan oversætte det direkte til andre sprog.

I praksis handler det slet ikke om at bruge mange penge. \da{Det vigtigste} er bare at sætte tid af til hinanden og skabe en rolig stemning. Både \dabold{de unge} på mit studie og \dabold{de ældre} i min familie elsker at mødes over god mad, brætspil eller en kop kaffe. \da{Det gode ved} hygge er, at alle bidrager med lidt, så ingen bliver stresset.

Næste fredag inviterer jeg nogle venner på aftensmad. Har du lyst til at komme og opleve \da{det hyggelige} sammen med os?

\textbf{Kærlig hilsen}\\
\textbf{Emma}
\end{emailbox}

\begin{tipbox}{本课发音与吞音要点（Udtale \& Reduktioner）}
\begin{itemize}
  \item \textbf{核心文化词的喉塞音与静音字母}：\textit{hygge} \ipa{ˈhyɡə}（短元音 \textit{y}，\textit{gg} 发清软腭塞音，不吞掉！）；\textit{hyggelig} \ipa{ˈhyɡəli}（词尾 \textit{g} 不发音）；\textit{egentlig} \ipa{ˈeːˀənli}（\textbf{\textit{g} 和 \textit{t} 均不发音}！）。
  \item \textbf{名词化词组的重音节奏}：\textit{det} 与 \textit{de} 在泛指词组中通常轻读，重音落在后面的形容词词根上：\textit{det vigtigste} \ipa{de ˈveɡdisdə}、\textit{det afslappede} \ipa{de ˈɑʊ̯ˌslɑbəðə}、\textit{de ældre} \ipa{di ˈɛlˀdʁɐ}、\textit{de unge} \ipa{di ˈɔŋə}。
  \item \textbf{高频词汇读音}：\textit{stearinlys} \ipa{sdɛɑˈʁiːˀnˌlyːˀs}（重音在 \textit{-rin-}）；\textit{fællesskab} \ipa{ˈfɛləsˌsgæːˀb}；\textit{tryg} \ipa{ˈtˢʁøɡ} / \textit{trygge} \ipa{ˈtˢʁøɡə}（\textit{y} 在 \textit{g} 前发音偏圆唇 \ipa{ø}）。
\end{itemize}
\end{tipbox}

\begin{transbox}
\small
\textbf{\color{NCDTeal}[Tekst B 中文对照]}\\
\textbf{马尔科来信}：嗨，艾玛！我一直听同事们谈起“hygge”。这在实际生活中到底是什么意思？要营造一个惬意的夜晚什么最重要？祝好，马尔科。

\textbf{艾玛回信}：嗨，马尔科！谢谢你的来信！我完全理解你为什么会问，因为“hygge”这个词难就难在无法直接翻译成其他语言。其实在生活中，它根本不关乎花多少钱。最重要的是留出彼此陪伴的时间，营造一种宁静的氛围。无论是我们系里的年轻人，还是我家的长辈，都很喜欢聚在一起吃顿好饭、玩桌游或者喝杯咖啡。Hygge 的好处在于每个人都略微出一份力，这样谁都不会觉得有压力。下周五我要请几位朋友吃晚饭，你想不想一起来亲身体验一下这种惬意的氛围？亲切的问候，艾玛。
\end{transbox}

% ---------------------------------------------------------------------
\ncdsubsec{三、Gloser（核心生词表）}

\begin{tipbox}{词表重点}
本课重点积累用于文化阐述与形容词名词化的核心 B1 词汇。特别留意：\textit{at opstå}（产生，出现）完成时必须使用助动词 \dabold{[er]}；\textit{tryg} 变复数或名词化时必须双写 \textit{g}（\textit{trygge}）；\textit{enkel} 加 \textit{-e} 时去掉中间的 \textit{e}（\textit{enkle}）。
\end{tipbox}

\begin{glosertable}
动词 [har] & at skabe & \textit{(skaber, skabte, har skabt)} & 创造，营造 \\
动词 [har] & at opleve & \textit{(oplever, oplevede, har oplevet)} & 体验，经历 \\
动词 [har] & at slappe af & \textit{(slapper af, slappede af, har slappet af)} & 放松，休息 \\
动词 [har] & at lægge vægt på & \textit{(lægger, lagde, har lagt)} & 重视，看重 \\
动词 [har] & at oversætte & \textit{(oversætter, oversatte, har oversat)} & 翻译 \\
动词 \dabold{[er]} & at opstå & \textit{(opstår, opstod, \dabold{er opstået})} & 产生，出现 \\
\midrule
名词 & et fællesskab & \textit{(fællesskabet, fællesskaber, fællesskaberne)} & 集体，社群归属感 \\
名词 & en stemning & \textit{(stemningen, stemninger, stemningerne)} & 氛围，气氛 \\
名词 & en tryghed & \textit{(trygheden, ingen flertal)} & 安全感，安心感 \\
名词 & et stearinlys & \textit{(stearinlyset, stearinlys, stearinlysene)} & 蜡烛 \\
名词 & en tilflytter & \textit{(tilflytteren, tilflyttere, tilflytterne)} & 新搬入者，新移民 \\
名词 & et begreb & \textit{(begrebet, begreber, begreberne)} & 概念，观念 \\
名词 & et tilhørsforhold & \textit{(tilhørsforholdet, tilhørsforhold, tilhørsforholdene)} & 归属感 \\
\midrule
形容词 & hyggelig & \textit{(hyggeligt, hyggelige)} & 惬意温馨的（\textit{det hyggelige}） \\
形容词 & enkel & \textit{(enkelt, enkle)} & 简单的，质朴的（\textit{det enkle}） \\
形容词 & afslappet & \textit{(afslappet, afslappede)} & 放松的（\textit{det afslappede}） \\
形容词 & tryg & \textit{(trygt, trygge)} & 安心的，有安全感的 \\
\end{glosertable}

% ---------------------------------------------------------------------
\ncdsubsec{四、Grammatik \& Faste vendinger（核心语法与地道表达精讲）}

\begin{grammatikbox}{语法点 1：形容词名词化的两大黄金公式（\textit{det + 形容词-e} vs \textit{de + 形容词-e}）}
当形容词不再修饰具体的名词，而是直接充当句子的主语、宾语或介词宾语时，就发生了“形容词名词化”。丹麦语只需在形容词前加上定冠词 \da{det} 或 \dabold{de}，并把形容词变为\textbf{定指形式（词尾加 \textit{-e}）}：
\begin{itemize}
  \item \textbf{公式 \ding{172}｜\da{det} + 形容词-\textit{e}（表抽象事物、本质或“……的那一点”）}：对应中文的“……的东西 / ……之处 / 最……的是”。无论原形还是最高级，只要泛指抽象概念，一律用中性单数 \da{det}！例：\textit{\da{Det vigtigste} er at have tid.}（最重要的是有时间。）/ \textit{Jeg elsker \da{det enkle} ved livet på landet.}（我喜欢乡村生活中质朴的一面。）
  \item \textbf{公式 \ding{173}｜\dabold{de} + 形容词-\textit{e}（表某一类人群全体）}：对应中文的“年轻人 / 老年人 / 富人 / 孤独的人 / 失业者”。例：\textit{\dabold{De unge} bruger ofte sociale medier, mens \dabold{de ældre} foretrækker aviser.}
  \item \textbf{词形拼写细则}：注意 \textit{-et} 结尾的分词形容词加 \textit{-e} 时词尾变为 \textit{-ede}（\textit{afslappet $\rightarrow$ det afslappede}）；以 \textit{-el/-en/-er} 结尾的双音节词加 \textit{-e} 时常省去中间的 \textit{e}（\textit{enkel $\rightarrow$ det enkle}, \textit{gammel $\rightarrow$ de gamle}）；以 \textit{-e} 或 \textit{-ende} 结尾的词不变（\textit{spændende $\rightarrow$ det spændende}，\textit{de studerende}）。
\end{itemize}
\end{grammatikbox}

{\small
\renewcommand{\arraystretch}{1.25}
\begin{xltabular}{\linewidth}{@{} >{\bfseries\color{NCDNavy}\raggedright\arraybackslash}p{3.0cm} >{\raggedright\arraybackslash}p{3.9cm} >{\raggedright\arraybackslash}X @{}}
\toprule
\textbf{名词化类别} & \textbf{高频词形结构} & \textbf{地道 B1 例句与语义} \\
\midrule
\endfirsthead
\toprule
\textbf{名词化类别} & \textbf{高频词形结构} & \textbf{地道 B1 例句与语义} \\
\midrule
\endhead
\bottomrule
\endfoot
\textbf{1. 抽象概念/本质}\newline （\textit{det} + 原级-\textit{e}） & \textit{det hyggelige, det gode, det svære, det sociale, det nye} & \textit{\da{Det gode} ved Danmark er den korte arbejdsdag.}\newline （丹麦的好之处在于较短的工作日。） \\
\midrule
\textbf{2. 最高级焦点}\newline （\textit{det} + 最高级-\textit{e}） & \textit{det vigtigste, det bedste, det sværeste, det mest spændende} & \textit{\da{Det sværeste} ved dansk er udtalen.}\newline （丹麦语最难的是发音。） \\
\midrule
\textbf{3. 社会群体泛指}\newline （\textit{de} + 原级/比较级-\textit{e}） & \textit{de unge, de ældre, de rige, de fattige, de ensomme, de arbejdsløse} & \textit{\dabold{De unge} og \dabold{de ældre} mødes ofte i lokale klubber.}\newline （年轻人和老年人经常在当地俱乐部相聚。） \\
\midrule
\textbf{4. 语言/文化/国籍}\newline （\textit{det} + 国籍形容词-\textit{e}） & \textit{det danske, det kinesiske, det typiske} & \textit{For mig er tillid \da{det mest typiske} ved det danske samfund.}\newline （对我而言，信任是丹麦社会最典型的一点。） \\
\end{xltabular}
}

\begin{grammatikbox}{语法点 2：中英文对比排雷——绝对不能踩的三个“假朋友”陷阱}
\begin{enumerate}
  \item \textbf{“最重要的是……”绝不用 \textit{den}，必须用 \da{det}！}：受英语 \textit{the most important thing} 或“一个重要的点（\textit{en vigtig ting}）”影响，学习者常误写成 *\textit{\dabold{Den} vigtigste er...}。记住：没有具体指代前面某一个 \textit{en}-名词时，抽象泛指一律用中性 \da{det}（\textit{\da{Det} vigtigste er...}）！
  \item \textbf{绝不可漏掉定指词尾 \textit{-e}！}：英语说 \textit{the young, the poor}，形容词后面没有词尾，但丹麦语如果漏掉 \textit{-e} 写成 *\textit{de ung} 或 *\textit{det vigtig} 就是严重语法错误！最高级同样必须加 \textit{-e}：\textit{det bedst\dabold{e}}、\textit{det vigtigst\dabold{e}}。
  \item \textbf{“一群/几个年轻人”不能用 \textit{de unge}！}：\textit{de unge} 只用来泛指“整个年轻人群体”（等同于英文 \textit{young people in general}）。如果要说“两个年轻人”或“我遇到了几个年轻人”，必须补回名词 \textit{mennesker} 或具体名词：\textit{to unge mennesker} / \textit{nogle unge studerende}。
\end{enumerate}
\end{grammatikbox}

\begin{tipbox}{Modul 3 写作高频框架句（Faste vendinger）}
\begin{itemize}
  \item \textit{\da{Det vigtigste} for mig er, at...}（对我而言最重要的是……）
  \item \textit{\da{Det gode / svære ved} [名词/at + 原形] er, at...}（……的好之处 / 困难之处在于……）
  \item \textit{Blandt \dabold{de unge} ser man ofte, at... mens \dabold{de ældre} derimod...}（在年轻人中常见……，而老年人相比之下……）
\end{itemize}
\end{tipbox}

% ---------------------------------------------------------------------
\ncdsubsec{五、Mønsterøvelser（句型操练）}

\textbf{练习 A：填空练习——填入 \textit{det} 或 \textit{de} 以及括号内形容词的正确名词化形式}
\begin{enumerate}
  \item \underline{\hspace{2.6cm}} (\textit{vigtig}, 最高级) ved en god fest er, at gæsterne slapper af.
  \item Både \underline{\hspace{2.4cm}} (\textit{ung}) og \underline{\hspace{2.4cm}} (\textit{ældre}) sætter pris på hygge.
  \item \underline{\hspace{2.6cm}} (\textit{svær}, 最高级) ved at flytte til et nyt land er at finde nære venner.
  \item Kommunen gør meget for at hjælpe \underline{\hspace{2.6cm}} (\textit{ensom}) i byen.
  \item Vi skal lægge \underline{\hspace{2.4cm}} (\textit{travl}) til side og nyde \underline{\hspace{2.4cm}} (\textit{enkel}) i hverdagen.
  \item \underline{\hspace{2.6cm}} (\textit{god}) ved at arbejde i Danmark er balancen mellem job og fritid.
  \item I samfundet diskuterer man ofte forskellen mellem \underline{\hspace{2.4cm}} (\textit{rig}) og \underline{\hspace{2.4cm}} (\textit{fattig}).
  \item Efter eksamen var \underline{\hspace{2.6cm}} (\textit{afslappet}) ved vores aftale i parken helt fantastisk.
\end{enumerate}

\textbf{练习 B：句型升级——将下列 A2 冗长表达改写为 B1 形容词名词化结构}
\begin{enumerate}
  \item \textit{Unge mennesker bruger meget tid på internettet.} $\rightarrow$ \underline{\hspace{2.0cm}} \textit{bruger meget tid på internettet.}
  \item \textit{Den vigtigste ting ved hygge er fællesskabet.} $\rightarrow$ \underline{\hspace{2.4cm}} \textit{ved hygge er fællesskabet.}
  \item \textit{Ældre mennesker foretrækker ofte en rolig eftermiddag.} $\rightarrow$ \underline{\hspace{2.0cm}} \textit{foretrækker ofte en rolig eftermiddag.}
  \item \textit{Det er svært at forstå den danske humor i starten.} $\rightarrow$ \underline{\hspace{2.2cm}} \textit{er at forstå den danske humor i starten.}
  \item \textit{Jeg kan godt lide de enkle ting ved en skovtur.} $\rightarrow$ \textit{Jeg kan godt lide} \underline{\hspace{2.2cm}} \textit{ved en skovtur.}
  \item \textit{Det er bedst, hvis alle kan være sig selv.} $\rightarrow$ \underline{\hspace{2.2cm}} \textit{er, hvis alle kan være sig selv.}
\end{enumerate}

\textbf{练习 C：中译丹（灵活运用 \textit{det / de + 形容词-e}）}
\begin{enumerate}
  \item 关于 hygge 最重要的是安静的氛围，而不是昂贵的东西。
  \item 无论是年轻人还是老年人，都喜欢在周日聚在一起喝咖啡。
  \item 搬到丹麦最难的是学习如何在冬天保持积极。
  \item 与朋友们共度时光好在（det gode ved）大家都可以完全放松。
  \item 这座城市有很多帮助孤独群体的免费活动。
\end{enumerate}

% ---------------------------------------------------------------------
\ncdsubsec{六、DU3 Modul 3 实战微任务与参考答案}

\begin{emailbox}{Skrivning Opgave 2 实战微任务（要求 100 词以上）}
\textbf{Situation}: Du har fået en mail fra din ven, Lucas, som bor i Kina. Han har læst en artikel om dansk ”hygge” og er nysgerrig.\\
\textbf{Opgave}: Skriv et svar til Lucas. Du skal svare på alle hans spørgsmål og bruge mindst tre substantiverede adjektiver (som \textit{det vigtigste, de unge, de ældre, det gode ved...}). Skriv \textbf{mindst 100 ord}.

\vspace{0.3em}
\textit{Hej!\\
Jeg har hørt så meget om ”hygge” i Danmark. Hvad betyder det egentlig for dig, nu hvor du bor der? Hvordan hygger folk i forskellige aldre sig, og hvordan ligner eller adskiller det sig fra den måde, vi slapper af på i vores hjemland?\\
Mange hilsner, Lucas}
\end{emailbox}

\begin{answerbox}
\textbf{Før du læser 参考回答}
\begin{enumerate}
  \item \textit{Det vigtigste ved hygge er det trygge og afslappede fællesskab uden stress. De unge skaber ofte hygge med mad og brætspil i små lejligheder, mens de ældre ofte foretrækker en rolig eftermiddag med kaffe, kage og lange samtaler.}
  \item \textit{Emma synes, at det svære er, at ordet ikke kan oversættes direkte til andre sprog. Hun giver Marco det råd, at man bare skal sætte tid af til hinanden, og hun inviterer ham på aftensmad næste fredag.}
\end{enumerate}

\textbf{练习 A 答案}
\begin{enumerate}
  \item \da{Det vigtigste}（抽象最高级，用 \textit{det} + \textit{-e}）
  \item \da{de unge}；\da{de ældre}（群体泛指用 \textit{de} + \textit{-e}）
  \item \da{Det sværeste}（抽象最高级，用 \textit{det} + \textit{-e}）
  \item \da{de ensomme}（社会群体“孤独的人们”）
  \item \da{det travle}；\da{det enkle}（注意 \textit{enkel $\rightarrow$ enkle} 省略中间的 \textit{e}）
  \item \da{Det gode}（抽象正面特质）
  \item \da{de rige}；\da{de fattige}（群体泛指）
  \item \da{det afslappede}（\textit{afslappet $\rightarrow$ afslappede}）
\end{enumerate}

\textbf{练习 B 答案}
\begin{enumerate}
  \item \da{De unge}（泛指年轻人整体）
  \item \da{Det vigtigste}（代替繁琐的 \textit{Den vigtigste ting}）
  \item \da{De ældre}（泛指老年人整体）
  \item \da{Det svære}（抽象难点）
  \item \da{det enkle}（抽象特质“质朴之处”）
  \item \da{Det bedste}（代替 \textit{Det er bedst} 作主语焦点）
\end{enumerate}

\textbf{练习 C 答案}
\begin{enumerate}
  \item \textit{\da{Det vigtigste} ved hygge er en rolig stemning og ikke \da{det dyre} / dyre ting.}
  \item \textit{Både \da{de unge} og \da{de ældre} kan lide at mødes til kaffe om søndagen.}
  \item \textit{\da{Det sværeste} ved at flytte til Danmark er at lære at holde sig positiv om vinteren.}
  \item \textit{\da{Det gode ved} at være sammen med venner er, at alle kan slappe helt af.}
  \item \textit{Byen har mange gratis aktiviteter, som hjælper \da{de ensomme}.}
\end{enumerate}

\textbf{Modul 3 实战参考范文（124 词）}\\
\textit{Til: Lucas \quad|\quad Emne: Hvordan jeg oplever hygge i Danmark}

\textit{Hej Lucas}

\textit{Tak for din mail! Det er sjovt, at du spørger om hygge. Nu hvor jeg har boet her i et år, har jeg opdaget, at \da{det vigtigste} ikke er stearinlys, men \da{det trygge} og rolige samvær med andre mennesker.}

\textit{Jeg ser lidt forskel på generationerne: \dabold{De unge} mødes tit spontant til fællesspisning eller brætspil, mens \dabold{de ældre} ofte inviterer til kaffe og kage i weekenden. \da{Det gode ved} den danske tradition er, at ingen behøver at imponere hinanden med dyre retter. Det ligner faktisk vores egne hyggelige måltider derhjemme med familie og venner, selvom man i Danmark måske gør endnu mere ud af \da{det enkle} i hverdagen.}

\textit{Hvornår kommer du selv på besøg, så du kan opleve det?}

\textit{Mange hilsner\\
\mbox{[Dit navn]}}

\small\textbf{点评}：全篇正文 124 词，完整回应了来信的三个提问；自然嵌入四处名词化结构（\textit{det vigtigste, det trygge, de unge, de ældre, det gode ved, det enkle}），并用 \textit{mens} 实现两代人的流畅对比，句法复杂度完全达到 Modul 3 B1 优秀水平。
\end{answerbox}

\ncdlesson{7}{Foreningsdanmark: Hvorfor er alle med i en klub?}{社团丹麦：为什么人人都参加俱乐部？}
{DU3 Modul 3（B1）｜对标 NCE 2 后半 \& Modul 3 \textit{Læsning Opgave 2A/2B (Chat)} \& \textit{Samfundsliv}}
{\textbf{因果表达进阶}：结果主语 \textit{skyldes (, at)}（归因于）vs 原因主语 \textit{føre til (, at)}（导致）· 介词 \textit{på grund af} + 名词 vs 连词 \textit{fordi / da} + 从句 · 因果句双向转换}
{丹麦全国拥有超过十万个民间社团。谈论社会现象时，因果逻辑方向绝对不能写反：\textbf{结果 \textit{skyldes} 原因}，而\textbf{原因 \textit{fører til} 结果}！此外，\textit{på grund af} 是介词，后面只能跟名词短语，想接完整从句请用 \textit{fordi} 或 \textit{skyldes, at...}！}

% ---------------------------------------------------------------------
\begin{foerduleaeser}
带着下面两个问题阅读 Tekst A 和 Tekst B：
\begin{enumerate}
  \item \textit{Hvad skyldes det ifølge Tekst A, at priserne på medlemskab i danske foreninger kan holdes nede?}（根据 Tekst A，丹麦社团的会员价格能保持低廉归因于什么？）
  \item \textit{Hvorfor tøver Kasper med at melde sig ind i løbeklubben i Tekst B, og hvad fører Sofies forslag til?}（在 Tekst B 的聊天中，卡斯珀为什么犹豫要不要加入跑步俱乐部？索菲的建议带来了什么结果？）
\end{enumerate}
\end{foerduleaeser}

% ---------------------------------------------------------------------
\ncdsubsec{二、Tekst A \& Tekst B（文化深读与 Modul 3 真题场景）}

\subsection*{Tekst A：\textit{Foreningsdanmark: Hvorfor er alle med i en klub?}（文化深读短文 · 约 240 词）}

\begin{tipbox}{语篇说明}
本篇剖析丹麦极具特色的“社团文化”（\textit{Foreningsdanmark}）。阅读时请特别区分\textbf{\color{NCDRed}红色标记的因果动词方向}（\textit{skyldes} 结果$\leftarrow$原因 vs \textit{føre til} 原因$\rightarrow$结果）以及\textbf{\color{NCDNavy}深蓝色标记的介词原因短语}（\textit{på grund af} + 名词）。
\end{tipbox}

I Danmark er næsten halvdelen af befolkningen aktiv i en forening eller en klub i både byer og små landsbyer. Det kan være alt fra en lokal fodboldklub og et kor til en spejdergruppe eller en haveforening. Denne særlige tradition, som ofte kaldes ”Foreningsdanmark”, spiller en stor rolle for både tillid og integration i hverdagen.

Men hvad \dabold{skyldes} den store interesse for foreningslivet egentlig? Det \dabold{skyldes først og fremmest, at} kontingentet som regel er lavt, så de fleste familier har råd til at være med. At priserne kan holdes nede, \dabold{skyldes} blandt andet økonomisk støtte fra kommunen og det enorme arbejde fra frivillige trænere. Derudover melder mange danskere sig ind \da{på grund af} ønsket om at komme væk fra arbejdsstress og møde nye mennesker.

Det frivillige engagement \dabold{fører til, at} børn og voksne kan dyrke hundredvis af forskellige aktiviteter tæt på deres hjem. Samtidig \dabold{fører} det faste samarbejde i bestyrelsen \dabold{til} en stærk demokratisk forståelse: Ved den årlige generalforsamling lærer medlemmerne at diskutere, lytte og træffe beslutninger i fællesskab. Flere undersøgelser viser endda, at aktiv deltagelse i foreninger \dabold{fører til} større tillid mellem borgerne i hele landet.

For nye tilflyttere kan foreningslivet dog virke lidt uoverskueligt i starten. Nogle tøver \da{på grund af} sproget, eller fordi de ikke kender de skrevne og uskrevne regler i klubben. Alligevel peger eksperter på, at et medlemskab ofte \dabold{fører til} bedre sprogkundskaber og dybere venskaber, fordi man mødes omkring en fælles passion.

\begin{transbox}
\small
\textbf{\color{NCDTeal}[Tekst A 中文精译]}\\
在丹麦，无论城市还是小村庄，将近一半的人口都活跃在某个社团或俱乐部里。项目从当地足球俱乐部、合唱团，到童子军或花园协会应有尽有。这项常被称为“社团丹麦”的独特传统，在日常生活中对社会信任和融合起着重要作用。

但人们对社团生活如此热衷，究竟归因于什么呢？首先是因为会费通常很低，大多数家庭都负担得起。价格之所以能压低，部分原因在于市政厅的经费资助以及志愿教练们付出的大量心血。此外，许多丹麦人报名入会，也是出于想暂时脱离工作压力、结识新朋友的愿望。

这种志愿参与使得儿童和成年人能在家附近参加数百种不同活动。同时，理事会中的常态化协作也带来了强烈的民主意识：在一年一度的会员全体大会上，成员们学会共同讨论、倾听并做出决策。多项调查甚至表明，积极参与社团还能增进全国公民之间的信任。

不过，对于新搬来的人而言，社团生活一开始可能让人觉得头绪繁杂。有些人因为语言原因而犹豫，或者因为不了解俱乐部里的成文及不成文规定。即便如此，专家指出，加入会员往往能带来更好的语言能力和更深厚的友谊，因为大家因共同的热爱而聚在一起。
\end{transbox}

\subsection*{Tekst B：\textit{Chat om at melde sig ind i en klub}（对标 Modul 3 Læsning Opgave 2A/2B）}

\textbf{阅读微任务（Læsning Opgave 2 聊天补全）：}\textit{Læs chatten mellem Sofie og Kasper. Vælg den sætning (A--E), der passer i hvert hul (1--3). Der er to sætninger for meget.}

\begin{emailbox}{社团报名聊天｜Sofie \& Kasper}
\textbf{Sofie:} Hej Kasper! Skal du med til prøvetræning i Aarhus Løbeklub på torsdag kl. 17?\\
\textbf{Kasper:} Hej Sofie! \textit{(Z) Jeg vil egentlig gerne, men jeg er lidt usikker på, om jeg kan nå det.}\\
\textbf{Sofie:} Okay, hvorfor det? Er det stadig \da{på grund af} dit nye projekt på arbejdet?\\
\textbf{Kasper:} \underline{\hspace{2.8cm}} \textbf{(1)}\\
\textbf{Sofie:} Nå, det forstår jeg godt. Men vores træning \dabold{skyldes ikke} kun motion --- vi drikker også kaffe bagefter, og det \dabold{fører tit til}, at man lærer nye naboer at kende.\\
\textbf{Kasper:} \underline{\hspace{2.8cm}} \textbf{(2)}\\
\textbf{Sofie:} Det koster kun 150 kr. for et halvt år, fordi vi får kommunalt tilskud.\\
\textbf{Kasper:} \underline{\hspace{2.8cm}} \textbf{(3)}\\
\textbf{Sofie:} Super! Så skriver jeg os begge på listen til på torsdag.
\end{emailbox}

\noindent\small\textbf{备选句子（A--E）：}\\
\textbf{A.} \textit{Hvor meget koster kontingentet egentlig, hvis man melder sig ind?}\\
\textbf{B.} \textit{Hold da op, det er billigt! Så prøver jeg gerne på torsdag.}\\
\textbf{C.} \textit{Ja, og det fører desværre til, at jeg tit først har fri kl. 17.}\\
\textbf{D.} \textit{Jeg har allerede løbet tre maratoner i år på grund af min træner.}\\
\textbf{E.} \textit{Det lyder hyggeligt! Men jeg skal også tænke på prisen.}

\begin{tipbox}{本课发音与吞音要点（Udtale \& Reduktioner）}
\begin{itemize}
  \item \textbf{极高频“静音 \textit{d}”陷阱}：\textit{skyldes} 必须读成 \ipa{ˈsgyləs}（\textbf{\textit{d} 完全不发音}，重音在第一音节且无长音！）；同样 \textit{på grund af} \ipa{pʰʌ ˈɡʁɔnˀ ˈæːˀ}（口语常弱读为 \ipa{pʰʌ ˈɡʁɔnˀ ɑ}，\textit{grund} 的 \textit{d} 不发音，\textit{af} 的 \textit{f} 经常脱落！）。
  \item \textbf{社团机构词重音}：\textit{forening} \ipa{fʌˈeːˀneŋ}（重音在第二音节 \textit{-e-}）；\textit{kontingent} \ipa{kʰʌnteŋˈɡɛnˀd}（重音在最后 \textit{-gent}）；\textit{bestyrelse} \ipa{beˈsdyːˀɐlsə}；\textit{generalforsamling} \ipa{ɡenəˈʁɑːˀlfʌˌsɑmˀleŋ}。
  \item \textbf{不发音 \textit{g} 与软 \textit{d}}：\textit{frivillig} \ipa{ˈfʁiˌvili}（词尾 \textit{g} 不发音）；\textit{tilskud} \ipa{ˈtˢelsɡuð}（末尾软 \textit{d}）；\textit{fører til} \ipa{ˈføːɐ tˢe}。
\end{itemize}
\end{tipbox}

\begin{transbox}
\small
\textbf{\color{NCDTeal}[Tekst B 中文对照]}\\
\textbf{索菲}：嗨，卡斯珀！这周四 17 点你一起去奥胡斯跑步俱乐部试训吗？\\
\textbf{卡斯珀}：（例 \textit{Z}）我其实挺想去的，但我不太确定能不能赶得上。\\
\textbf{索菲}：好吧，为什么呢？还是因为你工作上的新项目吗？\\
\textbf{卡斯珀}：（1）是的，而且这遗憾地导致我经常到 17 点才下班。\\
\textbf{索菲}：啊，我完全理解。不过我们训练不只是为了运动——之后还会一起喝咖啡，这经常会让人认识新邻居呢。\\
\textbf{卡斯珀}：（2）听起来很惬意！不过我也得考虑一下价格。\\
\textbf{索菲}：半年只需要 150 克朗，因为我们拿到了市政补贴。\\
\textbf{卡斯珀}：（3）天哪，这也太便宜了！那我周四很乐意去试试。\\
\textbf{索菲}：太棒了！那我把我们俩都写进周四的名单里。
\end{transbox}

% ---------------------------------------------------------------------
\ncdsubsec{三、Gloser（核心生词表）}

\begin{tipbox}{词表重点}
动词 \textit{at skyldes} 是“异态动词（deponent verbum）”，原形和现在时永远自带词尾 \textit{-s}（\textit{skyldes}），过去时为 \textit{skyldtes}，但表达的是主动语义“归因于”。注意它后面直接接名词或 \textit{, at} 从句，千万不要加多余的介词 \textit{af} 或 \textit{til}！
\end{tipbox}

\begin{glosertable}
动词 [har] & at skyldes & \textit{(skyldes, skyldtes, har skyldtes)} & 归因于，由……引起（结果在前） \\
动词 [har] & at føre til & \textit{(fører til, førte til, har ført til)} & 导致，带来（原因在前） \\
动词 [har] & at melde sig ind (i) & \textit{(melder, meldte, har meldt sig ind)} & 报名加入（俱乐部/协会） \\
动词 [har] & at deltage (i) & \textit{(deltager, deltog, har deltaget)} & 参加，参与 \\
动词 [har] & at træffe & \textit{(træffer, traf, har truffet)} & 做出（决定 \textit{beslutninger}）；遇见 \\
动词 [har] & at tøve & \textit{(tøver, tøvede, har tøvet)} & 犹豫，迟疑 \\
动词 [har] & at pege (på) & \textit{(peger, pegede, har peget)} & 指出，表明 \\
\midrule
名词 & en forening & \textit{(foreningen, foreninger, foreningerne)} & 社团，协会 \\
名词 & et kontingent & \textit{(kontingentet, kontingenter, kontingenterne)} & 会费，成员费 \\
名词 & en støtte & \textit{(støtten, ingen flertal)} & 支持，资助 \\
名词 & en bestyrelse & \textit{(bestyrelsen, bestyrelser, bestyrelserne)} & 理事会，管理委员会 \\
名词 & en generalforsamling & \textit{(generalforsamlingen, generalforsamlinger, generalforsamlingerne)} & 会员全体大会 \\
名词 & en deltagelse & \textit{(deltagelsen, ingen flertal)} & 参与，参加 \\
名词 & en sprogkundskab & \textit{(sprogkundskaben, sprogkundskaber, sprogkundskaberne)} & 语言能力，语言知识 \\
\midrule
形容词 & frivillig & \textit{(frivilligt, frivillige)} & 志愿的，义务的 \\
形容词 & uoverskuelig & \textit{(uoverskueligt, uoverskuelige)} & 头绪繁杂的，难以驾驭的 \\
介词短语 & på grund af & \textit{(固定介词短语 + 名词短语)} & 由于，因为 \\
\end{glosertable}

% ---------------------------------------------------------------------
\ncdsubsec{四、Grammatik \& Faste vendinger（核心语法与地道表达精讲）}

\begin{grammatikbox}{语法点 1：因果方向的黄金天平——\textit{skyldes} vs \textit{føre til} vs \textit{på grund af}}
在 B1 议论文和说明文中，我们需要熟练切换“由果溯因”和“由因推果”两种视角，确保主语与宾语方向精准无误：
\begin{itemize}
  \item \textbf{视角 \ding{172}｜由果溯因：\da{结果} + \dabold{skyldes (, at)} + \da{原因}}：主语是已经发生的结果，后面接导致它的原因。既可直接接名词，也可加逗号接 \textit{at} 从句。例：\textit{De lave priser \dabold{skyldes} kommunens støtte.} / \textit{Det \dabold{skyldes, at} alle arbejder frivilligt.}
  \item \textbf{视角 \ding{173}｜由因推果：\da{原因} + \dabold{fører til (, at)} + \da{结果}}：主语是起因、举措或条件，后面引出产生的结果。同样可直接接名词，或加逗号接 \textit{at} 从句。例：\textit{Frivilligt arbejde \dabold{fører til} billigere klubber.} / \textit{Et medlemskab \dabold{fører ofte til, at} man får nye venner.}
  \item \textbf{视角 \ding{174}｜介词原因：\dabold{på grund af} + \da{名词短语}}：相当于英语的 \textit{because of}，后面\textbf{绝对只能跟名词或动名词化短语}（如 \textit{på grund af sproget / på grund af det dårlige vejr}）。若放句首，别忘了触发\textbf{主句 V2 倒装}：\textit{På grund af arbejdet \dabold{har} han ikke tid.}
\end{itemize}
\end{grammatikbox}

{\small
\renewcommand{\arraystretch}{1.25}
\begin{xltabular}{\linewidth}{@{} >{\bfseries\color{NCDNavy}\raggedright\arraybackslash}p{2.9cm} >{\raggedright\arraybackslash}p{3.3cm} >{\raggedright\arraybackslash}p{3.2cm} >{\raggedright\arraybackslash}X @{}}
\toprule
\textbf{因果表达} & \textbf{语法词性与结构} & \textbf{逻辑方向} & \textbf{标准 B1 例句} \\
\midrule
\endfirsthead
\toprule
\textbf{因果表达} & \textbf{语法词性与结构} & \textbf{逻辑方向} & \textbf{标准 B1 例句} \\
\midrule
\endhead
\bottomrule
\endfoot
\textit{skyldes} / \textit{skyldtes} & 动词（不带介词）\newline + 名词 / \textit{, at} 从句 & 结果 $\leftarrow$ 原因 & \textit{Succesen \dabold{skyldes, at} medlemmerne hjælper til.} \\
\midrule
\textit{føre til} / \textit{førte til} & 动词 + 介词 \textit{til}\newline + 名词 / \textit{, at} 从句 & 原因 $\rightarrow$ 结果 & \textit{Samarbejdet \dabold{fører til} større tillid i byen.} \\
\midrule
\textit{på grund af} & 介词短语\newline + \textbf{仅限名词短语} & 说明原因状语 & \textit{Han tøver \da{på grund af} sit manglende dansk.} \\
\midrule
\textit{fordi} / \textit{da} & 从句连词\newline + \textbf{完整主谓从句} & 说明原因（从句不倒装） & \textit{Han tøver, \da{fordi} han \dabold{ikke taler} dansk endnu.} \\
\end{xltabular}
}

\begin{grammatikbox}{语法点 2：英汉对比排雷——四大高频考场扣分点}
\begin{enumerate}
  \item \textbf{严禁写成 *\textit{skyldes af / skyldes til}！}：英语说 \textit{is due to}，导致很多同学套用成 *\textit{Det skyldes til...} 或 *\textit{Det skyldes af...}。大错特错！丹麦语的 \textit{skyldes} 本身已经包含“归因于”的意思，后面\textbf{直接接原因名词或 \textit{, at}}，不加任何介词！
  \item \textbf{\textit{på grund af} 后禁止直接加句子！}：书面语里绝不能写 *\textit{på grund af han arbejder meget}。需要接主谓结构时，要么换回连词 \textit{\dabold{fordi} han arbejder meget}，要么把后面浓缩成名词短语 \textit{på grund af \dabold{sit arbejde}}！
  \item \textbf{因果方向别倒置}：中文“工作压力导致他没时间”或“他没时间是因为工作压力”，在丹麦语中不能混淆主语。请检查：\textit{Hans travlhed \dabold{fører til} manglende tid.}（压力$\rightarrow$没时间）；\textit{Den manglende tid \dabold{skyldes} hans travlhed.}（没时间$\leftarrow$压力）。
  \item \textbf{\textit{føre til} 后接动作要用 \textit{at + 原形} 或 \textit{, at + 从句}}：受英语 \textit{lead to + V-ing} 影响，不能只写名词外加光杆动词，要写 \textit{fører til \dabold{at man lærer}} 或 \textit{fører til \dabold{bedre sprog}}.
\end{enumerate}
\end{grammatikbox}

% ---------------------------------------------------------------------
\ncdsubsec{五、Mønsterøvelser（句型操练）}

\textbf{练习 A：用 \textit{skyldes, skyldtes, fører til, førte til, på grund af} 或 \textit{fordi} 填空}
\begin{enumerate}
  \item Den lave pris på kontingentet \underline{\hspace{2.2cm}} økonomisk støtte fra kommunen. (现在时)
  \item Mange udlændinge tøver med at melde sig ind \underline{\hspace{2.2cm}} sproget.
  \item Et aktivt foreningsliv \underline{\hspace{2.4cm}}, at borgerne får et stærkere netværk.
  \item Klubben blev lukket sidste år \underline{\hspace{2.2cm}} mangel på frivillige trænere.
  \item At mødet blev aflyst i går, \underline{\hspace{2.2cm}} en stor snestorm. (过去时)
  \item Han har meldt sig ind i et kor, \underline{\hspace{2.2cm}} han gerne vil øve sin danske udtale.
  \item Manglen på nye medlemmer \underline{\hspace{2.4cm}} i sidste måned, at klubben måtte hæve prisen. (过去时)
  \item Det gode samarbejde i bestyrelsen \underline{\hspace{2.2cm}}, at alle tager ansvar.
\end{enumerate}

\textbf{练习 B：句型双向转换——按提示翻转因果视角或互换“介词短语 vs 从句”}
\begin{enumerate}
  \item 用 \textit{skyldes, at} 改写（结果放前面）：\textit{At trænerne arbejder frivilligt, fører til, at priserne er lave.} $\rightarrow$ \textit{At priserne er lave,...}
  \item 用 \textit{fører til, at} 改写（原因放前面）：\textit{De mange nye venner skyldes, at han spiller fodbold.} $\rightarrow$ \textit{At han spiller fodbold,...}
  \item 用 \textit{på grund af} + 名词替换 \textit{fordi} 从句：\textit{Hun meldte afbud, fordi det regnede kraftigt.} $\rightarrow$ \textit{Hun meldte afbud...}
  \item 把 \textit{På grund af} 短语放到句首（注意 V2 倒装！）：\textit{Vi får ikke råd til en ny hal på grund af de høje priser.} $\rightarrow$ \textit{På grund af de høje priser...}
  \item 用 \textit{fordi} 改写（注意从句副词位置）：\textit{Han kom ikke til træning på grund af sin sygdom.} $\rightarrow$ \textit{Han kom ikke til træning, fordi han...} (\textit{ikke var rask / var syg})
  \item 纠错改写：\textit{Den store interesse skyldes af, at klubben har mange tilbud.} $\rightarrow$
\end{enumerate}

\textbf{练习 C：中译丹（精准运用本课因果结构）}
\begin{enumerate}
  \item 会员数量的增长归因于该俱乐部拥有友好的氛围。
  \item 由于工作太多，她暂时没有时间参加理事会的工作。
  \item 在当地社团做志愿者工作往往会带来更强的地方归属感。
  \item 我报名参加羽毛球俱乐部，是因为我想在工作之余结识新朋友。
  \item 高昂的房租导致许多年轻人不得不搬到郊区。
\end{enumerate}

% ---------------------------------------------------------------------
\ncdsubsec{六、DU3 Modul 3 实战微任务与参考答案}

\begin{emailbox}{Skrivning Opgave 1 推荐与评价微任务（要求 100 词以上）}
\textbf{Situation}: Din veninde eller kollega overvejer at melde sig ind i en lokal forening, men hun er lidt i tvivl, fordi hun ikke taler perfekt dansk endnu.\\
\textbf{Opgave}: Skriv en mail til hende, hvor du anbefaler en lokal forening eller klub. Du skal skrive om:
\begin{itemize}
  \item hvordan foreningslivet fungerer i Danmark, og hvad de lave priser skyldes;
  \item hvad du selv har fået ud af at være med;
  \item hvad et medlemskab kan føre til for hendes sprog og netværk;
  \item et konkret forslag til næste skridt.
\end{itemize}
Skriv \textbf{mindst 100 ord}.
\end{emailbox}

\begin{answerbox}
\textbf{Før du læser 参考回答}
\begin{enumerate}
  \item \textit{Det skyldes blandt andet økonomisk støtte fra kommunen og det store arbejde fra frivillige trænere.}
  \item \textit{Kasper tøver, fordi hans nye projekt på arbejdet tit fører til, at han først har fri kl. 17. Sofies forslag og oplysning om det lave kontingent fører til, at Kasper gerne vil prøve træningen på torsdag.}
\end{enumerate}

\textbf{Tekst B 聊天补全答案（Modul 3 Læsning Opgave 2 风格）}
\begin{enumerate}
  \item[(1)] \da{C}（呼应前文询问“是否因为工作项目”，接续说明这导致经常 17 点才下班）
  \item[(2)] \da{E}（接续“听起来很惬意”，但转折引出“价格”问题，刚好引出 Sofie 下一句回答 \textit{Det koster kun 150 kr.}）
  \item[(3)] \da{B}（得知半年只要 150 克朗后惊喜接受周四试训，与结尾 Sofie 说“我把我们俩都写上名单”无缝衔接；\textit{A} 为多余干扰项，因为下文索菲没等他问就已经主动报出了价格，\textit{D} 与语境无关）
\end{enumerate}

\textbf{练习 A 答案}
\begin{enumerate}
  \item \da{skyldes}（价格低归因于市政支持）
  \item \da{på grund af}（后接名词 \textit{sproget}）
  \item \da{fører til}（原因主语 \textit{foreningsliv} 带来结果）
  \item \da{på grund af}（后接名词短语 \textit{mangel på...}）
  \item \da{skyldtes}（昨天发生，过去时）
  \item \da{fordi}（后接完整主谓从句 \textit{han gerne vil...}）
  \item \da{førte til}（上个月发生，过去时）
  \item \da{fører til}（原因 $\rightarrow$ 结果）
\end{enumerate}

\textbf{练习 B 答案}
\begin{enumerate}
  \item \textit{At priserne er lave, \da{skyldes, at} der er frivillige trænere / frivillige trænere.}
  \item \textit{At han spiller fodbold, \da{fører til, at} han har mange nye venner.}
  \item \textit{Hun meldte afbud \da{på grund af} det kraftige regnvejr / regnen.}
  \item \textit{På grund af de høje priser \dabold{får vi} ikke råd til en ny hal.}（V2 倒装）
  \item \textit{Han kom ikke til træning, fordi han \da{var syg / ikke var rask}.}（注意从句中 \textit{ikke} 在动词前）
  \item \textit{Den store interesse \dabold{skyldes, at} klubben har mange tilbud.}（删掉多余的 \textit{af}）
\end{enumerate}

\textbf{练习 C 答案}
\begin{enumerate}
  \item \textit{Stigningen i antallet af medlemmer \da{skyldes, at} klubben har en venlig stemning.}
  \item \textit{\da{På grund af} for meget arbejde har hun foreløbig ikke tid til at arbejde i bestyrelsen.}
  \item \textit{Frivilligt arbejde i en lokal forening \da{fører ofte til} et stærkere tilhørsforhold til området.}
  \item \textit{Jeg har meldt mig ind i en badmintonklub, \da{fordi} jeg gerne vil møde nye venner efter arbejde.}
  \item \textit{Den høje husleje \da{fører til, at} mange unge må flytte ud til forstæderne.}
\end{enumerate}

\textbf{Modul 3 实战参考范文（120 词）}\\
\textit{Til: Sara \quad|\quad Emne: Skal du ikke med i min lokale forening?}

\textit{Hej Sara}

\textit{Jeg hørte, at du overvejer at starte i en klub her i byen, og jeg vil varmt anbefale min egen badmintonforening. I Danmark er foreningslivet åbent for alle, og de lave priser \dabold{skyldes}, at kommunen giver støtte, mens trænerne arbejder frivilligt.}

\textit{Jeg meldte mig selv ind for et år siden \da{på grund af} ønsket om at få mere motion og lære naboerne at kende. I starten var jeg også nervøs for sproget, men alle var meget tålmodige. Det faste fællesskab \dabold{fører hurtigt til, at} man taler mere afslappet dansk og udvider sit netværk.}

\textit{Hvad siger du til at tage med mig som gæst næste tirsdag kl. 18? Så kan du prøve det helt gratis!}

\textit{Kærlig hilsen\\
\mbox{[Dit navn]}}

\small\textbf{点评}：全篇正文 120 词，精准覆盖题干四个要求；自然运用 \textit{skyldes}, \textit{på grund af} 以及 \textit{fører hurtigt til, at}，完美展示了因果方向的精准操控能力，并给出清晰的行动建议。
\end{answerbox}

\ncdlesson{8}{Janteloven i det moderne samfund}{现代社会中的《詹代法则》}
{DU3 Modul 3（B1）｜对标 NCE 2 后半 \& Modul 3 \textit{Læsning Opgave 3 (Tekstudfyldning)} \& \textit{Samfund og Værdier}}
{\textbf{情态动词的主观推测与传闻用法（Epistemisk modalitet）}：高度肯定推测 \textit{må (have)} · 可能性推测 \textit{kan (ikke)} · 传闻转述 \textit{skal / skulle (efter sigende) være} · 理应推断 \textit{bør / burde}}
{在 Modul 3 谈论社会现象与职场文化时，直接下绝对断言往往不够客观。学会用 \textbf{\textit{må være}}（一定是）、\textbf{\textit{kan være}}（可能是）和 \textbf{\textit{skulle efter sigende være}}（据说是）来表达推断与审慎转述，既能显著提升语体成熟度，也是阅读理解判断作者立场的必考线索！}

% ---------------------------------------------------------------------
\begin{foerduleaeser}
带着下面两个问题阅读 Tekst A 和 Tekst B：
\begin{enumerate}
  \item \textit{Hvilke fordele og ulemper kan Janteloven have på arbejdspladsen ifølge Tekst A?}（根据 Tekst A，《詹代法则》在职场上可能带来哪些利与弊？）
  \item \textit{Hvordan oplevede skribenten i Tekst B sit første projektmøde i et dansk team, og hvad skyldtes den rolige stemning?}（在 Tekst B 中，作者在丹麦团队的第一次项目会议体验如何？那种平静的氛围可能源于什么？）
\end{enumerate}
\end{foerduleaeser}

% ---------------------------------------------------------------------
\ncdsubsec{二、Tekst A \& Tekst B（文化深读与 Modul 3 真题场景）}

\subsection*{Tekst A：\textit{Janteloven i det moderne samfund}（文化深读短文 · 约 240 词）}

\begin{tipbox}{语篇说明}
1933 年作家 Aksel Sandemose 写下著名的《詹代法则》（\textit{Janteloven}），深刻影响了北欧的平等与低调文化。阅读时请重点留意\textbf{\color{NCDNavy}深蓝色标出的主观推测情态动词}（\textit{må, kan, vil, bør}）与\textbf{\color{NCDRed}深红色标出的传闻与否定推测用法}（\textit{skulle efter sigende, kan ikke}）。
\end{tipbox}

I 1933 skrev den dansk-norske forfatter Aksel Sandemose romanen \textit{En flygtning krydser sit spor}, hvor han formulerede ti regler kendt som ”Janteloven”. Den første og mest berømte regel lyder: ”Du skal ikke tro, at du er noget.” Selvom bogen blev skrevet for næsten hundrede år siden, \da{kan man stadig mærke} dens indflydelse i dagens Danmark.

Når udlændinge hører om Janteloven første gang, \da{må det være} svært helt at forstå, hvorfor beskedenhed fylder så meget. For mange danskere \da{må lighed være} en af samfundets vigtigste værdier, og man ønsker ikke, at nogen praler af sin succes eller ser ned på andre. På den ene side \da{kan denne holdning skabe} et trygt arbejdsmiljø, hvor chefen og medarbejderne behandler hinanden ligeværdigt. På den anden side \da{kan Janteloven også være} en bremse, hvis dygtige unge bliver bange for at skille sig ud.

I dag diskuterer man meget, om de gamle regler stadig gælder. Janteloven \dabold{skulle efter sigende være} blevet svagere i de store byer, især blandt unge iværksættere og på sociale medier. Flere sociologer peger på, at kulturen \da{må have ændret sig}, fordi man nu i højere grad opmuntrer hinanden til at være stolte af gode resultater. Alligevel \dabold{kan det ikke passe}, at Janteloven er helt forsvundet. Hvis man praler for meget, \da{vil} reaktionen ofte \da{være} kølig stilhed. Derfor \da{bør} man som ny i Danmark både turde vise sine evner og samtidig huske respekten bag den danske ligeværdighed.

\begin{transbox}
\small
\textbf{\color{NCDTeal}[Tekst A 中文精译]}\\
1933 年，丹麦—挪威作家阿克塞尔·桑德莫斯写下了小说《逃亡者跨越自己的足迹》，在书中提出了被称为“詹代法则”的十条规则。第一条也是最著名的一条是：“不要以为你有什么了不起。”尽管这本书写于将近一百年前，人们在今天的丹麦依然能感受到它的影响。

外国人初次听说《詹代法则》时，一定会觉得很难完全理解为什么谦虚会占据这么大的分量。对许多丹麦人来说，平等必然是社会最重要的价值观之一，大家不希望有人炫耀自己的成功或轻视他人。一方面，这种态度可以营造出安全融洽的工作氛围，让老板和员工彼此平等相待；另一方面，如果优秀的年轻人因此不敢脱颖而出，《詹代法则》也可能成为一种阻力。

如今人们经常讨论这些老规矩是否依然适用。据说《詹代法则》在大城市里已经变弱了，尤其是在年轻创业者群体和社交媒体上。多位社会学家指出，这种文化一定已经发生了变化，因为现在人们越来越鼓励彼此为好成绩感到自豪。尽管如此，说《詹代法则》已经彻底消失是不可能成立的。如果一个人吹嘘过头，周围的反应往往会是一片冷淡的沉默。因此，初来丹麦的人既应该敢于展示自己的能力，同时也要记住这份丹麦式平等背后的相互尊重。
\end{transbox}

\subsection*{Tekst B：\textit{Debatindlæg: Mit første projektmøde i Danmark}（对标 Modul 3 Læsning Opgave 3）}

\textbf{阅读微任务（Læsning Opgave 3 段落逻辑补全）：}\textit{Læs den korte kommentar om dansk arbejdskultur, og vælg den sætning (A--D), der passer bedst i hullet [1].}

\begin{emailbox}{职场文化短评｜Ny i et dansk team}
Da jeg startede i mit danske it-team, troede jeg, at jeg skulle sælge mine idéer meget højlydt for at blive hørt. Jeg havde nemlig læst, at konkurrence er vigtig i internationale firmaer. \underline{\hspace{2.8cm}} \textbf{[1]} På mit allerførste projektmøde lagde jeg mærke til, at min chef talte meget roligt og lyttede lige så meget til praktikanterne som til de erfarne udviklere. Derudover drak alle kaffe sammen i pausen og snakkede om weekendens planer uden at vise sig frem. Det \da{må være} et tydeligt eksempel på, hvordan Janteloven stadig præger samarbejdet. Selvom loven \dabold{skulle efter sigende være} mindre synlig i dag, trives beskedenhed stadig på kontoret.
\end{emailbox}

\noindent\small\textbf{备选句子（A--D）：}\\
\textbf{A.} \textit{Derimod må man aldrig stille spørgsmål til sin chefs beslutninger.}\\
\textbf{B.} \textit{Men sådan fungerer et dansk kontor altså ikke altid.}\\
\textbf{C.} \textit{For vores afdeling skulle efter sigende lukke næste måned.}\\
\textbf{D.} \textit{Derfor pralede alle mine nye kolleger meget af deres løn.}

\begin{tipbox}{本课发音与吞音要点（Udtale \& Reduktioner）}
\begin{itemize}
  \item \textbf{传闻表达 \textit{skulle efter sigende} 的隐藏吞音}：\textit{sigende} 拼写里有 \textit{g} 和 \textit{d}，但实际标准读音是 \ipa{ˈsiːənə}（\textbf{\textit{g} 完全不发音，\textit{nd} 中的 \textit{d} 也不发音}！）；整串搭配在口语里常读作 \ipa{sgu ˈɛfdɐ ˈsiːənə}。
  \item \textbf{法则与价值观核心词}：\textit{Janteloven} \ipa{ˈjandəˌlʌʊ̯ˀən}（\textit{-lov-} 读合口双元音 \ipa{ʌʊ̯ˀ}）；\textit{samfund} \ipa{ˈsɑmˌfɔnˀ}（词尾 \textit{d} 不发音）；\textit{beskedenhed} \ipa{beˈsgeːˀðənˌheːˀð}（软 d）。
  \item \textbf{推测情态词的弱读与重音}：作主观推测时，情态动词常轻读或次重读，重音落在后面的动词或表语上：\textit{Det må \textbf{være}} \ipa{de mʌ ˈvɛːɐ}、\textit{Det kan \textbf{ikke}} \ipa{de kʰa ˈeɡə}（口语极快时 \textit{kan} 缩为 \ipa{kʰa}）；\textit{oprindeligt} \ipa{ʌbˈʁenˀəli}（\textit{d} 与词尾 \textit{gt} 均不发音）。
\end{itemize}
\end{tipbox}

\begin{transbox}
\small
\textbf{\color{NCDTeal}[Tekst B 中文对照]}\\
\textbf{初入丹麦团队}：刚加入我的丹麦 IT 团队时，我原以为必须大声推销自己的想法才能被听见，因为我以前读到过竞争在跨国公司里很重要。[1]（\textbf{但丹麦办公室其实并不总是那样运作的。}）在我的第一次项目会议上，我注意到老板说话非常平和，并且听取实习生意见的认真程度不亚于对待资深开发者。此外，大家在休息时一起喝咖啡、聊周末计划，毫不炫耀。这一定是《詹代法则》至今仍在影响团队协作的一个明显例子。尽管如今这条法则据说已经没那么明显，但谦逊低调在办公室里依然很有生命力。
\end{transbox}

% ---------------------------------------------------------------------
\ncdsubsec{三、Gloser（核心生词表）}

\begin{tipbox}{词表重点}
情态动词 \textit{at burde} 的现在时是特殊的 \textit{bør}（不规则无 \textit{-r} 以外的词干变音！），过去时是 \textit{burde}。另外注意 \textit{at skille sig ud}（脱颖而出）是反身动词短语，不要漏掉反身代词 \textit{sig / mig / dig}；\textit{grundlæggende} 作定语和表语时词形始终不变。
\end{tipbox}

\begin{glosertable}
动词 [har] & at prale (af) & \textit{(praler, pralede, har pralet)} & 炫耀，吹嘘（接介词 \textit{af}） \\
动词 [har] & at skille sig ud & \textit{(skiller, skilte, har skilt sig ud)} & 脱颖而出，显得与众不同 \\
动词 [har] & at opmuntre (til) & \textit{(opmuntrer, opmuntrede, har opmuntret)} & 鼓励，激励 \\
动词 [har] & at behandle & \textit{(behandler, behandlede, har behandlet)} & 对待，处理 \\
动词 [har] & at se ned på & \textit{(ser, så, har set ned på)} & 看不起，轻视 \\
动词 [har] & at ændre (sig) & \textit{(ændrer, ændrede, har ændret)} & 改变，变化 \\
情态动词 [har] & at burde & \textit{(bør, burde, har burdet)} & 应当，按理说 \\
\midrule
名词 & et samfund & \textit{(samfundet, samfund, samfundene)} & 社会 \\
名词 & en lighed & \textit{(ligheden, ligheder, lighederne)} & 平等；相似之处 \\
名词 & en værdi & \textit{(værdien, værdier, værdierne)} & 价值观；价值 \\
名词 & en beskedenhed & \textit{(beskedenheden, ingen flertal)} & 谦虚，低调 \\
名词 & en succes & \textit{(succesen, succeser, succeserne)} & 成功 \\
名词 & en iværksætter & \textit{(iværksætteren, iværksættere, iværksætterne)} & 创业者，企业家 \\
名词 & en bremse & \textit{(bremsen, bremser, bremserne)} & 阻碍，刹车 \\
名词 & en indflydelse & \textit{(indflydelsen, ingen flertal)} & 影响，影响力 \\
\midrule
形容词 & grundlæggende & \textit{(grundlæggende, 不变)} & 根本的，基础的 \\
副词短语 & efter sigende & \textit{(固定副词短语)} & 据说，传闻中 \\
\end{glosertable}

% ---------------------------------------------------------------------
\ncdsubsec{四、Grammatik \& Faste vendinger（核心语法与地道表达精讲）}

\begin{grammatikbox}{语法点 1：情态动词的“主观推测”与“传闻转述”功能}
在 A2 阶段我们学过情态动词的“能力/义务”本义（\textit{kan} 会、\textit{skal} 必须、\textit{må} 允许）。进入 B1 阶段，同一批情态动词用来表达说话人对事实真伪的\textbf{推断把握度}或\textbf{消息来源}（接不带 \textit{at} 的动词原形）：
\begin{itemize}
  \item \textbf{\da{må} + 原形（必然/高度确定，90\%+ 把握）}：根据现有线索合乎逻辑地推断“一定是……”。例：\textit{Hun har boet her i ti år, så hun \da{må kende} Janteloven.}（她在这儿住了十年，肯定了解詹代法则。）
  \item \textbf{\da{kan} + 原形（可能性推测，约 50\% 把握）}：表示“有可能会……”。例：\textit{Reglerne \da{kan virke} strenge for udlændinge.}
  \item \textbf{\dabold{kan ikke} + 原形（绝对不可能）}：断然否定某种可能性。例：\textit{Det \dabold{kan ikke passe}, at alle tænker ens!}
  \item \textbf{\dabold{skal / skulle (efter sigende)} + 原形（传闻转述“据说”）}：表明信息来自他人或舆论，\textbf{说话人自己不作担保}！用过去时形式 \dabold{skulle} 比用 \textit{skal} 更有“姑且一听、保持距离”的客观感。例：\textit{Byen \dabold{skulle efter sigende være} meget åben i dag.}
  \item \textbf{对“过去事件”的推测：\da{må / kan / skulle} + \dabold{have} + 过去分词}：表示“一定已经…… / 据说曾经……”。例：\textit{Kulturen \da{må have ændret sig} gennem de sidste år.}
\end{itemize}
\end{grammatikbox}

{\small
\renewcommand{\arraystretch}{1.25}
\begin{xltabular}{\linewidth}{@{} >{\bfseries\color{NCDNavy}\raggedright\arraybackslash}p{2.9cm} >{\raggedright\arraybackslash}p{3.2cm} >{\raggedright\arraybackslash}p{3.2cm} >{\raggedright\arraybackslash}X @{}}
\toprule
\textbf{情态结构} & \textbf{推测/语义色彩} & \textbf{把握度 / 立场} & \textbf{B1 例句} \\
\midrule
\endfirsthead
\toprule
\textbf{情态结构} & \textbf{推测/语义色彩} & \textbf{把握度 / 立场} & \textbf{B1 例句} \\
\midrule
\endhead
\bottomrule
\endfoot
\textit{må} + 原形 & 逻辑必然推测\newline “一定，肯定” & 高度肯定（约 90\%+） & \textit{Det \da{må være} svært at begynde i et nyt firma.} \\
\midrule
\textit{kan} + 原形 & 开放性可能\newline “可能，也许” & 半信半疑（约 50\%） & \textit{Det \da{kan skabe} en god balance på kontoret.} \\
\midrule
\textit{kan ikke} + 原形 & 逻辑排除\newline “不可能” & 坚决排除 & \textit{Det \dabold{kan ikke være} hele forklaringen.} \\
\midrule
\textit{bør / burde} + 原形 & 合理推断或建议\newline “理应，按说应该” & 期待/规范判断 & \textit{Man \da{bør} også rose unge iværksættere.} \\
\midrule
\textit{skulle (efter sigende)} + 原形 & 传闻转述\newline “据说，听说” & 客观转述，不担保真假 & \textit{Den nye chef \dabold{skulle efter sigende være} meget direkte.} \\
\end{xltabular}
}

\begin{grammatikbox}{语法点 2：中英文对比排雷——三大推测致命误区}
\begin{enumerate}
  \item \textbf{“一定不是……”绝不能用 *\textit{må ikke}！}：很多同学把“一定是 \textit{må være}”直接加 \textit{ikke} 变成 *\textit{Det må ikke være sandt.}——错！在丹麦语中，\dabold{\textit{må ikke}} 永远只表示\textbf{“禁止/不许”}（等同于英语 \textit{must not}）。表达推测“一定不是 / 不可能是”，必须彻底换成 \dabold{\textit{kan ikke være}}！
  \item \textbf{传闻“据说”绝不逐字翻译成 *\textit{det er sagt, at...}！}：英语有 \textit{is said to be}，但丹麦语最地道的对应表达就是情态动词 \dabold{\textit{skal / skulle (efter sigende) være}}（也可以用 \textit{man siger, at...}，但 \textit{skulle efter sigende} 是标准 B1/B2 书面高分标志）。
  \item \textbf{情态动词后面绝不加 \textit{at}！}：哪怕中间插入了 \textit{efter sigende}，后面的不定式依然是光杆原形：\textit{Han skulle efter sigende \dabold{bo}} i Aarhus（绝不是 *\textit{at bo}！）。
\end{enumerate}
\end{grammatikbox}

% ---------------------------------------------------------------------
\ncdsubsec{五、Mønsterøvelser（句型操练）}

\textbf{练习 A：语境选词——用 \textit{må, kan, kan ikke, bør} 或 \textit{skulle} 填空（注意推测与传闻语义）}
\begin{enumerate}
  \item Når han arbejder 60 timer om ugen, \underline{\hspace{1.8cm}} han være meget træt om aftenen.
  \item Det \underline{\hspace{2.2cm}} passe, at hun allerede har glemt aftalen --- hun ringede jo for ti minutter siden! (否定)
  \item Den nye restaurant i midtbyen \underline{\hspace{2.0cm}} efter sigende have den bedste fisk i byen.
  \item Janteloven \underline{\hspace{1.8cm}} både have positive og negative sider for unge medarbejdere.
  \item Vi \underline{\hspace{1.8cm}} give nye kolleger mere plads til at fortælle om deres idéer.
  \item Hun svarer ikke på telefonen; hun \underline{\hspace{1.8cm}} allerede være taget hjem fra arbejde.
  \item Kulturen på arbejdspladser \underline{\hspace{2.0cm}} efter sigende have ændret sig meget de sidste ti år.
  \item Du \underline{\hspace{1.8cm}} ikke prale for meget på dit første jobmøde, hvis du vil skabe et godt indtryk.
\end{enumerate}

\textbf{练习 B：语体改写——将下列直接陈述改写为带情态推测或传闻转述的 B1 句子}
\begin{enumerate}
  \item 用 \textit{må være} 改写：\textit{Jeg er helt sikker på, at det er svært at starte i et nyt land.} $\rightarrow$ \textit{Det...}
  \item 用 \textit{skulle efter sigende være} 改写：\textit{Folk siger, at kulturen i København er mere international i dag.} $\rightarrow$ \textit{Kulturen i København...}
  \item 用 \textit{kan ikke passe} 改写：\textit{Det er helt umuligt, at Janteloven er forsvundet fuldstændigt.} $\rightarrow$ \textit{Det...}
  \item 用 \textit{må have + 分词} 改写（过去推测）：\textit{Jeg er sikker på, at hun har glemt sin paraply på kontoret.} $\rightarrow$ \textit{Hun...}
  \item 用 \textit{kan skabe} 改写：\textit{Måske skaber for meget konkurrence stress på arbejdspladsen.} $\rightarrow$ \textit{For meget konkurrence...}
  \item 纠错改写：\textit{Det må ikke være rigtigt, at han tjener så mange penge i et studiejob!} $\rightarrow$
\end{enumerate}

\textbf{练习 C：中译丹（精准使用 \textit{må, kan (ikke), skulle efter sigende}）}
\begin{enumerate}
  \item 初来乍到时，这一定很难习惯，因为大家都很低调。
  \item 据说这所大学在创新领域非常开放。
  \item 他不可能已经回家了，因为他的包还放在桌子上呢。
  \item 一点点谦虚可以营造出一个更加平等的合作氛围。
  \item 他们昨天一定在办公室里聊得很愉快。
\end{enumerate}

% ---------------------------------------------------------------------
\ncdsubsec{六、DU3 Modul 3 实战微任务与参考答案}

\begin{emailbox}{Skrivning Opgave 2 实战微任务（要求 100 词以上）}
\textbf{Situation}: Din ven Daniel overvejer at søge et job i en dansk virksomhed, men han har hørt om ”Janteloven” og er bange for, at han ikke må vise sine kompetencer.\\
\textbf{Opgave}: Skriv en mail til Daniel. Forklar ham, hvordan Janteloven fungerer på en moderne arbejdsplads, hvad han bør passe på, og hvordan han alligevel kan vise sin erfaring. Brug mindst tre forskellige modale udtryk (f.eks. \textit{må være, kan være, skulle efter sigende, bør}). Skriv \textbf{mindst 100 ord}.
\end{emailbox}

\begin{answerbox}
\textbf{Før du læser 参考回答}
\begin{enumerate}
  \item \textit{På den ene side kan Janteloven skabe et trygt arbejdsmiljø, hvor chefen og medarbejderne behandler hinanden ligeværdigt. På den anden side kan den være en bremse, hvis unge bliver bange for at skille sig ud.}
  \item \textit{Skribenten blev overrasket over, at chefen talte meget roligt og lyttede lige så meget til praktikanterne som til de erfarne udviklere. Det skyldtes blandt andet, at Janteloven stadig præger samarbejdet på kontoret.}
\end{enumerate}

\textbf{Tekst B 段落逻辑补全答案（Modul 3 Læsning Opgave 3 风格）}
\begin{enumerate}
  \item[] \da{B}（前文说“我原以为必须大声推销自己的想法”，下文突然转折描写丹麦老板很平静、倾听所有人意见，因此必须用转折句 \textit{Men sådan fungerer et dansk kontor altså ikke altid.} 承上启下；\textit{A} 语义完全违背事实，\textit{C} 和 \textit{D} 与上下文逻辑脱节）
\end{enumerate}

\textbf{练习 A 答案}
\begin{enumerate}
  \item \da{må}（每周工作 60 小时，逻辑必然“一定很累”）
  \item \da{kan ikke}（十分钟前才打过电话，“不可能”忘记）
  \item \da{skulle / skal}（\textit{skulle efter sigende} 据说/传闻，用过去时 \textit{skulle} 更显客观存疑）
  \item \da{kan}（表示客观可能性“可能有正反两面”）
  \item \da{bør}（提出合理规范建议“应当”）
  \item \da{må}（不接电话，高度推测“一定是已经回家了”）
  \item \da{skulle / skal}（与 \textit{efter sigende} 搭配表示传闻）
  \item \da{bør / må}（劝告/提醒“不应该/不许”）
\end{enumerate}

\textbf{练习 B 答案}
\begin{enumerate}
  \item \textit{Det \da{må være} svært at starte i et nyt land.}
  \item \textit{Kulturen i København \da{skulle efter sigende være} mere international i dag.}
  \item \textit{Det \da{kan ikke passe}, at Janteloven er forsvundet fuldstændigt.}
  \item \textit{Hun \da{må have glemt} sin paraply på kontoret.}
  \item \textit{For meget konkurrence \da{kan skabe} stress på arbejdspladsen.}
  \item \textit{Det \dabold{kan ikke være} rigtigt, at han tjener så mange penge i et studiejob!}（否定推测禁区：必须把 \textit{må ikke} 改为 \textit{kan ikke}）
\end{enumerate}

\textbf{练习 C 答案}
\begin{enumerate}
  \item \textit{Når man er ny, \da{må det være} svært at vænne sig til det, fordi alle er så beskedne.}
  \item \textit{Dette universitet \da{skulle efter sigende være} meget åbent for innovation.}
  \item \textit{Han \da{kan ikke} allerede være taget hjem, for hans taske står stadig på bordet.}
  \item \textit{En lille smule beskedenhed \da{kan skabe} et mere ligeværdigt samarbejde.}
  \item \textit{De \da{må have haft} en rigtig god samtale på kontoret i går.}
\end{enumerate}

\textbf{Modul 3 实战参考范文（123 词）}\\
\textit{Til: Daniel \quad|\quad Emne: Janteloven og dit nye job i Danmark}

\textit{Hej Daniel}

\textit{Jeg forstår godt din bekymring! Når man læser om Janteloven online, \da{må det næsten lyde}, som om man slet ikke må være ambitiøs i Danmark. Men det \dabold{kan ikke passe} i praksis på en moderne arbejdsplads.}

\textit{Janteloven \dabold{skulle efter sigende være} blevet meget svagere i dag, især i internationale firmaer. Her \da{kan} det tværtimod \da{være} et problem, hvis man er for stille omkring sine idéer. Det vigtige er bare måden, man gør det på: Du \da{bør} fokusere på, hvordan din viden kan hjælpe hele teamet i stedet for kun at prale af dig selv. Danske chefer sætter nemlig stor pris på ærlighed og godt samarbejde.}

\textit{Held og lykke med ansøgningen --- det skal nok gå rigtig godt!}

\textit{Mange hilsner\\
\mbox{[Dit navn]}}

\small\textbf{点评}：全篇正文 123 词，全面回应丹尼尔的顾虑并给出实用建议；恰当运用 \textit{må næsten lyde}、\textit{kan ikke passe}、\textit{skulle efter sigende være}、\textit{kan være} 与 \textit{bør} 五种情态用法，语气自然克制、逻辑清晰。
\end{answerbox}

\ncdlesson{9}{Fra vikingetid til velfærdsstat}{从维京时代到福利国家}
{DU3 Modul 3（B1）｜对标 NCE 2 后半 \& Modul 3 \textit{Læsning Opgave 4 (Tre perspektiver)} \& \textit{Skrivning Opgave 1}}
{\textbf{历史叙事中的时态规范与长跨度时间衔接词}：一般过去时（\textit{Datid}）锚定已结束历史阶段 · 过去完成时（\textit{Førdatid}）追溯先发生背景 · 现在完成时（\textit{Førnutid}）承接至今 · 长跨度时间短语（\textit{i vikingetiden, op gennem 1900-tallet, i efterkrigstiden, siden da}）}
{讲述丹麦历史或时代变迁时，时态一旦乱跳就容易失分。\textbf{只要时代已彻底结束（如 \textit{i 1849, i 1900-tallet, i efterkrigstiden}），无论跨度多长都用一般过去时}；只有连接“从那时起一直延续到今天”的结果（配 \textit{siden da, helt frem til i dag}），才切回现在完成时。此外切记：\textbf{\textit{1900-tallet} 指的是 20 世纪（1900--1999 年），不是 19 世纪！}}

% ---------------------------------------------------------------------
\begin{foerduleaeser}
带着下面两个问题阅读 Tekst A 和 Tekst B：
\begin{enumerate}
  \item \textit{Hvornår fik Danmark sin første grundlov, og hvad byggede politikerne op i efterkrigstiden?}（丹麦何时颁布了第一部宪法？政治家们在战后时期建设了什么？）
  \item \textit{Hvilke tre historiske steder i Danmark har Jonas, Mette og Lukas besøgt ifølge Tekst B?}（根据 Tekst B，约纳斯、梅特和卢卡斯分别参观了丹麦哪三个历史场所？）
\end{enumerate}
\end{foerduleaeser}

% ---------------------------------------------------------------------
\ncdsubsec{二、Tekst A \& Tekst B（文化深读与 Modul 3 真题场景）}

\subsection*{Tekst A：\textit{Fra vikingetid til velfærdsstat}（历史文化短文 · 约 255 词）}

\begin{tipbox}{语篇说明}
本篇用一条清晰的时间轴串联丹麦千年演变。阅读时请特别关注\textbf{\color{NCDNavy}深蓝色标记的长跨度时间衔接短语}（注意它们放在句首时引发的 V2 倒装），以及\textbf{\color{NCDRed}深红色标记的三种时态分工}（一般过去时 vs 过去完成时 vs 现在完成时）。
\end{tipbox}

Hvordan udviklede Danmark sig fra et krigerisk vikingesamfund til en af verdens mest kendte velfærdsstater? Svaret strækker sig over mere end tusind års historie.

\da{I vikingetiden}, fra omkring år 800 til 1050, \dabold{sejlede} danske vikinger ud på havene for at handle og erobre nye områder. Kong Harald Blåtand samlede riget og indførte kristendommen omkring år 965. \da{Gennem middelalderen} og helt frem til 1800-tallet \dabold{havde} kongen meget stor magt. Men efter at borgerne længe \dabold{havde krævet} politiske reformer, underskrev kongen Danmarks første grundlov i 1849. Dermed sluttede enevælden, og folkestyret begyndte. Samtidig ændrede landbruget sig \da{i slutningen af 1800-tallet}, da bønderne dannede andelsmejerier og folkehøjskoler, hvilket styrkede samarbejdet i landsbyerne.

\da{Op gennem 1900-tallet} \dabold{ændrede} Danmark sig hurtigt fra et fattigt landbrugssamfund til et moderne industriland. \da{I 1930'erne} oplevede mange familier stor fattigdom og arbejdsløshed under den økonomiske krise. Netop fordi befolkningen \dabold{havde mærket} den hårde tid på egen krop, ønskede de fleste et stærkere socialt sikkerhedsnet. \da{I efterkrigstiden}, især fra 1950'erne og frem, \dabold{byggede} politikerne den danske velfærdsstat op med gratis uddannelse, sygehuse og støtte til ældre. Kvinderne kom også i stort tal ud på arbejdsmarkedet \da{op gennem 1960'erne og 1970'erne}, fordi der var brug for mere arbejdskraft, og kommunerne oprettede daginstitutioner til børnene.

\da{Siden da har} velfærdsmodellen \dabold{været} en grundpille i dansk politik, og mange historiske idéer om fællesskab \dabold{har overlevet} helt frem til i dag. Selvom samfundet nu står over for en voksende ældregruppe, diskuterer danskerne fortsat, hvordan man bedst bevarer den model, som tidligere generationer skabte.

\begin{transbox}
\small
\textbf{\color{NCDTeal}[Tekst A 中文精译]}\\
丹麦是如何从一个好战的维京社会发展成世界上最著名的福利国家之一的？答案跨越了一千多年的历史。

在维京时代，大约公元 800 年至 1050 年间，丹麦维京人扬帆出海，开展贸易并征服新的土地。国王蓝牙哈拉尔德（Harald Blåtand）统一了王国，并于公元 965 年左右引入了基督教。贯穿整个中世纪一直到 19 世纪，国王都掌握着极大的权力。但在市民们长期要求政治改革之后，国王于 1849 年签署了丹麦第一部宪法。君主专制就此终结，民主制度由此开启。与此同时，19 世纪末农业也发生了变革，农民们创办了合作乳品厂和民众高等学院，增强了乡村之间的协作。

整个 20 世纪期间，丹麦迅速从一个贫困的农业社会转变为现代工业国家。20 世纪 30 年代，许多家庭在经济危机中经历了严重的贫困与失业。正是因为民众亲身感受过那段艰难岁月，大多数人都渴望建立一张更强大的社会安全网。在二战后的时期，尤其是从 20 世纪 50 年代起，政治家们建起了丹麦福利国家，提供免费教育、公立医院和养老支持。整个 60 至 70 年代期间，女性也大规模进入劳动力市场，因为当时需要更多劳动力，市政厅也随之建起了托儿所等日间照料机构。

从那以后，福利模式一直是丹麦政治的基石，许多关于共同体的历史理念一直延续至今。尽管当今社会正面临老龄人口不断增加的新挑战，丹麦人依然在持续探讨如何最好地守护前几代人创造出的这一模式。
\end{transbox}

\subsection*{Tekst B：\textit{Tre danskere om historieoplevelser}（对标 Modul 3 Læsning Opgave 4）}

\textbf{阅读微任务（Læsning Opgave 4 三人物信息匹配）：}\textit{Læs teksterne om Jonas (A), Mette (B) og Lukas (C), og match spørgsmålene (1--4) med den rigtige person.}

\begin{emailbox}{Tre stemmer om dansk historie（A / B / C）}
\textbf{A. Jonas:} I sommer besøgte jeg Vikingeskibsmuseet i Roskilde. Jeg blev imponeret over de fem originale skibe fra 1000-tallet, som arkæologer \dabold{havde fundet} på bunden af fjorden i 1962. Man forstår virkelig, hvor dygtige håndværkere vikingerne \da{var}.\\
\textbf{B. Mette:} Min yndlingsudstilling er Arbejdermuseet i København. Her ser man, hvordan almindelige familier \da{levede} op gennem 1900-tallet, og hvordan kvinder og mænd kæmpede for bedre arbejdsvilkår i efterkrigstiden. Museet \dabold{har bevaret} stemningen helt frem til i dag.\\
\textbf{C. Lukas:} Sidste weekend tog jeg til Den Gamle By i Aarhus. Der kan man gå gennem gader fra både 1800-tallet og 1970'erne, og skuespillere fortæller historier fra hverdagen, før man \dabold{havde fået} moderne teknologi.
\end{emailbox}

{\small
\renewcommand{\arraystretch}{1.24}
\begin{xltabular}{\linewidth}{@{} X c c c @{}}
\toprule
\textbf{Hvem...（哪位人物符合以下描述？请选择 \textit{A}、\textit{B} 或 \textit{C}）} & \textbf{Jonas (\textit{A})} & \textbf{Mette (\textit{B})} & \textbf{Lukas (\textit{C})} \\
\midrule
\endfirsthead
\toprule
\textbf{Hvem...} & \textbf{Jonas (\textit{A})} & \textbf{Mette (\textit{B})} & \textbf{Lukas (\textit{C})} \\
\midrule
\endhead
\bottomrule
\endfoot
1. Hvem så noget, som var blevet fundet under vandet i 1960'erne? & A & B & C \\
2. Hvem oplevede gader fra to forskellige århundreder samme sted? & A & B & C \\
3. Hvem fokuserede mest på arbejdernes kamp og hverdagsliv i 1900-tallet? & A & B & C \\
4. Hvem mødte udklædte formidlere, der fortalte om livet før moderne teknik? & A & B & C \\
\bottomrule
\end{xltabular}
}

\begin{tipbox}{本课发音与吞音要点（Udtale \& Reduktioner）}
\begin{itemize}
  \item \textbf{历史复合词里的静音 \textit{d} 与 \textit{g}}：\textit{velfærdsstat} \ipa{ˈvɛlfɛɐ̯sˌsdæːˀd}（\textbf{\textit{-færds-} 中的 \textit{d} 完全不发音}！）；\textit{grundlov} \ipa{ˈɡʁɔnˌlʌʊ̯ˀ}（\textit{d} 不发音，\textit{lov} 读 \ipa{lʌʊ̯ˀ}）；\textit{efterkrigstiden} \ipa{ˈɛfdɐˌkʰʁiːˀsˌtˢiðˀən}（\textit{krig} 的 \textit{g} 脱落）；\textit{middelalderen} \ipa{ˈmiðəlˌalˀɐn}（前软 \textit{d}，后 \textit{-ld-} 的 \textit{d} 不发音）。
  \item \textbf{维京与年代数字读法}：\textit{1849} 年读 \textit{atten hundrede niogfyrre}；\textit{vikingetiden} \ipa{ˈviːɡeŋəˌtˢiðˀən}；\textit{1900-tallet}（\textit{nittenhundrede-tallet}）读作 \ipa{ˈnedənˌhunʁəðəˌtˢalˀəð}。
  \item \textbf{介词短语连读}：\textit{op gennem} \ipa{ʌb ˈɡɛnˀəm}（注意 \textit{gennem} 的 \textit{g} 是硬音 \ipa{ɡ}，绝不读成 \textit{j}！）；\textit{siden da} \ipa{ˈsiːðən ˈda}（口语常弱化为 \ipa{ˈsiːn ˈda}）。
\end{itemize}
\end{tipbox}

\begin{transbox}
\small
\textbf{\color{NCDTeal}[Tekst B 中文对照]}\\
\textbf{A. 约纳斯}：今年夏天我参观了罗斯基勒的维京船博物馆。那五艘来自 11 世纪的原装古船让我惊叹不已，考古学家们早在 1962 年就在峡湾底部发现了它们。在那里你真切地体会到维京人曾是多么高超的工匠。\\
\textbf{B. 梅特}：我最喜欢的展览是哥本哈根的工人博物馆。在这里可以看到普通家庭在整个 20 世纪是如何生活的，以及男女劳动者在战后时期如何争取更好的工作条件。这座博物馆把当年的氛围一直保存到了今天。\\
\textbf{C. 卢卡斯}：上周末我去了奥胡斯的“老城”（Den Gamle By）。在那里你可以穿梭于 19 世纪和 20 世纪 70 年代的街道之间，演员们还会讲述在拥有现代科技之前人们日常生活的故事。
\end{transbox}

% ---------------------------------------------------------------------
\ncdsubsec{三、Gloser（核心生词表）}

\begin{tipbox}{词表重点}
本课词汇聚焦丹麦历史与制度演进。\textit{at kræve}（要求）与 \textit{at indføre}（引入，推行）是撰写社会背景的黄金动词；注意 \textit{en grundlov} 复数变音或加 \textit{-e}（\textit{grundlove}），而 \textit{vikingetiden, middelalderen, efterkrigstiden, enevælden, arbejdsløsheden} 等抽象时代名词通常只用单数（\textit{ingen flertal}）。
\end{tipbox}

\begin{glosertable}
动词 [har] & at udvikle (sig) & \textit{(udvikler, udviklede, har udviklet)} & 发展，演变 \\
动词 [har] & at erobre & \textit{(erobrer, erobrede, har erobret)} & 征服，占领 \\
动词 [har] & at samle & \textit{(samler, samlede, har samlet)} & 统一，汇集 \\
动词 [har] & at indføre & \textit{(indfører, indførte, har indført)} & 引入，推行（制度/信仰） \\
动词 [har] & at kræve & \textit{(kræver, krævede, har krævet)} & 要求，争取 \\
动词 [har] & at opbygge & \textit{(opbygger, opbyggede, har opbygget)} & 建设，构建（制度/国家） \\
动词 [har] & at overleve & \textit{(overlever, overlevede, har overlevet)} & 幸存，延续至今 \\
动词 [har] & at bevare & \textit{(bevarer, bevarede, har bevaret)} & 保留，保护，维持 \\
\midrule
名词 & en vikingetid & \textit{(vikingetiden, ingen flertal)} & 维京时代 \\
名词 & en middelalder & \textit{(middelalderen, ingen flertal)} & 中世纪 \\
名词 & en grundlov & \textit{(grundloven, grundlove, grundlovene)} & 宪法 \\
名词 & en enevælde & \textit{(enevælden, ingen flertal)} & 君主专制 \\
名词 & en velfærdsstat & \textit{(velfærdsstaten, velfærdsstater, velfærdsstaterne)} & 福利国家 \\
名词 & et landbrugssamfund & \textit{(landbrugssamfundet, landbrugssamfund, landbrugssamfundene)} & 农业社会 \\
名词 & en efterkrigstid & \textit{(efterkrigstiden, ingen flertal)} & 战后时期（二战后） \\
名词 & en arbejdsløshed & \textit{(arbejdsløsheden, ingen flertal)} & 失业，失业率 \\
名词 & en grundpille & \textit{(grundpillen, grundpiller, grundpillerne)} & 基石，核心支柱 \\
\end{glosertable}

% ---------------------------------------------------------------------
\ncdsubsec{四、Grammatik \& Faste vendinger（核心语法与地道表达精讲）}

\begin{grammatikbox}{语法点 1：历史叙事的三重时态锚定法则}
无论写国家历史、公司沿革还是个人成长经历，只要按时间推进，丹麦语严格遵循三层时态分工：
\begin{enumerate}
  \item \textbf{锚定历史时代与过程 $\rightarrow$ 一般过去时（\da{Datid}）}：只要时间段\textbf{已经彻底结束}（哪怕持续了几百年！），主干叙事一律用一般过去时！常见标志词：\textit{i vikingetiden, gennem middelalderen, i 1849, op gennem 1900-tallet, i efterkrigstiden}。例：\textit{Op gennem 1900-tallet \da{blev} Danmark et industriland.}（不可用完成时 *har blevet！）
  \item \textbf{回溯“过去的过去”原因 $\rightarrow$ 过去完成时（\dabold{Førdatid}: \textit{havde/var + 分词}）}：当在讲述某个过去事件 \textit{A} 时，需要补充说明在此\textbf{之前就已经发生}的铺垫 \textit{B}，必须使用 \textit{havde/var} + 过去分词（常由 \textit{efter at, før, fordi} 引导）。例：\textit{I 1849 fik man en grundlov, efter at folket længe \dabold{havde krævet} reformer.}
  \item \textbf{延续至今的遗产或现状 $\rightarrow$ 现在完成时（\da{Førnutid}: \textit{har/er + 分词}）}：只有当历史成果“一直持续或影响到今天”（标志词：\textit{siden da, helt frem til i dag, gennem de sidste 70 år}），才切换回现在完成时！例：\textit{Siden da \da{har} velfærdsmodellen \da{spillet} en stor rolle.}
\end{enumerate}
\end{grammatikbox}

{\small
\renewcommand{\arraystretch}{1.25}
\begin{xltabular}{\linewidth}{@{} >{\bfseries\color{NCDNavy}\raggedright\arraybackslash}p{3.2cm} >{\raggedright\arraybackslash}p{3.1cm} >{\raggedright\arraybackslash}p{2.5cm} >{\raggedright\arraybackslash}X @{}}
\toprule
\textbf{长跨度时间衔接词} & \textbf{时间区间含义} & \textbf{配套时态} & \textbf{句首 V2 倒装例句} \\
\midrule
\endfirsthead
\toprule
\textbf{长跨度时间衔接词} & \textbf{时间区间含义} & \textbf{配套时态} & \textbf{句首 V2 倒装例句} \\
\midrule
\endhead
\bottomrule
\endfoot
\textit{I vikingetiden / I middelalderen} & 在维京时代 / 中世纪 & \textbf{一般过去时} & \textit{I vikingetiden \da{sejlede} vikingerne ud.} \\
\midrule
\textit{I 1800-tallet / I slutningen af 1800-tallet} & 在 19 世纪 / 19 世纪末 & \textbf{一般过去时} & \textit{I slutningen af 1800-tallet \da{ændrede} landbruget sig.} \\
\midrule
\textit{Op gennem 1900-tallet} & 贯穿整个 20 世纪期间 & \textbf{一般过去时} & \textit{Op gennem 1900-tallet \da{voksede} byerne hurtigt.} \\
\midrule
\textit{I efterkrigstiden} & 在二战后时期（1945 后） & \textbf{一般过去时} & \textit{I efterkrigstiden \da{byggede} man velfærdsstaten op.} \\
\midrule
\textit{Siden da / Helt frem til i dag} & 从那时起至今 & \textbf{现在完成时} & \textit{Siden da \dabold{har} modellen \dabold{været} meget stabil.} \\
\end{xltabular}
}

\begin{grammatikbox}{语法点 2：中英文对比排雷——年份“差一百年”与时态误用}
\begin{enumerate}
  \item \textbf{超级陷阱：\textit{1900-tallet} 究竟是几世纪？}：中文说“20 世纪”，英语说 \textit{the 20th century}，但丹麦语最常用的表达是 \dabold{\textit{1900-tallet}}（字面是“1900 年代段”，即 \textbf{1900--1999 年 = 20 世纪}）！同理，\dabold{\textit{1800-tallet}} 是 \textbf{19 世纪（1800--1899 年）}，而 \textit{1960'erne} 则专指“20 世纪 60 年代”。Modul 3 阅读理解极爱用“文本写 \textit{1800-tallet}，题干写 \textit{det 19. århundrede}”来考查同义替换，千万别算差一百年！
  \item \textbf{有明确结束跨度，绝不能用现在完成时！}：受“持续了一百年”的中文习惯影响，学生常写出 *\textit{Op gennem 1900-tallet \dabold{har} Danmark udviklet sig...}——错！因为 1900--1999 年已经结束，必须用一般过去时 \textit{Op gennem 1900-tallet \da{udviklede} Danmark sig...}！
  \item \textbf{历史时间短语放句首必触发 V2 倒装}：如 \textit{I efterkrigstiden \dabold{fik} alle borgere ret til gratis behandling.}（动词 \textit{fik} 紧跟时间短语之后排第 2 位！）。
\end{enumerate}
\end{grammatikbox}

% ---------------------------------------------------------------------
\ncdsubsec{五、Mønsterøvelser（句型操练）}

\textbf{练习 A：用括号内动词的正确历史时态填空（一般过去时 \textit{Datid}、过去完成时 \textit{Førdatid} 或现在完成时 \textit{Førnutid}）}
\begin{enumerate}
  \item I vikingetiden \underline{\hspace{2.2cm}} (\textit{sejle}) danske skibe helt til England og Frankrig.
  \item Efter at folket længe \underline{\hspace{2.6cm}} (\textit{kræve}) forandring, fik Danmark en grundlov i 1849.
  \item Op gennem 1900-tallet \underline{\hspace{2.4cm}} (\textit{blive}) Danmark gradvist et rigt samfund.
  \item Siden 1950'erne \underline{\hspace{1.4cm}} velfærdsstaten \underline{\hspace{2.2cm}} (\textit{spille}) en central rolle i danskernes liv.
  \item I efterkrigstiden \underline{\hspace{2.2cm}} (\textit{opbygge}) politikerne nye sygehuse og universiteter.
  \item Arkæologerne fandt i 1962 fem skibe, som vikingerne \underline{\hspace{2.6cm}} (\textit{sænke}, 沉没) i fjorden næsten tusind år tidligere.
  \item Helt frem til i dag \underline{\hspace{1.4cm}} mange gamle traditioner \underline{\hspace{2.2cm}} (\textit{overleve}) i landsbyerne.
  \item I slutningen af 1800-tallet \underline{\hspace{2.2cm}} (\textit{danne}) danske bønder mange fælles andelsmejerier.
\end{enumerate}

\textbf{练习 B：句型转换——时间状语提前倒装与跨世纪表达精准转换}
\begin{enumerate}
  \item 把时间短语提前到句首（注意 V2 倒装）：\textit{Danske vikinger erobrede store områder i vikingetiden.} $\rightarrow$ \textit{I vikingetiden...}
  \item 把时间短语提前到句首（注意现在完成时倒装）：\textit{Velfærdsstaten har skabt stor tryghed siden da.} $\rightarrow$ \textit{Siden da...}
  \item 用 \textit{efter at} + 过去完成时合并两句（先把较早发生的事改过去完成时）：\textit{Danskerne oplevede stor krise i 1930'erne. Derefter byggede de velfærdsstaten.} $\rightarrow$ \textit{Efter at danskerne...}
  \item 将汉语时间转换为地道丹麦语表达：\textit{在 20 世纪期间（贯穿 1900--1999 年）} $\rightarrow$ \underline{\hspace{3.0cm}}
  \item 将汉语时间转换为地道丹麦语表达：\textit{在 19 世纪末（1880--1899 年左右）} $\rightarrow$ \underline{\hspace{3.2cm}}
  \item 纠错改写：\textit{Op gennem 1900-tallet har flere kvinder kommet ud på arbejdsmarkedet.} $\rightarrow$
\end{enumerate}

\textbf{练习 C：中译丹（精准把控历史时态与衔接词）}
\begin{enumerate}
  \item 在维京时代，丹麦维京人建起了庞大的贸易网络（et netværk）。
  \item 在 1849 年宪法签署之后，君主专制就结束了。
  \item 贯穿整个 20 世纪，丹麦从一个农业国转变为现代社会。
  \item 在二战后时期，政府建设了免费学校和公立医院。
  \item 从那时起，平等的价值观一直延续到了今天。
\end{enumerate}

% ---------------------------------------------------------------------
\ncdsubsec{六、DU3 Modul 3 实战微任务与参考答案}

\begin{emailbox}{Skrivning Opgave 1 文化体验推荐信（Anbefaling · 要求 100 词以上）}
\textbf{Situation}: Du har besøgt et historisk museum i Danmark (f.eks. \textit{Den Gamle By} i Aarhus, \textit{Nationalmuseet} eller \textit{Arbejdermuseet} i København) og synes, det var en rigtig god oplevelse.\\
\textbf{Opgave}: Skriv en kort anbefaling til andre kursister på din sprogskole. Du skal skrive om:
\begin{itemize}
  \item hvilket museum du har besøgt, og hvor det ligger;
  \item hvilke historiske perioder man kan lære om der;
  \item hvad du selv syntes var mest spændende ved besøget;
  \item hvem du især vil anbefale museet til, og hvorfor.
\end{itemize}
Skriv \textbf{mindst 100 ord}.
\end{emailbox}

\begin{answerbox}
\textbf{Før du læser 参考回答}
\begin{enumerate}
  \item \textit{Danmark fik sin første grundlov i 1849. I efterkrigstiden byggede politikerne den danske velfærdsstat op med gratis uddannelse, sygehuse og støtte til ældre.}
  \item \textit{Jonas har besøgt Vikingeskibsmuseet i Roskilde, Mette har besøgt Arbejdermuseet i København, og Lukas har besøgt Den Gamle By i Aarhus.}
\end{enumerate}

\textbf{Tekst B 三人物匹配答案（Modul 3 Læsning Opgave 4 风格）}
\begin{enumerate}
  \item \da{A}（Jonas：\textit{på bunden af fjorden i 1962} 同义对应“1960 年代在水下发现”）
  \item \da{C}（Lukas：\textit{gader fra både 1800-tallet og 1970'erne} 同义对应“两个不同世纪的街道”）
  \item \da{B}（Mette：提到工人博物馆呈现 20 世纪普通家庭和战后劳动者争取更好工作条件）
  \item \da{C}（Lukas：提到演员们 \textit{skuespillere} 讲述现代科技出现前的日常生活）
\end{enumerate}

\textbf{练习 A 答案}
\begin{enumerate}
  \item \da{sejlede}（维京时代已彻底结束，一般过去时）
  \item \da{havde krævet}（在 1849 年签署宪法之前，过去完成时）
  \item \da{blev}（\textit{Op gennem 1900-tallet} 已结束，一般过去时）
  \item \da{har ... spillet}（\textit{Siden 1950'erne} 延续至今，现在完成时）
  \item \da{opbyggede}（战后时期已结束，一般过去时）
  \item \da{havde sænket}（早于 1962 年发掘时间近千年，过去完成时）
  \item \da{har ... overlevet}（\textit{Helt frem til i dag} 延续至今，现在完成时）
  \item \da{dannede}（19 世纪末已结束，一般过去时）
\end{enumerate}

\textbf{练习 B 答案}
\begin{enumerate}
  \item \textit{I vikingetiden \dabold{erobrede} danske vikinger store områder.}（V2 倒装）
  \item \textit{Siden da \dabold{har} velfærdsstaten \da{skabt} stor tryghed.}（助动词 \textit{har} 倒装到第 2 位）
  \item \textit{Efter at danskerne \dabold{havde oplevet} stor krise i 1930'erne, \dabold{byggede} de velfærdsstaten.}（从句过去完成时 + 前置触发主句倒装）
  \item \textit{\da{op gennem 1900-tallet} / i 1900-tallet}（记住是 1900 不是 2000！）
  \item \textit{\da{i slutningen af 1800-tallet}}
  \item \textit{Op gennem 1900-tallet \dabold{kom} flere kvinder ud på arbejdsmarkedet.}（1900 年代已结束，不能用完成时，必须用一般过去时 \textit{kom}）
\end{enumerate}

\textbf{练习 C 答案}
\begin{enumerate}
  \item \textit{I vikingetiden byggede danske vikinger et stort handelsnetværk op.}
  \item \textit{Efter at grundloven \dabold{var blevet} underskrevet i 1849, sluttede enevælden.}
  \item \textit{Op gennem 1900-tallet ændrede Danmark sig fra et landbrugsland til et moderne samfund.}
  \item \textit{I efterkrigstiden byggede regeringen gratis skoler og offentlige sygehuse.}
  \item \textit{Siden da \da{har} værdierne om lighed \da{overlevet} helt frem til i dag.}
\end{enumerate}

\textbf{Modul 3 实战参考范文（119 词）}\\
\textit{Anbefaling: Besøg Arbejdermuseet i hjertet af København!}

\textit{Hvis du gerne vil forstå det moderne Danmark bedre, kan jeg varmt anbefale Arbejdermuseet, som ligger i Rømersgade tæt på Nørreport Station i København.}

\textit{Udstillingerne viser, hvordan almindelige familier levede og arbejdede fra slutningen af 1800-tallet og op gennem 1900-tallet. Man lærer især meget om, hvordan arbejderbevægelsen kæmpede for højere løn og bedre vilkår, efter at mange \dabold{havde oplevet} stor fattigdom i 1930'erne. I efterkrigstiden skabte denne udvikling grundlaget for den velfærdsstat, vi kender i dag.}

\textit{Det mest spændende ved mit besøg var de gamle lejligheder fra 1950'erne og den autentiske kaffebar i kælderen. Jeg vil især anbefale museet til nye tilflyttere og kursister på Modul 3, fordi man både får et stærkt historisk indblik og øver dansk kulturforståelse.}

\small\textbf{点评}：全篇正文 119 词，紧扣推荐信四大要点；精准调动 \textit{fra slutningen af 1800-tallet, op gennem 1900-tallet, havde oplevet, i efterkrigstiden} 等时间坐标与历史时态，完美复现 Modul 3 Opgave 1 高分推荐文体。
\end{answerbox}

\ncdlesson{10}{Tillid: Når barnevognen står ude på gaden}{信任：当婴儿车停在街边}
{DU3 Modul 3（B1）｜对标 NCE 2 后半 \& Modul 3 \textit{Kulturkontrast, Skrivning Opgave 2 \& Mundtlig Opgave 2}}
{\textbf{高级让步与对比结构}：让步从句 \textit{til trods for, at...}（尽管……）vs 让步介词 \textit{på trods af / til trods for} + 名词短语 · 对比转折 \textit{mens ... derimod}（而相比之下）与 \textit{i modsætning til} + 名词}
{在 Modul 3 文化对比写作与口试中，只用简单的 \textit{men} 和 \textit{selvom} 很难拿高分。必须严格区分：\textbf{\textit{på trods af} 后只能跟名词短语，\textit{til trods for, at} 后才能跟完整句子}；做两地习惯对比时，使用 \textbf{\textit{mens} + 主语 + \textit{derimod} + 动词} 可以把反差写得精准而有说服力！}

% ---------------------------------------------------------------------
\begin{foerduleaeser}
带着下面两个问题阅读 Tekst A 和 Tekst B：
\begin{enumerate}
  \item \textit{Hvorfor lader danske forældre ofte deres børn sove i barnevognen udenfor, og hvor i samfundet ser man ellers den høje tillid?}（为什么丹麦父母经常让孩子睡在户外的婴儿车里？在社会其他哪些地方还能看到这种高度信任？）
  \item \textit{Hvad lagde skribenten i Tekst B mærke til på sin første arbejdsdag og i landsbyen i forhold til sit hjemland?}（在 Tekst B 中，作者在第一天上班以及在丹麦乡村时注意到了哪些与自己祖国不同的地方？）
\end{enumerate}
\end{foerduleaeser}

% ---------------------------------------------------------------------
\ncdsubsec{二、Tekst A \& Tekst B（文化深读与 Modul 3 真题场景）}

\subsection*{Tekst A：\textit{Tillid: Når barnevognen står ude på gaden}（文化深读短文 · 约 245 词）}

\begin{tipbox}{语篇说明}
在丹麦街头看到婴儿车独自停在咖啡馆门外，几乎是每一位新移民印象最深的文化冲击之一。阅读时请特别关注\textbf{\color{NCDNavy}深蓝色标记的高级让步结构}（\textit{på trods af} + 名词 vs \textit{til trods for, at} + 从句）与\textbf{\color{NCDRed}深红色标记的对比结构}（\textit{mens ... derimod}）及从句中 \textit{ikke} 的位置。
\end{tipbox}

En helt almindelig formiddag i Aarhus eller København kan turister pludselig få øje på et overraskende syn: Uden for en travl café står tre eller fire barnevogne på fortovet, mens forældrene sidder inde i varmen og drikker kaffe. Små babyer sover roligt i den friske luft med en babyalarm i vognen. \da{Til trods for, at} dette billede kan virke chokerende på udlændinge, er det en helt normal del af dansk hverdag.

\da{På trods af} det omskiftelige vejr mener danske forældre og sundhedsplejersker, at børn sover bedre og bliver sundere udendørs. Og \da{på trods af} de mange fremmede mennesker, der går forbi på gaden, er forældrene som regel rolige. Danmark har nemlig en af verdens højeste grader af social tillid. I mange andre lande holder man altid sit barn helt tæt ved sig i det offentlige rum, \dabold{mens} danskerne \dabold{derimod} opfatter tillid som en fælles tryghed, man passer på sammen.

Tilliden ses ikke kun ved caféerne, men overalt i samfundet. Langs de danske landeveje står der ofte ubemandede vejboder med friske jordbær, kartofler eller honning. Kunderne tager selv varerne og lægger penge i en åben kasse eller betaler via MobilePay. I andre kulturer sikrer man varer med kameraer og låse, \dabold{mens} man \dabold{derimod} stoler på, at folk på landet betaler ærligt. \da{Til trods for, at} tilliden selvfølgelig \dabold{ikke er} blind, og man stadig skal bruge sin sunde fornuft, sparer denne høje tillid hele samfundet for meget kontrol og bekymring.

\begin{transbox}
\small
\textbf{\color{NCDTeal}[Tekst A 中文精译]}\\
在奥胡斯或哥本哈根一个再普通不过的上午，游客们可能会突然看到一幕令人吃惊的场景：一家繁忙的咖啡馆门外人行道上停着三四辆婴儿车，而父母们正坐在暖和的室内喝咖啡。小宝宝们车里放着婴儿监护器，在新鲜空气中安然熟睡。尽管这一幕在外国人看来颇具冲击力，但在丹麦却是日常生活的常态。

尽管天气变化无常，丹麦父母和儿童保健护士仍认为孩子在户外睡得更香、身体更健康。而且尽管大街上走过许多陌生人，父母们通常也十分坦然。原因在于丹麦拥有世界上最高水平的社会信任度之一。在许多其他国家，人们在公共场所总是把孩子紧紧带在身边；而相比之下，丹麦人则把信任视为一种大家共同守护的公共安全感。

这种信任不仅体现在咖啡馆外，还遍布社会的方方面面。在丹麦的乡间公路旁，经常能看到摆放着新鲜草莓、土豆或蜂蜜的无人售卖摊位。顾客自己拿取商品，然后把钱放进敞开的钱箱或用 MobilePay 转账。在其他文化里，人们会用监控摄像头和锁具来保障货物安全；而相比之下，丹麦乡间则相信大家会诚实付款。尽管这种信任当然不是盲目的，人们也依然要保持基本的常识，但高度的社会信任为整个社会省去了大量的监控成本与焦虑。
\end{transbox}

\subsection*{Tekst B：\textit{Kulturkontrast: Tillid i hverdagen}（对标 Modul 3 文化对比短评）}

\begin{emailbox}{新移民随笔｜Fra min første måned på landet i Danmark}
\da{I modsætning til} mit travle hjemland, hvor man næsten altid kontrollerer alting grundigt, oplevede jeg hurtigt den danske tillid på min nye arbejdsplads. Allerede den første dag fik jeg nøglen til kontoret, \da{på trods af} min korte tid i firmaet. \dabold{Derimod} skal man selv huske at passe på sin cykel ved stationen, fordi cykler desværre ofte bliver stjålet. I weekenden købte jeg friske æbler i en lille vejbod uden personale. \da{Til trods for, at} jeg \dabold{ikke mødte} sælgeren, virkede hele oplevelsen utrolig varm. I min hjemby sætter folk store hegn om deres haver, \dabold{mens} man her \dabold{derimod} ofte nøjes med en lav hæk.
\end{emailbox}

\begin{tipbox}{本课发音与吞音要点（Udtale \& Reduktioner）}
\begin{itemize}
  \item \textbf{本课核心词“静音 \textit{d} 与 \textit{g}”大集合}：\textit{tillid} \ipa{ˈtˢeliðˀ}（首音节元音偏 \ipa{e}，末尾软 \textit{d} + 喉塞音）；\textit{trods} \ipa{ˈtˢʁʌs}（\textbf{\textit{d} 完全不发音}！）；\textit{på trods af} \ipa{pʰʌ ˈtˢʁʌs ˈæːˀ}（口语常弱化为 \ipa{pʰʌ ˈtˢʁʌs ɑ}）；\textit{barnevogn} \ipa{ˈbɑːnəˌvʌʊ̯ˀn}（\textit{g} 不发音）；\textit{offentlig} \ipa{ˈʌfənli}（\textit{t} 与 \textit{g} 均不发音）。
  \item \textbf{对比副词与环境词读音}：\textit{derimod} \ipa{ˈdɛɐ̯iˌmoːˀð}（重音常在 \textit{der-} 或 \textit{-mod}，末尾软 \textit{d}）；\textit{i modsætning til} \ipa{i ˈmoðˌsɛdneŋ tˢe}；\textit{fremmede} \ipa{ˈfʁɑməðə}（注意第一音节是 \ipa{ɑ} 而不是 \ipa{ɛ}！）。
  \item \textbf{街头与生活高频词}：\textit{fortov} \ipa{ˈfɒːˌtˢʌʊ̯ˀ}（\textit{-ov} 读 \ipa{ʌʊ̯ˀ}）；\textit{vejbod} \ipa{ˈvɑi̯ˌboːˀð}（软 \textit{d}）；\textit{ubemandet} \ipa{ˈubəˌmanˀəð}。
\end{itemize}
\end{tipbox}

\begin{transbox}
\small
\textbf{\color{NCDTeal}[Tekst B 中文对照]}\\
\textbf{我在丹麦乡村的第一个月}：与我那繁忙的祖国不同——在那里人们几乎总是把一切检查得严严实实——我在新工作单位很快就体会到了丹麦式的信任。尽管进公司没几天，我在上班第一天就拿到了办公室钥匙。不过反过来，在火车站大家还是得自己看好自行车，因为自行车遗憾地经常被偷。周末我在一个没有店员的小路边摊买了新鲜苹果。尽管我根本没见到卖家，整个购物体验却让人觉得无比温暖。在我的家乡，人们会给院子装上高高的围栏，而在这里相比之下，大家往往只用一排低矮的树篱就足够了。
\end{transbox}

% ---------------------------------------------------------------------
\ncdsubsec{三、Gloser（核心生词表）}

\begin{tipbox}{词表重点}
注意 \textit{en tillid}（信任）为不可数共性名词（\textit{stor tillid}, \textit{ingen flertal}），搭配固定动词 \textit{at stole på} 或 \textit{at have tillid til}（注意两个介词不同：\textit{stole \dabold{på}} vs \textit{tillid \dabold{til}}！）。\textit{chokerende} 与现在分词同形，词尾始终不变。
\end{tipbox}

\begin{glosertable}
动词 [har] & at stole (på) & \textit{(stoler, stolede, har stolet)} & 信任，信赖（固定接 \textit{på}） \\
动词 [har] & at få øje på & \textit{(får, fik, har fået øje på)} & 瞥见，注意到 \\
动词 [har] & at virke & \textit{(virker, virkede, har virket)} & 显得，看起来；起作用 \\
动词 [har] & at opfatte & \textit{(opfatter, opfattede, har opfattet)} & 看待，理解 \\
动词 [har] & at sikre & \textit{(sikrer, sikrede, har sikret)} & 保护，保障安全 \\
动词 [har] & at spare (for) & \textit{(sparer, sparede, har sparet)} & 省去，使免于（麻烦/成本） \\
\midrule
名词 & en tillid & \textit{(tilliden, ingen flertal)} & 信任，信赖 \\
名词 & et fortov & \textit{(fortovet, fortove, fortovene)} & 人行道 \\
名词 & en babyalarm & \textit{(babyalarmen, babyalarmer, babyalarmerne)} & 婴儿监护器 \\
名词 & en sundhedsplejerske & \textit{(sundhedsplejersken, sundhedsplejersker, -ne)} & 儿童保健护士 \\
名词 & en vejbod & \textit{(vejboden, vejboder, vejboderne)} & 路边无人售卖摊位 \\
名词 & en fornuft & \textit{(fornuften, ingen flertal)} & 理智，常识（\textit{sund fornuft}） \\
名词 & en modsætning & \textit{(modsætningen, modsætninger, modsætningerne)} & 对比，对立（\textit{i modsætning til}） \\
\midrule
形容词 & omskiftelig & \textit{(omskifteligt, omskiftelige)} & 多变的，变化无常的 \\
形容词 & ubemandet & \textit{(ubemandet, ubemandede)} & 无人值守的 \\
介词短语 & på trods af & \textit{(介词短语 + 名词短语)} & 尽管，不顾（后接名词） \\
从句连词 & til trods for, at & \textit{(从句连词 + 从句)} & 尽管（事实是）…… \\
\end{glosertable}

% ---------------------------------------------------------------------
\ncdsubsec{四、Grammatik \& Faste vendinger（核心语法与地道表达精讲）}

\begin{grammatikbox}{语法点 1：高级让步结构（\textit{på trods af} vs \textit{til trods for, at}）与对比结构（\textit{mens ... derimod}）}
在 B1 书面语和讨论题中，“虽然……但是……”和“……而相比之下……”是拉开表达层次的关键骨架：
\begin{enumerate}
  \item \textbf{让步介词：\da{på trods af}（或 \textit{til trods for}）+ \dabold{名词短语}}：等同于英语的 \textit{despite / in spite of}。后面\textbf{绝对只能接名词或名词短语}（如 \textit{det kolde vejr, de mange fremmede}）。放在句首时触发\textbf{主句 V2 倒装}！例：\textit{\da{På trods af} kulden \dabold{sover} barnet godt udenfor.}
  \item \textbf{让步从句连词：\da{til trods for, at} + \dabold{完整主谓从句}}：语气比 \textit{selvom} 更书面、更有说服力，内部必须严格遵守\textbf{从句语序（\textit{ikke/ofte} 放动词前）}！若整个从句前置，后方主句一样发生\textbf{V2 倒装}。例：\textit{\da{Til trods for, at} det \dabold{ikke er} varmt, \dabold{sover} børnene ude.}
  \item \textbf{对比双保险：\da{mens} + 主语 + \dabold{derimod} + 动词}：在 \textit{mens} 引导的对比从句中插入副词 \dabold{derimod}（“相反，相比之下”），用来凸显两类人或两国习惯的反差。注意在 \textit{mens} 从句里，\dabold{derimod} 属于句中副词，放在主语后面、变位动词前面！而如果不用 \textit{mens}，直接把 \dabold{Derimod} 放在新主句句首，则后面必须\textbf{V2 倒装}（\textit{Derimod \dabold{stoler} man på folk.}）。
\end{enumerate}
\end{grammatikbox}

{\small
\renewcommand{\arraystretch}{1.25}
\begin{xltabular}{\linewidth}{@{} >{\bfseries\color{NCDNavy}\raggedright\arraybackslash}p{3.0cm} >{\raggedright\arraybackslash}p{3.4cm} >{\raggedright\arraybackslash}p{3.1cm} >{\raggedright\arraybackslash}X @{}}
\toprule
\textbf{让步 / 对比结构} & \textbf{后续成分与语序} & \textbf{语体与功能} & \textbf{标准 B1 例句} \\
\midrule
\endfirsthead
\toprule
\textbf{让步 / 对比结构} & \textbf{后续成分与语序} & \textbf{语体与功能} & \textbf{标准 B1 例句} \\
\midrule
\endhead
\bottomrule
\endfoot
\textit{på trods af} + 名词 & 接名词短语；\newline 句首引发主句 V2 倒装 & 紧凑客观让步\newline （相当于 \textit{despite}） & \textit{\da{På trods af} mørket \dabold{cykler} mange stadig til arbejde.} \\
\midrule
\textit{til trods for, at} + 从句 & 接完整从句（\textit{ikke} 在动词前）；前置引发主句倒装 & 正式书面让步\newline （比 \textit{selvom} 更正式） & \textit{\da{Til trods for, at} man \dabold{ikke kender} ejeren, betaler man.} \\
\midrule
\textit{mens} + 主语 + \textit{derimod} + 动词 & 从句语序（\textit{derimod} 在动词前，不倒装） & 鲜明对照反差 & \textit{Nogle låser alt, \da{mens} danskerne \dabold{derimod} har åbne haver.} \\
\midrule
\textit{Derimod} + 动词 + 主语 & 独立主句副词，\newline 必须 V2 倒装 & 另起一句转折对比 & \textit{\dabold{Derimod} \da{skal} man altid låse sin cykel.} \\
\midrule
\textit{I modsætning til} + 名词 & 介词短语，\newline 句首引发主句 V2 倒装 & 句首直接点出对比对象 & \textit{\da{I modsætning til} mange storbyer \dabold{er} her meget trygt.} \\
\end{xltabular}
}

\begin{grammatikbox}{语法点 2：中英文对比排雷——三大高频失误点}
\begin{enumerate}
  \item \textbf{介词后面不能直接套英语 \textit{-ing} 动名词或光杆句子！}：英语写 \textit{Despite having cold weather...} 或中文直接连动词，但丹麦语\textbf{绝不能写 *\textit{på trods af det regner} 或 *\textit{på trods af regner}}！要么把后面变成名词（\textit{på trods af \dabold{regnen}}），要么把结构升级成连词（\textit{til trods for, \dabold{at} det regner}）。
  \item \textbf{别把 \textit{mens ... derimod} 写成从句倒装！}：\textit{mens} 引导的是标准从句，千万别模仿主句写成 *\textit{mens \dabold{ser man} derimod...}！正确语序永远是：\textit{mens man \dabold{derimod ser}...}。
  \item \textbf{介词区分：\textit{på trods af} vs \textit{til trods for} vs \textit{i modsætning til}}：记住每个短语的结尾介词不同——\textit{på trods \dabold{af}}、\textit{til trods \dabold{for} (, at)}、\textit{i modsætning \dabold{til}}，一个字母都别写串！
\end{enumerate}
\end{grammatikbox}

% ---------------------------------------------------------------------
\ncdsubsec{五、Mønsterøvelser（句型操练）}

\textbf{练习 A：用 \textit{på trods af, til trods for, mens, derimod} 或 \textit{i modsætning til} 填空}
\begin{enumerate}
  \item \underline{\hspace{2.4cm}} det kolde vejr sover mange danske babyer ude i barnevognen.
  \item \underline{\hspace{2.6cm}}, at vejboden er ubemandet, betaler næsten alle kunder ærligt.
  \item I min hjemby bruger man høje hegn, \underline{\hspace{1.6cm}} man i Danmark \underline{\hspace{1.8cm}} ofte nøjes med en lav hæk.
  \item \underline{\hspace{2.8cm}} mange andre storbyer har København en meget høj grad af tillid.
  \item Han cyklede hele vejen til arbejde \underline{\hspace{2.4cm}} den kraftige blæst.
  \item I caféen kan man slappe helt af. \underline{\hspace{2.0cm}} skal man passe på sin cykel ved stationen. (注意主句 V2 倒装位置)
  \item \underline{\hspace{2.6cm}}, at hun ikke kendte naboerne så godt, fik hun hurtigt hjælp til at flytte.
  \item Nogle lande satser på meget kontrol, \underline{\hspace{1.6cm}} det danske samfund \underline{\hspace{1.8cm}} bygger på tillid.
\end{enumerate}

\textbf{练习 B：句型转换——让步互换与对比句合并}
\begin{enumerate}
  \item 用 \textit{på trods af} + 名词替换 \textit{selvom} 从句：\textit{Selvom det regner meget, går de en tur med barnevognen.} (\textit{regnen}) $\rightarrow$ \textit{På trods af...}
  \item 用 \textit{til trods for, at} 改写下列 \textit{på trods af} 句子：\textit{På trods af den høje pris købte han friske jordbær.} (\textit{prisen er/var høj}) $\rightarrow$ \textit{Til trods for, at...}
  \item 用 \textit{mens ... derimod} 合并两句：\textit{I nogle kulturer låser man alt. I Danmark stoler man ofte på fremmede.} $\rightarrow$
  \item 用 \textit{I modsætning til} 开头改写（注意 V2 倒装）：\textit{I andre lande er vejboder sjældne, men her er de meget almindelige.} $\rightarrow$ \textit{I modsætning til andre lande...}
  \item 用 \textit{Til trods for, at} 合并两句（注意从句中 \textit{ikke} 的位置）：\textit{Man kender ikke sælgeren. Man betaler alligevel via MobilePay.} $\rightarrow$
  \item 纠错改写：\textit{På trods af at gaden er meget travl, så passer forældrene aldrig på barnevognen, mens i andre lande holder man derimod barnet tæt.}（找出 \textit{på trods af} 与 \textit{mens} 从句中的两处语法错误并改正）$\rightarrow$
\end{enumerate}

\textbf{练习 C：中译丹（综合运用高级让步与对比结构）}
\begin{enumerate}
  \item 尽管天气很冷，很多丹麦婴儿仍然在户外睡得很好。
  \item 尽管大家彼此不认识，但人们通常还是互相信任。
  \item 在一些国家人们很少和陌生人说话，而相比之下在丹麦社团里很容易交朋友。
  \item 与许多其他大城市不同，这里的人把婴儿车停在咖啡馆外面。
  \item 尽管存在语言挑战，我还是在当地俱乐部里感受到了温暖的氛围。
\end{enumerate}

% ---------------------------------------------------------------------
\ncdsubsec{六、DU3 Modul 3 实战微任务与参考答案}

\begin{emailbox}{Skrivning Opgave 2 跨文化讨论回信（Unit 2 收官任务 · 要求 100 词以上）}
\textbf{Situation}: Din veninde Astrid bor i udlandet og har lige set et foto på nettet af tre barnevogne foran en café i København. Hun er meget overrasket og skriver til dig.\\
\textbf{Opgave}: Skriv et svar til Astrid. Du skal svare på alle hendes spørgsmål og bruge mindst tre strukturer fra denne lektion (\textit{på trods af, til trods for at, mens... derimod, i modsætning til}). Skriv \textbf{mindst 100 ord}.

\vspace{0.3em}
\textit{Hej!\\
Jeg så lige et billede fra Danmark, hvor små børn sov i barnevogne ude på gaden, mens forældrene sad inde på en café! Hvorfor gør danske forældre det? Har du selv oplevet den danske tillid andre steder, og hvad er forskellen i forhold til vores eget hjemland?\\
Mange hilsner, Astrid}
\end{emailbox}

\begin{answerbox}
\textbf{Før du læser 参考回答}
\begin{enumerate}
  \item \textit{Danske forældre gør det, fordi de og sundhedsplejerskerne mener, at børn sover bedre i frisk luft, og fordi tilliden i samfundet er høj. Man ser også tilliden i ubemandede vejboder på landet.}
  \item \textit{Skribenten lagde mærke til, at man allerede den første dag fik nøglen til kontoret, og at der var ubemandede vejboder samt lave hække mellem haverne i stedet for høje hegn.}
\end{enumerate}

\textbf{练习 A 答案}
\begin{enumerate}
  \item \da{På trods af}（后接名词短语 \textit{det kolde vejr}）
  \item \da{Til trods for}（与后面的 \textit{, at} 组成从句连词）
  \item \da{mens}；\da{derimod}（从句中 \textit{derimod} 置于主语 \textit{man} 之后、动词之前）
  \item \da{I modsætning til}（后接名词短语表示“与……不同”）
  \item \da{på trods af}（后接名词短语 \textit{den kraftige blæst}）
  \item \da{Derimod}（位于新主句句首，引发后面的 \textit{skal man} 倒装）
  \item \da{Til trods for}（后接 \textit{, at} 从句）
  \item \da{mens}；\da{derimod}（对比从句，\textit{derimod} 在动词 \textit{bygger} 前面）
\end{enumerate}

\textbf{练习 B 答案}
\begin{enumerate}
  \item \textit{\da{På trods af} regnen \dabold{går de} en tur med barnevognen.}（介词短语放句首，主句 V2 倒装）
  \item \textit{\da{Til trods for, at} prisen var høj, \dabold{købte han} friske jordbær.}
  \item \textit{I nogle kulturer låser man alt, \da{mens} man i Danmark \dabold{derimod ofte} stoler på fremmede / \da{mens} danskerne \dabold{derimod ofte} stoler på fremmede.}
  \item \textit{\da{I modsætning til} andre lande \dabold{er} vejboder meget almindelige her.}（V2 倒装）
  \item \textit{\da{Til trods for, at} man \dabold{ikke kender} sælgeren, \dabold{betaler} man alligevel via MobilePay.}（注意从句 \textit{ikke} 在动词前）
  \item \textit{\da{Til trods for, at} gaden er meget travl, stoler forældrene på omgivelserne, \da{mens} man i andre lande \dabold{derimod} holder barnet tæt.}（修正两点：\ding{172} 加从句连词须用 \textit{Til trods for, at} 而非 \textit{På trods af}；\ding{173} \textit{mens} 从句不倒装，语序须改为 \textit{mens man ... derimod holder}）
\end{enumerate}

\textbf{练习 C 答案}
\begin{enumerate}
  \item \textit{\da{På trods af} det kolde vejr sover mange danske babyer stadig godt udenfor.}
  \item \textit{\da{Til trods for, at} man \dabold{ikke kender} hinanden, stoler folk som regel på hinanden.}
  \item \textit{I nogle lande taler man sjældent med fremmede, \da{mens} det \dabold{derimod} er nemt at få venner i danske foreninger.}
  \item \textit{\da{I modsætning til} mange andre storbyer stiller folk her barnevognen uden for caféen.}
  \item \textit{\da{På trods af} de sproglige udfordringer oplevede jeg en varm stemning i den lokale klub.}
\end{enumerate}

\textbf{Modul 3 实战参考范文（129 词）}\\
\textit{Til: Astrid \quad|\quad Emne: Sv: Barnevogne og tillid i Danmark}

\textit{Hej Astrid}

\textit{Jeg forstår godt, at billedet ser vildt ud! \da{På trods af} det kolde vejr tror danskerne på, at børn sover sundere i den friske luft. \da{Til trods for, at} forældrene \dabold{ikke sidder} lige ved siden af vognen, har de en babyalarm med ind på caféen.}

\textit{Jeg har også oplevet den samme tillid ude på landet, hvor folk sælger jordbær og honning i ubemandede vejboder. Her lægger man bare pengene i en kasse eller bruger MobilePay. \da{I modsætning til} vores hjemland, hvor man typisk har meget kontrol og høje hegn af sikkerhedshensyn, ser danskerne \dabold{derimod} tillid som noget, alle passer på sammen. Det gør hverdagen meget mere afslappet, selvom man selvfølgelig stadig bruger sin sunde fornuft.}

\textit{Hvad synes du selv om den tradition?}

\textit{Kærlig hilsen\\
\mbox{[Dit navn]}}

\small\textbf{点评}：全篇 129 词，完整覆盖来信三大问题；精准联动本课四大核心表达（\textit{på trods af} + 名词短语、\textit{til trods for, at} + 带 \textit{ikke} 从句、\textit{i modsætning til} 以及 \textit{derimod}），同时呼应 Unit 2 前几课的 \textit{afslappet} 与社会文化视角，是一篇极具说服力的 Modul 3 高分范文。
\end{answerbox}


% =====================================================================
% Unit 3
% =====================================================================
\ncdunit{3}{Familie, Forbrug og Digitalisering}{家庭、消费与数字化}{Lektion 11--15：介词 + 形式宾语 det + 从句 · jo... desto... · 舆论引述 · 趋势动词 · 分数与群体搭配}

\ncdlesson{11}{MitID, e-Boks og det papirløse samfund}{MitID、电子邮箱与无纸化社会}
{DU3 Modul 3（B1）｜对标 NCE 2 后半--NCE 3 前半 \& Modul 3 \textit{Skrivning Opgave 1（半正式求助/投诉信）}}
{丹麦语必考句法：\textbf{介词 +（形式宾语 \textit{det}）+ \textit{at / hv-} 从句}（\textit{regne med, at... / være i tvivl om, hvorvidt... / sætte pris på (det), at...}）}
{本课聚焦丹麦高数字化社会中的 MitID 与 e-Boks。进入 Modul 3 后，阅读和写作大量出现“动词/形容词 + 介词 + 从句”的长句。请牢记：丹麦语中\textbf{介词绝不能省略}，后面接 \textit{at / hv- / om} 从句时必须保留介词，并在需要占位、强调或接 \textit{hvis/når} 从句时插入形式宾语 \textit{det}！}

% ---------------------------------------------------------------------
\begin{foerduleaeser}
带着下面两个问题阅读 Tekst A 和 Tekst B：
\begin{enumerate}
  \item \textit{Hvad kan man bruge MitID og e-Boks til i Danmark, og hvilke ulemper kan det papirløse samfund give ifølge Tekst A?}（在丹麦可以用 MitID 和 e-Boks 做什么？根据 Tekst A，无纸化社会可能带来哪些弊端？）
  \item \textit{Hvilket problem har Freja med sit MitID i Tekst B, og hvad beder hun borgerservice om at gøre?}（Tekst B 中芙蕾雅的 MitID 遇到了什么问题？她请求市民服务中心做什么？）
\end{enumerate}
\end{foerduleaeser}

% ---------------------------------------------------------------------
\ncdsubsec{二、Tekst A \& Tekst B（数字社会与 Modul 3 真题场景）}

\subsection*{Tekst A：\textit{Livet med MitID og e-Boks}（社会观察短文 · 约 240 词）}

\begin{tipbox}{语篇说明}
本篇介绍丹麦“无纸化数字政务”的日常运作与社会争议。阅读时请重点标出文中\textbf{\color{NCDNavy}深蓝色加粗的“动词/形容词 + 介词 + 从句”}以及\textbf{\color{NCDRed}深红色加粗的“介词 + 形式宾语 \textit{det} + 从句”}结构。
\end{tipbox}

I Danmark foregår næsten al kommunikation mellem borgere og det offentlige digitalt. Så snart man flytter til landet og får sit CPR-nummer, skal man oprette MitID og lære at tjekke sin digitale post i e-Boks. Myndighederne regner nemlig \da{med, at} alle borgere læser deres breve elektronisk.

For mange tilflyttere kan overgangen til det papirløse samfund virke overvældende i begyndelsen. De fleste bliver hurtigt overraskede \da{over, hvor} mange ting der kræver et digitalt login: lægebesøg, skat, bankoverførsler og endda biblioteksbøger. Efter nogle måneder sætter næsten alle dog stor pris \da{på, at} man kan klare hverdagens praktiske opgaver hjemmefra på få minutter. Man slipper for lange køer på kommunen, og vigtige dokumenter bliver aldrig væk i postkassen.

Samtidig må man være opmærksom \da{på, at} systemet også stiller store krav. Hvis man glemmer at tjekke e-Boks i flere uger, risikerer man at overse en vigtig frist eller en regning. Derudover er der stadig en del borgere, især ældre, som er i tvivl \da{om, hvorvidt} den hurtige digitalisering nu også gavner alle. Mange har svært ved at vænne sig til \dabold{det, at} fysiske kontorer lukker, og at man bliver helt afhængig af sin smartphone.

Hvis telefonen pludselig går i stykker eller bliver stjålet, bliver man mindet \da{om, hvor} sårbar hverdagen kan være. Heldigvis kan man få en fysisk kodeviser på borgerservice eller i særlige tilfælde blive fritaget for Digital Post. Alligevel er de fleste eksperter enige \da{om, at} det papirløse samfund er kommet for at blive.

\begin{transbox}
\small
\textbf{\color{NCDTeal}[Tekst A 中文精译]}\\
在丹麦，公民与公共部门之间几乎所有的沟通都在网上进行。人一搬到丹麦并拿到个人身份号（CPR 号），就必须开通 MitID，并学会在 e-Boks 里查看电子政务信件。因为政府部门默认所有公民都会以电子方式阅读信件。

对许多新移民来说，向无纸化社会的过渡起初可能会让人有些招架不住。大多数人很快会惊讶地发现，竟然有这么多事情需要数字登录：看医生、报税、银行转账，甚至连借阅图书馆书籍也不例外。不过几个月后，几乎所有人都会非常欣赏这种便利——人们在家花几分钟就能搞定日常琐事，既不用在市政厅排长队，重要文件也不会在信箱里弄丢。

与此同时，人们也必须留意到这套系统同样提出了很高的要求。如果一连几周忘记查看 e-Boks，就可能错过重要截止日期或账单。此外，仍然有一部分公民——尤其是老年人——对如此快速的数字化是否真的惠及所有人持怀疑态度。许多人很难习惯实体办事处纷纷关闭、而且人变得完全依赖智能手机这一事实。

如果手机突然坏了或被偷了，人们就会立刻意识到日常生活能有多脆弱。幸好，人们可以在市民服务中心领取实体密码器（kodeviser），在特殊情况下也可以申请豁免接收电子政务邮件。尽管如此，大多数专家都一致认为，无纸化社会已成定局。
\end{transbox}

\subsection*{Tekst B：\textit{Henvendelse til borgerservice om MitID}（对标 DU3 Modul 3 Skrivning Opgave 1）}

\textbf{写作场景对照：}\textit{Freja har skiftet telefon og kan ikke længere bruge MitID. Hun skriver en halvformel besked til kommunens borgerservice og dækker alle fire punkter: problemet, konsekvensen, hvad hun har prøvet, og hvad hun ønsker.}

\vspace{0.3em}
\begin{emailbox}{市民服务求助信（125 词）｜Til: Borgerservice \quad Emne: Hjælp til MitID efter telefonskift}
\textbf{Kære Borgerservice}

Jeg skriver til jer, fordi jeg desværre har fået et problem med mit MitID efter at have skiftet telefon i går.

Jeg regnede \da{med, at} jeg selv kunne flytte appen til min nye mobil, men systemet viser en fejlmeddelelse hele tiden. Nu kan jeg hverken logge ind i e-Boks eller i min netbank. Jeg er derfor i tvivl \da{om, hvorvidt} jeg har fået et vigtigt brev fra kommunen om min boligstøtte.

Jeg har allerede prøvet at genstarte telefonen og ringe til support, men desværre uden held.

Jeg vil sætte stor pris \dabold{på det, hvis} I kan give mig en hurtig tid på rådhuset, så jeg kan få en ny aktiveringskode. Kan jeg også bestille en fysisk kodeviser som backup?

\textbf{Venlig hilsen}\\
\textbf{Freja Hansen}
\end{emailbox}

\begin{transbox}
\small
\textbf{\color{NCDTeal}[Tekst B 中文对照]}\\
\textbf{尊敬的市民服务中心}：我写信给你们，是因为我昨天换了手机后，很遗憾 MitID 出了问题。我原本以为自己能把应用迁移到新手机上，但系统一直显示报错信息。现在我既登录不了 e-Boks，也登录不了网银。因此我也不确定自己是否收到了市政厅关于住房补贴的重要信件。我已经试过重启手机和拨打客服热线，但遗憾的是没有成功。如果你们能在市政厅给我安排一个较快的预约时间，让我获取新的激活码，我将非常感激。我也可以顺便申请一个实体密码器作为备用吗？诚挚的问候，芙蕾雅·汉森。
\end{transbox}

\begin{tipbox}{本课发音与吞音要点（Udtale \& Reduktioner）}
\begin{itemize}
  \item \textbf{数字政务高频词（软 \textit{d} 与不发音 \textit{d/g}）}：\textit{et samfund} \ipa{ˈsɑmˌfʌnˀ}（词尾 \textit{nd} 中的 \textit{d} 不发音，带喉塞音 \ipa{ˀ}）；\textit{en myndighed} \ipa{ˈmøndiˌheðˀ}（中间 \textit{g} 不发音，词尾 \textit{-hed} 读轻软 \textit{d} 加喉塞音）；\textit{digitalisering} \ipa{diɡitˢaliˈseːˀɐeŋ}（重音在 \textit{-se-}）。
  \item \textbf{介词短语弱读与连音}：\textit{regner med} \ipa{ˈʁɑi̯nɐ mɛ}（\textit{egn} 发滑音 \ipa{ɑi̯}，\textit{med} 在句中经常弱读省掉软 \textit{d}）；\textit{i tvivl om} \ipa{i ˈtˢviʊ̯ˀl ʌm}（\textit{tvivl} 中第二个 \textit{v} 滑音化为 \ipa{ʊ̯}）；\textit{hvorvidt} \ipa{vɒˈvid}（词首 \textit{hv-} 的 \textit{h} 不发音，词尾 \textit{dt} 合并读 \ipa{d}）。
  \item \textbf{其他高频词}：\textit{afhængig} \ipa{ɑʊ̯ˈhɛŋˀi}（前缀 \textit{af-} 读 \ipa{ɑʊ̯}，词尾 \textit{-ig} 的 \textit{g} 不发音）；\textit{opmærksom} \ipa{ʌbˈmɛɐ̯ɡsʌmˀ}；\textit{sætte pris på} \ipa{ˈsɛdə ˈpʰʁiːˀs pʰʌ}。
\end{itemize}
\end{tipbox}

% ---------------------------------------------------------------------
\ncdsubsec{三、Gloser（核心生词表）}

\begin{tipbox}{词表重点}
本课重点掌握“动词/形容词 + 固定介词”的捆绑记忆：例如 \textit{at regne \textbf{med}}（预计/指望）、\textit{at sætte pris \textbf{på}}（欣赏/感激）、\textit{at minde nogen \textbf{om}}（提醒某人……）、\textit{at være i tvivl \textbf{om}}（对……存疑）。另外注意 \textit{en myndighed} 等以 \textit{-hed} 结尾的名词均为共性（\textit{en}）。
\end{tipbox}

\begin{glosertable}
名词 & en borger & \textit{(borgeren, borgere, borgerne)} & 公民，市民 \\
名词 & en myndighed & \textit{(myndigheden, myndigheder, myndighederne)} & 官方机构，政府部门 \\
名词 & et samfund & \textit{(samfundet, samfund, samfundene)} & 社会 \\
名词 & en tilflytter & \textit{(tilflytteren, tilflyttere, tilflytterne)} & 新移居者，新住民 \\
名词 & en overgang & \textit{(overgangen, overgange, overgangene)} & 过渡，转变 \\
名词 & en ulempe & \textit{(ulempen, ulemper, ulemperne)} & 缺点，弊端 \\
名词 & en frist & \textit{(fristen, frister, fristerne)} & 截止期限 \\
名词 & en kodeviser & \textit{(kodeviseren, kodevisere, kodeviserne)} & 实体密码器（MitID 硬件） \\
\midrule
动词 [har] & at regne med & \textit{(regner, regnede, har regnet med)} & 指望，预计，认为 \\
动词 [har] & at sætte pris på & \textit{(sætter, satte, har sat pris på)} & 欣赏，珍惜，感激 \\
动词 [har] & at minde (nogen) om & \textit{(minder, mindede, har mindet om)} & 提醒（某人）想起某事 \\
动词 [har] & at overse & \textit{(overser, overså, har overset)} & 忽略，漏看 \\
动词 [har] & at fritage (for) & \textit{(fritager, fritog, har fritaget)} & 豁免，免除（接 \textit{for}） \\
动词 [har] & at genstarte & \textit{(genstarter, genstartede, har genstartet)} & 重启（设备/系统） \\
\midrule
形容词 & papirløs & \textit{(papirløst, papirløse)} & 无纸化的 \\
形容词 & afhængig (af) & \textit{(afhængigt, afhængige)} & 依赖……的（接 \textit{af}） \\
形容词 & sårbar / utryg (ved) & \textit{(sårbart, sårbare) / (utrygt, utrygge)} & 脆弱的 / 不安的，没把握的 \\
\end{glosertable}

% ---------------------------------------------------------------------
\ncdsubsec{四、Grammatik \& Faste vendinger（核心语法与地道表达精讲）}

\begin{grammatikbox}{语法点 1：丹麦语核心句法——介词 +（形式宾语 \textit{det}）+ \textit{at / hv-} 从句}
在丹麦语中，大量核心动词、形容词和名词短语都自带固定介词。当后面需要接一整个句子（从句）作为宾语时，\textbf{绝对不能把介词删掉}，必须按以下三种路径连接：
\begin{enumerate}
  \item \textbf{路径 A｜介词 + \textit{, at...}（陈述一个事实或预期）}：\\
    公式：\textbf{动词/形容词 + 介词 + \textit{, at} + 从句}（注意：逗号永远打在\textbf{介词后面}！）。\\
    例：\textit{Myndighederne regner \da{med, at} alle læser e-Boks.} / \textit{Vi er enige \da{om, at} det er smart.}
  \item \textbf{路径 B｜介词 + \textit{, om / hvorvidt / hv-}（表达不确定、问题或“是否”）}：\\
    公式：\textbf{动词/形容词 + 介词 + \textit{, om / hvorvidt / hv-} + 从句}（这类从句前一般\textbf{不加} \textit{det}）。\\
    例：\textit{De er i tvivl \da{om, hvorvidt} det hjælper.} / \textit{Man bliver mindet \da{om, hvor} sårbar man er.}
  \item \textbf{路径 C｜介词 + 形式宾语 \textit{det} + \textit{, at / hvis / når...}（占位、强调或防歧义）}：\\
    什么时候必须在介词与从句之间插入形式宾语 \dabold{det}？主要看三大场景：
    \begin{itemize}
      \item \ding{172} \textbf{后面接 \textit{hvis} 或 \textit{når} 从句时（必须加 \textit{det}！）}：\textit{Jeg sætter stor pris \dabold{på det, hvis} I kan hjælpe mig.}
      \item \ding{173} \textbf{强调“……这件事本身”（相当于“……这一事实”）}：\textit{Mange har svært ved at vænne sig \dabold{til det, at} fysiske kontorer lukker.}
      \item \ding{174} \textbf{介词短语接 \textit{at} 易与连词混淆时（如 \textit{trods, på grund af, ved, efter}）}：\textit{Trods \dabold{det, at} systemet er nyt, fungerer det godt.}
    \end{itemize}
\end{enumerate}
\end{grammatikbox}

\vspace{0.25em}
\renewcommand{\arraystretch}{1.24}
{\small
\begin{xltabular}{\linewidth}{@{} >{\bfseries\color{NCDNavy}\raggedright\arraybackslash}p{3.8cm} >{\raggedright\arraybackslash}p{3.9cm} >{\raggedright\arraybackslash}X @{}}
\toprule
\textbf{高频“词 + 介词”搭配} & \textbf{中文语义} & \textbf{标准 Modul 3 造句示范} \\
\midrule
\endfirsthead
\toprule
\textbf{高频“词 + 介词”搭配} & \textbf{中文语义} & \textbf{标准 Modul 3 造句示范（续表）} \\
\midrule
\endhead
\bottomrule
\endfoot
regne \textbf{med} & 预计，指望，默认 & \textit{Jeg regner \da{med, at} jeg får svar i morgen.} \\
sætte pris \textbf{på} & 欣赏，重视，感激 & \textit{Jeg sætter pris \da{på, at} det går hurtigt.} / \textit{... \dabold{på det, hvis}...} \\
være i tvivl \textbf{om} & 对……不确定/存疑 & \textit{Han er i tvivl \da{om, hvorvidt} brevet er sendt.} \\
være opmærksom \textbf{på} & 留意到，注意到 & \textit{Vær opmærksom \da{på, at} fristen udløber på fredag.} \\
være overrasket \textbf{over} & 对……感到惊讶 & \textit{Jeg blev overrasket \da{over, hvor} nemt det var.} \\
være enig \textbf{om} & 对……达成一致 & \textit{Alle er enige \da{om, at} vi skal hjælpe ældre.} \\
være utryg \textbf{ved} & 对……感到不安/没底 & \textit{Nogle er utrygge \da{ved, at} alt er digitalt.} \\
minde (nogen) \textbf{om} & 提醒（某人）…… & \textit{Appen minder os \da{om, hvornår} vi skal betale.} \\
komme an \textbf{på} & 取决于…… & \textit{Det kommer an \da{på, hvordan} man bruger systemet.} \\
vænne sig \textbf{til} & 习惯于…… & \textit{Man vænner sig hurtigt \da{til (det), at} bruge MitID.} \\
\end{xltabular}
}

\begin{grammatikbox}{英丹对比避坑指南（English vs Danish Pitfalls）}
\begin{enumerate}
  \item \textbf{英语删介词 vs 丹麦语死守介词}：英语不能说 \textit{*I count on that...}（往往省略介词改用 \textit{expect that} 或动名词 \textit{count on doing}）；但丹麦语\textbf{绝不能丢掉介词}，必须说 \ding{51}\ \textit{Jeg regner \dabold{med, at}...}，不能写成 \ding{55}\ \textit{*Jeg regner, at...}！
  \item \textbf{礼貌请求别漏掉 \textit{det}}：写邮件表达“如果你们能……我会很感激”时，英语说 \textit{I would appreciate it if...}，丹麦语同样必须带 \dabold{det}：\ding{51}\ \textit{Jeg vil sætte pris \dabold{på det, hvis} I kan svare snart.}（绝不能写成 \ding{55}\ \textit{*sætte pris på, hvis...}）。
  \item \textbf{表“是否”用 \textit{om / hvorvidt}，严禁用 \textit{hvis}}：在 \textit{i tvivl om, ...} 后面表示“是否”，只能用 \da{om} 或更正式的 \da{hvorvidt}，千万不要受英语 \textit{if} 影响误用 \textit{hvis}！
  \item \textbf{从句内部语序 + \textit{der} 补位}：从句中副词 \textit{ikke / ofte / stadig} 放在动词前；若 \textit{hv-} 词充当从句主语（如 \textit{over, hvor mange ting...}），记住要补主语占位词 \dabold{der}：\textit{overrasket over, hvor mange der bruger MitID}。
\end{enumerate}
\end{grammatikbox}

% ---------------------------------------------------------------------
\ncdsubsec{五、Mønsterøvelser（句型操练）}

\textbf{练习 A：填入正确的固定介词（\textit{med, på, om, over, ved, til, af}）或形式宾语/连词（\textit{det, at, om, hvorvidt}）}
\begin{enumerate}
  \item Kommunen regner \underline{\hspace{1.4cm}}, \underline{\hspace{1.2cm}} alle borgere tjekker deres digitale post.
  \item Jeg vil sætte stor pris på \underline{\hspace{1.4cm}}, hvis du kan ringe tilbage i eftermiddag.
  \item Mange nye studerende er i tvivl \underline{\hspace{1.4cm}}, \underline{\hspace{1.8cm}} de skal bruge MitID til eksamen.
  \item Vi blev meget overraskede \underline{\hspace{1.4cm}}, hvor hurtigt vi fik svar fra banken.
  \item Vær opmærksom \underline{\hspace{1.4cm}}, at du ikke deler din personlige kode med andre.
  \item Flere ældre borgere er utrygge \underline{\hspace{1.4cm}}, at de ikke længere kan betale med girokort på posthuset.
  \item Systemet minder mig altid \underline{\hspace{1.4cm}}, hvornår jeg skal forny mit pas.
  \item Det kommer an \underline{\hspace{1.4cm}}, hvordan man sikrer sin telefon med en stærk kode.
\end{enumerate}

\textbf{练习 B：句型外科手术（合并两句或修正下列病句，注意保留介词与形式宾语 \textit{det}）}
\begin{enumerate}
  \item 合并句子：\textit{Jeg regner med det. Min nye kodeviser kommer med posten på fredag.} $\rightarrow$ \textit{Jeg regner med, at...}
  \item 合并句子：\textit{Hun er i tvivl om det. Har hun fået en besked i e-Boks?}（用 \textit{hvorvidt}）$\rightarrow$ \textit{Hun er i tvivl om, hvorvidt...}
  \item 合并句子：\textit{Vi sætter stor pris på det. Hvis I kan løse problemet inden weekenden.} $\rightarrow$ \textit{Vi sætter stor pris...}
  \item 纠错改写（漏了介词）：\textit{*Jeg er meget glad, at du kunne hjælpe mig med MitID.} $\rightarrow$
  \item 纠错改写（误用 \textit{hvis} 表“是否”）：\textit{*Han er i tvivl om, hvis han skal ringe til borgerservice.} $\rightarrow$
  \item 纠错改写（逗号位置与从句否定语序错误）：\textit{*De regner, med at systemet virker ikke i nat.} $\rightarrow$
\end{enumerate}

\textbf{练习 C：中译丹（灵活运用“介词 +（\textit{det}）+ 从句”）}
\begin{enumerate}
  \item 我预计市政厅会在两天内给我回复。
  \item 我们非常感激大家能按时查看自己的电子邮箱。
  \item 我不确定我是否已经收到了来自银行的信件。
  \item 这完全取决于你在新手机上如何使用这个应用。
  \item 如果您能尽快查看我的申请，我将非常感激。（用 \textit{på det, hvis}）
\end{enumerate}

% ---------------------------------------------------------------------
\ncdsubsec{六、DU3 Modul 3 实战微任务与参考答案}

\begin{emailbox}{Modul 3 Skrivning Opgave 1 实战微任务（银行数字化故障投诉/求助信）}
\textbf{Situation}: Du er kunde i NordBank. I tre dage har du ikke kunnet betale dine regninger i netbanken, fordi appen fryser, når du skal godkende med MitID. Du har ringet til kundeservice to gange, men ingen har ringet tilbage.\\
\textbf{Opgave}: Skriv en høflig, men bestemt klage til bankens kundeservice. Du skal komme ind på alle fire punkter:
\begin{enumerate}
  \item Hvad problemet er, og hvor længe det har varet.
  \item Hvilke konsekvenser problemet har for dig.
  \item Hvad du allerede har prøvet at gøre.
  \item Hvad du forventer, at banken gør nu (brug mindst to sætninger med \textit{præposition + (det) + at/hv-}).
\end{enumerate}
\end{emailbox}

\begin{answerbox}
\textbf{Før du læser 参考回答}
\begin{enumerate}
  \item \textit{Man kan bruge MitID og e-Boks til at læse post fra det offentlige, bestille lægetid, tjekke skat, bruge netbank og låne biblioteksbøger. Ulemperne er, at man kan overse vigtige frister, at nogle ældre bliver utrygge, og at hverdagen er sårbar, hvis telefonen går i stykker.}
  \item \textit{Freja kan ikke flytte MitID-appen til sin nye telefon og kan derfor ikke logge ind i e-Boks eller netbank. Hun beder borgerservice om en hurtig tid på rådhuset for at få en ny aktiveringskode og spørger, om hun kan bestille en fysisk kodeviser.}
\end{enumerate}

\textbf{练习 A 答案}
\begin{enumerate}
  \item \da{med}；\da{at}（固定搭配：\textit{regne med, at...}）
  \item \da{det}（后接 \textit{hvis} 从句，介词 \textit{på} 后必须加形式宾语 \textit{det}）
  \item \da{om}；\da{om} / \da{hvorvidt}（\textit{være i tvivl om, om/hvorvidt...}）
  \item \da{over}（\textit{overrasket over, hvor...}）
  \item \da{på}（\textit{opmærksom på, at...}）
  \item \da{ved}（\textit{utryg ved, at...}）
  \item \da{om}（\textit{minde nogen om, hvornår...}）
  \item \da{på}（\textit{komme an på, hvordan...}）
\end{enumerate}

\textbf{练习 B 答案（句型手术参考答案）}
\begin{enumerate}
  \item \textit{Jeg regner med, at min nye kodeviser kommer med posten på fredag.}
  \item \textit{Hun er i tvivl om, hvorvidt hun har fået en besked i e-Boks.}（注意一般疑问句语序还原为主语 \textit{hun} 在动词 \textit{har} 前！）
  \item \textit{Vi sætter stor pris på det, hvis I kan løse problemet inden weekenden.}
  \item \textit{Jeg er meget glad \dabold{for}, at du kunne hjælpe mig med MitID.}（\textit{glad for, at...} 不可丢掉介词 \textit{for}）
  \item \textit{Han er i tvivl om, \dabold{om / hvorvidt} han skal ringe til borgerservice.}（不用 \textit{hvis}）
  \item \textit{De regner \dabold{med, at} systemet \dabold{ikke virker} i nat.}（逗号在 \textit{med} 后；从句里 \textit{ikke} 放在动词 \textit{virker} 前）
\end{enumerate}

\textbf{练习 C 答案}
\begin{enumerate}
  \item \textit{Jeg regner med, at kommunen svarer mig inden for to dage.}
  \item \textit{Vi sætter stor pris på, at alle tjekker deres e-Boks til tiden.}
  \item \textit{Jeg er i tvivl om, hvorvidt jeg allerede har fået brevet fra banken.}
  \item \textit{Det kommer helt an på, hvordan du bruger appen på den nye telefon.}
  \item \textit{Jeg vil sætte stor pris på det, hvis I kan se på min ansøgning snarest.}
\end{enumerate}

\textbf{Modul 3 实战高分范文（109 词）}\\
\textit{Til: NordBank Kundeservice \quad|\quad Emne: Klage over fejl i netbank og MitID}

\textit{Kære NordBank}

\textit{Jeg skriver til jer, fordi jeg i tre dage desværre ikke har kunnet betale mine regninger i jeres netbank. Hver gang jeg skal godkende en betaling med MitID, fryser appen og lukker ned.}

\textit{Det er et stort problem for mig, fordi flere faste regninger har frist i morgen. Jeg er nu meget bekymret \da{for, at} jeg skal betale gebyrer for for sen betaling. Jeg har allerede ringet til jeres kundeservice to gange og genstartet appen, men desværre uden svar.}

\textit{Jeg regner \da{med, at} I hurtigst muligt kan undersøge fejlen, og jeg vil sætte stor pris \dabold{på det, hvis} en medarbejder ringer til mig i dag.}

\textit{Venlig hilsen\\
\mbox{[Dit navn]}}

\small\textbf{点评}：严格覆盖 Modul 3 \textit{Skrivning Opgave 1} 全部四个要点；对银行团队统一使用复数 \textit{jer / jeres / I}；自然融入了 \textit{bekymret for, at...}、\textit{regner med, at...} 和 \textit{sætte stor pris på det, hvis...} 三组本课核心句型，逻辑清晰、礼貌而坚定，完全达到 B1 高分标准。
\end{answerbox}
\ncdlesson{12}{Delingsøkonomi og genbrugstøj}{共享经济与二手服装}
{DU3 Modul 3（B1）｜对标 NCE 2 后半--NCE 3 前半 \& Modul 3 \textit{Skrivning Opgave 1（二手交易维权沟通）}}
{\textbf{递进比较句型（\textit{jo} + 比较级..., \textit{desto/jo} + 比较级...）}与“前从后主”V2 倒装铁律 · \textbf{等比比较（\textit{lige så... som} vs \textit{ikke så... som}）}}
{丹麦的二手商店（\textit{genbrugsbutik}）、跳蚤市场与共享出行极度普及。本课核心句型 \textit{Jo..., desto...} 是 B1 阶段表达因果递进的必杀技，但中文与英语学习者极易在“后半句忘记倒装”上丢分——请牢记：\textbf{前半句 \textit{jo} 是从句（不倒装），后半句 \textit{desto/jo} 是主句（动词必须提前）！}}

% ---------------------------------------------------------------------
\begin{foerduleaeser}
带着下面两个问题阅读 Tekst A 和 Tekst B：
\begin{enumerate}
  \item \textit{Hvorfor vælger flere danskere genbrugstøj og delingsøkonomi i dag ifølge Tekst A?}（根据 Tekst A，为什么如今越来越多丹麦人选择二手衣物和共享经济？）
  \item \textit{Hvad er Sofie utilfreds med i sin besked til Jonas i Tekst B, og hvilke to løsninger foreslår hun?}（在 Tekst B 中，索菲对发给乔纳斯的交易商品有什么不满？她提出了哪两种解决方案？）
\end{enumerate}
\end{foerduleaeser}

% ---------------------------------------------------------------------
\ncdsubsec{二、Tekst A \& Tekst B（数字社会与 Modul 3 真题场景）}

\subsection*{Tekst A：\textit{Fra ”køb og smid væk” til genbrug og deling}（社会观察短文 · 约 245 词）}

\begin{tipbox}{语篇说明}
本篇探讨丹麦从“快消弃置文化”向循环共享生活方式的转变。阅读时请特别留意文中\textbf{\color{NCDNavy}深蓝色加粗的 \textit{jo..., desto...} 与 \textit{lige så... som} 比较结构}，以及\textbf{\color{NCDRed}深红色标注的 \textit{desto} 后半句 V2 倒装}。
\end{tipbox}

I dag tænker mange danskere mere over klima og forbrug end tidligere. Førhen købte man ofte billigt tøj og smed det ud efter få måneder, men nu vælger flere at købe brugt i genbrugsbutikker eller via apps. \da{Jo} mere opmærksomme forbrugerne bliver på miljøet, \dabold{desto større bliver} interessen for genbrugstøj.

Kvaliteten er tit overraskende høj. En varm vinterjakke fra et kendt mærke er ofte \da{lige så pæn og holdbar som} en helt ny model, men den koster måske kun halvt så meget. \da{Jo} flere gange en trøje bliver brugt, \dabold{desto mindre belaster} den kloden. Mange unge finder også sjældne jakker og støvler på loppemarkeder i weekenden, så brugt tøj føles ikke længere som en nødløsning, men som et bevidst valg.

Samme tendens ser man i resten af delingsøkonomien. I storbyer som København og Aarhus er det blevet \da{lige så nemt} at leje en nabos bil via en app \da{som} at tage bussen. Flere boligforeninger har endda et fælles værksted, hvor beboerne kan låne boremaskiner og symaskiner i stedet for at eje alting selv. Selvom det kan kræve lidt mere planlægning at dele ressourcer med andre, oplever de fleste, at fordelene vejer tungere. \da{Jo} bedre vi bliver til at dele og genbruge, \dabold{desto færre penge bruger vi}, og \dabold{desto bedre passer vi} på fremtiden.

Forbrugerne har dog også lært, at man skal se sig godt for på nettet. En vare er nemlig ikke altid \da{lige så fejlfri som} på billederne i annoncen, og derfor er tydelige beskrivelser afgørende.

\begin{transbox}
\small
\textbf{\color{NCDTeal}[Tekst A 中文精译]}\\
如今，许多丹麦人比过去更加关注气候与消费问题。从前人们常常买便宜衣服，穿几个月就扔掉，但现在越来越多人选择在二手店或通过手机应用购买二手衣物。消费者对环境越关注，对二手衣物的兴趣就越大。

二手衣物的质量往往高得令人意外。一件知名品牌的保暖冬衣通常和全新款一样美观耐用，价格却可能只有一半。一件毛衣被穿着的次数越多，对地球的负担就越小。许多年轻人周末也会在跳蚤市场上淘到罕见的外套和靴子，因此二手服饰不再被视为无奈之举，而是一种有意识的生活选择。

在共享经济的其他领域也能看到同样的趋势。在哥本哈根和奥胡斯这样的大城市，通过应用租用邻居的汽车已经变得和乘公交车一样方便。不少住房协会甚至设有公共工具室，住户可以借用电钻和缝纫机，而不必样样都自己买。尽管与他人共享资源可能需要多花一点心思规划，但大多数人感到利远大于弊。我们越擅长分享和循环利用，花的钱就越少，也就能越好好地守护未来。

不过，消费者也学到了在网上交易必须擦亮眼睛。因为商品并不总是像广告照片上看起来那样完美无瑕，所以清晰准确的成色描述至关重要。
\end{transbox}

\subsection*{Tekst B：\textit{Reklamation efter et genbrugskøb}（对标 DU3 Modul 3 二手交易维权沟通）}

\textbf{沟通场景：}\textit{Sofie har købt en brugt uldtrøje via en app, men den har en skjult fejl. Læs beskederne mellem køber og sælger.}

\vspace{0.3em}
\begin{emailbox}{平台交易消息（122 词）｜Køber: Sofie \quad $\leftrightarrow$ \quad Sælger: Jonas}
\textbf{Sofie:} Hej Jonas! Jeg har lige hentet pakken med den grå uldtrøje, jeg købte af dig. Tak for hurtig forsendelse. Desværre er trøjen ikke \da{lige så pæn som} beskrevet i annoncen: der er et lille hul på højre ærme, og en knap mangler. \da{Jo} hurtigere vi finder en god løsning, \dabold{desto bedre er} det for os begge. Kan vi aftale et afslag på 100 kroner via MobilePay, eller skal jeg sende den retur?

\vspace{0.35em}
\hrule
\vspace{0.35em}

\textbf{Jonas:} Hej Sofie! Det er jeg virkelig ked af. Jeg havde slet ikke set hullet, da jeg pakkede trøjen. Selvfølgelig skal prisen være \da{lige så rimelig som} standen. Jeg overfører gerne 100 kroner tilbage med det samme, hvis du stadig vil beholde den --- ellers betaler jeg selvfølgelig for returfragten!
\end{emailbox}

\begin{transbox}
\small
\textbf{\color{NCDTeal}[Tekst B 中文对照]}\\
\textbf{索菲}：你好，乔纳斯！我刚取到从你这里买的灰色羊毛衫包裹。谢谢你发货这么快。很遗憾，这件毛衣并不像广告里描述的那么完好：右袖子上有一个小洞，还少了一颗纽扣。我们越快找到一个好办法，对我们俩就越好。我们能协商通过 MobilePay 退 100 克朗差价吗，还是说我应该把衣服退回去？

\textbf{乔纳斯}：你好，索菲！我真的很抱歉。我打包毛衣时完全没注意到那个洞。价格当然应该和实际成色一样合理。如果你还想留着它，我很乐意立刻退回 100 克朗——不然的话，退货运费当然由我来承担！
\end{transbox}

\begin{tipbox}{本课发音与吞音要点（Udtale \& Reduktioner）}
\begin{itemize}
  \item \textbf{二手与环保词群（注意硬 \textit{g} 与不发音 \textit{d/g}）}：\textit{genbrugstøj} \ipa{ˈɡɛnˌbʁuːˀsˌtˢʌi̯}（词首 \textit{gen-} 读硬 \ipa{ɡ}，不可读成 \textit{j}！\textit{-øj} 为滑音 \ipa{ʌi̯}）；\textit{holdbar} \ipa{ˈhʌlˌbɑːˀ}（\textit{-ld-} 中的 \textit{d} 不发音）；\textit{bæredygtig} \ipa{ˈbɛːɐˌdøɡdi}（词尾 \textit{-ig} 的 \textit{g} 不发音）；\textit{uldtrøje} \ipa{ˈulˌtˢʁʌi̯ə}（\textit{uld} 中的 \textit{d} 不发音）。
  \item \textbf{比较句连读与弱读}：\textit{lige så pæn som} \ipa{ˈliːə sʌ ˈpʰɛːˀn sʌm}（\textit{lige} 中的 \textit{g} 不发音，\textit{så} 轻读）；\textit{halvt så meget} \ipa{ˈhalˀd sʌ ˈmɑːð}（\textit{halv-} 的 \textit{v} 常脱落，\textit{meget} 吞音为 \ipa{mɑːð}）；\textit{desto bedre} \ipa{desˈtˢo ˈbɛðʁɐ}（软 \textit{d} \ipa{ð}）。
  \item \textbf{其他高频词}：\textit{delingsøkonomi} \ipa{ˈdeːleŋsˌøkʰonoˌmiːˀ}；\textit{et miljø} \ipa{milˈjøːˀ}；\textit{returfragt} \ipa{ʁeˈtˢuɐ̯ˀˌfʁɑɡd}。
\end{itemize}
\end{tipbox}

% ---------------------------------------------------------------------
\ncdsubsec{三、Gloser（核心生词表）}

\begin{tipbox}{词表重点}
注意 \textit{at beholde}（保留，留下）的过去时与过去分词是不规则短音变化 \textit{beholdt / har beholdt}；\textit{at smide ud}（扔掉）为 \textit{smed ud / har smidt ud}。\textit{en stand}（成色、状态）在描述二手物品时固定用单数定指 \textit{standen}（如 \textit{i god stand} 成色良好）。
\end{tipbox}

\begin{glosertable}
名词 & en delingsøkonomi & \textit{(delingsøkonomien, ingen flertal)} & 共享经济 \\
名词 & genbrugstøj (utælleligt) & \textit{(genbrugstøjet, ingen flertal)} & 二手服装（不可数） \\
名词 & en forbruger & \textit{(forbrugeren, forbrugere, forbrugerne)} & 消费者 \\
名词 & et miljø & \textit{(miljøet, miljøer, miljøerne)} & 环境 \\
名词 & en nødløsning & \textit{(nødløsningen, nødløsninger, nødløsningerne)} & 权宜之计，无奈的选择 \\
名词 & en vare & \textit{(varen, varer, varerne)} & 商品，货物 \\
名词 & en stand & \textit{(standen, ingen flertal)} & （物品的）成色，状况 \\
名词 & et ærme & \textit{(ærmet, ærmer, ærmerne)} & 袖子 \\
名词 & en knap & \textit{(knappen, knapper, knapperne)} & 纽扣；按钮 \\
名词 & et afslag & \textit{(afslaget, afslag, afslagene)} & 折扣，减价，退差价 \\
名词 & en returfragt & \textit{(returfragten, ingen flertal)} & 退货运费 \\
\midrule
动词 [har] & at dele & \textit{(deler, delte, har delt)} & 分享，共享 \\
动词 [har] & at leje & \textit{(lejer, lejede, har lejet)} & 租用，租赁 \\
动词 [har] & at belaste & \textit{(belaster, belastede, har belastet)} & 给……造成负担，消耗 \\
动词 [har] & at beholde & \textit{(beholder, beholdt, har beholdt)} & 保留，留下 \\
动词 [har] & at smide ud & \textit{(smider ud, smed ud, har smidt ud)} & 扔掉，丢弃 \\
\midrule
形容词 & holdbar / bæredygtig & \textit{(holdbart, holdbare) / (bæredygtigt, bæredygtige)} & 耐用的 / 可持续的 \\
形容词 & fejlfri / rimelig & \textit{(fejlfrit, fejlfrie) / (rimeligt, rimelige)} & 无瑕疵的 / 合理的 \\
\end{glosertable}

% ---------------------------------------------------------------------
\ncdsubsec{四、Grammatik \& Faste vendinger（核心语法与地道表达精讲）}

\begin{grammatikbox}{语法点 1：递进比较句 \textit{jo + 比较级..., desto / jo + 比较级...}（“越……，就越……”）}
表示两个变化量同步递增或递减时，丹麦语使用成对关联结构 \da{jo..., desto...}（口语中也常用 \da{jo..., jo...}）。最核心的考核点在于\textbf{前半句与后半句截然相反的语序规则}：
\begin{enumerate}
  \item \textbf{前半句（\textit{jo}-从句）：严格遵循从句语序（不倒装！）}
    \begin{itemize}
      \item 公式：\textbf{\textit{Jo} + 比较级短语 + 主语 +（副词 \textit{ikke/ofte}）+ 动词}
      \item 例：\textit{\da{Jo} flere penge \textbf{man sparer}, ...} / \textit{\da{Jo} oftere \textbf{vi ikke køber} nyt, ...}
    \end{itemize}
  \item \textbf{后半句（\textit{desto/jo}-主句）：必须触发 V2 主谓倒装！}
    \begin{itemize}
      \item 公式：\textbf{\textit{desto / jo} + 比较级短语 + \dabold{变位动词 + 主语}}
      \item 例：\textit{..., \dabold{desto bedre bliver} økonomien.} / \textit{..., \dabold{desto mindre belaster den} miljøet.}
    \end{itemize}
\end{enumerate}
\end{grammatikbox}

\vspace{0.25em}
\renewcommand{\arraystretch}{1.24}
{\small
\begin{xltabular}{\linewidth}{@{} >{\bfseries\color{NCDNavy}\raggedright\arraybackslash}p{3.4cm} >{\raggedright\arraybackslash}p{4.3cm} >{\raggedright\arraybackslash}X @{}}
\toprule
\textbf{比较级成分类型} & \textbf{结构公式} & \textbf{地道 B1 例句（留意后半句倒装）} \\
\midrule
\endfirsthead
\toprule
\textbf{比较级成分类型} & \textbf{结构公式} & \textbf{地道 B1 例句（续表）} \\
\midrule
\endhead
\bottomrule
\endfoot
形容词/副词比较级 & \textit{Jo + -ere/mere..., desto + -ere/mere...} & \textit{\da{Jo} hurtigere vi handler, \dabold{desto nemmere bliver} løsningen.} \\
不规则比较级 & \textit{bedre, værre, større, mindre, længere} & \textit{\da{Jo} længere tid man venter, \dabold{desto værre bliver} problemet.} \\
带名词（可数复数） & \textit{Jo flere / færre + 复数名词...} & \textit{\da{Jo} flere der deler bil, \dabold{desto færre biler kører} på gaden.} \\
带名词（不可数/程度） & \textit{Jo mere / mindre (+ 不可数名词)...} & \textit{\da{Jo} mindre tøj vi smider ud, \dabold{desto mere sparer vi}.} \\
\end{xltabular}
}

\begin{grammatikbox}{语法点 2：等比与倍数比较（\textit{lige så... som}）与英丹对比避坑指南}
表达“和……一样……”或“一半/两倍……”时，用原级形容词或副词夹在 \da{lige så} 与 \da{som} 之间：
\begin{itemize}
  \item \textbf{肯定等比}：\textit{lige så + 原级 + som} $\rightarrow$ \textit{Jakken er \da{lige så varm som} en ny model.}
  \item \textbf{否定不等比}：\textit{ikke (lige) så + 原级 + som} $\rightarrow$ \textit{Trøjen er \da{ikke så pæn som} på billedet.}
  \item \textbf{倍数对比}：\textit{halvt så + 原级 + som}（一半……）、\textit{dobbelt så + 原级 + som}（两倍……）：\textit{Den koster kun \da{halvt så meget som} en ny.}
\end{itemize}

\vspace{0.3em}
\textbf{英丹对比避坑指南（English vs Danish Pitfalls）}：
\begin{enumerate}
  \item \textbf{英语前后均不倒装 vs 丹麦语后半句必倒装}：英语说 \textit{The more you share, the more you save}（前后都是主语在前）；丹麦语后半句必须倒装成 \ding{51}\ \textit{Jo mere man deler, desto mere \dabold{sparer man}}，绝不能写成 \ding{55}\ \textit{*desto mere man sparer}！
  \item \textbf{可数 \textit{flere/færre} vs 不可数 \textit{mere/mindre}}：英语的 \textit{more} 与 \textit{less} 在口语中常混用，但丹麦语在可数复数名词前（如 \textit{penge, biler, gange}）必须用 \da{flere / færre}；在不可数名词或单独修饰动词时必须用 \da{mere / mindre}。
  \item \textbf{短形容词比较级别套英语 \textit{more}}：不能把 \textit{cheaper} 译成 \textit{*jo mere billig}，而必须写规则比较级 \ding{51}\ \textit{jo billigere}。
  \item \textbf{肯定句不能省略 \textit{lige}}：英语说 \textit{as good as}，丹麦语肯定等比不能只写 \textit{*så god som}，必须完整写出 \ding{51}\ \textit{lige så god som}（仅在否定句 \textit{ikke så god som} 中可省去 \textit{lige}）。
\end{enumerate}
\end{grammatikbox}

% ---------------------------------------------------------------------
\ncdsubsec{五、Mønsterøvelser（句型操练）}

\textbf{练习 A：用 \textit{jo, desto, lige, så, som, mere, flere, mindre, færre} 填空（注意句意与可数/不可数搭配）}
\begin{enumerate}
  \item \underline{\hspace{1.4cm}} mere man køber i genbrugsbutikker, \underline{\hspace{1.6cm}} mere skåner man miljøet.
  \item Min brugte cykel er \underline{\hspace{1.4cm}} så hurtig \underline{\hspace{1.4cm}} min vens helt nye cykel.
  \item Jo \underline{\hspace{1.6cm}} folk bruger delebiler, desto \underline{\hspace{1.6cm}} biler er der i centrum.
  \item Denne frakke er desværre ikke \underline{\hspace{1.4cm}} så fejlfri \underline{\hspace{1.4cm}} på billedet i annoncen.
  \item Jo \underline{\hspace{1.6cm}} tid vi bruger på at reparere gamle ting, desto \underline{\hspace{1.6cm}} affald laver vi.
  \item En brugt sofa koster tit kun halvt \underline{\hspace{1.2cm}} meget \underline{\hspace{1.2cm}} en ny fra butikken.
  \item Jo \underline{\hspace{1.8cm}} (tidligt) du skriver til sælgeren, desto \underline{\hspace{1.8cm}} (godt) er det.
  \item Jeg synes ikke, at fast fashion er \underline{\hspace{1.4cm}} så bæredygtigt \underline{\hspace{1.4cm}} genbrugstøj.
\end{enumerate}

\textbf{练习 B：句型手术——按要求用 \textit{Jo..., desto...} 或 \textit{lige så... som} 合并或改错}
\begin{enumerate}
  \item 合并句子：\textit{Vi genbruger meget tøj. Vi sparer mange penge.} $\rightarrow$ \textit{Jo mere tøj vi genbruger, desto...}
  \item 合并句子：\textit{Man bor længe i storbyen. Man bruger delingsøkonomi ofte.} $\rightarrow$ \textit{Jo længere...}
  \item 合并句子：\textit{Den gamle symaskine er god. Den nye symaskine er også god.} $\rightarrow$ \textit{Den gamle symaskine er...}
  \item 纠错改写（后半句漏了倒装）：\textit{*Jo billigere billetten bliver, desto hurtigere vi køber den.} $\rightarrow$
  \item 纠错改写（比较级与可数词错配）：\textit{*Jo mere billig trøjen er, desto flere tid sparer jeg.} $\rightarrow$
  \item 纠错改写（肯定句缺少 \textit{lige}）：\textit{*Den brugte frakke er så varm som en ny vinterjakke.} $\rightarrow$
\end{enumerate}

\textbf{练习 C：中译丹（灵活运用递进比较与等比比较）}
\begin{enumerate}
  \item 我们买越多的二手衣服，对环境的负担就越小。
  \item 你越早给我回信，我们就越容易找到解决办法。
  \item 这件二手外套和全新的一样暖和，却只花一半的钱。
  \item 越多人共享汽车，街上的交通问题就越少。
  \item 遗憾的是，这双鞋并不像图片上看起来那样完好无损。
\end{enumerate}

% ---------------------------------------------------------------------
\ncdsubsec{六、DU3 Modul 3 实战微任务与参考答案}

\begin{emailbox}{Modul 3 Skrivning Opgave 1 实战微任务（二手平台交易维权信）}
\textbf{Situation}: Du har købt en brugt vinterjakke af en privat sælger, Mikkel, på en dansk genbrugsapp. I annoncen stod der, at jakken var ”som ny”, men da du åbnede pakken, opdagede du, at lynlåsen (拉链) ikke virker, og at der er en stor plet på lommen.\\
\textbf{Opgave}: Skriv en høflig, men tydelig besked til Mikkel. Du skal dække alle fire punkter:
\begin{enumerate}
  \item Hvad du har købt, og hvad der stod i annoncen.
  \item Hvilke to fejl du har fundet på jakken.
  \item Hvordan det opleves for dig (brug mindst én sætning med \textit{ikke lige så... som}).
  \item Hvilken løsning du foreslår (brug mindst én sætning med \textit{jo..., desto...}).
\end{enumerate}
\end{emailbox}

\begin{answerbox}
\textbf{Før du læser 参考回答}
\begin{enumerate}
  \item \textit{Flere danskere vælger genbrugstøj og delingsøkonomi for at passe på klimaet og miljøet, og fordi brugt tøj ofte er lige så pænt og holdbart som nyt tøj til halv pris, samtidig med at man sparer penge.}
  \item \textit{Sofie er utilfreds med, at den grå uldtrøje har et lille hul på højre ærme og mangler en knap. Hun foreslår enten et afslag på 100 kroner via MobilePay eller at sende trøjen retur.}
\end{enumerate}

\textbf{练习 A 答案}
\begin{enumerate}
  \item \da{Jo}；\da{desto}
  \item \da{lige}；\da{som}（肯定等比：\textit{lige så... som}）
  \item \da{flere}；\da{færre}（\textit{folk} 与 \textit{biler} 均为复数可数概念）
  \item \da{lige}；\da{som}（也可省略 \textit{lige}，但保留 \textit{ikke lige så... som} 语气更精准）
  \item \da{mere}；\da{mindre}（\textit{tid} 和 \textit{affald} 为不可数名词）
  \item \da{så}；\da{som}（\textit{halvt så meget som} 一半那么多）
  \item \da{tidligere}；\da{bedre}（不规则比较级：\textit{godt $\rightarrow$ bedre}）
  \item \da{lige}；\da{som}（\textit{ikke lige så... som}）
\end{enumerate}

\textbf{练习 B 答案（句型手术参考答案）}
\begin{enumerate}
  \item \textit{Jo mere tøj vi genbruger, desto \dabold{flere penge sparer vi}.}（后半句触发 V2 倒装 \textit{sparer vi}！）
  \item \textit{Jo længere tid man bor i storbyen, desto \dabold{oftere bruger man} delingsøkonomi.}
  \item \textit{Den gamle symaskine er lige så god som den nye.}
  \item \textit{Jo billigere billetten bliver, desto \dabold{hurtigere køber vi} den.}（注意动词 \textit{køber} 必须抢在第 2 位，置于主语 \textit{vi} 之前！）
  \item \textit{Jo \dabold{billigere} trøjen er, desto \dabold{mere} tid sparer jeg.}（短形容词用 \textit{-ere}，不可数 \textit{tid} 用 \textit{mere}）
  \item \textit{Den brugte frakke er \dabold{lige} så varm som en ny vinterjakke.}（肯定等比必须有 \textit{lige}）
\end{enumerate}

\textbf{练习 C 答案}
\begin{enumerate}
  \item \textit{Jo mere genbrugstøj vi køber, desto mindre belaster vi miljøet.}
  \item \textit{Jo tidligere du svarer mig, desto nemmere finder vi en løsning.}
  \item \textit{Denne brugte jakke er lige så varm som en helt ny, men den koster kun halvt så meget.}
  \item \textit{Jo flere mennesker der deler bil, desto færre trafikproblemer er der på gaderne.}
  \item \textit{Desværre er disse sko ikke lige så fejlfrie, som de ser ud på billederne.}
\end{enumerate}

\textbf{Modul 3 实战高分范文（126 词）}\\
\textit{Til: Mikkel \quad|\quad Emne: Problem med den brugte vinterjakke}

\textit{Hej Mikkel}

\textit{Jeg har i dag modtaget den vinterjakke, som jeg købte af dig i mandags. I din annonce stod der, at jakken var ”som ny”, så jeg glædede mig meget til at bruge den.}

\textit{Desværre viste det sig, da jeg åbnede pakken, at jakken ikke er \da{lige så pæn som} beskrevet. For det første virker lynlåsen slet ikke, og for det andet er der en stor mørk plet på den ene lomme. Det er jeg selvfølgelig ret skuffet over, fordi jeg havde regnet med, at den var i god stand.}

\textit{\da{Jo} hurtigere vi kan finde en fair løsning, \dabold{desto bedre er} det for os begge. Kunne jeg få et afslag på 150 kroner til reparation, eller skal jeg sende jakken tilbage til dig?}

\textit{Venlig hilsen\\
\mbox{[Dit navn]}}

\small\textbf{点评}：流畅运用“先肯定后指瑕”的丹麦协商礼仪；准确嵌入 \textit{ikke lige så pæn som} 和 \textit{Jo hurtigere..., desto bedre er det...}，同时自然复现了第 11 课的 \textit{regnet med, at...} 句型。126 词结构完整，达到 B1 优秀水平。
\end{answerbox}
\ncdlesson{13}{Børneopdragelse og ”curlingforældre”}{育儿观念与“冰壶父母”}
{DU3 Modul 3（B1）｜对标 NCE 2 后半--NCE 3 前半 \& Modul 3 \textit{Læsning Opgave 4（多方观点对比）}}
{\textbf{客观引述与社会舆论表达}（\textit{Ifølge eksperter..., Det diskuteres ofte, om..., Mange peger på, at..., Undersøgelser viser, at...}）· \textit{På den ene side... på den anden side...} 立场平衡}
{“冰壶父母”（\textit{curlingforældre}）是丹麦社会关于家庭教育的热门词，也是 Modul 3 阅读与口试的高频思辨主题。本课帮你摆脱千篇一律的 \textit{Jeg synes}，掌握 B1 考官最看重的\textbf{客观引述与舆论对比框架}，并同步演练 \textit{Læsning Opgave 4} 三人观点匹配法！}

% ---------------------------------------------------------------------
\begin{foerduleaeser}
带着下面两个问题阅读 Tekst A 和 Tekst B：
\begin{enumerate}
  \item \textit{Hvad betyder ordet ”curlingforældre”, og hvorfor advarer pædagoger imod denne opdragelse ifølge Tekst A?}（“冰壶父母”是什么意思？根据 Tekst A，教育工作者为什么对这种育儿方式提出警示？）
  \item \textit{Hvilke tre forskellige syn på børns selvstændighed og forældrekontrol møder vi hos Henrik, Maja og Sofie i Tekst B?}（在 Tekst B 中，亨里克、玛雅和索菲对孩子独立性与父母干预分别持哪三种不同看法？）
\end{enumerate}
\end{foerduleaeser}

% ---------------------------------------------------------------------
\ncdsubsec{二、Tekst A \& Tekst B（数字社会与 Modul 3 真题场景）}

\subsection*{Tekst A：\textit{Når forældre fejer vejen foran børnene}（社会文化短文 · 约 235 词）}

\begin{tipbox}{语篇说明}
本篇剖析丹麦社会关于“冰壶父母”的代际育儿讨论。阅读时请特别关注文中\textbf{\color{NCDNavy}深蓝色加粗的客观引述与舆论表达}，以及\textbf{\color{NCDRed}深红色标注的句首状语 V2 倒装}。
\end{tipbox}

I Danmark lægger man traditionelt stor vægt på, at børn tidligt bliver selvstændige. Allerede i børnehaven øver børnene sig i at tage tøj på selv, og mange skolebørn cykler alene til skole. Alligevel hører man i dag ofte ordet ”curlingforældre” i den offentlige debat. Udtrykket beskriver forældre, der hele tiden fejer alle problemer af vejen foran deres børn, ligesom i sporten curling.

\da{Ifølge flere psykologer} \dabold{er fænomenet blevet} mere tydeligt i de seneste år. Nogle forældre pakker stadig deres store børns skoletaske, kører dem til fritidsaktiviteter hver dag eller skriver lange mails til læreren om små konflikter. \da{Det diskuteres ofte, om} denne form for omsorg i virkeligheden hjælper eller skader børnene.

\da{Mange eksperter peger på, at} børn har brug for at begå små fejl og selv finde løsninger. Hvis mor og far altid rydder op efter konflikter i skolen eller på fritidshjemmet, lærer de unge ikke at håndtere modgang senere i livet. \da{En ny undersøgelse viser desuden, at} næsten hver femte studerende oplever, at forældrene blander sig i deres studie eller første job.

\da{På den ene side} ønsker forældre naturligvis bare at skabe tryghed og sikre, at deres børn klarer sig godt i et konkurrencepræget samfund. \da{På den anden side} \dabold{advarer pædagoger om, at} overbeskyttede børn ofte bliver mere usikre som voksne. \da{Derfor foreslår mange rådgivere, at} voksne tør træde et skridt tilbage og give børnene plads til selv at vokse med ansvaret.

\begin{transbox}
\small
\textbf{\color{NCDTeal}[Tekst A 中文精译]}\\
在丹麦，人们传统上非常重视让孩子早早学会独立。早在幼儿园阶段，孩子们就练习自己穿衣服，许多学龄儿童也会独自骑车上学。尽管如此，如今人们在公共讨论中仍常听到“冰壶父母”（curlingforældre）一词。这个词用来形容那些像玩冰壶一样，时刻替孩子扫清前方一切障碍的父母。

据多位心理学家称，这一现象近年来变得越来越明显。有些父母甚至还替大孩子收拾书包，每天开车接送他们参加课外活动，或者因为一点小矛盾就给老师写长邮件。人们经常讨论，这种关怀究竟是在帮助还是在伤害孩子。

许多专家指出，孩子需要经历犯小错并自己寻找解决办法的过程。如果爸爸妈妈总是替孩子收拾学校或课后托管中心里的矛盾，年轻人就学不会在日后生活中应对挫折。此外，一项新调查显示，近五分之一的大学生表示父母会干预他们的学业或第一份工作。

一方面，父母当然只是想营造安全感，确保孩子在竞争激烈的社会中表现良好；另一方面，教育工作者警告说，受到过度保护的孩子成年后往往会变得更加缺乏自信。因此，许多顾问建议成年人应当敢于退后一步，给孩子留出伴随责任成长的空间。
\end{transbox}

\subsection*{Tekst B：\textit{Tre stemmer i debatten om opdragelse}（对标 DU3 Modul 3 Læsning Opgave 4）}

\textbf{阅读任务（Læsning Opgave 4 模拟）：}\textit{Læs de tre synspunkter fra Henrik (A), Maja (B) og Sofie (C), og match spørgsmålene (1)--(4) med den rigtige person.}

\vspace{0.3em}
\begin{emailbox}{三方观点对比（135 词）｜Børneopdragelse og digitale vaner}
\textbf{A. Henrik (far til to på 9 og 12 år)}\\
Jeg mener, at børn lærer bedst, når de selv får ansvar. Mine børn cykler selv til fodbold, også når det regner. Glemmer de deres sportstøj, kører jeg ikke efter dem, for de bliver stærkere af at mærke konsekvensen.

\vspace{0.3em}
\textbf{B. Maja (mor til en dreng på 11 år)}\\
Jeg tjekker tit min søns lokation på mobilen, når han er på vej hjem. Det gør jeg ikke for at styre ham, men fordi trafikken er tæt. God opdragelse handler også om tryghed, indtil børnene er helt klar.

\vspace{0.3em}
\textbf{C. Sofie (pædagog i en fritidsklub)}\\
Mange forældre skriver til os om helt små konflikter mellem børnene. Men børn har brug for selv at øve sig i at tale sammen. Hvis voksne altid ordner alting, bliver de unge usikre senere.
\end{emailbox}

\vspace{0.25em}
\textbf{Hvem siger hvad?（请选择最符合该描述的发言人 A、B 或 C）：}
\begin{enumerate}
  \item \textit{Hvem bruger teknologi til at holde øje med, hvor barnet befinder sig?} \quad \textbf{A / B / C}
  \item \textit{Hvem synes, at børn må tage konsekvensen, hvis de glemmer deres ting?} \quad \textbf{A / B / C}
  \item \textit{Hvem oplever på sit arbejde, at forældre blander sig i børnenes små uenigheder?} \quad \textbf{A / B / C}
  \item \textit{Hvem peger på, at travl trafik gør forældrekontrol nødvendig for trygheden?} \quad \textbf{A / B / C}
\end{enumerate}

\begin{transbox}
\small
\textbf{\color{NCDTeal}[Tekst B 中文对照]}\\
\textbf{A. 亨里克（9 岁和 12 岁两个孩子的父亲）}：我认为孩子在自己承担责任时学得最好。我的孩子哪怕下雨也是自己骑车去踢足球。如果他们忘了带运动服，我不会开车送过去，因为亲身感受后果会让他们变得更坚强。

\textbf{B. 玛雅（11 岁男孩的母亲）}：当儿子在回家的路上时，我经常在手机上查看他的定位。我这么做不是为了控制他，而是因为交通很拥堵。在孩子完全准备好之前，良好的育儿同样关乎安全感。

\textbf{C. 索菲（课后俱乐部教育员）}：很多父母会因为孩子们之间非常小的摩擦给我们写信。但孩子需要自己练习如何沟通。如果成年人总是包办一切，年轻人以后就会变得不自信。
\end{transbox}

\begin{tipbox}{本课发音与吞音要点（Udtale \& Reduktioner）}
\begin{itemize}
  \item \textbf{外来词与静音字母}：\textit{psykolog} \ipa{sykʰoˈloːˀ}（词首 \textit{p} 完全不发音！重音在末音节带喉塞音）；\textit{curlingforældre} \ipa{ˈkʰøɐ̯leŋˌfɒˌɛlˀdʁə}（丹麦语把英语 \textit{curling} 元音圆唇化为 \ipa{øɐ̯}）。
  \item \textbf{长复合词与软 \textit{d} / 不发音 \textit{g}}：\textit{børneopdragelse} \ipa{ˈbɶɐ̯nəˌʌbˌdʁɑːˀəlsə}（中间的 \textit{g} 不发音）；\textit{selvstændighed} \ipa{sɛlˈsdɛnˀdiˌheðˀ}（\textit{selv} 常省掉 \textit{v} 读 \ipa{sɛl}，\textit{-ig} 中 \textit{g} 不发音）；\textit{undersøgelse} \ipa{ˈʌnɐˌsøːˀəlsə}（\textit{g} 不发音）；\textit{pædagog} \ipa{pʰɛðɑˈɡoːˀ}（前一个 \textit{d} 为软 \textit{d} \ipa{ð}）。
  \item \textbf{引述高频词读音}：\textit{ifølge} \ipa{iˈføljə}（\textit{lg} 中的 \textit{g} 发成半元音 \ipa{j}！）；\textit{modgang} \ipa{ˈmoðˌɡɑŋˀ}；\textit{ansvar} \ipa{ˈanˌsvɑːˀ}。
\end{itemize}
\end{tipbox}

% ---------------------------------------------------------------------
\ncdsubsec{三、Gloser（核心生词表）}

\begin{tipbox}{词表重点}
注意情态动词 \textit{at turde}（敢于）的不规则变化：\textit{tør, turde, har turdet}，后面直接加不带 \textit{at} 的动词原形（如 \textit{tør træde tilbage}）。动词 \textit{at begå}（犯错/实施）常与 \textit{en fejl} 搭配，而“应对逆境”常用 \textit{at håndtere modgang}。
\end{tipbox}

\begin{glosertable}
名词 & en opdragelse & \textit{(opdragelsen, opdragelser, opdragelserne)} & 教育，养育方式 \\
名词 & en selvstændighed & \textit{(selvstændigheden, ingen flertal)} & 独立性，自主能力 \\
名词 & et ansvar & \textit{(ansvaret, ingen flertal)} & 责任，职责 \\
名词 & en omsorg & \textit{(omsorgen, ingen flertal)} & 关怀，照料 \\
名词 & en tryghed & \textit{(trygheden, ingen flertal)} & 安全感 \\
名词 & en modgang & \textit{(modgangen, ingen flertal)} & 逆境，挫折 \\
名词 & en konflikt & \textit{(konflikten, konflikter, konflikterne)} & 冲突，矛盾 \\
名词 & en undersøgelse & \textit{(undersøgelsen, undersøgelser, undersøgelserne)} & 调查，研究 \\
名词 & en pædagog & \textit{(pædagogen, pædagoger, pædagogerne)} & 教育员，幼教老师 \\
名词 & et fænomen & \textit{(fænomenet, fænomener, fænomenerne)} & 现象 \\
\midrule
动词 [har] & at feje & \textit{(fejer, fejede, har fejet)} & 清扫，扫除 \\
动词 [har] & at pege på & \textit{(peger, pegede, har peget på)} & 指出，表明 \\
动词 [har] & at advare (om) & \textit{(advarer, advarede, har advaret om)} & 警告，提醒注意 \\
动词 [har] & at begå (en fejl) & \textit{(begår, begik, har begået)} & 犯（错），做出 \\
动词 [har] & at blande sig (i) & \textit{(blander, blandede, har blandet sig i)} & 干预，插手 \\
动词 [har] & at turde & \textit{(tør, turde, har turdet)} & 敢于（情态动词，接原形） \\
\midrule
形容词 & selvstændig & \textit{(selvstændigt, selvstændige)} & 独立的，自主的 \\
形容词 & overbeskyttet & \textit{(overbeskyttet, overbeskyttede)} & 被过度保护的 \\
\end{glosertable}

% ---------------------------------------------------------------------
\ncdsubsec{四、Grammatik \& Faste vendinger（核心语法与地道表达精讲）}

\begin{grammatikbox}{语法点 1：客观引述与社会舆论四大句型骨架}
在 DU3 Modul 3 的辩论、阅读匹配和书信写作中，不能通篇只用 \textit{Jeg synes}。把个人观点升级为“客观引述 + 正反平衡”，只需掌握下面四套句法模板：
\end{grammatikbox}

\vspace{0.25em}
\renewcommand{\arraystretch}{1.24}
{\small
\begin{xltabular}{\linewidth}{@{} >{\bfseries\color{NCDNavy}\raggedright\arraybackslash}p{3.6cm} >{\raggedright\arraybackslash}p{4.3cm} >{\raggedright\arraybackslash}X @{}}
\toprule
\textbf{语用功能} & \textbf{核心句型公式} & \textbf{课文与实战示范} \\
\midrule
\endfirsthead
\toprule
\textbf{语用功能} & \textbf{核心句型公式} & \textbf{课文与实战示范（续表）} \\
\midrule
\endhead
\bottomrule
\endfoot
引述来源（据……） & \textit{Ifølge + [专家/调查], + \dabold{V2 倒装}} & \textit{\da{Ifølge flere psykologer} \dabold{er} fænomenet blevet mere tydeligt.} \\
引述研究/观点 & \textit{[专家/研究] + peger på / viser / advarer om, at + 从句} & \textit{\da{En ny undersøgelse viser, at} næsten hver femte oplever det.} \\
无人称舆论引入 & \textit{Det diskuteres ofte / Man spørger, \dabold{om / hvorvidt} + 从句} & \textit{\da{Det diskuteres ofte, om} denne omsorg skader børnene.} \\
双面平衡过渡 & \textit{På den ene side + \dabold{V2}... På den anden side + \dabold{V2}...} & \textit{\da{På den ene side} ønsker de tryghed. \da{På den anden side} advarer eksperter...} \\
\end{xltabular}
}

\begin{grammatikbox}{语法点 2：引述动词的立场光谱 \& 英丹对比避坑指南}
不同引述动词传达的作者态度并不一样，B1 写作可以灵活切换：
\begin{itemize}
  \item \textbf{客观中立陈述}：\textit{pege på, at...}（指出）、\textit{påpege, at...}（强调指出）、\textit{vise, at...}（表明/显示）、\textit{mener, at...}（认为）。
  \item \textbf{警示或建议立场}：\textit{advare om, at...}（警告说……）、\textit{foreslå, at...}（建议……）、\textit{fremhæve, at...}（突出强调）。
\end{itemize}

\vspace{0.3em}
\textbf{英丹对比避坑指南（English vs Danish Pitfalls）}：
\begin{enumerate}
  \item \textbf{\textit{Ifølge} 后面千万不能不倒装，也不能加 \textit{til}}：英语说 \textit{According to experts, children need...}（有逗号且主语在动词前）；丹麦语 \da{Ifølge eksperter} \textbf{本身就是整个第 1 位}，后面必须直接 V2 倒装且不打逗号：\ding{51}\ \textit{Ifølge eksperter \dabold{har børn} brug for...}（绝不能写成 \ding{55}\ \textit{*Ifølge til eksperter, børn har...}！）。
  \item \textbf{\textit{om / hvorvidt} vs \textit{hvis}}：在 \textit{Det diskuteres, om...}（人们讨论“是否……”）中，对应英语的 \textit{whether/if}，丹麦语\textbf{只能用 \da{om} 或 \da{hvorvidt}}，绝对不能受英语 \textit{if} 影响误写成 \ding{55}\ \textit{*Det diskuteres, hvis...}！
  \item \textbf{\textit{pege på} vs \textit{påpege}}：英语说 \textit{point out that}，丹麦语不能照搬写成 \textit{*pege ud, at}！正确形式要么是两个词 \ding{51}\ \textit{pege \dabold{på}, at...}，要么是一体词 \ding{51}\ \textit{\dabold{påpege}, at...}（注意前重后轻差异）。
  \item \textbf{从句副词前移}：在 \textit{Det diskuteres, om...} 或 \textit{Mange peger på, at...} 后面的从句中，\textit{ikke / ofte / altid} 依然必须放到动词前面：\ding{51}\ \textit{... at unge \dabold{ikke lærer} at håndtere modgang.}
\end{enumerate}
\end{grammatikbox}

% ---------------------------------------------------------------------
\ncdsubsec{五、Mønsterøvelser（句型操练）}

\textbf{练习 A：用本课客观引述与舆论表达词填空（\textit{ifølge, på, om, hvorvidt, viser, advarer, ene, anden}）}
\begin{enumerate}
  \item \underline{\hspace{1.6cm}} nye tal fra kommunen cykler flere børn til skole i dag end sidste år.
  \item Det diskuteres ofte i medierne, \underline{\hspace{1.8cm}} unge mennesker bliver for afhængige af deres forældre.
  \item Mange pædagoger peger \underline{\hspace{1.2cm}}, at børn skal have lov til at løse små konflikter selv.
  \item En ny rapport \underline{\hspace{1.6cm}}, at hver fjerde elev får hjælp til lektierne hver aften.
  \item På den \underline{\hspace{1.2cm}} side vil forældre beskytte deres børn; på den \underline{\hspace{1.4cm}} side skal børnene lære ansvar.
  \item Psykologer \underline{\hspace{1.8cm}} om, at for meget kontrol kan gøre børnene usikre.
  \item Mange er i tvivl om, \underline{\hspace{1.8cm}} det er en god idé at bruge GPS-apps til store børn.
  \item Rådgiveren \underline{\hspace{1.8cm}} (foreslå), at forældre lader børnene pakke deres egen taske.
\end{enumerate}

\textbf{练习 B：语体升级与纠错——把“个人口语句”升级为“客观社会引述句”，或修正语法错误}
\begin{enumerate}
  \item 升级为 \textit{Ifølge...}：\textit{Eksperterne siger, at børn lærer meget af deres fejl.} $\rightarrow$ \textit{Ifølge eksperterne...}
  \item 升级为 \textit{Det diskuteres ofte, om...}：\textit{Folk spørger hele tiden: Hjælper forældre deres børn for meget?} $\rightarrow$
  \item 升级为 \textit{Mange peger på, at...}：\textit{Jeg tror, at unge bliver mere selvstændige uden kontrol.} $\rightarrow$
  \item 纠错改写（受英语倒装与介词干扰）：\textit{*Ifølge til en undersøgelse, børn bruger ikke mindre tid med forældre.} $\rightarrow$
  \item 纠错改写（误用 \textit{hvis} 和直译 \textit{pege ud}）：\textit{*Læreren peger ud, hvis forældre skriver for mange mails.} $\rightarrow$ （注：把句意“老师指出父母写太多邮件”修正为地道丹麦语）
  \item 纠错改写（\textit{På den ene side} 漏了倒装）：\textit{*På den anden side eksperter mener, at debatten er vigtig.} $\rightarrow$
\end{enumerate}

\textbf{练习 C：中译丹（运用客观引述、\textit{om/at} 从句与正反对比）}
\begin{enumerate}
  \item 据专家称，如今许多父母帮孩子做得太多了。
  \item 人们经常讨论，电子设备是否会影响家庭里的沟通。
  \item 许多老师指出，孩子需要学会自己面对小挫折。
  \item 一项新的调查显示，大部分年轻人其实想自己做决定。
  \item 一方面家长想创造安全感，另一方面孩子也需要空间去独立成长。
\end{enumerate}

% ---------------------------------------------------------------------
\ncdsubsec{六、DU3 Modul 3 实战微任务与参考答案}

\begin{emailbox}{Modul 3 Skrivning Opgave 2 / 口试 Opgave 2 实战微任务（回应朋友关于育儿观的提问）}
\textbf{Situation}: Din ven, Thomas, har læst en artikel om ”curlingforældre” i Danmark og skriver til dig for at høre din mening.\\
\textbf{Opgave}: Skriv en mail til Thomas på \textbf{mindst 100 ord}, hvor du svarer på alle spørgsmål og bruger mindst tre objektive vendinger fra denne lektion (\textit{Ifølge..., Det diskuteres..., Mange peger på..., På den ene side... på den anden side...}):
\begin{enumerate}
  \item Hvad synes du om, at nogle forældre hjælper deres børn med næsten alting?
  \item Hvordan var børneopdragelsen i din egen familie eller i dit hjemland?
  \item Hvad er det vigtigste, som børn skal lære for at blive selvstændige?
\end{enumerate}
\end{emailbox}

\begin{answerbox}
\textbf{Før du læser 参考回答}
\begin{enumerate}
  \item \textit{”Curlingforældre” er forældre, der fejer alle problemer af vejen foran deres børn. Pædagoger advarer imod det, fordi overbeskyttede børn ikke lærer at håndtere modgang og ofte bliver mere usikre som voksne.}
  \item \textit{Henrik mener, at børn lærer bedst ved selv at få ansvar og mærke konsekvensen; Maja bruger mobilens lokation for tryghed i tæt trafik; Sofie oplever, at forældre blander sig i små konflikter, men mener, at børn selv skal øve sig i at løse dem.}
\end{enumerate}

\textbf{Tekst B 阅读答案（Læsning Opgave 4）}
\begin{enumerate}
  \item \textbf{B (Maja)} —— \textit{Jeg tjekker tit min søns lokation på mobilen, når han er på vej hjem.}
  \item \textbf{A (Henrik)} —— \textit{Glemmer de deres sportstøj, kører jeg ikke efter dem, for de bliver stærkere af at mærke konsekvensen.}
  \item \textbf{C (Sofie)} —— \textit{Mange forældre skriver til os om helt små konflikter mellem børnene.}
  \item \textbf{B (Maja)} —— \textit{... fordi trafikken er tæt. God opdragelse handler også om tryghed...}
\end{enumerate}

\textbf{练习 A 答案}
\begin{enumerate}
  \item \da{Ifølge}（据……数据称）
  \item \da{om} / \da{hvorvidt}（\textit{Det diskuteres ofte, om/hvorvidt...}）
  \item \da{på}（固定搭配：\textit{pege på, at...}）
  \item \da{viser}（研究显示：\textit{viser, at...}）
  \item \da{ene}；\da{anden}（\textit{På den ene side..., på den anden side...}）
  \item \da{advarer}（\textit{advare om, at...}）
  \item \da{hvorvidt} / \da{om}（\textit{i tvivl om, hvorvidt/om...}）
  \item \da{foreslår}（单数主语现在时）
\end{enumerate}

\textbf{练习 B 答案（语体升级参考答案）}
\begin{enumerate}
  \item \textit{Ifølge eksperterne \dabold{lærer børn} meget af deres fejl.}（V2 倒装：\textit{lærer børn}）
  \item \textit{Det diskuteres ofte, \dabold{om} forældre hjælper deres børn for meget.}
  \item \textit{Mange peger \dabold{på, at} unge bliver mere selvstændige uden kontrol.}
  \item \textit{Ifølge en undersøgelse \dabold{bruger børn ikke} mindre tid med forældre.}（删去 \textit{til} 和逗号，触发 V2 倒装 \textit{bruger børn}）
  \item \textit{Læreren \dabold{peger på, at} (eller \dabold{påpeger, at}) forældre skriver for mange mails.}（不用 \textit{*pege ud}，陈述事实用 \textit{at}）
  \item \textit{På den anden side \dabold{mener eksperter}, at debatten er vigtig.}（V2 倒装：\textit{mener eksperter}）
\end{enumerate}

\textbf{练习 C 答案}
\begin{enumerate}
  \item \textit{Ifølge eksperter hjælper mange forældre deres børn for meget i dag.}
  \item \textit{Det diskuteres ofte, om digitale medier påvirker kommunikationen i familien.}
  \item \textit{Mange lærere peger på, at børn skal lære selv at håndtere modgang.}（或：\textit{klare små nederlag}）
  \item \textit{En ny undersøgelse viser, at de fleste unge faktisk gerne vil træffe egne beslutninger.}
  \item \textit{På den ene side vil forældre skabe tryghed, men på den anden side har børn også brug for plads til at vokse selvstændigt.}
\end{enumerate}

\textbf{Modul 3 实战高分范文（128 词）}\\
\textit{Til: Thomas \quad|\quad Emne: Sv: Tanker om ”curlingforældre”}

\textit{Hej Thomas}

\textit{Tak for din spændende mail! Jeg har også læst meget om ”curlingforældre” her i Danmark, og det er et virkelig interessant emne.}

\textit{\da{Det diskuteres ofte, om} forældre i dag hjælper deres børn for meget. \da{På den ene side} kan jeg godt forstå, at mor og far vil beskytte deres børn mod stress. \da{På den anden side peger mange eksperter på, at} børn kun bliver selvstændige, hvis de får lov til at begå fejl og selv løse små problemer.}

\textit{I min egen familie lærte jeg tidligt at hjælpe til derhjemme og tage ansvar for mine egne lektier. \da{Ifølge min erfaring} \dabold{giver det} stor selvtillid senere i livet. Jeg synes derfor, at det vigtigste er at give børn både tryghed og frihed.}

\textit{Hvad tænker du selv?}

\textit{Mange hilsner\\
\mbox{[Dit navn]}}

\small\textbf{点评}：恰当地把个人经验与社会舆论结合，流畅使用 \textit{Det diskuteres ofte, om...}、\textit{På den ene side... På den anden side...}、\textit{peger mange eksperter på, at...} 以及 \textit{Ifølge min erfaring giver det...}（V2 倒装零失误），共 128 词，完整覆盖邮件三问并自然回问收尾，完全符合 B1 优秀水平。
\end{answerbox}
\ncdlesson{14}{Når kontanterne forsvinder}{当现金逐渐消失}
{DU3 Modul 3（B1）｜对标 NCE 2 后半--NCE 3 前半 \& Modul 3 \textit{Læsning Opgave 3（段落补句）\& Skrivning Opgave 2}}
{\textbf{趋势变化动词与程度副词}（\textit{stige, falde, vokse, forsvinde} + \dabold{\textit{er}} 助动词；\textit{i stigende grad, stadigt flere/færre, markant, kraftigt}）}
{丹麦是全球无现金化程度最高的国家之一，现金使用率逐年走低，MobilePay 与刷卡早已成为日常标配。本课集中攻克 B1 社会报道与未来 PD3 图表作文最核心的\textbf{趋势变化动词、助动词 \textit{er} vs \textit{har} 区分，以及程度副词的精确搭配}！}

% ---------------------------------------------------------------------
\begin{foerduleaeser}
带着下面两个问题阅读 Tekst A 和 Tekst B：
\begin{enumerate}
  \item \textit{Hvordan har brugen af kontanter og digitale betalinger udviklet sig i Danmark ifølge Tekst A, og hvilke bekymringer skaber udviklingen?}（根据 Tekst A，丹麦的现金使用和数字支付是如何演变的？这一趋势引发了哪些担忧？）
  \item \textit{Hvad har Folketinget besluttet om fysiske butikker og kontanter i Tekst B, og hvorfor?}（在 Tekst B 中，丹麦议会对实体店和现金作出了什么决定？为什么？）
\end{enumerate}
\end{foerduleaeser}

% ---------------------------------------------------------------------
\ncdsubsec{二、Tekst A \& Tekst B（数字社会与 Modul 3 真题场景）}

\subsection*{Tekst A：\textit{Farvel til mønter og sedler?}（社会观察短文 · 约 245 词）}

\begin{tipbox}{语篇说明}
本篇追踪丹麦从现金向数字支付转型的社会变化与潜在隐患。请重点留意文中\textbf{\color{NCDRed}深红色加粗的 \textit{er} 助动词趋势变化（如 \textit{er faldet, er steget}）}与\textbf{\color{NCDNavy}深蓝色加粗的趋势副词与数量词（\textit{i stigende grad, stadigt flere/færre}）}。
\end{tipbox}

For bare tyve år siden havde næsten alle danskere en pung fuld af mønter og sedler i lommen. I dag er billedet et helt andet. Ifølge tal fra Nationalbanken \dabold{er brugen af kontanter faldet} kraftigt i de seneste ti år, og i mange butikker betaler mindre end hver tiende kunde med fysiske penge.

Samtidig \dabold{er antallet af} digitale betalinger \dabold{steget} markant. Dankort, kontaktløs betaling og især MobilePay har ændret danskernes vaner fuldstændigt. \da{Stadigt flere} caféer, frisører og loppemarkeder tager slet ikke imod kontanter længere, fordi det både er hurtigere og billigere at håndtere penge digitalt. Forretningerne sparer tid på at tælle kassen op, og risikoen for tyveri bliver mindre.

Udviklingen ses også i gadebilledet. I mange mindre byer \dabold{falder} antallet af hæveautomater år for år, mens omkostningerne ved at transportere sedler og mønter \da{stiger}. Derfor oplever nogle borgere, at det er blevet besværligt at få fat i kontanter, selv når de faktisk har brug for dem.

Alligevel vækker det kontantløse samfund bekymring hos flere grupper. \da{I stigende grad} \dabold{peger eksperter på}, at børn og unge har sværere ved at forstå værdien af penge, når de kun ser tal på en skærm. Desuden bliver ældre borgere og turister let udelukket, hvis en café kun accepterer MobilePay. Endelig stiller mange sig kritiske over for, hvad der sker ved et stort IT-nedbrud eller en strømafbrydelse. \da{Stadigt færre} danskere bærer kontanter til daglig, men de fleste er enige om, at mønter og sedler næppe \dabold{forsvinder} helt som backup.

\begin{transbox}
\small
\textbf{\color{NCDTeal}[Tekst A 中文精译]}\\
就在二十年前，几乎所有丹麦人的口袋里都揣着装满硬币和纸币的钱包。而如今情况已截然不同。根据丹麦国家银行的数据，过去十年现金的使用量急剧下降，在许多商店里，用实体货币付款的顾客还不到十分之一。

与此同时，数字支付的数量大幅上升。Dankort 银行卡、非接触式支付，尤其是 MobilePay，彻底改变了丹麦人的习惯。越来越多的咖啡馆、理发店和跳蚤市场不再接收现金，因为数字化处理资金既更快又更省钱。商家省下了清点收银机的时间，失窃的风险也随之降低。

这一趋势同样体现在街头景象中。在许多中小城镇，自动取款机的数量逐年减少，而运输纸币和硬币的成本却在不断攀升。因此，一些市民发现，即便真有需要，取到现金也变得麻烦起来。

尽管如此，无现金社会仍在多个群体中引发担忧。专家们越来越频繁地指出，当孩子和年轻人只在屏幕上看到数字时，他们更难理解金钱的实际价值。此外，如果一家咖啡馆只接受 MobilePay，老年市民和游客很容易被排除在外。最后，许多人也对发生大规模 IT 故障或停电时会怎样表示质疑。虽然日常携带现金的丹麦人越来越少，但大多数人一致认为，硬币和纸币作为备用手段几乎不会彻底消失。
\end{transbox}

\subsection*{Tekst B：\textit{Skal vi afskaffe kontanterne helt?}（对标 DU3 Modul 3 Læsning Opgave 3）}

\textbf{阅读任务（Læsning Opgave 3 模拟）：}\textit{Læs de tre korte afsnit, og vælg den sætning (A, B, C eller D), der passer bedst i hvert af de tre huller (1)--(3).}

\vspace{0.3em}
\begin{emailbox}{段落补句训练（107 词）｜Kontantloven i Danmark}
Danske supermarkeder og caféer har i mange år arbejdet på at gøre betalingen ved kassen så hurtig som muligt. \textbf{(1) \underline{\hspace{2.0cm}}} I dag bruger størstedelen af kunderne nemlig betalingskort eller telefon, hvilket sparer meget tid for både medarbejdere og handlende.

Alligevel har Folketinget besluttet, at almindelige fysiske butikker fortsat har pligt til at tage imod kontanter i dagtimerne. \textbf{(2) \underline{\hspace{2.0cm}}} Mange ældre og andre borgere uden smartphone kan ellers let blive udelukket fra at købe mad og medicin i hverdagen.

Også i udlandet følger mange økonomer den hurtige udvikling i Danmark tæt. \textbf{(3) \underline{\hspace{2.0cm}}} Spørgsmålet er dog stadig, om samfundet bliver for sårbart, hvis der pludselig sker et stort IT-nedbrud eller en strømafbrydelse.
\end{emailbox}

\vspace{0.25em}
\textbf{Vælg den rigtige sætning til hvert hul（请为每个空选出逻辑衔接最自然的一句）：}
\begin{enumerate}
  \item \textbf{Hul (1)：}\\
    \textbf{A.} \textit{Derfor er brugen af kontanter faldet kraftigt i hele landet.}\\
    \textbf{B.} \textit{Desværre er priserne på mobiltelefoner steget markant.}\\
    \textbf{C.} \textit{På den anden side foretrækker unge stadig at betale med mønter.}\\
    \textbf{D.} \textit{Derudover må butikkerne ikke længere sælge varer om aftenen.}
  \item \textbf{Hul (2)：}\\
    \textbf{A.} \textit{Det skyldes, at ingen danskere bruger Dankort længere.}\\
    \textbf{B.} \textit{Formålet er at sikre, at alle borgere kan betale for nødvendige varer.}\\
    \textbf{C.} \textit{Derfor skal alle banker lukke deres lokale filialer næste år.}\\
    \textbf{D.} \textit{Samtidig stiger antallet af hæveautomater i de mindste landsbyer.}
  \item \textbf{Hul (3)：}\\
    \textbf{A.} \textit{I stigende grad mener eksperter, at kontanter er den eneste løsning.}\\
    \textbf{B.} \textit{Stadigt færre lande vil nemlig bruge internettet til handel.}\\
    \textbf{C.} \textit{Stadigt flere europæiske lande overvejer at følge samme digitale vej.}\\
    \textbf{D.} \textit{Til gengæld har man aldrig brug for strøm i moderne supermarkeder.}
\end{enumerate}

\begin{transbox}
\small
\textbf{\color{NCDTeal}[Tekst B 中文对照]}\\
丹麦的超市和咖啡馆多年来一直致力于让收银台付款尽可能快捷。\textbf{(1) 填入正确句。} 因为如今绝大多数顾客都使用银行卡或手机，这为店员和消费者都节省了大量时间。

尽管如此，丹麦议会依然决定，普通实体商店在日间仍有义务接收现金。\textbf{(2) 填入正确句。} 否则，许多老年人以及其他没有智能手机的公民很容易在日常生活中被排除在购买食物和药品之外。

在国外，许多经济学家也在密切关注丹麦的快速发展。\textbf{(3) 填入正确句。} 不过问题依然在于，如果突然发生大规模 IT 故障或停电，社会是否会变得过于脆弱。
\end{transbox}

\begin{tipbox}{本课发音与吞音要点（Udtale \& Reduktioner）}
\begin{itemize}
  \item \textbf{趋势核心词组（软 \textit{d} 与不发音 \textit{g/d}）}：\textit{i stigende grad} \ipa{i ˈsdiːˀənə ˈɡʁɑːˀð}（\textit{stigende} 中间的 \textit{g} 完全不发音，\textit{grad} 词尾发软 \textit{d} 加喉塞音）；\textit{faldet} \ipa{ˈfɑləð}（\textit{-ld-} 中的 \textit{d} 不发音）；\textit{forsvinder} \ipa{fʌˈsvenˀɐ}（\textit{-nd-} 中的 \textit{d} 不发音，带喉塞音）。
  \item \textbf{货币与支付高频词}：\textit{kontanter} \ipa{kʰʌnˈtˢanˀdɐ}（重音在第二音节）；\textit{mønter} \ipa{ˈmøndɐ}；\textit{sedler} \ipa{ˈsɛðˀlɐ}（\textit{d} 发软音 \ipa{ð} 且带喉塞音）；\textit{hæveautomat} \ipa{ˈhɛːvəˌɑʊ̯tˢoˌmæːˀd}。
  \item \textbf{特殊读音提示}：\textit{efterspørgsel} \ipa{ˈɛfdɐˌsbɶɐ̯ˀsəl}（\textit{g} 不发音）；\textit{stadigt} \ipa{ˈsdæːði}（口语中词尾 \textit{-gt} 常脱落）；\textit{bekymring} \ipa{beˈkʰømˀʁeŋ}（\textit{y} 在 \textit{m} 前圆唇微降）。
\end{itemize}
\end{tipbox}

% ---------------------------------------------------------------------
\ncdsubsec{三、Gloser（核心生词表）}

\begin{tipbox}{词表重点}
本课三大趋势动词 \textit{at stige}（上升）、\textit{at falde}（下降）、\textit{at forsvinde}（消失）均为不规则强变化，且构成完成时\textbf{全部必须用 \dabold{[er]}}！记住对比：\textit{Priserne \dabold{er steget}}（物价涨了——自动变化，用 \textit{er}）vs \textit{Butikken \da{har øget} priserne}（商店提高了物价——人为动作，用 \textit{har}）。
\end{tipbox}

\begin{glosertable}
名词 & kontanter (flertal) & \textit{(kontanterne, kun flertal)} & 现金（通常只用复数） \\
名词 & en mønt & \textit{(mønten, mønter, mønterne)} & 硬币 \\
名词 & en seddel & \textit{(sedlen, sedler, sedlerne)} & 纸币；便条 \\
名词 & et betalingskort & \textit{(betalingskortet, betalingskort, betalingskortene)} & 银行卡，支付卡 \\
名词 & en hæveautomat & \textit{(hæveautomaten, hæveautomater, hæveautomaterne)} & 自动取款机（ATM） \\
名词 & et IT-nedbrud & \textit{(IT-nedbruddet, IT-nedbrud, IT-nedbruddene)} & IT 系统崩溃/故障 \\
名词 & en strømafbrydelse & \textit{(strømafbrydelsen, strømafbrydelser, strømafbrydelserne)} & 停电，电力中断 \\
名词 & en værdi & \textit{(værdien, værdier, værdierne)} & 价值；价值观 \\
名词 & en udvikling & \textit{(udviklingen, udviklinger, udviklingerne)} & 发展，演变，趋势 \\
名词 & en omkostning & \textit{(omkostningen, omkostninger, omkostningerne)} & 成本，费用 \\
\midrule
动词 \dabold{[er]} & at stige & \textit{(stiger, steg, \dabold{er steget})} & 上升，增长 \\
动词 \dabold{[er]} & at falde & \textit{(falder, faldt, \dabold{er faldet})} & 下降，减少；跌落 \\
动词 \dabold{[er]} & at forsvinde & \textit{(forsvinder, forsvandt, \dabold{er forsvundet})} & 消失 \\
动词 \dabold{[er]} & at vokse & \textit{(vokser, voksede, \dabold{er vokset})} & 增长，成长，扩大 \\
动词 [har] & at øge & \textit{(øger, øgede, har øget)} & （人为）增加，提高 \\
动词 [har] & at håndtere & \textit{(håndterer, håndterede, har håndteret)} & 处理，应对，操作 \\
\midrule
形容词 & markant & \textit{(markant, markante)} & 显著的，明显的（副词同形） \\
形容词 & kraftig & \textit{(kraftigt, kraftige)} & 剧烈的，大幅度的 \\
\end{glosertable}

% ---------------------------------------------------------------------
\ncdsubsec{四、Grammatik \& Faste vendinger（核心语法与地道表达精讲）}

\begin{grammatikbox}{语法点 1：趋势变化动词与完成时助动词铁律（\textit{er} vs \textit{har}）}
描述社会统计变化、物价起伏或习惯演变时，丹麦语严格区分“\textbf{事物自身状态变化}（不及物 $\rightarrow$ \dabold{er}）”与“\textbf{人为施加影响}（及物 $\rightarrow$ \da{har}）”：
\begin{itemize}
  \item \textbf{不及物趋势动词（无宾语，主语自己在变）$\rightarrow$ 完成时必用 \dabold{er}}：\\
    \textit{stige}（上升）$\rightarrow$ \textit{er steget}；\textit{falde}（下降）$\rightarrow$ \textit{er faldet}；\textit{forsvinde}（消失）$\rightarrow$ \textit{er forsvundet}；\textit{vokse}（增长）$\rightarrow$ \textit{er vokset}。\\
    例：\textit{Antallet af brugere \dabold{er steget} kraftigt.} / \textit{Brugen af kontanter \dabold{er faldet} markant.}
  \item \textbf{及物动词（后面带宾语，“把某物提高/降低”）$\rightarrow$ 完成时必用 \da{har}}：\\
    \textit{øge}（增加）、\textit{sænke}（降低）、\textit{mindske}（减少）。\\
    例：\textit{Butikken \da{har øget} brugen af digitale kasser.}
\end{itemize}
\end{grammatikbox}

\vspace{0.25em}
\renewcommand{\arraystretch}{1.24}
{\small
\begin{xltabular}{\linewidth}{@{} >{\bfseries\color{NCDNavy}\raggedright\arraybackslash}p{3.6cm} >{\raggedright\arraybackslash}p{4.1cm} >{\raggedright\arraybackslash}X @{}}
\toprule
\textbf{趋势与程度高频模块} & \textbf{语法功能与搭配规则} & \textbf{地道 B1 句型示范} \\
\midrule
\endfirsthead
\toprule
\textbf{趋势与程度高频模块} & \textbf{语法功能与搭配规则} & \textbf{地道 B1 句型示范（续表）} \\
\midrule
\endhead
\bottomrule
\endfoot
\textit{kraftigt / markant / støt} & 修饰动词“大幅/显著/稳步地” & \textit{Priserne er steget \da{markant / kraftigt} i år.} \\
\textit{langsomt / gradvist} & 修饰动词“缓慢/逐渐地” & \textit{Kontanterne forsvinder \da{gradvist} fra bybilledet.} \\
\textit{i stigende grad} & 副词短语“越来越、日益”（\textbf{放句首必倒装}） & \textit{\da{I stigende grad} \dabold{bruger danskerne} mobilen til at betale.} \\
\textit{stadigt flere / færre} & 接\textbf{可数复数名词}“越来越多/少……” & \textit{\da{Stadigt flere} butikker er lukket om søndagen.} \\
\textit{stadigt mere / mindre} & 接\textbf{不可数名词或动词}“越来越多/少” & \textit{Man bruger \da{stadigt mindre} papir på kontoret.} \\
\end{xltabular}
}

\begin{grammatikbox}{语法点 2：动态变化表达与英丹对比避坑指南}
\textbf{英丹对比避坑指南（English vs Danish Pitfalls）}：
\begin{enumerate}
  \item \textbf{完成时助动词：\textit{have} 陷阱！}：英语说 \textit{Cash usage \textbf{has} dropped} 和 \textit{Prices \textbf{have} risen}，全部用 \textit{have}；但丹麦语中 \textit{falde / stige / forsvinde} 属于状态改变，绝对不能说 \textit{*Brugen har faldet}，必须用 \dabold{er}：\ding{51}\ \textit{Brugen \dabold{er faldet}.} / \textit{Antallet \dabold{er steget}.}
  \item \textbf{\textit{more and more / fewer and fewer} 不能逐字翻！}：英语可以说 \textit{more and more people}，但丹麦语不能照搬成 \textit{*mere og mere mennesker}！正确对应如下：
    \begin{itemize}
      \item 可数复数名词（如 \textit{mennesker, caféer, sedler}）$\rightarrow$ \ding{51}\ \da{stadigt flere / stadigt færre}（或 \textit{flere og flere / færre og færre}）。
      \item 不可数名词/动作程度（如 \textit{tid, kontanter, bekymret}）$\rightarrow$ \ding{51}\ \da{stadigt mere / stadigt mindre} 或 \da{i stigende grad}。
    \end{itemize}
  \item \textbf{形容词变副词时注意 \textit{-t} 与不变形词}：\textit{kraftig} 变副词必须加 \textit{-t}（\textit{falde kraftig\dabold{t}}）；但以 \textit{-t} 结尾的 \textit{markant} 变副词\textbf{不再加第二个 \textit{t}}（\textit{stige markant}，不是 \textit{*markantt}！）。
  \item \textbf{\textit{I stigende grad} 句首倒装}：放在句首时算第 1 位，动词紧跟在后面：\ding{51}\ \textit{I stigende grad \dabold{betaler vi} med telefonen.}（绝不能写成 \textit{*I stigende grad vi betaler...}）。
\end{enumerate}
\end{grammatikbox}

% ---------------------------------------------------------------------
\ncdsubsec{五、Mønsterøvelser（句型操练）}

\textbf{练习 A：用括号内动词的正确形式（注意 \textit{er / har}）或本课趋势词（\textit{flere, færre, mere, mindre, grad, kraftigt, markant}）填空}
\begin{enumerate}
  \item Brugen af MobilePay \underline{\hspace{1.2cm}} \underline{\hspace{1.8cm}} (\textit{stige}, 现在完成时) meget hurtigt i Danmark.
  \item Siden 2015 \underline{\hspace{1.2cm}} antallet af kontantbetalinger \underline{\hspace{1.8cm}} (\textit{falde}, 现在完成时) hvert eneste år.
  \item Stadigt \underline{\hspace{1.6cm}} unge betaler næsten altid med mobilen i kantinen.
  \item I mindre landsbyer findes der desværre stadigt \underline{\hspace{1.6cm}} hæveautomater end tidligere.
  \item I stigende \underline{\hspace{1.4cm}} handler danskerne ind på nettet i stedet for i fysiske butikker.
  \item Man bruger stadigt \underline{\hspace{1.6cm}} tid på at hæve kontanter i banken.
  \item Flere eksperter mener ikke, at sedler og mønter \underline{\hspace{1.2cm}} \underline{\hspace{2.0cm}} (\textit{forsvinde}, 现在完成时) helt endnu.
  \item Kommunen \underline{\hspace{1.2cm}} \underline{\hspace{1.6cm}} (\textit{øge}, 现在完成时) digitale tilbud til ældre borgere.
\end{enumerate}

\textbf{练习 B：句型手术——将下列普通句改写为高分趋势句，或纠正错误}
\begin{enumerate}
  \item 把 \textit{I stigende grad} 移到句首（注意 V2 倒装）：\textit{Danskerne bruger i stigende grad digitale løsninger i hverdagen.} $\rightarrow$ \textit{I stigende grad...}
  \item 改为现在完成时并加程度副词 \textit{markant}：\textit{Antallet af brugere stiger i år.} $\rightarrow$ \textit{Antallet af brugere...}
  \item 合并并用 \textit{mens} 对比两条相反趋势：\textit{Antallet af hæveautomater falder. Brugen af betalingskort vokser.} $\rightarrow$ \textit{Antallet af hæveautomater...}
  \item 纠错改写（助动词错误）：\textit{*Siden sidste sommer har brugen af kontanter faldet kraftig.} $\rightarrow$
  \item 纠错改写（可数与不可数搭配错误）：\textit{*Stadigt mere butikker nægter at tage imod mønter, og vi får stadigt mindre sedler.} $\rightarrow$
  \item 纠错改写（句首副词短语漏了倒装）：\textit{*I stigende grad folk bekymrer sig for IT-nedbrud.} $\rightarrow$
\end{enumerate}

\textbf{练习 C：中译丹（运用趋势动词、\textit{er} 完成时与程度副词）}
\begin{enumerate}
  \item 近年来，现金的使用量大幅下降了。
  \item 越来越少的年轻人会在钱包里放硬币和纸币。
  \item 人们越来越频繁地用手机支付一杯咖啡的钱。
  \item 数字支付的数量在过去五年里显著增长了。
  \item 随着越来越多人用卡付款，街头取款机的数量也在减少。
\end{enumerate}

% ---------------------------------------------------------------------
\ncdsubsec{六、DU3 Modul 3 实战微任务与参考答案}

\begin{emailbox}{Modul 3 Skrivning Opgave 2 实战微任务（回复朋友关于“无现金社会”的邮件）}
\textbf{Situation}: Din ven, Henrik, bor i udlandet og overvejer at rejse til Danmark næste måned. Han skriver til dig for at høre om betalingsvaner i Danmark.\\
\textbf{Opgave}: Skriv et svar til Henrik på \textbf{mindst 100 ord}, hvor du svarer på alle hans spørgsmål og bruger mindst tre trendudtryk fra denne lektion (\textit{er steget/faldet, stadigt flere/færre, i stigende grad}):
\begin{enumerate}
  \item Hvordan betaler man typisk i Danmark i dag, og hvordan har det ændret sig?
  \item Behøver man at have kontanter med på ferie i Danmark?
  \item Hvilke fordele og ulemper ser du selv ved et næsten kontantløst samfund?
\end{enumerate}
\end{emailbox}

\begin{answerbox}
\textbf{Før du læser 参考回答}
\begin{enumerate}
  \item \textit{Brugen af kontanter er faldet kraftigt, mens antallet af digitale betalinger som Dankort og MobilePay er steget markant. Udviklingen vækker bekymring for, om børn forstår penges værdi, om ældre og turister bliver udelukket, og hvad der sker ved IT-nedbrud eller strømafbrydelse.}
  \item \textit{Folketinget har besluttet, at almindelige fysiske butikker fortsat skal tage imod kontanter i dagtimerne, for at ældre og borgere uden smartphone ikke bliver udelukket fra at købe nødvendige varer.}
\end{enumerate}

\textbf{Tekst B 阅读答案（Læsning Opgave 3）}
\begin{enumerate}
  \item \textbf{A} —— 前句说“努力让结账尽可能快”，后句紧接“因为如今多数人用卡或手机”，中间选 \textit{Derfor er brugen af kontanter faldet kraftigt...} 构成完整因果链。
  \item \textbf{B} —— 前句说议会规定实体店白天必须收现金，本句解释“目的是确保所有公民都能购买必需品”（\textit{Formålet er at sikre...}），与后句 \textit{nemlig} 严丝合缝。
  \item \textbf{C} —— 前句说国外经济学家密切关注丹麦趋势，本句顺接“越来越多欧洲国家考虑走同样的数字化道路”，随后用 \textit{dog} 转折引出电网与 IT 故障隐患。
\end{enumerate}

\textbf{练习 A 答案}
\begin{enumerate}
  \item \da{er}；\da{steget}（状态上升，必用 \textit{er}）
  \item \da{er}；\da{faldet}（注意 \textit{Siden 2015} 放句首引发 V2 倒装 \textit{er antallet... faldet}）
  \item \da{flere}（\textit{unge} 复数可数）
  \item \da{færre}（\textit{hæveautomater} 复数可数）
  \item \da{grad}（固定短语：\textit{i stigende grad}）
  \item \da{mindre}（\textit{tid} 不可数）
  \item \da{er}；\da{forsvundet}（\textit{forsvinde} 完成时用 \textit{er}）
  \item \da{har}；\da{øget}（及物动词“增加某物”，用 \textit{har}）
\end{enumerate}

\textbf{练习 B 答案（句型手术参考答案）}
\begin{enumerate}
  \item \textit{I stigende grad \dabold{bruger danskerne} digitale løsninger i hverdagen.}（\textit{bruger} 必须位于主语 \textit{danskerne} 前！）
  \item \textit{Antallet af brugere \dabold{er steget markant} i år.}
  \item \textit{Antallet af hæveautomater falder, \da{mens} brugen af betalingskort vokser.}
  \item \textit{Siden sidste sommer \dabold{er} brugen af kontanter faldet \dabold{kraftigt}.}（助动词用 \textit{er}，副词 \textit{kraftigt} 带 \textit{-t}）
  \item \textit{Stadigt \dabold{flere} butikker nægter at tage imod mønter, og vi får stadigt \dabold{færre} sedler.}（\textit{butikker} 与 \textit{sedler} 均为可数复数！）
  \item \textit{I stigende grad \dabold{bekymrer folk sig} for IT-nedbrud.}（句首副词短语引发倒装：\textit{bekymrer folk sig}）
\end{enumerate}

\textbf{练习 C 答案}
\begin{enumerate}
  \item \textit{I de seneste år er brugen af kontanter faldet kraftigt.}
  \item \textit{Stadigt færre unge har mønter og sedler i pungen.}
  \item \textit{Folk betaler i stigende grad for en kop kaffe med telefonen.}
  \item \textit{Antallet af digitale betalinger er steget markant i de sidste fem år.}
  \item \textit{Mens stadigt flere betaler med kort, falder antallet af hæveautomater på gaden.}
\end{enumerate}

\textbf{Modul 3 实战高分范文（121 词）}\\
\textit{Til: Henrik \quad|\quad Emne: Sv: Betaling og kontanter i Danmark}

\textit{Hej Henrik}

\textit{Hvor dejligt at høre, at du snart kommer til Danmark! Du behøver faktisk næsten ikke at veksle mange kontanter inden turen. I de seneste ti år \dabold{er brugen af kontanter faldet} kraftigt her i landet, og næsten alle danskere betaler med kort eller telefon.}

\textit{Det er dog en god idé at tage et internationalt betalingskort med, fordi \da{stadigt flere} små caféer og boder slet ikke tager imod sedler længere. På den ene side er det meget hurtigt og nemt for os som kunder. På den anden side bekymrer jeg mig \da{i stigende grad} for, hvad man gør, hvis netværket pludselig går ned. Alligevel går hverdagen utrolig glat med kortet!}

\textit{Sig til, når du ved, hvornår du lander!}

\textit{Mange hilsner\\
\mbox{[Dit navn]}}

\small\textbf{点评}：亲切回应朋友关切；自然嵌入 \textit{er brugen af kontanter faldet kraftigt}、\textit{stadigt flere} 和 \textit{i stigende grad}，同时巧妙复用前面单元学到的 \textit{På den ene side... På den anden side...} 和 \textit{bekymrer mig for, hvad...}，全篇 121 词，句法准确丰富，符合 Modul 3 满分回信标准。
\end{answerbox}
\ncdlesson{15}{Alene, sammen eller i regnbuefamilie?}{独居、同居还是多元家庭？}
{DU3 Modul 3（B1）｜对标 NCE 2 后半--NCE 3 前半 \& Modul 3 \textit{Mundtlig Opgave 2（偏好对比互动）\& 社会调查阅读}}
{\textbf{分数、比例与群体主谓/表语搭配}（\textit{hver tredje dansker, halvdelen af befolkningen, størstedelen af, flertallet af, de fleste}）· \textbf{表语形容词单复数一致}}
{丹麦的家庭与居住形态非常多元：独居者、未婚同居、重组家庭与彩虹家庭在日常生活中平等共存。在 Modul 3 口语对比与写作引用数据时，中国与英语学习者最常在\textbf{“\textit{hver tredje} 后名词单复数”}以及\textbf{“比例短语后的形容词词尾”}上出错，本课一网打尽！}

% ---------------------------------------------------------------------
\begin{foerduleaeser}
带着下面两个问题阅读 Tekst A 和 Tekst B：
\begin{enumerate}
  \item \textit{Hvilke forskellige familie- og boformer findes der i Danmark i dag ifølge Tekst A, og hvordan ser befolkningen på regnbuefamilier?}（根据 Tekst A，当今丹麦有哪些不同的家庭与居住形态？民众如何看待彩虹家庭？）
  \item \textit{Hvad foretrækker Lukas og Maja i deres diskussion om boformer i Tekst B, og hvilke argumenter giver de?}（在 Tekst B 关于居住方式的讨论中，卢卡斯和玛雅分别倾向于哪种选择？他们给出了什么理由？）
\end{enumerate}
\end{foerduleaeser}

% ---------------------------------------------------------------------
\ncdsubsec{二、Tekst A \& Tekst B（数字社会与 Modul 3 真题场景）}

\subsection*{Tekst A：\textit{Den moderne danske familie i tal}（社会统计与文化短文 · 约 245 词）}

\begin{tipbox}{语篇说明}
本篇通过丹麦统计局视角呈现当代家庭结构的变化。阅读时请特别关注文中\textbf{\color{NCDNavy}深蓝色加粗的比例与群体表达}，以及\textbf{\color{NCDRed}深红色加粗的表语形容词与代词单复数配合}。
\end{tipbox}

Hvordan ser en typisk familie ud i Danmark i dag? Tidligere tænkte mange på en mor, en far og to børn i et parcelhus, men virkeligheden er blevet langt mere mangfoldig. Ifølge Danmarks Statistik består \da{mere end hver tredje} husstand i landet i dag af én enkelt person. I København bor næsten \da{halvdelen af befolkningen} i lejligheder med kun én voksen, fordi mange unge studerer, arbejder eller venter længere med at flytte sammen med en partner.

Samtidig vælger \da{størstedelen af} de unge par at bo sammen uden at være gift --- det kaldes på dansk et ”papirløst samliv”. Mange får også børn, før de eventuelt holder bryllup, og \da{flertallet af} danskerne \dabold{er enige} om, at kærlighed og tryghed er vigtigere end papir og vielsesring.

Der findes desuden mange sammenbragte familier med ”dine, mine og vores børn”. Omkring \da{en femtedel af} alle danske børn lever en del af deres barndom i en familie med halvsøskende eller stedsøskende. Endelig er regnbuefamilier --- for eksempel familier med to mødre eller to fædre --- blevet en naturlig del af hverdagen i børnehaver og skoler. \da{Over tre fjerdedele af} befolkningen \dabold{er positiv} over for, at forskellige familieformer har de samme juridiske rettigheder.

Selvom flere bor alene end før, betyder det ikke nødvendigvis, at folk er ensomme. \da{Hver fjerde} enlige er meget \dabold{tilfreds} med \dabold{sit} sociale liv og mødes med nære venner flere gange om ugen, og mange ældre vælger at flytte ind i bofællesskaber for at dele både måltider og fællesskab.

\begin{transbox}
\small
\textbf{\color{NCDTeal}[Tekst A 中文精译]}\\
如今丹麦的一个典型家庭是什么样的呢？过去很多人会想到独栋住宅里的一位妈妈、一位爸爸和两个孩子，但现实早已变得更加多元。根据丹麦统计局的数据，如今全国超过三分之一的家庭住户由单身一人组成。在哥本哈根，近一半的人口住在只有一位成年人的公寓里，因为许多年轻人还在读书、工作，或者推迟了与伴侣同居的时间。

与此同时，大多数年轻伴侣选择不结婚而同居——这在丹麦语中被称为“无纸化同居生活”（papirløst samliv）。许多人甚至在举办婚礼之前就生儿育女，而且大多数丹麦人都一致认为，爱与安全感比一纸婚书和结婚戒指更重要。

此外，还有许多由“你的、我的和我们的孩子”组成的重组家庭。大约五分之一的丹麦儿童在童年的一部分时间里，生活在拥有同父异母/同母异父兄弟姐妹或继兄弟姐妹的家庭中。最后，彩虹家庭——例如由两位母亲或两位父亲组成的家庭——也已经成为幼儿园和学校日常生活中的自然组成部分。超过四分之三的民众支持不同家庭形态享有同等的法律权利。

尽管如今独居的人比过去更多，但这并不一定意味着人们感到孤独。每四位独居者中就有一位对自己的社交生活非常满意，每周与亲密朋友聚会数次；许多老年人也选择搬进共居社区（bofællesskab），一起分享三餐与集体陪伴感。
\end{transbox}

\subsection*{Tekst B：\textit{To kursister diskuterer bolig- og familieformer}（对标 DU3 Modul 3 Mundtlig Opgave 2）}

\textbf{口语互动场景：}\textit{Lukas og Maja øver sig til Modul 3 Mundtlig Opgave 2. De sammenligner forskellige boformer og fortæller, hvad de foretrækker.}

\vspace{0.3em}
\begin{emailbox}{双人话题讨论（124 词）｜Lukas \& Maja om at bo alene eller sammen}
\textbf{Lukas:} Hvis jeg skulle vælge mellem de forskellige måder at bo på, ville jeg helt klart foretrække at bo sammen med en partner eller i en lille familie. \da{Størstedelen af} mine venner i Aarhus bor stadig alene i små lejligheder, men efter arbejdet synes jeg, det er rart at have nogen at dele hverdagen med.

\vspace{0.35em}
\textbf{Maja:} Det forstår jeg godt! Men for mig ser det lidt anderledes ud. Jeg har boet alene i tre år, og i dag er \da{flertallet af} mine kolleger faktisk meget \dabold{glade} for deres frihed som enlige. På den anden side lyder et bofællesskab også spændende, fordi man har både privatliv og fælles middage et par gange om ugen. \da{Hver tredje} beboer i vores gade bor i et bofællesskab nu.
\end{emailbox}

\begin{transbox}
\small
\textbf{\color{NCDTeal}[Tekst B 中文对照]}\\
\textbf{卢卡斯}：如果让我在不同的居住方式之间选择，我肯定更喜欢和伴侣或者在一个小家庭里一起生活。我在奥胡斯的大部分朋友现在依然一个人住在小公寓里，但下班后我觉得有人一起分享日常生活很惬意。

\textbf{玛雅}：我很理解！不过对我来说情况有点不同。我已经独居三年了，如今我大多数同事其实都非常享受单身的自由。另一方面，共居社区听起来也很吸引人，因为既能拥有私人空间，每周又能吃几次公共晚餐。现在我们这条街上每三位住户里就有一位住在共居社区里。
\end{transbox}

\begin{tipbox}{本课发音与吞音要点（Udtale \& Reduktioner）}
\begin{itemize}
  \item \textbf{比例与序数词核心发音}：\textit{hver tredje} \ipa{vɛɐ̯ ˈtˢʁɛðjə}（\textit{hver} 词首 \textit{h} 不发音，\textit{tredje} 中的 \textit{d} 读软音 \ipa{ð}）；\textit{halvdelen} \ipa{ˈhalˌdeːˀlən}（\textit{halv} 中的 \textit{v} 不发音，\textit{-delen} 读软 \textit{d}）；\textit{størstedelen} \ipa{ˈsdɶɐ̯sdəˌdeːˀlən}；\textit{flertallet} \ipa{ˈfleɐ̯ˌtˢalˀəð}。
  \item \textbf{家庭与社会词群}：\textit{en befolkning} \ipa{beˈfʌlˀɡneŋ}（重音在第二音节）；\textit{en husstand} \ipa{ˈhuːsˌsdanˀ}（词尾 \textit{nd} 中的 \textit{d} 不发音）；\textit{regnbuefamilie} \ipa{ˈʁɑi̯nˌbuːəfaˌmiːljə}（\textit{egn} 读滑音 \ipa{ɑi̯}）；\textit{rettighed} \ipa{ˈʁadiˌheðˀ}（中间 \textit{g} 不发音，结尾带软 \textit{d} 与喉塞音）。
  \item \textbf{高频率词尾弱化}：\textit{sammenbragt} \ipa{ˈsɑmənˌbʁɑɡd}；\textit{enlige} \ipa{ˈeːnliə}（\textit{-ig-} 中 \textit{g} 不发音）。
\end{itemize}
\end{tipbox}

% ---------------------------------------------------------------------
\ncdsubsec{三、Gloser（核心生词表）}

\begin{tipbox}{词表重点}
注意 \textit{halvsøskende}（同父异母或同母异父兄弟姐妹）与普通 \textit{søskende} 一样，通常只有复数形式；而“搬到一起同居”用 \textit{at flytte sammen}，完成时必须用 \dabold{[er]}（状态和居所发生位移改变）。另外注意 \textit{at foretrække}（更喜欢、偏爱）的不规则过去时 \textit{foretrak / har foretrukket}。
\end{tipbox}

\begin{glosertable}
名词 & en befolkning & \textit{(befolkningen, befolkninger, befolkningerne)} & 人口，民众 \\
名词 & en husstand & \textit{(husstanden, husstande, husstandene)} & 家庭住户，一户人家 \\
名词 & en enlig & \textit{(den enlige, enlige, de enlige)} & 单身人士，独居者 \\
名词 & et samliv & \textit{(samlivet, ingen flertal)} & 共同生活，同居关系 \\
名词 & en regnbuefamilie & \textit{(regnbuefamilien, regnbuefamilier, regnbuefamilierne)} & 彩虹家庭（多元/同性家庭） \\
名词 & halvsøskende (flertal) & \textit{(halvsøskendene, kun flertal)} & 半血缘兄弟姐妹 \\
名词 & et bofællesskab & \textit{(bofællesskabet, bofællesskaber, bofællesskaberne)} & 共居社区，合居公寓 \\
名词 & et flertal & \textit{(flertallet, flertal, flertallene)} & 多数，大多数 \\
名词 & en størstedel & \textit{(størstedelen, ingen flertal)} & 绝大部分，大多数 \\
名词 & en halvdel & \textit{(halvdelen, halvdele, halvdelene)} & 一半 \\
名词 & en tredjedel & \textit{(tredjedelen, tredjedele, tredjedelene)} & 三分之一 \\
名词 & en fjerdedel & \textit{(fjerdedelen, fjerdedele, fjerdedelene)} & 四分之一 \\
名词 & en rettighed & \textit{(rettigheden, rettigheder, rettighederne)} & 权利 \\
\midrule
动词 [har] & at bestå af & \textit{(består, bestod, har bestået af)} & 由……组成 \\
动词 \dabold{[er]} & at flytte sammen & \textit{(flytter, flyttede, \dabold{er flyttet} sammen)} & 搬到一起住，同居 \\
动词 [har] & at foretrække & \textit{(foretrækker, foretrak, har foretrukket)} & 更喜欢，宁愿选择 \\
\midrule
形容词 & mangfoldig & \textit{(mangfoldigt, mangfoldige)} & 多元化的，丰富多样的 \\
形容词 & ensom & \textit{(ensomt, ensomme)} & 孤独的，寂寞的 \\
\end{glosertable}

% ---------------------------------------------------------------------
\ncdsubsec{四、Grammatik \& Faste vendinger（核心语法与地道表达精讲）}

\begin{grammatikbox}{语法点 1：分数、比例与“每第 \textit{N} 个”精准构词表}
在阅读调研报告以及参加 Modul 3 口语与写作任务时，需要精准表达部分与整体的比例关系。丹麦语最核心的三大工具组如下：
\end{grammatikbox}

\vspace{0.25em}
\renewcommand{\arraystretch}{1.24}
{\small
\begin{xltabular}{\linewidth}{@{} >{\bfseries\color{NCDNavy}\raggedright\arraybackslash}p{3.4cm} >{\raggedright\arraybackslash}p{4.3cm} >{\raggedright\arraybackslash}X @{}}
\toprule
\textbf{比例结构类型} & \textbf{构词与核心公式} & \textbf{地道 B1 例句（留意单复数！）} \\
\midrule
\endfirsthead
\toprule
\textbf{比例结构类型} & \textbf{构词与核心公式} & \textbf{地道 B1 例句（续表）} \\
\midrule
\endhead
\bottomrule
\endfoot
\textit{hver / hvert + 序数词} & \textit{hver anden/tredje/fjerde/tiende} + \textbf{单数无冠词名词} & \textit{\da{Hver tredje dansker} bor alene.} / \textit{\da{Hvert fjerde barn} har halvsøskende.} \\
\textit{X ud af Y}（几分之几） & \textit{én/to ud af tre/fire} + \textbf{复数名词} & \textit{\da{To ud af tre danskere} mener, at tryghed er vigtigst.} \\
分数词 + \textit{af} & \textit{halvdelen, en tredjedel, tre fjerdedele, en femtedel} + \textit{af} & \textit{Næsten \da{halvdelen af befolkningen} bor i storbyer.} \\
百分比 + \textit{af} & 数字 + \textit{procent af}（\textit{procent} 永不加 \textit{-s/-er}） & \textit{\da{35 procent af de adspurgte} er enlige.} \\
群体总称词 & \textit{størstedelen af, flertallet af, mindretallet af, de fleste (af)} & \textit{\da{Flertallet af} danskerne / \da{De fleste} unge bor sammen før brylluppet.} \\
\end{xltabular}
}

\begin{grammatikbox}{语法点 2：动词不变形，但“表语形容词与代词”必须对号入座！}
很多同学以为丹麦语动词不区分单复数（如 \textit{er, bor, mener} 永远同形），就可以不管单复数。错了！只要句子里出现\textbf{表语形容词（\textit{glad/glade, enig/enige, tilfreds/tilfredse}）或反身物主代词（\textit{sin/sit/sine} vs \textit{deres}）}，就必须严格分清真实主语是单数还是复数：
\begin{enumerate}
  \item \textbf{\textit{hver tredje + 单数名词} $\rightarrow$ 视作单数！配单数形容词与 \textit{sin/sit/sine}}：\\
    \ding{51}\ \textit{Hver tredje dansker er \dabold{tilfreds} med \dabold{sit} arbejde.}（绝对不能写 \textit{*tilfredse / *deres}！）
  \item \textbf{\textit{halvdelen / størstedelen / flertallet / X procent af + 名词} $\rightarrow$ 单复数看 \textit{af} 后面的核心名词！}：
    \begin{itemize}
      \item \textit{af} 后接\textbf{不可数或集合单数名词}（如 \textit{befolkningen, tiden, landet}）$\rightarrow$ \textbf{单数形容词}：\\
        \ding{51}\ \textit{Halvdelen af \textbf{befolkningen} er \dabold{positiv}.}
      \item \textit{af} 后接\textbf{复数名词}（如 \textit{danskerne, kollegerne, børnene}）$\rightarrow$ \textbf{复数形容词（\textit{+e}）}：\\
        \ding{51}\ \textit{Flertallet af \textbf{danskerne} er \dabold{enige}.} / \textit{Halvdelen af \textbf{mine venner} er \dabold{glade} for at bo alene.}
    \end{itemize}
\end{enumerate}

\vspace{0.3em}
\textbf{英丹对比避坑指南（English vs Danish Pitfalls）}：
\begin{enumerate}
  \item \textbf{英语 \textit{one in three Danes}（复数）vs 丹麦语 \textit{hver tredje dansker}（单数！）}：英语说 \textit{every third Dane} 较少而常说 \textit{one in three Danes}（带复数 \textit{-s}）；丹麦语高频句式 \da{hver tredje} 后面\textbf{绝对只能接单数无冠词名词}：\ding{51}\ \textit{hver tredje \dabold{dansker}}（不能写 \ding{55}\ \textit{*hver tredje danskere}！若想用复数名词，请换用 \ding{51}\ \textit{én ud af tre danskere}）。
  \item \textbf{中性名词必须用 \textit{hvert}}：\textit{et barn, et år, et hus} 属于 \textit{et}-词，前面不能用 \textit{hver}，必须变成 \dabold{hvert}：\ding{51}\ \textit{hvert femte barn}（每五个孩子中的一个）、\textit{hvert andet år}（每两年）。
  \item \textbf{\textit{de fleste} vs \textit{mest}（不要混淆“大多数人”和“最多”）}：英语 \textit{most people} 绝对不能直译成 \textit{*mest folk}！复数概念“大多数……”要用 \da{de fleste}：泛指直接接复数名词不加 \textit{af}（\textit{de fleste danskere}），特指（带定冠词或物主代词）才加 \textit{af}（\textit{de fleste af mine venner}）。
  \item \textbf{\textit{procent} 永远不加复数词尾}：不能受英语 \textit{percentages} 或其它名词复数影响写成 \textit{*20 procenter}，只能写 \ding{51}\ \textit{20 procent af borgerne}。
\end{enumerate}
\end{grammatikbox}

% ---------------------------------------------------------------------
\ncdsubsec{五、Mønsterøvelser（句型操练）}

\textbf{练习 A：填入正确的比例表达（\textit{hver, hvert, halvdelen, størstedelen, procent, fleste}）或形容词单复数形式}
\begin{enumerate}
  \item \underline{\hspace{1.4cm}} fjerde dansker over 65 år bor alene i dag.
  \item \underline{\hspace{1.4cm}} femte barn i klassen vokser op i en sammenbragt familie.
  \item Næsten \underline{\hspace{2.0cm}} (\textit{一半}) af befolkningen i storbyen bor i små lejligheder.
  \item Flertallet af mine venner er meget \underline{\hspace{1.8cm}} (\textit{tilfreds}) med at bo sammen med en partner.
  \item Størstedelen af befolkningen er \underline{\hspace{1.8cm}} (\textit{positiv}) over for forskellige familieformer.
  \item Over 70 \underline{\hspace{1.8cm}} af de unge voksne flytter sammen før ægteskabet.
  \item De \underline{\hspace{1.6cm}} (\textit{大多数}) studerende har desværre ikke råd til en stor lejlighed i centrum.
  \item Hver tredje deltager var \underline{\hspace{1.6cm}} (\textit{enig}) i konklusionen og tog \underline{\hspace{1.2cm}} (\textit{sin / deres}) egen beslutning.
\end{enumerate}

\textbf{练习 B：句型手术——按要求改写句子或纠正下列典型错误}
\begin{enumerate}
  \item 用 \textit{hver tredje} 改写（注意名词单数）：\textit{Én ud af tre lejligheder i gaden har kun én beboer.} $\rightarrow$ \textit{Hver tredje...}
  \item 用 \textit{halvdelen af} 改写：\textit{50 procent af alle nye boliger er bygget til singler.} $\rightarrow$ \textit{Halvdelen af...}
  \item 合并句子并用 \textit{de fleste} 开头：\textit{Næsten alle danske par bor sammen først. Derefter bliver de gift.} $\rightarrow$ \textit{De fleste...}
  \item 纠错改写（名词单复数与形容词错误）：\textit{*Hver tredje danskere er glade for deres boform.} $\rightarrow$
  \item 纠错改写（集合单数 vs 复数词尾错配）：\textit{*Halvdelen af befolkningen er utilfredse med huslejen, mens størstedelen af de unge er enig.} $\rightarrow$
  \item 纠错改写（误用 \textit{mest} 与 \textit{hver} 的性别）：\textit{*Mest familier har en bil, og hver andet barn har en cykel.} $\rightarrow$
\end{enumerate}

\textbf{练习 C：中译丹（灵活运用分数、比例与群体搭配）}
\begin{enumerate}
  \item 如今在哥本哈根，超过三分之一的家庭住户由一个人组成。
  \item 大多数丹麦人认为，安全感比一纸婚书更重要。
  \item 调查中一半的年轻人对自己的居住条件非常满意。
  \item 每五个孩子中就有一个在重组家庭里长大。
  \item 我大多数的同事都更喜欢住在城市中心附近。
\end{enumerate}

% ---------------------------------------------------------------------
\ncdsubsec{六、DU3 Modul 3 实战微任务与参考答案}

\begin{emailbox}{Modul 3 Mundtlig Opgave 2 / Skrivning Opgave 2 实战微任务（对比居住与家庭形式）}
\textbf{Situation}: Til DU3 Modul 3 mundtlig opgave 2 (eller i en mail til din ven, Sara) skal du sammenligne forskellige måder at bo på i Danmark: at bo alene, at bo i en kernefamilie/med en partner, og at bo i et bofællesskab.\\
\textbf{Opgave}: Skriv et sammenhængende svar på \textbf{mindst 100 ord}, hvor du kommer ind på følgende punkter og bruger mindst tre mængde-/procentudtryk fra denne lektion:
\begin{enumerate}
  \item Hvordan bor de fleste mennesker, du kender i Danmark eller i dit hjemland?
  \item Hvilke fordele og ulemper er der ved at bo alene vs. at bo i et bofællesskab?
  \item Hvilken boform vil du helst selv vælge i fremtiden, og hvorfor?
\end{enumerate}
\end{emailbox}

\begin{answerbox}
\textbf{Før du læser 参考回答}
\begin{enumerate}
  \item \textit{I Danmark findes der enlige husstande, par i papirløst samliv, kernefamilier, sammenbragte familier med halvsøskende og regnbuefamilier. Over tre fjerdedele af befolkningen er positiv over for, at forskellige familieformer har de samme juridiske rettigheder.}
  \item \textit{Lukas foretrækker at bo sammen med en partner eller i en lille familie, fordi han synes, det er rart at dele hverdagen efter arbejde. Maja er glad for sin frihed som enlig, men synes også, at et bofællesskab lyder spændende på grund af kombinationen af privatliv og fælles middage.}
\end{enumerate}

\textbf{练习 A 答案}
\begin{enumerate}
  \item \da{Hver}（\textit{dansker} 为共性单数名词）
  \item \da{Hvert}（\textit{et barn} 为中性词，故用 \textit{hvert}）
  \item \da{halvdelen}（\textit{halvdelen af befolkningen}）
  \item \da{tilfredse}（核心名词 \textit{venner} 为复数 $\rightarrow$ \textit{-e}）
  \item \da{positiv}（核心名词 \textit{befolkningen} 为集合单数 $\rightarrow$ 不加 \textit{-e}）
  \item \da{procent}（\textit{procent} 不加复数词尾）
  \item \da{fleste}（\textit{de fleste studerende}）
  \item \da{enig}；\da{sin}（\textit{Hver tredje deltager} 为单数 $\rightarrow$ 单数形容词 \textit{enig} + 反身单数 \textit{sin}）
\end{enumerate}

\textbf{练习 B 答案（句型手术参考答案）}
\begin{enumerate}
  \item \textit{Hver tredje \dabold{lejlighed} i gaden har kun én beboer.}（\textit{lejlighed} 用单数无冠词形式！）
  \item \textit{Halvdelen af alle nye boliger er bygget til singler.}
  \item \textit{De fleste danske par bor sammen, før de bliver gift.}
  \item \textit{Hver tredje \dabold{dansker} er \dabold{glad} for \dabold{sin} boform.}（\textit{dansker}、\textit{glad} 和 \textit{sin} 全部用单数！）
  \item \textit{Halvdelen af befolkningen er \dabold{utilfreds} med huslejen, mens størstedelen af de unge er \dabold{enige}.}（\textit{befolkningen} 单数配 \textit{utilfreds}；\textit{de unge} 复数配 \textit{enige}）
  \item \textit{\dabold{De fleste} familier har en bil, og \dabold{hvert} andet barn har en cykel.}
\end{enumerate}

\textbf{练习 C 答案}
\begin{enumerate}
  \item \textit{I København består mere end hver tredje husstand i dag af én enkelt person.}
  \item \textit{Flertallet af danskerne mener, at tryghed er vigtigere end et stykke papir.}（或：\textit{De fleste danskere mener, at tryghed er vigtigere end papir.}）
  \item \textit{Halvdelen af de unge i undersøgelsen er meget tilfredse med deres boligforhold.}
  \item \textit{Hvert femte barn vokser op i en sammenbragt familie.}
  \item \textit{Størstedelen af mine kolleger foretrækker at bo tæt på byens centrum.}
\end{enumerate}

\textbf{Modul 3 实战高分范文（125 词）}\\
\textit{Til: Sara \quad|\quad Emne: Tanker om forskellige måder at bo på}

\textit{Hej Sara}

\textit{Tak for dit spørgsmål om boligformer! I mit hjemland bor \da{størstedelen af} befolkningen sammen med familien, indtil de bliver gift, men her i Danmark ser jeg noget helt andet. \da{Hver tredje} husstand består jo af én person, og \da{halvdelen af} mine danske studiekammerater \dabold{er meget glade} for at bo alene.}

\textit{På den ene side giver det stor frihed at bo alene, fordi man selv bestemmer over sin tid. På den anden side kan man hurtigt føle sig lidt ensom om vinteren. Et bofællesskab lyder derimod rigtig hyggeligt, fordi man deler fælles middage, men stadig har sit eget værelse.}

\textit{I fremtiden vil jeg helst bo sammen med en partner eller i et lille bofællesskab, da fællesskab i hverdagen betyder meget for mig.}

\textit{Kærlig hilsen\\
\mbox{[Dit navn]}}

\small\textbf{点评}：清晰完成了跨文化对比、利弊分析与个人偏好论证（同时胜任书面 \textit{Skrivning Opgave 2} 与口试 \textit{Opgave 2}）；精准使用 \textit{størstedelen af befolkningen}、\textit{Hver tredje husstand består} 和 \textit{halvdelen af mine danske studiekammerater er meget glade}（单复数词尾 \textit{glade} 准确无误），共 125 词，逻辑连贯自然。
\end{answerbox}

% =====================================================================
% Unit 4
% =====================================================================
\ncdunit{4}{Miljø, Byer og Sundhed}{环境、城市与公共医疗}{Lektion 16--20：复合名词构词 fuge-s/e · 分裂排比句 · 书面复合介词 · 城乡对比结构 · 第三册总汇演}

\ncdlesson{16}{Vindmøller og drømmen om grøn energi}{风车与绿色能源之梦}
{DU3 Modul 3（B1）｜对标 NCE 2 后半 \& Modul 3 \textit{Klima, Energi og Samfund}}
{\textbf{复合名词构词法（\textit{Sammensatte navneord}）}：从右往左拆解长复合词词性与词义 · \textbf{三大连接规则（零连接 / \textit{fuge-s} / \textit{fuge-e}）} · 英丹复合词“分与合”避坑指南}
{丹麦语报刊与社会话题文章中充满长达十几、二十个字母的“火车厢式复合名词”。千万不要被长度吓到！只要掌握“\textbf{最右端决定词性}”和“\textbf{连接字母 \textit{fuge-s / fuge-e}}”规律，你不仅能秒懂长词，还能在 Modul 3 写作中自造地道复合词，避免被英语空格习惯拖累丢分。}

% ---------------------------------------------------------------------
\begin{foerduleaeser}
带着下面两个问题阅读 Tekst A 和 Tekst B：
\begin{enumerate}
  \item \textit{Hvordan har almindelige borgere været med til at skabe opbakning omkring vindkraft i Danmark, og hvilke udfordringer er der stadig?}（普通民众是如何推动丹麦风电支持的？目前仍然存在哪些挑战？）
  \item \textit{Hvordan forandrede øen Samsø sit energisystem fra 1997 og frem til i dag ifølge Tekst B?}（根据 Tekst B，萨姆索岛是如何从 1997 年起改变其能源系统的？）
\end{enumerate}
\end{foerduleaeser}

% ---------------------------------------------------------------------
\ncdsubsec{二、Tekst A \& Tekst B（社会深读与 Modul 3 真题场景）}

\subsection*{Tekst A：\textit{Vindmøller og drømmen om grøn energi}（社会深读 · 约 245 词）}

\begin{tipbox}{语篇说明}
本篇探讨丹麦风电发展、市民风车合作社与绿色能源转型。请重点观察文中\textbf{\color{NCDNavy}深蓝色标注的基础复合词}，以及\textbf{\color{NCDRed}深红色标注的含连接字母 \textit{fuge-s} 复合词}。
\end{tipbox}

Når man rejser gennem Danmark med tog eller bil, ser man næsten altid hvide \da{vindmøller} i landskabet eller langt ude på havet. Danmark har i mange årtier været et \dabold{foregangsland} inden for \da{vindkraft}, og drømmen om en renere \da{elforsyning} begyndte allerede i 1970'erne, da \da{oliekrisen} ramte Europa.

I dag spiller både små lokale \da{vindmøllelaug} og store \da{havvindmølleparker} en afgørende rolle for den grønne omstilling. Mange borgere ejer stadig andele i de lokale møller, hvilket skaber bred opbakning i \da{lokalsamfundet}. Samtidig giver denne nye \da{energisektor} tusindvis af grønne \dabold{arbejdspladser} til ingeniører, teknikere og projektledere.

Den teknologiske udvikling går hurtigt. Nye, enorme møller kan producere strøm nok til hele byer, og danske \dabold{forskningsprojekter} arbejder tæt sammen med universiteterne om at gøre energien endnu mere \da{bæredygtig}. Når vinden blæser kraftigt i Nordsøen, dækker \da{vindstrømmen} ofte hele landets forbrug, og overskydende strøm sendes videre til \da{nabolandene} gennem et fælles \da{elnet}. Desuden bruges den grønne strøm i stigende grad til at producere \da{fjernvarme} og drive nye \da{elbiler}, så både boliger og daglig transport kan fungere uden kul, olie og gas.

Alligevel er vejen til de nationale \da{klimamål} ikke uden udfordringer. Hvad gør man på vindstille vinterdage, hvor efterspørgslen efter varme og lys er størst? Derfor investerer Danmark massivt i nye \dabold{lagringsmuligheder}, varmepumper og grøn brint. Hvis denne \da{energiomstilling} lykkes gennem tætte \dabold{samarbejdsaftaler} med andre lande, kan Danmark fortsat sætte \dabold{verdensrekord} i \da{miljøvenlig} energi og inspirere næste generation til nye \dabold{fremtidsdrømme}.

\begin{transbox}
\small
\textbf{\color{NCDTeal}[Tekst A 中文精译]}\\
当人们乘火车或汽车穿越丹麦时，几乎总能看到田野里或远海上的白色风力涡轮机。几十年来，丹麦一直是风电领域的先行国家，而追求更清洁电力供应的梦想早在 20 世纪 70 年代石油危机席卷欧洲时就已经开始。

如今，无论是小规模的本地风电合作社，还是大型海上风电场，都在绿色转型中发挥着决定性作用。许多市民至今仍持有本地风车的股份，这在地方社区中营造了广泛的支持。与此同时，这个新兴的能源行业为工程师、技术人员和项目经理创造了数以千计的绿色工作岗位。

技术发展日新月异。新型巨型风机能够生产足以供应整座城市的电力，丹麦的科研项目也与高校紧密合作，致力于让能源变得更加可持续。当北海风力强劲时，风电往往能满足全国的用电需求，多余的电力还会通过共享电网输送给邻国。此外，绿色电力正越来越多地被用于生产区域供暖和驱动新型电动汽车，从而让住宅和日常出行在不使用煤炭、石油和天然气的情况下也能正常运转。

尽管如此，通往国家气候目标的道路并非毫无挑战。在风平浪静、而供暖与照明需求又最大的冬日里该怎么办呢？为此，丹麦正大力投资新型储能方案、热泵以及绿氢。如果这一能源转型能够通过与其他国家的紧密合作协议取得成功，丹麦就能继续在环保能源领域创造世界纪录，并激励下一代怀抱新的未来梦想。
\end{transbox}

\subsection*{Tekst B：\textit{Energiøen Samsø}（对标 DU3 Modul 3 \textit{Læsning Opgave 3}：段落补句）}

\textbf{阅读任务（Opgave 3 模拟 · 约 130 词）：}\textit{Læs de tre afsnit om øen Samsø. I hvert afsnit mangler der én sætning. Vælg den sætning (A--D), der passer bedst ind i sammenhængen.}

\begin{emailbox}{Læsning Opgave 3｜Energiøen Samsø（萨姆索绿色能源岛）}
\textbf{1.}\ I 1997 vandt den lille danske ø Samsø en national konkurrence om at blive selvforsynende med vedvarende energi på kun ti år. Dengang brugte øboerne mest olie og kul til varme og transport. \textbf{(1) \underline{\hspace{2.8cm}}} For mange mente, at projektet både var for dyrt og alt for ambitiøst for et lille øsamfund.
\begin{itemize}[label={}]
  \item \textbf{A.} Derfor var alle øens borgere begejstrede fra første dag.
  \item \textbf{B.} I begyndelsen var flere lokale borgere dog ret skeptiske.
  \item \textbf{C.} Desuden producerede øen allerede al sin strøm fra havvindmøller.
  \item \textbf{D.} Men konkurrencen blev aflyst året efter på grund af manglende støtte.
\end{itemize}

\textbf{2.}\ Med hjælp fra lokale landmænd, håndværkere og energirådgivere lykkedes planen alligevel hurtigere end forventet. Øens borgere investerede sammen i både landvindmøller og store solvarmeanlæg til landsbyerne. \textbf{(2) \underline{\hspace{2.8cm}}} Samtidig skabte omstillingen nye lokale arbejdspladser og styrkede sammenholdet på øen.
\begin{itemize}[label={}]
  \item \textbf{A.} Det betød, at øboerne selv fik glæde af indtægterne fra den grønne strøm.
  \item \textbf{B.} Derfor måtte mange unge familier desværre flytte væk fra Samsø.
  \item \textbf{C.} Alligevel nægtede landmændene at deltage i de fælles møder.
  \item \textbf{D.} Til gengæld steg øens udledning af CO2 meget hurtigt.
\end{itemize}

\textbf{3.}\ I dag producerer Samsø mere grøn energi, end øen selv forbruger, og overskuddet sendes til fastlandet. Hvert år besøger tusindvis af forskere, politikere og studerende fra hele verden øens Energiakademi. \textbf{(3) \underline{\hspace{2.8cm}}} Erfaringerne fra Samsø viser nemlig, at lokale fællesskaber kan skabe store forandringer.
\begin{itemize}[label={}]
  \item \textbf{A.} Men akademiet er desværre lukket for udenlandske gæster om sommeren.
  \item \textbf{B.} De kommer for at lære, hvordan en grøn omstilling kan gennemføres i praksis.
  \item \textbf{C.} Ingen ved dog endnu, om vindmøller overhovedet virker på små øer.
  \item \textbf{D.} Derfor har Samsø besluttet at vende tilbage til olie og kul næste år.
\end{itemize}
\end{emailbox}

\begin{tipbox}{本课发音与长复合词重音要点（Udtale \& Sammensatte ord）}
\begin{itemize}
  \item \textbf{复合词重音铁律（主重音在第一成分）}：不管复合词多长，\textbf{主重音永远落在第一个词根上}，后面的词根降为次重音（\ipa{ˌ}）；且第一成分原有的喉塞音 \textit{stød}（\ipa{ˀ}）往往脱落！例如：\textit{vind} \ipa{ˈvenˀ} $\rightarrow$ \textit{vindmølle} \ipa{ˈvenˌmølə}（\textit{vind} 的 \textit{d} 也不发音！）；\textit{hav} \ipa{ˈhɑʊ̯ˀ} $\rightarrow$ \textit{havvindmøllepark} \ipa{ˈhɑʊ̯venˌmøləˌpʰɑːɡ}。
  \item \textbf{外来重音词干作第一成分}：以 \textit{-i} 结尾的词干保持原有重音，例如：\textit{energiforsyning} \ipa{enɐˈɡiːfʌˌsyːˀneŋ}（重音在 \textit{-gi}）。
  \item \textbf{本课高频词与不发音辅音}：\textit{bæredygtig} \ipa{ˈbɛːɐˌdyɡdi}（词尾 \textit{g} 不发音）；\textit{miljøvenlig} \ipa{milˈjøːˌvɛnli}（\textit{g} 不发音）；\textit{arbejdsplads} \ipa{ˈɑːbɑɪ̯ðsˌpʰlas}（含软 d \ipa{ð} 与 \textit{fuge-s}）；\textit{fremtid} \ipa{ˈfʁɑmˌtˢiðˀ}；\textit{samarbejde} \ipa{ˈsɑmɑˌbɑɪ̯ˀðə}。
\end{itemize}
\end{tipbox}

\begin{transbox}
\small
\textbf{\color{NCDTeal}[Tekst B 中文对照（含正确补全线索）]}\\
\textbf{1.} 1997 年，丹麦小岛萨姆索赢得了一项国家级竞赛，目标是在短短十年内实现可再生能源自给自足。当时岛民主要使用石油和煤炭来供暖和交通。\textbf{[空 1]} 因为许多人认为这个项目对于一个小型海岛社区来说既太昂贵又过于激进。

\textbf{2.} 在当地农民、手工业者和能源顾问的帮助下，这项计划还是比预期更快地取得了成功。岛民们共同投资了陆上风车和用于村庄的大型太阳能供热站。\textbf{[空 2]} 与此同时，这次转型创造了新的本地就业机会，并增强了岛上的凝聚力。

\textbf{3.} 如今，萨姆索生产的绿色能源超过了自身消耗，多余的电力被输送往丹麦本土。每年都有来自世界各地的数千名学者、政治家和学生造访岛上的能源学院。\textbf{[空 3]} 来自萨姆索的经验表明，地方社区同样能够带来巨大的变革。
\end{transbox}

% ---------------------------------------------------------------------
\ncdsubsec{三、Gloser（核心生词表）}

\begin{tipbox}{词表重点}
注意观察复合词的\textbf{词性继承规律}：\textit{en vind} + \textit{en mølle} $\rightarrow$ \textit{en vindmølle}，但 \textit{en forskning}（共性）+ \textit{et projekt}（中性）$\rightarrow$ \dabold{et} \textit{forskningsprojekt}（以右端词 \textit{projekt} 为准）！
\end{tipbox}

\begin{glosertable}
名词 & en vindmølle & \textit{(vindmøllen, vindmøller, vindmøllerne)} & 风力涡轮机，风车 \\
名词 & en havvindmøllepark & \textit{(havvindmølleparken, -parker, -erne)} & 海上风电场 \\
名词 & en omstilling & \textit{(omstillingen, omstillinger, -erne)} & 转型，转变（\textit{grøn omstilling} 绿色转型） \\
名词 & en energikilde & \textit{(energikilden, energikilder, -rne)} & 能源（来源） \\
名词 & en elforsyning & \textit{(elforsyningen, elforsyninger, -erne)} & 电力供应 \\
名词 & en arbejdsplads & \textit{(arbejdspladsen, arbejdspladser, -erne)} & 工作岗位；工作场所 \\
名词 & et forskningsprojekt & \textit{(forskningsprojektet, -projekter, -erne)} & 科研项目 \\
名词 & et samarbejde & \textit{(samarbejdet, samarbejder, -ne)} & 合作，协作 \\
名词 & et klimamål & \textit{(klimamålet, klimamål, klimamålene)} & 气候目标 \\
名词 & en udfordring & \textit{(udfordringen, udfordringer, -erne)} & 挑战，难题 \\
名词 & en borger & \textit{(borgeren, borgere, borgerne)} & 市民，公民 \\
\midrule
动词 [har] & at producere & \textit{(producerer, producerede, har produceret)} & 生产，制造 \\
动词 [har] & at udlede & \textit{(udleder, udledte, har udledt)} & 排放（气体/污染物） \\
动词 [har] & at investere (i) & \textit{(investerer, investerede, har investeret)} & 投资（于……） \\
动词 [har] & at bidrage (til) & \textit{(bidrager, bidrog, har bidraget)} & 贡献，促成 \\
\midrule
形容词 & bæredygtig & \textit{(bæredygtigt, bæredygtige)} & 可持续的 \\
形容词 & miljøvenlig & \textit{(miljøvenligt, miljøvenlige)} & 环保的，环境友好的 \\
\end{glosertable}

% ---------------------------------------------------------------------
\ncdsubsec{四、Grammatik \& Faste vendinger（核心语法与地道表达精讲）}

\begin{grammatikbox}{语法点 1：长复合名词“黄金拆解法”——从右往左定词性、定单复数}
丹麦语的复合词就像把几节车厢焊成一列火车：\textbf{前面的词修饰后面的词，最右端的“火车头”决定整个词的词性（\textit{en/et}）、定指后缀与复数形式！}

\vspace{0.3em}
\renewcommand{\arraystretch}{1.24}
{\small
\begin{tabularx}{\linewidth}{@{} >{\bfseries\raggedright\arraybackslash}p{3.4cm} >{\raggedright\arraybackslash}p{4.8cm} >{\raggedright\arraybackslash}X @{}}
\toprule
\textbf{复合名词} & \textbf{从右往左逐层拆解} & \textbf{词性与屈折判定} \\
\midrule
\textit{et vindmølleprojekt} & \textit{vind}（风）+ \textit{mølle}（磨坊）+ \dabold{projekt}（项目） & 由最右端 \textit{et projekt} 决定 $\rightarrow$ \dabold{et} \\
\textit{en havvindmøllepark} & \textit{hav}（海）+ \textit{vindmølle}（风机）+ \dabold{park}（园区） & 由最右端 \textit{en park} 决定 $\rightarrow$ \dabold{en} \\
\textit{en samarbejdsaftale} & \textit{sam} + \textit{arbejde} + \textit{s} + \dabold{aftale}（协议） & 由最右端 \textit{en aftale} 决定 $\rightarrow$ \dabold{en} \\
\textit{en klimaudfordring} & \textit{klima}（气候）+ \dabold{udfordring}（挑战） & 由最右端 \textit{en udfordring} 决定 $\rightarrow$ \dabold{en} \\
\bottomrule
\end{tabularx}
}
\end{grammatikbox}

\begin{grammatikbox}{语法点 2：三大连接方式（零连接 vs 连接字母 \textit{fuge-s} vs \textit{fuge-e}）}
两个词拼在一起时，中间需不需要加“胶水字母”？掌握下面三组实用规律即可覆盖 90\% 的日常与考场词汇：

\vspace{0.3em}
\renewcommand{\arraystretch}{1.24}
{\small
\begin{tabularx}{\linewidth}{@{} >{\bfseries\raggedright\arraybackslash}p{2.7cm} >{\raggedright\arraybackslash}p{5.5cm} >{\raggedright\arraybackslash}X @{}}
\toprule
\textbf{连接类型} & \textbf{适用条件与高频线索} & \textbf{典型例词（注意红色连接字母）} \\
\midrule
\textbf{1. 零连接}\newline （直接拼接） & 绝大多数双名词基础组合、以元音或 \textit{-el/-er} 结尾的普通词干 & \textit{vind + mølle} $\rightarrow$ \textit{vindmølle}\newline \textit{energi + kilde} $\rightarrow$ \textit{energikilde}\newline \textit{cykel + sti} $\rightarrow$ \textit{cykelsti} \\
\midrule
\textbf{2. 连接字母 \textit{fuge-s}}\newline （最高频考点） &
\ding{172} 第一成分以派生后缀 \textbf{\textit{-ning, -hed, -skab, -else, -tion, -tet}} 结尾；\newline
\ding{173} 第一成分\textbf{本身已经是复合词}（如 \textit{sam+arbejde, fore+gang, frem+tid}）；\newline
\ding{174} 高频单音节习惯词（\textit{arbejds-, livs-, tids-, lands-, verdens-}） &
\textit{forskning\dabold{s}projekt}（科研项目）\newline
\textit{sundhed\dabold{s}væsen}（医疗体系）\newline
\textit{uddannelse\dabold{s}sted}（教育机构）\newline
\textit{samarbejd\dabold{s}aftale}（注意省 \textit{e} 加 \textit{s}）\newline
\textit{arbejd\dabold{s}plads} / \textit{verden\dabold{s}rekord} \\
\midrule
\textbf{3. 连接字母 \textit{fuge-e}}\newline （及动词词干 \textit{-e}） & 部分带传统复数感或单音节词（如 \textit{hund, barn $\rightarrow$ børn, jul, gæst}），以及部分动词“动作用途”复合词 & \textit{hund\dabold{e}hus}（狗窝）\newline \textit{børn\dabold{e}have}（幼儿园）\newline \textit{jul\dabold{e}aften} / \textit{spis\dabold{e}bord} \\
\bottomrule
\end{tabularx}
}
\end{grammatikbox}

\begin{grammatikbox}{英丹对比排雷 \& 地道搭配（Faste vendinger）}
\begin{enumerate}
  \item \textbf{英语空格分开 $\neq$ 丹麦语分开！严禁拆词（\textit{Særskrivning}）}：英语写 \textit{wind turbine, green energy transition, job market}，但丹麦语名词复合\textbf{必须紧密合写}（\textit{vindmølle, energiomstilling, arbejdsmarked}）。拆开写不仅扣分，还常闹笑话：\textit{en dansk lærer}（一位丹麦籍老师）vs \textit{en dansklærer}（一位丹麦语老师）；\textit{en rød vin}（一杯红色的酒）vs \textit{en rødvin}（红葡萄酒）！
  \item \textbf{何时合法使用连字符（Bindestreg）}：当复合词包含缩写、数字或极其冗长影响阅读时，必须用连字符：\textit{CO2-udledning}（二氧化碳排放）、\textit{IT-arbejdsplads}、\textit{2030-klimamål}。
  \item \textbf{形容词修饰复合词不合写}：遇到“绿色转型”时，形容词 \textit{grøn} 不合写，只修饰整个复合词（\textit{den grønne energiomstilling}）。
\end{enumerate}

\textbf{本课高频地道搭配（Faste vendinger）：}
\begin{itemize}
  \item \da{den grønne omstilling}（绿色转型）/ \da{vedvarende energi}（可再生能源）
  \item \da{at spille en afgørende rolle for...}（对……发挥决定性作用）
  \item \da{at være selvforsynende med...}（在……方面实现自给自足）
  \item \da{at sætte rekord i...}（在……方面创下纪录）
\end{itemize}
\end{grammatikbox}

% ---------------------------------------------------------------------
\ncdsubsec{五、Mønsterøvelser（句型操练）}

\textbf{练习 A：复合词“手术台”——拆解词义、判定词性，并合成新复合词}
\begin{enumerate}
  \item 拆解并判定冠词（\textit{en} 或 \textit{et}）：\underline{\hspace{0.9cm}} \textit{vindmølleprojekt} = \underline{\hspace{1.8cm}} + \underline{\hspace{1.8cm}} + \underline{\hspace{1.8cm}}
  \item 拆解并判定冠词（\textit{en} 或 \textit{et}）：\underline{\hspace{0.9cm}} \textit{fjernvarmesystem} = \underline{\hspace{1.8cm}} + \underline{\hspace{1.8cm}} + \underline{\hspace{1.8cm}}
  \item 合成（判定零连接还是 \textit{fuge-s / fuge-e}）：\textit{en produktion + en pris} $\rightarrow$ \underline{\hspace{3.6cm}}
  \item 合成：\textit{en sol + en celle + et anlæg} $\rightarrow$ \underline{\hspace{3.6cm}}（太阳能电池板设施）
  \item 合成：\textit{en udvikling + et projekt} $\rightarrow$ \underline{\hspace{3.6cm}}（开发项目）
  \item 合成：\textit{et klima + en forandring} $\rightarrow$ \underline{\hspace{3.6cm}}（气候变化）
  \item 合成：\textit{en uddannelse + et system} $\rightarrow$ \underline{\hspace{3.6cm}}（教育体系）
  \item 找出并改正“英语式拆词错误”：\textit{Jeg arbejder på en stor vindmølle fabrik i Aarhus.} $\rightarrow$
\end{enumerate}

\textbf{练习 B：语篇选词填空（从方框中选出 6 个词填入短文，注意每词只用一次，有 2 个干扰词）}

\textit{备选词：vindkraft · arbejdspladser · bæredygtig · udfordringer · investerer · udleder · klimamål · borgere}

\vspace{0.2em}
\begin{emailbox}{Læsning Opgave 1 热身｜Danmarks grønne fremtid}
I Danmark er (1) \underline{\hspace{2.4cm}} blevet en helt central del af hverdagen. Fordi vindmøller ikke (2) \underline{\hspace{2.4cm}} CO2, når de drejer rundt, hjælper de landet med at nå sine ambitiøse (3) \underline{\hspace{2.4cm}}. Samtidig skaber den grønne sektor mange nye (4) \underline{\hspace{2.4cm}} i både store byer og mindre kystområder. Derfor (5) \underline{\hspace{2.4cm}} både staten og private virksomheder stadig flere penge i (6) \underline{\hspace{2.4cm}} teknologi og grøn strøm.
\end{emailbox}

\textbf{练习 C：中译丹（重点演练复合名词合写与 \textit{fuge-s} 连接）}
\begin{enumerate}
  \item 绿色转型在丹麦的能源政策（\textit{energipolitik}）中发挥着决定性作用。
  \item 许多科研项目致力于减少二氧化碳排放（\textit{CO2-udledning}）。
  \item 海上风电场能够为整座城市生产足够的可持续电力。
  \item 萨姆索岛的市民们共同投资了一个大型太阳能供热设施（\textit{solvarmeanlæg}）。
  \item 尽管存在许多挑战，可再生能源仍在创造新的工作岗位。
\end{enumerate}

% ---------------------------------------------------------------------
\ncdsubsec{六、DU3 Modul 3 实战微任务与参考答案}

\begin{emailbox}{Mundtlig Opgave 1 (Emneopgave) 备考微写作}
\textbf{Situation}: Til DU3 Modul 3 mundtlig prøve skal du holde et oplæg på 1--2 minutter ud fra et mindmap med stikord.\\
\textbf{Opgave}: Dit emne er \textit{”Grøn energi og miljø i min hverdag”}. Skriv først \textbf{5 stikord} (brug sammensatte ord), og skriv derefter et sammenhængende oplæg på \textbf{100--120 ord}, hvor du forklarer, hvordan man oplever grøn energi i Danmark, og hvad du selv gør for at leve mere miljøvenligt.
\end{emailbox}

\begin{answerbox}
\textbf{Før du læser 参考回答}
\begin{enumerate}
  \item \textit{Mange almindelige borgere ejer stadig andele i lokale vindmøllelaug, hvilket skaber bred opbakning. De største udfordringer er vindstille vinterdage, hvor efterspørgslen efter varme og lys er størst, og derfor investerer Danmark i lagringsmuligheder, varmepumper og grøn brint.}
  \item \textit{Fra 1997 blev Samsø selvforsynende med vedvarende energi ved at investere sammen i landvindmøller og solvarmeanlæg. I dag producerer øen mere grøn energi, end den selv bruger.}
\end{enumerate}

\textbf{Tekst B 阅读理解（Opgave 3）答案}
\begin{enumerate}
  \item \textbf{B} —— 后句 \textit{For mange mente, at projektet både var for dyrt...} 解释前文为什么“一开始当地居民颇为怀疑”。
  \item \textbf{A} —— 前句说岛民共同投资风车与太阳能，A 句 \textit{Det betød, at øboerne selv fik glæde af indtægterne...} 顺承说明共同投资带来的直接收益。
  \item \textbf{B} —— 前句说世界各地研究者和学生来访能源学院，B 句 \textit{De kommer for at lære, hvordan...} 精准衔接代词 \textit{De}。
\end{enumerate}

\textbf{练习 A 答案}
\begin{enumerate}
  \item \da{et} vindmølleprojekt = \textit{vind} + \textit{mølle} + \textit{projekt}（最右端 \textit{et projekt} 决定中性）
  \item \da{et} fjernvarmesystem = \textit{fjern} + \textit{varme} + \textit{system}（最右端 \textit{et system} 决定中性）
  \item \da{en produktionspris}（\textit{-tion} 结尾必加 \textit{fuge-s}）
  \item \da{et solcelleanlæg}（零连接；\textit{et anlæg} 决定中性）
  \item \da{et udviklingsprojekt}（\textit{-ing} 结尾加 \textit{fuge-s}）
  \item \da{en klimaforandring}（零连接）
  \item \da{et uddannelsessystem}（\textit{-else} 结尾加 \textit{fuge-s}）
  \item \textit{Jeg arbejder på en stor \dabold{vindmøllefabrik} i Aarhus.}（绝对不可分开写）
\end{enumerate}

\textbf{练习 B 答案}
(1) \da{vindkraft} \quad (2) \da{udleder} \quad (3) \da{klimamål} \quad (4) \da{arbejdspladser} \quad (5) \da{investerer} \quad (6) \da{bæredygtig}

\textbf{练习 C 答案}
\begin{enumerate}
  \item \textit{Den grønne omstilling spiller en afgørende rolle i dansk energipolitik.}
  \item \textit{Mange forskningsprojekter arbejder på at mindske CO2-udledningen.}
  \item \textit{Havvindmølleparker kan producere nok bæredygtig strøm til hele byer.}
  \item \textit{Borgerne på Samsø investerede sammen i et stort solvarmeanlæg.}
  \item \textit{På trods af mange udfordringer skaber vedvarende energi stadig nye arbejdspladser.}
\end{enumerate}

\textbf{Modul 3 Emneopgave 备考示范讲稿（115 词）}\\
\textbf{Stikord (5 个关键词)：} \textit{vindmølleparker · elforsyning · affaldssortering · elcykel · fremtidsdrømme}\\[3pt]
\textit{”Jeg har valgt at tale om grøn energi og miljø i min hverdag. Når man bor i Danmark, ser man hurtigt, hvor stor en rolle \da{vindkraft} og \da{havvindmølleparker} spiller for landets \da{elforsyning}. Det inspirerer også mig til at tænke over mine egne vaner. For det første tager jeg næsten altid min \da{elcykel} til arbejde i stedet for bil, så jeg sparer energi og undgår CO2-udledning. For det andet bruger jeg meget tid på \dabold{affaldssortering} i min opgang og prøver at købe færre nye ting. Derudover sparer vi på varmen i lejligheden om vinteren. Selvom én person ikke kan løse alle \da{klimaudfordringer} alene, tror jeg på, at mange små valg kan skabe en mere \da{bæredygtig} hverdag.”}

\small\textbf{考官点评}：讲稿自然嵌入了 \textit{vindmølleparker, elforsyning, elcykel, affaldssortering, klimaudfordringer} 五个地道复合词，结构采用 \textit{For det første... For det andet... Derudover...} 逻辑清晰，长度恰好符合 1.5 分钟陈述节奏。
\end{answerbox}

\ncdlesson{17}{Cykelstier i København: Mere end bare transport}{哥本哈根自行车道：不止于交通}
{DU3 Modul 3（B1）｜对标 NCE 2 后半 \& Modul 3 \textit{By, Transport og Hverdagskultur}}
{\textbf{分裂句进阶（\textit{Det er først og fremmest..., der/at...}）} · \textbf{排比递进句（\textit{Ikke alene..., men også...} 句首 V2 倒装）} · DU3 投诉信（\textit{Klage}）四要点模板}
{NCD 2 L22 已经学过基础分裂句，而在 Modul 3 书面与口语论证中，考官看重你能不能“先肯定一点，再层层递进拔高”。牢记：\da{Ikke alene} 放在句首时，变位动词必须立刻紧跟第 2 位（\dabold{Ikke alene + V + S}），千万别受英语语序误导！}

% ---------------------------------------------------------------------
\begin{foerduleaeser}
带着下面两个问题阅读 Tekst A 和 Tekst B：
\begin{enumerate}
  \item \textit{Hvad gør det trygt for både børn og ældre at cykle i København, og hvad er ”den grønne bølge”?}（是什么让哥本哈根的老人和孩子都敢放心骑车？什么是“绿波带”？）
  \item \textit{Hvorfor klager borgeren til Københavns Kommune i Tekst B, og hvad beder hun kommunen om at gøre?}（在 Tekst B 中，市民为什么向哥本哈根市政府投诉？她希望市政府采取什么措施？）
\end{enumerate}
\end{foerduleaeser}

% ---------------------------------------------------------------------
\ncdsubsec{二、Tekst A \& Tekst B（社会深读与 Modul 3 真题场景）}

\subsection*{Tekst A：\textit{Cykelstier i København: Mere end bare transport}（社会深读 · 约 245 词）}

\begin{tipbox}{语篇说明}
哥本哈根拥有世界闻名的骑行基础设施。请重点观察文中\textbf{\color{NCDNavy}深蓝色标注的进阶分裂句}与\textbf{\color{NCDRed}深红色标注的 \textit{Ikke alene / Ikke kun..., men også...} 排比句}，留意句首倒装与从句词序的配合。
\end{tipbox}

Hver morgen strømmer tusindvis af cyklister gennem Københavns gader. I mange storbyer verden over bliver cyklen mest brugt til fritid og sport, men i den danske hovedstad er den et helt naturligt hverdagsredskab. \dabold{Ikke alene sparer man} tid og penge i myldretiden, \dabold{men man får også} frisk luft og daglig motion på vej til arbejde eller studie.

\da{Det er først og fremmest} de brede, adskilte cykelstier, \da{der} gør det trygt for både skolebørn, pendlere og ældre at bevæge sig rundt i byen. Cykelstierne ligger nemlig hævet over kørebanen, så man slipper for at køre lige ved siden af biler og tunge lastbiler. Desuden bliver stierne ryddet for sne om vinteren, ofte før bilvejene. \da{Det er faktisk ikke} det danske vejr, \da{der} afgør, hvor mange der cykler --- det er derimod trygheden og den smarte byplanlægning.

Forbindelsen mellem lyskryds og cyklister spiller også en stor rolle. \da{Det er især} om morgenen langs travle hovedgader som Nørrebrogade, \da{at} man oplever den berømte ”grønne bølge”: hvis man holder en fast fart på tyve kilometer i timen, rammer man grønt lys hele vejen ind mod centrum. Samtidig forbinder nye cykelbroer over havnen bydelene hurtigt, så turen over vandet nu kun tager få minutter.

For mange familier har den praktiske ladcykel endda erstattet bil nummer to. \dabold{Ikke kun skaber} cykelkulturen mindre støj og renere byluft, \dabold{men den giver også} en særlig følelse af frihed og lighed i hverdagen, hvor både direktører og studerende sidder på cyklen side om side.

\begin{transbox}
\small
\textbf{\color{NCDTeal}[Tekst A 中文精译]}\\
每天早上，成千上万的骑行者穿行在哥本哈根的街道上。在世界上许多大城市，自行车主要用于休闲和运动，但在丹麦首都，它是一种完全自然的日常交通工具。在早晚高峰期，人们不仅节省了时间和金钱，而且还在上班或上学的路上获得了新鲜空气和每日锻炼。

首先正是宽阔且与机动车道隔离的自行车道，让学童、通勤者和老年人都能安心地在城市里出行。这是因为自行车道高于行车道，所以骑车人不必紧挨着小汽车和重型卡车骑行。此外，冬天道路积雪时，自行车道往往比机动车道更早被清扫。其实决定有多少人骑车的并不是丹麦的天气，而是安全感与巧妙的城市规划。

信号灯与骑行者之间的衔接也发挥着重要作用。尤其是在早晨沿着诺勒布罗大街等繁忙主干道骑行时，人们会体验到著名的“绿波带”：只要保持每小时二十公里的稳定速度，就能一路绿灯直达市中心。与此同时，横跨海港的新建自行车桥迅速连接起各个城区，使跨水通勤现在只需短短几分钟。

对许多家庭而言，实用的货运自行车（三轮亲子自行车）甚至取代了家里的第二辆汽车。骑行文化不仅营造了更少的噪音和更洁净的城市空气，而且还赋予日常生活一种独特的自由感与平等感——在这里，公司高管和大学生并肩骑在车座上。
\end{transbox}

\subsection*{Tekst B：\textit{Klage til kommunen over en cykelsti}（对标 DU3 Modul 3 \textit{Skrivning Opgave 1: Klage}）}

\textbf{写作题目（Opgave 1 规范）：}\textit{Du cykler hver dag på en cykelsti i din bydel, men der er store problemer med den. Skriv en klage til Københavns Kommune. Du skal fortælle: \ding{172} hvor cykelstien ligger; \ding{173} hvorfor du vil klage; \ding{174} hvordan det påvirker din hverdag; \ding{175} hvad du håber, kommunen vil gøre.}

\begin{emailbox}{高分投诉信范本（126 词）｜Til: Teknik- og Miljøforvaltningen, Københavns Kommune}
\textbf{Kære Københavns Kommune}

Jeg skriver til jer, fordi jeg hver dag pendler på cykelstien langs Amagerbrogade mellem nr. 40 og 90, og her har der desværre været store problemer i flere uger.

\da{Det er først og fremmest} de mange dybe huller i asfalten og manglende gadebelysning, \da{der} gør strækningen meget farlig i øjeblikket. Derudover står der ofte store varevogne parkeret delvist inde på cykelstien om morgenen. \dabold{Ikke alene må} vi cyklister bremse pludseligt i mørket, \dabold{men vi bliver også} tvunget ud på kørebanen mellem travle busser og biler. Sidste tirsdag faldt en anden cyklist næsten på grund af et hul i vejen.

Jeg håber derfor, at I hurtigt vil reparere asfalten, tjekke lyset og sætte tydelige skilte op mod ulovlig parkering.

\textbf{Med venlig hilsen}\\
\textbf{Sofie Holm}
\end{emailbox}

\begin{tipbox}{本课发音与语流重音要点（Udtale \& Rytme）}
\begin{itemize}
  \item \textbf{排比句与分裂句的停顿重音}：读 \textit{Ikke alene sparer man tid...} 时，重音落在 \textit{a\textbf{le}ne} 与动词 \textit{\textbf{spa}rer} 上：\ipa{ˈeɡə aˈleːnə ˈsbɑːɐ man ˈtˢiðˀ}；\textit{først og fremmest} 中的 \textit{og} 弱读为 \ipa{ʌ}：\ipa{ˈfœɐ̯sd ʌ ˈfʁɑməsd}。
  \item \textbf{城市交通核心词（不发音辅音与软 d）}：\textit{en cykelsti} \ipa{ˈsyɡəlˌsdiːˀ}（重音在前，词尾带 \textit{stød}）；\textit{København} \ipa{kʰøbənˈhɑʊ̯ˀn}（注意重音在后半截 \textit{-havn}！）；\textit{myldretid} \ipa{ˈmylɐˌtˢiðˀ}（\textit{-ld-} 中的 \textit{d} 不发音！）；\textit{ladcykel} \ipa{ˈlæðˌsyɡəl}（软 d）。
  \item \textbf{公共服务词读音}：\textit{offentlig} \ipa{ˈʌfənli}（词尾 \textit{g} 与词中 \textit{t} 均常不发音）；\textit{tryghed} \ipa{ˈtˢʁyɡˌheðˀ}；\textit{gadebelysning} \ipa{ˈɡæːðəbeˌlyːˀsneŋ}。
\end{itemize}
\end{tipbox}

\begin{transbox}
\small
\textbf{\color{NCDTeal}[Tekst B 中文对照]}\\
\textbf{致哥本哈根市政府的投诉信}：尊敬的哥本哈根市政府，我写信给你们，是因为我每天都在阿迈厄布罗大街 40 号到 90 号之间的自行车道上骑行通勤，而这里几周以来遗憾地一直存在严重问题。

目前让这段路变得非常危险的，首先正是沥青路面上众多深坑以及缺失的街道照明。此外，早上还经常有大型厢式货车半占着自行车道停放。我们骑行者不仅不得不在黑暗中急刹车，而且还被迫驶入繁忙的公交车和小汽车车道之间。上周二，另一位骑车人就差点因为路上的一个坑摔倒。

因此我希望你们能尽快修补沥青路面、检修路灯，并设置清晰的禁止违章停车标识。谨致问候，索菲·霍尔姆。
\end{transbox}

% ---------------------------------------------------------------------
\ncdsubsec{三、Gloser（核心生词表）}

\begin{tipbox}{词表重点}
分清 \textit{sikker} 与 \textit{tryg} 是 B1 写作的加分细节：\textit{sikker} 侧重客观“没有事故/危险”（\textit{trafiksikkerhed} 交通安全），而 \textit{tryg} 侧重人的主观“感到安心、有安全感”（形容词复数和定指必须双写 \textit{g}：\dabold{trygge}）。
\end{tipbox}

\begin{glosertable}
名词 & en cykelsti & \textit{(cykelstien, cykelstier, cykelstierne)} & 自行车专用道 \\
名词 & en ladcykel & \textit{(ladcyklen, ladcykler, ladcyklerne)} & 货运/亲子三轮自行车 \\
名词 & en myldretid & \textit{(myldretiden, ingen flertal)} & 上下班高峰期 \\
名词 & en trafik & \textit{(trafikken, ingen flertal)} & 交通，车流 \\
名词 & en infrastruktur & \textit{(infrastrukturen, infrastrukturer, -erne)} & 基础设施 \\
名词 & en sikkerhed & \textit{(sikkerheden, ingen flertal)} & 安全，安全性 \\
名词 & en tryghed & \textit{(trygheden, ingen flertal)} & 安全感，安心感 \\
名词 & et lyskryds & \textit{(lyskrydset, lyskryds, lyskrydsene)} & 红绿灯路口，十字路口 \\
名词 & en pendler & \textit{(pendleren, pendlere, pendlerne)} & 通勤者 \\
名词 & en kommune & \textit{(kommunen, kommuner, kommunerne)} & 市政当局，市政府 \\
\midrule
动词 [har] & at pendle & \textit{(pendler, pendlede, har pendlet)} & 通勤，往返上班/上学 \\
动词 [har] & at prioritere & \textit{(prioriterer, prioriterede, har prioriteret)} & 优先考虑，重点投入 \\
动词 [har] & at skabe & \textit{(skaber, skabte, har skabt)} & 创造，营造 \\
动词 [har] & at forbedre & \textit{(forbedrer, forbedrede, har forbedret)} & 改善，改进 \\
动词 [har] & at undgå & \textit{(undgår, undgik, har undgået)} & 避免，避开 \\
\midrule
形容词 & tryg & \textit{(trygt, trygge)} & 安心的，感到安全的（复数双写 \textit{g}） \\
形容词 & offentlig & \textit{(offentligt, offentlige)} & 公共的，公共交通的 \\
\end{glosertable}

% ---------------------------------------------------------------------
\ncdsubsec{四、Grammatik \& Faste vendinger（核心语法与地道表达精讲）}

\begin{grammatikbox}{语法点 1：分裂句进阶（\textit{Det er først og fremmest..., der/som/at...}）}
在 NCD 2 L22 学过基础型 \textit{Det er X, der...} 后，B1 议论文与投诉信常在焦点前加上\textbf{限定副词}（\textit{først og fremmest} 首先、\textit{især} 尤其、\textit{faktisk} 其实、\textit{netop} 恰好）来精确控制论证重心，或用 \textit{ikke..., men derimod...} 做纠偏对比：

\vspace{0.3em}
\renewcommand{\arraystretch}{1.24}
{\small
\begin{tabularx}{\linewidth}{@{} >{\bfseries\raggedright\arraybackslash}p{3.0cm} >{\raggedright\arraybackslash}p{4.8cm} >{\raggedright\arraybackslash}X @{}}
\toprule
\textbf{进阶分裂句型} & \textbf{关系词选择铁律} & \textbf{地道 B1 例句} \\
\midrule
\textbf{首要聚焦型}\newline \textit{Det er først og fremmest...} & 强调主语用 \da{der}（或 \textit{som}）；强调宾语只用 \dabold{som} & \textit{Det er først og fremmest de brede cykelstier, \da{der} gør byen tryg.} \\
\midrule
\textbf{时间/地点聚焦型}\newline \textit{Det er især i/om/langs...} & 强调时间、地点或方式状语，必须用 \da{(at)}，\textbf{禁填 \textit{der/som}} & \textit{Det er især om morgenen, \da{(at)} man oplever den grønne bølge.} \\
\midrule
\textbf{纠偏反驳型}\newline \textit{Det er ikke A, der..., men derimod B} & 先否定大众误解 A，再用 \textit{men derimod}（相反地，而是）引出核心原因 B & \textit{Det er faktisk ikke vejret, \da{der} afgør det, men derimod planlægningen.} \\
\bottomrule
\end{tabularx}
}
\end{grammatikbox}

\begin{grammatikbox}{语法点 2：排比递进句（\textit{Ikke alene..., men også...} 的双轨语序）}
想把“不仅 A，而且 B”写出高级书面感，可以用 \da{Ikke alene..., men også...}（或 \da{Ikke kun..., men også...}）。根据它在句中的位置，语序规则截然不同：
\begin{enumerate}
  \item \textbf{轨 A｜放在句首（最推荐，强烈突出对比递进）$\rightarrow$ 前半句必须立刻 V2 倒装！}
    \begin{itemize}
      \item 公式：\textbf{\textit{Ikke alene} + \dabold{变位动词} + 主语 + ..., \textit{men} + 主语 + 变位动词 + \da{også} + ...}
      \item 例句：\textit{\dabold{Ikke alene sparer man} tid, \dabold{men man får også} daglig motion.}
      \item 例句：\textit{\dabold{Ikke alene må vi} bremse, \dabold{men vi bliver også} tvunget ud på kørebanen.}
    \end{itemize}
  \item \textbf{轨 B｜放在句中（结构较平缓，共享主语）}：
    \begin{itemize}
      \item 公式：\textbf{主语 + 变位动词 + \textit{ikke alene} + ..., \textit{men også} + ...}
      \item 例句：\textit{Cykelkulturen giver \da{ikke alene} motion, \da{men også} renere luft.}
    \end{itemize}
  \item \textbf{加分变体}：若两半句共用主语且只对动词或名词作简单并列，\textit{men} 后也可省略重复主语，但写完整主谓（\textit{men man får også...}）在 Modul 3 考场中最不易出错！
\end{enumerate}
\end{grammatikbox}

\begin{grammatikbox}{英丹对比排雷 \& 地道搭配（Faste vendinger）}
\begin{enumerate}
  \item \textbf{\textit{men} 后面的 \textit{også} 不要放错位置！}：英语写 \textit{but you also get exercise}（\textit{also} 紧跟在主语后、动词前），但丹麦语 \textit{men} 是并列连词，后面是完整主句，中心状语 \textit{også} 必须放在\textbf{变位动词之后}！
    \begin{itemize}
      \item \ding{55}\ \textit{Ikke alene er det billigt, men det \dabold{også sparer} tid.}（英语语序污染！）
      \item \ding{51}\ \textit{Ikke alene er det billigt, men det \da{sparer også} tid.}
    \end{itemize}
  \item \textbf{\textit{Ikke alene} 放句首绝不能忘记倒装}：\ding{55}\ \textit{Ikke alene det er sundt...} $\rightarrow$ \ding{51}\ \textit{Ikke alene \dabold{er det} sundt...}
  \item \textbf{分裂句中别写英语 \textit{which}}：英语可以说 \textit{It is the bike lanes, which...}，但丹麦语分裂句绝不能用 \textit{hvilket}，只能用 \da{der}（主语）、\da{som}（宾语）或 \da{at}（状语）。
\end{enumerate}

\textbf{本课高频地道搭配（Faste vendinger）：}
\begin{itemize}
  \item \da{i myldretiden}（在高峰期）/ \da{fra A til B}（从 A 地到 B 地）
  \item \da{at holde sig i form}（保持体形/锻炼身体）/ \da{at få daglig motion}（获得每日运动）
  \item \da{at komme hurtigt frem}（快速顺利到达）/ \da{at tage hensyn til hinanden i trafikken}（在交通中互相礼让）
\end{itemize}
\end{grammatikbox}

% ---------------------------------------------------------------------
\ncdsubsec{五、Mønsterøvelser（句型操练）}

\textbf{练习 A：句法“手术台”——按要求升级为排比递进句或进阶分裂句}
\begin{enumerate}
  \item 用 \textit{Ikke alene} 放句首改写：\textit{Cyklen er billig, og den er også god for miljøet.} $\rightarrow$ \textit{Ikke alene...}
  \item 用 \textit{Ikke alene} 放句首改写：\textit{Man sparer penge, og man undgår også lange bilkøer.} $\rightarrow$ \textit{Ikke alene...}
  \item 用 \textit{Det er først og fremmest..., der...} 强调划线成分：\textit{\underline{De adskilte cykelstier} skaber tryghed i byen.} $\rightarrow$
  \item 用 \textit{Det er især..., at...} 强调划线成分：\textit{Der opstår farlige situationer \underline{i myldretiden}.} $\rightarrow$
  \item 用纠偏分裂句改写（不是 A，而是 B）：\textit{Regnen stopper ikke folk, men dårlige stier gør.} $\rightarrow$ \textit{Det er faktisk ikke regnen, der...}
  \item 找出并改正 1 处典型语序错误：\textit{Ikke alene cyklisterne bliver glade, men kommunen også sparer penge på vejene.} $\rightarrow$
\end{enumerate}

\textbf{练习 B：段落句法补全（对标 DU3 Modul 3 \textit{Læsning Opgave 3}）}

\textit{Læs de tre korte afsnit om ”Supercykelstier” i hovedstadsområdet, og vælg den rigtige sætning (A--D) til hvert hul.}

\begin{enumerate}
  \item Rundt om København har man i de senere år bygget særlige ”supercykelstier”, som forbinder forstæderne med centrum. \textbf{(1) \underline{\hspace{2.6cm}}} Derfor vælger stadig flere pendlere nu elcyklen frem for bilen til arbejde.
  \newline\textbf{A.} Ikke alene er ruterne brede og lige, men de har også færre stop ved lyskryds.
  \newline\textbf{B.} Ikke alene stierne er smalle, men bilerne kører også alt for hurtigt der.
  \newline\textbf{C.} Det er især om natten, der ingen mennesker bruger cykelstierne.
  \newline\textbf{D.} Men desværre er det forbudt at bruge elcykel på de nye stier.
  \item På nogle strækninger er der sat små fodhvilere og luftpumper op langs ruten. \textbf{(2) \underline{\hspace{2.6cm}}} Det viser, at kommunerne tænker meget over de små detaljer for cyklisterne.
  \newline\textbf{A.} Alligevel vil pendlerne hellere vente i bilkøen på motorvejen.
  \newline\textbf{B.} Det er netop sådanne praktiske løsninger, der gør turen meget mere behagelig.
  \newline\textbf{C.} Derfor har næsten alle borgere solgt deres gamle cykler i år.
  \newline\textbf{D.} Det er faktisk ikke cyklisterne, at man har bygget pumperne til.
  \item Selvom Danmark ofte har blæst og regn om efteråret, fortsætter mange med at pendle på to hjul året rundt. \textbf{(3) \underline{\hspace{2.6cm}}} Hvis bare tøjet passer til vejret, kommer man hurtigt og frisk frem til kontoret.
  \newline\textbf{A.} For det er først og fremmest en god vane, som man hurtigt vænner sig til.
  \newline\textbf{B.} Men derfor lukker alle supercykelstierne allerede i oktober måned.
  \newline\textbf{C.} Ikke kun regner det altid, men også sner det hver uge om efteråret.
  \newline\textbf{D.} Det er først i maj, der man kan tage cyklen frem igen.
\end{enumerate}

\textbf{练习 C：中译丹（重点演练 \textit{Ikke alene..., men også...} 与进阶分裂句）}
\begin{enumerate}
  \item 骑车不仅有利于健康，而且还能减少城市噪音。（用 \textit{Ikke alene} 放句首）
  \item 首先正是安全的自行车道，让孩子们能够自己骑车去上学。
  \item 尤其是在冬天的早晨，人们需要更好的街道照明。
  \item 通勤者不仅节省了时间，而且还在路上获得了新鲜空气。
  \item 决定大家是否骑车的其实不是天气，而是城市的基础设施。
\end{enumerate}

% ---------------------------------------------------------------------
\ncdsubsec{六、DU3 Modul 3 实战微任务与参考答案}

\begin{emailbox}{Skrivning Opgave 1 (Klage) 实战模拟}
\textbf{Situation}: Du cykler ofte gennem en bestemt tunnel eller langs en cykelsti i din kommune, hvor belysningen har været ødelagt i lang tid, og hvor der ofte står affaldsspande eller parkerede cykler i vejen.\\
\textbf{Opgave}: Skriv en klage til kommunen. Du skal fortælle:
\begin{itemize}
  \item hvor problemet er
  \item hvad der er galt, og hvor længe det har stået på
  \item hvordan det påvirker din og andres sikkerhed
  \item hvad du håber, kommunen vil gøre.
\end{itemize}
\textit{（要求：使用至少一组 \textbf{Ikke alene..., men også...} 句首倒装句和一组 \textbf{Det er først og fremmest..., der...} 分裂句，建议写出 105--130 词。）}
\end{emailbox}

\begin{answerbox}
\textbf{Før du læser 参考回答}
\begin{enumerate}
  \item \textit{Det er først og fremmest de brede, adskilte cykelstier, der er hævet over kørebanen, og som bliver ryddet for sne om vinteren. ”Den grønne bølge” betyder, at hvis man holder en fast fart på 20 km/t om morgenen, rammer man grønt lys hele vejen ind mod centrum.}
  \item \textit{Hun klager over dybe huller i asfalten, manglende gadebelysning og ulovligt parkerede varevogne på Amagerbrogade. Hun beder kommunen om at reparere asfalten, tjekke lyset og sætte skilte op mod ulovlig parkering.}
\end{enumerate}

\textbf{练习 A 答案}
\begin{enumerate}
  \item \textit{\dabold{Ikke alene er} cyklen billig, \dabold{men den er også} god for miljøet.}
  \item \textit{\dabold{Ikke alene sparer} man penge, \dabold{men man undgår også} lange bilkøer.}
  \item \textit{\da{Det er først og fremmest} de adskilte cykelstier, \da{der} skaber tryghed i byen.}
  \item \textit{\da{Det er især} i myldretiden, \da{(at)} der opstår farlige situationer.}
  \item \textit{Det er faktisk ikke regnen, \da{der} stopper folk, men derimod dårlige stier.}
  \item 错误两处（缺倒装 + \textit{også} 错位）$\rightarrow$ \textit{Ikke alene \dabold{bliver cyklisterne} glade, men kommunen \dabold{sparer også} penge på vejene.}
\end{enumerate}

\textbf{练习 B 答案}
\begin{enumerate}
  \item \textbf{A}（语序唯一正确：\textit{Ikke alene er ruterne..., men de har også...}；B 项未倒装且语义不符；C 项时间状语误用 \textit{der}。）
  \item \textbf{B}（\textit{sådanne praktiske løsninger} 回指前句的“脚踏板和打气筒”，且分裂句使用正确；D 项状语/介词结构误用 \textit{at} 且语义荒谬。）
  \item \textbf{A}（\textit{god vane} 解释为什么风雨无阻；C 项逻辑荒谬，D 项时间状语误用 \textit{der}。）
\end{enumerate}

\textbf{练习 C 答案}
\begin{enumerate}
  \item \textit{Ikke alene er det sundt at cykle, men det mindsker også støjen i byen.}
  \item \textit{Det er først og fremmest de sikre cykelstier, der gør, at børn selv kan cykle i skole.}
  \item \textit{Det er især om vintermorgenen, (at) man har brug for bedre gadebelysning.}
  \item \textit{Ikke alene sparer pendlerne tid, men de får også frisk luft på vejen.}
  \item \textit{Det er faktisk ikke vejret, der afgør, om folk cykler, men derimod byens infrastruktur.}
\end{enumerate}

\textbf{Modul 3 Klage 高分范文（117 词）}\\
\textit{Til: Teknik- og Miljøforvaltningen \quad|\quad Emne: Manglende lys og spærret cykelsti ved Nørrebro Station}

\textit{Kære kommune}

\textit{Jeg skriver til jer for at klage over forholdene på cykelstien og i tunnelen lige ved Nørrebro Station mod Frederikssundsvej. I næsten tre uger har lyset i tunnelen desværre ikke virket, og samtidig står der hver dag overfyldte affaldsspande og ødelagte cykler midt på stien.}

\textit{\da{Det er først og fremmest} kombinationen af mørke og forhindringer, \da{der} skaber meget farlige situationer i myldretiden. \dabold{Ikke alene kan vi} cyklister dårligt se hinanden om morgenen og aftenen, \dabold{men vi bliver også} nødt til at køre op på fortovet for at komme forbi.}

\textit{Jeg håber meget, at I snarest vil få repareret lamperne og fjernet de gamle cykler, så alle kan færdes trygt igen.}

\textit{Med venlig hilsen\\
\mbox{[Dit navn]}}

\small\textbf{考官点评}：全信精准覆盖投诉四要点，严格使用复数机构称呼 \textit{jer / I}；主体段巧妙将 \textit{Det er først og fremmest..., der...} 与 \textit{Ikke alene kan vi..., men vi bliver også...} 并置，句式严谨、语气客观恳切，稳居 Modul 3 书面高分段。
\end{answerbox}

\ncdlesson{18}{Det danske sundhedsvæsen under lup}{放大镜下的丹麦医疗体系}
{DU3 Modul 3（B1）｜对标 NCE 3 前半 \& Modul 3 \textit{Sundhed, Velfærd og Offentlig Service}}
{\textbf{书面复合介词短语}（\textit{i forhold til, i forbindelse med, med hensyn til, på baggrund af}）· 介词后接名词 vs 不定式 vs 从句的“搭桥规则” · DU3 好评推荐信（\textit{Anbefaling}）写法}
{在丹麦就医、读政府公告或写 Modul 3 社会说明文时，这四个“介词 + 名词 + 介词”组成的四词块出现频率极高。它们不仅能帮你把零散小句浓缩成紧凑的 B1 书面句，放在句首时还能自然触发 V2 倒装，让考官一眼看出你的书面语体水平！}

% ---------------------------------------------------------------------
\begin{foerduleaeser}
带着下面两个问题阅读 Tekst A 和 Tekst B：
\begin{enumerate}
  \item \textit{Hvilken rolle spiller den praktiserende læge i det danske sundhedsvæsen, og hvad diskuteres i øjeblikket?}（在丹麦医疗体系中，全科家庭医生扮演什么角色？目前大家正在讨论什么问题？）
  \item \textit{Hvorfor vil skribenten anbefale Lægehuset i Midtbyen i Tekst B, og hvordan oplevede hun ventetiden?}（在 Tekst B 中，作者为什么推荐市中心的全科诊所？她对候诊时间有怎样的体验？）
\end{enumerate}
\end{foerduleaeser}

% ---------------------------------------------------------------------
\ncdsubsec{二、Tekst A \& Tekst B（社会深读与 Modul 3 真题场景）}

\subsection*{Tekst A：\textit{Det danske sundhedsvæsen under lup}（社会深读 · 约 250 词）}

\begin{tipbox}{语篇说明}
本篇以客观视角剖析丹麦公共医疗系统的运转逻辑与现实争议。请重点观察文中\textbf{\color{NCDNavy}深蓝色标注的四大复合介词短语}以及\textbf{\color{NCDRed}深红色标注的从句/不定式衔接方式}。
\end{tipbox}

Det danske sundhedsvæsen er finansieret over skatten, hvilket betyder, at alle borgere med adresse i landet har fri og lige adgang til lægehjælp. \da{I forbindelse med} flytning til Danmark modtager man et gult sundhedskort, hvor navnet på ens egen praktiserende læge står skrevet, og kortet skal altid vises frem ved hvert fysisk besøg eller gennem den digitale sundhedskort-app på telefonen.

Den praktiserende læge fungerer som den første indgang til systemet. Hvis man bliver syg, kontakter man først sin egen læge eller lægevagten uden for åbningstiden. \da{På baggrund af} en indledende samtale og eventuelle prøver vurderer lægen, om man kan behandles i klinikken, eller om der er behov for en henvisning til en speciallæge eller et hospital.

\da{Med hensyn til} digitalisering er Danmark kommet meget langt \da{i forhold til} mange andre europæiske lande. Gennem appen ”Min Læge” og hjemmesiden \textit{sundhed.dk} kan borgerne både bestille tid, forny recepter og læse deres egne prøvesvar. Dette giver tryghed \dabold{i forbindelse med, at} man skal følge et længere behandlingsforløb.

Alligevel er sundhedsvæsenet ofte under lup i den offentlige debat. Selvom kvaliteten af selve behandlingen generelt er høj, oplever mange patienter lange ventetider, især \dabold{i forhold til at} komme til hos en speciallæge. \dabold{På baggrund af, hvad} flere eksperter og sygeplejersker peger på, handler udfordringen ikke kun om penge, men også om mangel på personale på sygehusene. \da{Med hensyn til} fremtiden kræver det derfor nye løsninger, så velfærdsstaten fortsat kan sikre god hjælp til alle.

\begin{transbox}
\small
\textbf{\color{NCDTeal}[Tekst A 中文精译]}\\
丹麦的医疗体系由税收资助，这意味着所有在丹麦登记住址的公民都能平等、免费地获得医疗救助。在搬到丹麦办理落户之际，每个人都会收到一张黄色的健康卡（医保卡），上面印着自己专属全科家庭医生的名字，而且每次线下就诊时都必须出示该卡或使用手机上的电子医保卡应用。

全科家庭医生充当着进入医疗系统的第一道入口。如果生病，人们首先要联系自己的家庭医生，或在非营业时间联系急诊值班医生。基于初步问诊和相关检查，医生会评估病人是否可以在诊所内治疗，或者是否需要开具转诊单前往专科医生处或医院就诊。

在数字化方面，与许多其他欧洲国家相比，丹麦已经走得很远。通过“我的医生（Min Læge）”应用和 \textit{sundhed.dk} 网站，市民既可以预约挂号、续开处方，也能查看自己的化验结果。这在需要跟进长期治疗过程时给人们带来了安心感。

尽管如此，医疗体系依然常常在公共舆论中被置于放大镜下审视。虽然治疗本身的质量总体很高，但许多患者都会遇到漫长的候诊时间，尤其是相对于及时见到专科医生而言。根据多位专家和护士指出的情况，这项挑战不仅关乎资金，也关乎医院医护人员的短缺。因此，着眼于未来，这需要新的解决方案，才能让福利国家继续为所有人保障优质医疗。
\end{transbox}

\subsection*{Tekst B：\textit{Anbefaling af Lægehuset i Midtbyen}（对标 DU3 Modul 3 \textit{Skrivning Opgave 1: Anbefaling}）}

\textbf{写作题目（Opgave 1 规范 · 约 125 词）：}\textit{Du har lige skiftet til et nyt lægehus i din by, og du er meget tilfreds. Skriv en anbefaling til en lokal onlinegruppe for nye tilflyttere. Du skal fortælle: \ding{172} hvor lægehuset ligger, og hvornår du brugte det; \ding{173} hvad du oplevede i forbindelse med dit besøg; \ding{174} hvordan stedet er i forhold til tidligere oplevelser; \ding{175} hvem du vil anbefale lægehuset til.}

\begin{emailbox}{好评推荐信范本（133 词）｜Forum: Nye i Odense --- Anbefaling af lægehus}
\textbf{Hej alle sammen}

Jeg vil gerne anbefale ”Lægehuset i Midtbyen”, som ligger lige ved Kongensgade i Odense. Jeg skiftede hertil for to måneder siden, og i sidste uge havde jeg mit første besøg \da{i forbindelse med} en kraftig hoste og feber.

Allerede om morgenen fik jeg en akuttid samme dag. Lægen var utrolig venlig og tog sig god tid til at lytte, og \da{på baggrund af} en grundig undersøgelse fik jeg hurtigt den rigtige medicin. \da{Med hensyn til} kommunikation forklarede hun alt i et klart og roligt dansk.

\da{I forhold til} mit gamle lægehus var ventetiden i telefonen meget kortere, og hele personalet virkede slet ikke stresset. Jeg vil derfor varmt anbefale klinikken til alle --- især internationale studerende og børnefamilier, der sætter pris på tryghed og god service.

\textbf{Venlig hilsen}\\
\textbf{Linnea Chen}
\end{emailbox}

\begin{tipbox}{本课发音与复合介词节奏要点（Udtale \& Reduktioner）}
\begin{itemize}
  \item \textbf{复合介词的“两头轻、中间重”韵律}：前面和后面的小介词（\textit{i, på, med, til, af}）极度轻读，主重音全部扛在\textbf{中间名词}上！注意小介词中的大量不发音辅音与软化：
    \begin{itemize}
      \item \textit{i \textbf{for}hold til} \ipa{i ˈfɒːˌhʌlˀ tˢe}（\textit{forhold} 词尾 d 不发音且带 \textit{stød}；\textit{til} 弱读 \ipa{tˢe}）
      \item \textit{i for\textbf{bin}delse med} \ipa{i fʌˈbenˀəlsə mɛ}（\textit{-nd-} 中的 d 不发音；\textit{med} 词尾 d 常弱化脱落读 \ipa{mɛ}）
      \item \textit{med \textbf{hen}syn til} \ipa{mɛ ˈhɛnˌsyːˀn tˢe}
      \item \textit{på \textbf{bag}grund af} \ipa{pʰʌ ˈbɑʊ̯ˌɡʁɔnˀ ɑ}（\textit{bag} 读 \ipa{bɑʊ̯}，\textit{grund} 词尾 d 不发音，\textit{af} 弱读为 \ipa{ɑ}！）
    \end{itemize}
  \item \textbf{医疗词汇不发音辅音（-d / -g）}：\textit{sundhedsvæsen} \ipa{ˈsonˌheðsˌvɛːˀsən}（\textit{sund} 中的 d 不发音）；\textit{praktiserende læge} \ipa{pʰʁɑɡtˢiˈseːˀɐnə ˈlɛːə}（\textit{læge} 中的 g 不发音）；\textit{undersøgelse} \ipa{ˈʌnɐˌsøːˀəlsə}；\textit{henvisning} \ipa{ˈhɛnˌviːˀsneŋ}。
\end{itemize}
\end{tipbox}

\begin{transbox}
\small
\textbf{\color{NCDTeal}[Tekst B 中文对照]}\\
\textbf{市中心全科诊所推荐贴（欧登塞新居民群）}：大家好！我想推荐位于欧登塞国王街旁的“市中心全科诊所（Lægehuset i Midtbyen）”。我两个月前刚转来这家诊所，上周因为剧烈咳嗽和发烧进行了第一次就诊。

一大早我就约到了当天的急诊号。医生非常友善，耐心地听我讲述病情；基于一次细致的检查，我很快就拿到了对症药物。在沟通方面，她用清晰而平缓的丹麦语解释了一切。

与我以前的诊所相比，这里的电话等待时间短得多，全体医护人员也完全不显得匆忙焦虑。因此我热烈向大家推荐这家诊所——尤其是重视安心感与优质服务的国际学生和有孩子的家庭。诚挚的问候，林妮亚·陈。
\end{transbox}

% ---------------------------------------------------------------------
\ncdsubsec{三、Gloser（核心生词表）}

\begin{tipbox}{词表重点}
\textit{at henvise}（转诊）和 \textit{at undersøge}（检查）是丹麦就医最核心的两个及物动词。注意 \textit{gratis}（免费的）无论修饰什么名词或复数都\textbf{永远不变词尾}（\textit{et gratis besøg, gratis timer}）。
\end{tipbox}

\begin{glosertable}
名词 & et sundhedsvæsen & \textit{(sundhedsvæsenet, ingen flertal)} & 医疗卫生体系 \\
名词 & et sundhedskort & \textit{(sundhedskortet, sundhedskort, -ene)} & 医疗保险卡（黄卡） \\
名词短语 & en praktiserende læge & \textit{(den praktiserende læge, -læger, -ne)} & 全科家庭医生 \\
名词 & en henvisning & \textit{(henvisningen, henvisninger, -erne)} & 转诊单，转介 \\
名词 & en speciallæge & \textit{(speciallægen, speciallæger, -ne)} & 专科医生 \\
名词 & en ventetid & \textit{(ventetiden, ventetider, -erne)} & 等待时间，候诊时间 \\
名词 & en behandling & \textit{(behandlingen, behandlinger, -erne)} & 治疗，处理 \\
名词 & en undersøgelse & \textit{(undersøgelsen, undersøgelser, -erne)} & 检查，体检；调研 \\
名词 & et personale & \textit{(personalet, ingen flertal)} & 全体员工，医护人员 \\
名词 & en skat & \textit{(skatten, skatter, skatterne)} & 税，税收 \\
名词 & en recept & \textit{(recepten, recepter, recepterne)} & 处方，药方 \\
\midrule
动词 [har] & at henvise (til) & \textit{(henviser, henviste, har henvist)} & 转诊至；参考，指引 \\
动词 [har] & at vurdere & \textit{(vurderer, vurderede, har vurderet)} & 评估，判断 \\
动词 [har] & at undersøge & \textit{(undersøger, undersøgte, har undersøgt)} & 检查，调查 \\
动词 [har] & at forny & \textit{(fornyer, fornyede, har fornyet)} & 续签，续开（处方/卡） \\
\midrule
形容词 & grundig & \textit{(grundigt, grundige)} & 细致的，彻底的 \\
形容词 & gratis & \textit{(不变)} & 免费的 \\
\end{glosertable}

% ---------------------------------------------------------------------
\ncdsubsec{四、Grammatik \& Faste vendinger（核心语法与地道表达精讲）}

\begin{grammatikbox}{语法点 1：四大高频书面复合介词短语精确分工}
这四个“介词 + 名词 + 介词”组合是 B1 说明文与官方书信的灵魂胶水，分别负责对比、伴随场景、话题限定与依据推论：

\vspace{0.3em}
\renewcommand{\arraystretch}{1.24}
{\small
\begin{tabularx}{\linewidth}{@{} >{\bfseries\raggedright\arraybackslash}p{3.0cm} >{\raggedright\arraybackslash}p{4.3cm} >{\raggedright\arraybackslash}X @{}}
\toprule
\textbf{复合介词短语} & \textbf{核心语义与功能} & \textbf{典型书面例句} \\
\midrule
\da{i forhold til} & 与……相比；关于，针对 & \textit{\da{I forhold til} sidste år er ventetiden nu kortere.} \\
\midrule
\da{i forbindelse med} & 在……之际/过程中；与……有关 & \textit{\da{I forbindelse med} flytningen får man et nyt sundhedskort.} \\
\midrule
\da{med hensyn til} & 至于，在……方面（用于切换话题维度） & \textit{\da{Med hensyn til} digitalisering er Danmark kommet langt.} \\
\midrule
\da{på baggrund af} & 基于，根据（基于某事实/调查做评估） & \textit{\da{På baggrund af} prøverne vurderer lægen næste skridt.} \\
\bottomrule
\end{tabularx}
}
\end{grammatikbox}

\begin{grammatikbox}{语法点 2：复合介词后面接“名词”、“不定式”还是“完整从句”？}
把这四个短语用对只是第一步，能否在后面丝滑接上\textbf{动词动作或整句从句}，决定了你能不能拿 B1+ 高分：
\begin{enumerate}
  \item \textbf{接名词/名词短语（最稳妥、最地道）}：尽量将动作转化为抽象名词（\textit{flytning, undersøgelse, digitalisering, besøget}）：
    \begin{itemize}
      \item \textit{\da{I forbindelse med} \dabold{ansøgningen} skal du vedlægge dit pas.}（在申请之际……）
    \end{itemize}
  \item \textbf{接 \textit{at} + 动词不定式}：当动作主体就是主句主语时，可直接用 \textit{at + 原形}（千万不能像英语那样直接写分词！）：
    \begin{itemize}
      \item \textit{Der er lang ventetid \da{i forhold til at} få en tid hos en speciallæge.}
      \item \textit{\da{Med hensyn til at} bestille tid bruger de fleste appen.}
    \end{itemize}
  \item \textbf{接完整从句（复现 L11 介词从句机制）}：
    \begin{itemize}
      \item 接 \textbf{\textit{hv-} 疑问词从句}（\textit{hvad, hvordan, hvorvidt}）直接连接：\textit{\da{På baggrund af, hvad} eksperterne siger...} / \textit{\da{Med hensyn til, hvordan} systemet virker...}
      \item \textit{i forbindelse med} 与 \textit{på baggrund af} 接 \textbf{\textit{at} 从句}时可直接加逗号连接：\textit{Det giver tryghed \da{i forbindelse med, at} man skal behandles.}（注意从句内部不倒装，\textit{ikke} 放动词前！）
    \end{itemize}
  \item \textbf{放在句首必须触发 V2 倒装！}：\textit{\da{På baggrund af samtalen} \dabold{fik hun} en henvisning.}（介词短语整体占第 1 位，变位动词 \textit{fik} 紧跟第 2 位！）
\end{enumerate}
\end{grammatikbox}

\begin{grammatikbox}{英丹对比排雷 \& 地道搭配（Faste vendinger）}
\begin{enumerate}
  \item \textbf{介词后面严禁加 \textit{-ende}（英语 \textit{-ing} 动名词陷阱）}：英语说 \textit{in connection with moving / regarding booking an appointment}；丹麦语中 \textit{-ende} 只能作形容词或副词，\textbf{绝对不能作介词宾语}！
    \begin{itemize}
      \item \ding{55}\ \textit{I forbindelse med \dabold{flyttende} til Danmark...}（致命英式错误！）
      \item \ding{51}\ \textit{I forbindelse med \da{flytning} til Danmark...} / \textit{I forbindelse med, \da{at man flytter}...}
    \end{itemize}
  \item \textbf{分清 \textit{på baggrund af} 与 \textit{på grund af}}：\textit{på baggrund af} 表示“\textbf{基于（证据/经验/数据）而作出判断或决定}”；而 \textit{på grund af} 表示“\textbf{因为/由于（某原因导致某结果）}”（如 \textit{på grund af sygdom} 因为生病），二者不可随意互换！
  \item \textbf{段落转场别只会用 \textit{Om...}}：谈到“至于价格/交通/服务方面”，英语学生常直译 \textit{Om prisen...}，地道书面丹麦语应该用 \da{Med hensyn til prisen...}。
\end{enumerate}

\textbf{本课高频地道搭配（Faste vendinger）：}
\begin{itemize}
  \item \da{at bestille en tid hos lægen}（向医生预约就诊）/ \da{at få en henvisning til...}（获得去……的转诊单）
  \item \da{at være tilmeldt en læge}（在某位医生名下注册）/ \da{at forny en recept}（续开处方）
  \item \da{finansieret over skatten}（由税收资助的）/ \da{at sætte under lup}（置于放大镜下审视）
\end{itemize}
\end{grammatikbox}

% ---------------------------------------------------------------------
\ncdsubsec{五、Mønsterøvelser（句型操练）}

\textbf{练习 A：用四大复合介词或常用因果介词选填（\textit{i forhold til, i forbindelse med, med hensyn til, på baggrund af, på grund af}）}
\begin{enumerate}
  \item \underline{\hspace{3.2cm}} en grundig undersøgelse besluttede lægen, at jeg skulle have antibiotika.
  \item \underline{\hspace{3.2cm}} ventetiden i Danmark er folk generelt utilfredse med de lange lister hos speciallæger.
  \item Danmark bruger en stor del af budgettet på hospitaler \underline{\hspace{3.2cm}} mange andre lande.
  \item Alle nye borgere modtager vigtige breve i e-Boks \underline{\hspace{3.2cm}} deres flytning til kommunen.
  \item Lægeklinikken var desværre lukket hele fredagen \underline{\hspace{3.2cm}} sygdom blandt personalet.
  \item \underline{\hspace{3.2cm}}, hvordan man bruger ”Min Læge”-appen, kan man finde en video på nettet.
  \item \underline{\hspace{3.2cm}} hendes erfaring fra tre forskellige sygehuse fik hun hurtigt stillingen.
  \item Det er gratis at gå til egen læge i Danmark, men \underline{\hspace{3.2cm}} tandlægebesøg skal voksne selv betale.
\end{enumerate}

\textbf{练习 B：语体升级手术——将口语或“中式/英式直译”句改写为带复合介词的 B1 书面句（注意 V2 倒装）}
\begin{enumerate}
  \item 口语句：\textit{Når man snakker om prisen, er dansk medicin ikke gratis for voksne.} $\rightarrow$ \textit{Med hensyn til...}
  \item 英语污染改错：\textit{I forbindelse med bestillende en tid på nettet fik jeg en fejl.} $\rightarrow$
  \item 句首倒装合并：\textit{Vi læste rapporten. Derefter ændrede vi vores plan.}（用 \textit{På baggrund af} 合并）$\rightarrow$
  \item 对比改写：\textit{Mit nye lægehus er meget bedre, når man sammenligner med det gamle.} $\rightarrow$ \textit{I forhold til...}
  \item 从句衔接：\textit{I forbindelse med at patienter skal vente længe på operation, bliver de ofte utrygge.}（把从句移到句首，注意主句倒装）$\rightarrow$
\end{enumerate}

\textbf{练习 C：中译丹（灵活运用四大复合介词与医疗核心词汇）}
\begin{enumerate}
  \item 与去年相比，这家医院的候诊时间现在短了很多。
  \item 在搬到哥本哈根之际，我收到了一张新的黄色医保卡。
  \item 基于化验结果（\textit{prøvesvarene}），家庭医生给我开了一张去专科医生的转诊单。
  \item 至于数字医疗服务，丹麦已经比许多国家走得更远。
  \item 由于医护人员短缺，有时很难很快约到号（\textit{at få en tid hurtigt}）。
\end{enumerate}

% ---------------------------------------------------------------------
\ncdsubsec{六、DU3 Modul 3 实战微任务与参考答案}

\begin{emailbox}{Skrivning Opgave 1 (Anbefaling) 实战模拟}
\textbf{Situation}: Du har haft en rigtig god oplevelse på en ny tandlægeklinik i din by, og du vil gerne skrive en anbefaling på en dansk hjemmeside for lokale borgere.\\
\textbf{Opgave}: Skriv en anbefaling. Du skal fortælle:
\begin{itemize}
  \item hvor klinikken ligger, og hvorfor du tog derhen
  \item hvordan du oplevede personalet og selve behandlingen
  \item hvordan klinikken er i forhold til andre steder, du kender (eller priser/ventetid)
  \item hvem du især vil anbefale klinikken til.
\end{itemize}
\textit{（要求：使用至少 \textbf{3 个} 本课书面复合介词短语，建议写出 105--130 词。）}
\end{emailbox}

\begin{answerbox}
\textbf{Før du læser 参考回答}
\begin{enumerate}
  \item \textit{Den praktiserende læge fungerer som den første indgang til systemet og vurderer, om man kan behandles i klinikken eller skal have en henvisning til en speciallæge eller et hospital. I øjeblikket diskuteres især lange ventetider hos speciallæger og mangel på personale på sygehusene.}
  \item \textit{Hun anbefaler lægehuset, fordi hun hurtigt fik en akuttid, lægen var venlig og grundig, og hun kommunikerede klart. Ventetiden i telefonen var desuden meget kortere end i hendes gamle lægehus.}
\end{enumerate}

\textbf{练习 A 答案}
\begin{enumerate}
  \item \da{På baggrund af}（基于检查结果作出医疗判断）
  \item \da{Med hensyn til}（至于候诊时间方面）
  \item \da{i forhold til}（与许多其他国家相比）
  \item \da{i forbindelse med}（在搬迁登记之际）
  \item \da{på grund af}（因为生病导致关门——注意与 \textit{på baggrund af} 区分！）
  \item \da{Med hensyn til}（接 \textit{hv-} 从句：至于如何使用 app）
  \item \da{På baggrund af}（基于她的工作经验获得职位）
  \item \da{med hensyn til}（至于看牙医方面，成年人须自费）
\end{enumerate}

\textbf{练习 B 答案}
\begin{enumerate}
  \item \textit{Med hensyn til prisen er dansk medicin ikke gratis for voksne.}（注意 V2 倒装 \textit{er dansk medicin}）
  \item \textit{I forbindelse med \da{at bestille} en tid på nettet fik jeg en fejl.}（或 \textit{I forbindelse med \da{tidsbestillingen} på nettet...}，禁止用 \textit{-ende}）
  \item \textit{På baggrund af rapporten \dabold{ændrede vi} vores plan.}（句首介词短语触发主句倒装）
  \item \textit{I forhold til mit gamle lægehus \dabold{er det nye} meget bedre.}
  \item \textit{I forbindelse med, at patienter skal vente længe på operation, \dabold{bliver de} ofte utrygge.}
\end{enumerate}

\textbf{练习 C 答案}
\begin{enumerate}
  \item \textit{I forhold til sidste år er ventetiden på dette hospital nu blevet meget kortere.}
  \item \textit{I forbindelse med min flytning til København modtog jeg et nyt gult sundhedskort.}
  \item \textit{På baggrund af prøvesvarene gav min praktiserende læge mig en henvisning til en speciallæge.}
  \item \textit{Med hensyn til digitale sundhedstilbud er Danmark kommet længere end mange andre lande.}
  \item \textit{På grund af mangel på personale er det nogle gange svært at få en tid hurtigt.}
\end{enumerate}

\textbf{Modul 3 Anbefaling 高分范文（110 词）}\\
\textit{Anbefaling af ”Smil \& Tryghed Tandlægeklinik” i Aarhus C}

\textit{Jeg vil meget gerne anbefale ”Smil \& Tryghed Tandlægeklinik”, som ligger tæt på banegården i Aarhus. I sidste måned fik jeg pludselig stærk tandpine og måtte bestille en akuttid \da{i forbindelse med} en betændelse.}

\textit{Selvom jeg normalt er lidt nervøs hos tandlægen, følte jeg mig helt tryg fra første minut. Tandlægen lyttede roligt og lagde en klar plan \da{på baggrund af} nogle hurtige røntgenbilleder. \da{Med hensyn til} prisen fik jeg altid et præcist overslag, før behandlingen begyndte.}

\textit{\da{I forhold til} andre klinikker, jeg har prøvet før, var både ventetiden og atmosfæren her langt mere behagelig. Jeg anbefaler derfor varmt klinikken til alle, der søger en grundig og venlig tandlæge i centrum.}

\small\textbf{考官点评}：自然融入了全部四个目标复合介词短语（\textit{i forbindelse med, på baggrund af, med hensyn til, i forhold til}），句首倒装（\textit{I sidste måned fik jeg..., Med hensyn til prisen fik jeg..., I forhold til andre klinikker... var både...}）零失误，完美覆盖推荐信四要素。
\end{answerbox}

\ncdlesson{19}{Udkantsdanmark eller storbydrømme?}{边缘乡村还是大城市之梦？}
{DU3 Modul 3（B1）｜对标 NCE 3 前半 \& Modul 3 \textit{Bolig, Geografi og Livsvalg}}
{\textbf{城乡对比说明文的段落衔接架构（平行对比法 vs 逐点对比法）} · 高阶对比衔接词与语序矩阵（\textit{hvorimod} 从句 vs \textit{derimod / til gengæld / på den anden side} 倒装）}
{DU3 Modul 3 的阅读理解与写作常常让你在“两种生活模式”之间做对比与权衡。写对比段落最怕“一段城市、下一句突然跳乡村、最后又跳回城市”。学完本课的两种黄金段落架构，你就能用清晰的对比连词写出让考官眼前一亮的 B1 结构化文章！}

% ---------------------------------------------------------------------
\begin{foerduleaeser}
带着下面两个问题阅读 Tekst A 和 Tekst B：
\begin{enumerate}
  \item \textit{Hvilke fordele og ulemper nævnes i Tekst A ved at bo i henholdsvis storbyen og Udkantsdanmark?}（在 Tekst A 中，分别提到了住在大城市和边缘乡村地区的哪些优势与劣势？）
  \item \textit{Hvad overvejer Jonas og Sara i deres chat i Tekst B, og hvilken løsning vælger de til sidst?}（在 Tekst B 的聊天中，乔纳斯和萨拉在纠结什么？他们最后倾向哪种方案？）
\end{enumerate}
\end{foerduleaeser}

% ---------------------------------------------------------------------
\ncdsubsec{二、Tekst A \& Tekst B（社会深读与 Modul 3 真题场景）}

\subsection*{Tekst A：\textit{Udkantsdanmark eller storbydrømme?}（社会深读 · 约 250 词）}

\begin{tipbox}{语篇说明}
本篇采用\textbf{“逐点对比法”（Punkt-for-punkt-metoden）}依次从\ding{172}住房成本、\ding{173}就业通勤、\ding{174}生活节奏三个维度展开城乡对比。请重点观察\textbf{\color{NCDNavy}深蓝色标注的句首对比过渡短语}与\textbf{\color{NCDRed}深红色标注的从句/副词对比连接词}。
\end{tipbox}

Selvom Danmark geografisk er et lille land, er forskellen mellem de store byer og de såkaldte yderområder --- ofte kaldet ”Udkantsdanmark” --- blevet tydeligere i de seneste år. Mange unge drømmer om København eller Aarhus på grund af studier og et aktivt byliv, mens børnefamilier og pensionister oftere kigger mod mindre byer og landdistrikter, hvor hverdagen føles mere overskuelig.

\da{På den ene side} er huslejen og boligpriserne i storbyerne steget så kraftigt, at det kan være svært for almindelige lønmodtagere at købe en større bolig. \da{På den anden side} får man langt mere plads for pengene i provinsen eller på små øer, \dabold{hvorimod} udvalget af små, centrale lejligheder her kan være begrænset. Mange flytter netop ud for at få en privat have til børnene.

Når det gælder arbejde og uddannelse, ligger de store universiteter og internationale virksomheder stadig tæt i storbyerne, \da{mens} man i landdistrikterne \dabold{derimod} ofte må acceptere en længere pendlertid og færre offentlige busser. I de senere år har mere fleksibelt hjemmearbejde dog gjort det lettere for nogle pendlere at bo længere væk fra kontoret, uden at de mister deres karrieremuligheder.

\da{Sammenlignet med} storbyens travle tempo tilbyder Udkantsdanmark frisk luft, tæt natur og et stærkt lokalsamfund, hvor naboerne kender hinanden. \dabold{Til gengæld} savner nogle tilflyttere storbyens brede kulturliv, mange restauranter og internationale miljø. Valget mellem by og land handler derfor sjældent om, hvad der er ”bedst” for alle, men snarere om, hvilken livsfase man befinder sig i.

\begin{transbox}
\small
\textbf{\color{NCDTeal}[Tekst A 中文精译]}\\
尽管丹麦在地理上是个小国，但大城市与所谓外围地区——常被称为“边缘丹麦（Udkantsdanmark）”——之间的差异近年来变得越来越明显。许多年轻人因为求学和丰富活跃的城市生活而向往哥本哈根或奥胡斯，而有孩子的家庭和退休人员则更常把目光投向日常生活显得更从容可控的中小城镇和乡村地区。

一方面，大城市的房租和房价上涨得如此之快，以至于普通工薪阶层很难买到一套较大的住房。另一方面，在地方城镇或小岛上，同样的钱可以换来大得多的居住空间，不过当地市中心小户型的选择相对有限。许多家庭正是为了给孩子们换取一个私人花园才搬出城市。

至于工作和教育，大型高校和跨国企业依然密集分布在大城市，而在乡村地区则往往不得不接受更长的通勤时间和更少的公共交通班次。不过近年来，更灵活的居家办公方式让一些通勤者即使住得离办公室更远，也不会失去职业发展机会。

与大城市忙碌的节奏相比，丹麦乡村地区提供了清新的空气、亲近的自然以及邻里彼此相熟的紧密地方社区。不过相应地，也有一些新搬来的人会怀念大城市丰富的文化生活、众多餐馆和国际化氛围。因此，选择城市还是乡村，很少取决于哪种对所有人来说“最好”，而更多取决于一个人正处于哪个人生阶段。
\end{transbox}

\subsection*{Tekst B：\textit{Flytte ud af storbyen?}（对标 DU3 Modul 3 \textit{Læsning Opgave 2A}：对话气泡补全）}

\textbf{阅读任务（Opgave 2A 模拟）：}\textit{Situation: Jonas skriver til sin kæreste, Sara, som er på arbejde, om deres boligplaner. Læs Jonas' beskeder, og find de svar fra Sara (A--F), der passer til hullerne (1)--(3). Der er tre svar, du ikke skal bruge. Se eksemplet (0).}

\begin{emailbox}{Læsning Opgave 2A｜Chat mellem Jonas og Sara}
\textbf{Jonas (0):} Hej Sara! Har du haft tid til at kigge på det lille hus i Silkeborg, som jeg sendte i går?\\
\textbf{Sara (0):} \textbf{\underline{\ Z\ }} \textit{Ja, det har jeg lige set i frokostpausen. Det ser virkelig hyggeligt ud med den store have!}

\vspace{0.3em}
\textbf{Jonas:} Ikke sandt? Men jeg er stadig lidt nervøs for afstanden. Jeg skal jo pendle næsten en time til Aarhus hver dag.\\
\textbf{Sara (1):} \underline{\hspace{1.2cm}}

\vspace{0.3em}
\textbf{Jonas:} Det har du ret i. Og min chef sagde faktisk i morges, at jeg godt må arbejde hjemme to dage om ugen.\\
\textbf{Sara (2):} \underline{\hspace{1.2cm}}

\vspace{0.3em}
\textbf{Jonas:} Ja, bestemt! På den anden side vil vi jo savne vores venner og caféerne i midtbyen lidt.\\
\textbf{Sara (3):} \underline{\hspace{1.2cm}}

\vspace{0.3em}
\textbf{Jonas:} Super idé! Jeg skriver til ejendomsmægleren nu og bestiller en fremvisning på lørdag.

\vspace{0.4em}
\hrule
\vspace{0.3em}
\small
\textbf{Saras mulige svar（A--F，6 选 3，其中 3 项多余）：}
\begin{itemize}[label={}, leftmargin=0.8em]
  \item \textbf{A.} Til gengæld ligger mine forældres hus kun ti minutter derfra, og vores venner kan jo altid komme på besøg i weekenden!
  \item \textbf{B.} Ej, hvor er det ærgerligt, at du slet ikke kan arbejde hjemmefra i dit nye job.
  \item \textbf{C.} Det er rigtigt, men til gengæld er huslejen jo næsten halvt så dyr som i vores lille lejlighed her.
  \item \textbf{D.} Derfor synes jeg hellere, vi skal købe en dyrere lejlighed lige midt i Aarhus.
  \item \textbf{E.} Ej, hvor godt! Det gør jo en kæmpe forskel for pendlertiden. Skal vi så tage ud at se huset i weekenden?
  \item \textbf{F.} Sammenlignet med Silkeborg er vores have her i Aarhus desværre alt for stor at passe.
\end{itemize}
\end{emailbox}

\begin{tipbox}{本课发音与对比连接词读音要点（Udtale \& Reduktioner）}
\begin{itemize}
  \item \textbf{对比连词与副词的读音及不发音字母}：
    \begin{itemize}
      \item \textit{hvorimod} \ipa{vɒːiˈmoðˀ}（词首 \textit{h} 不发音！重音在最后的 \textit{-mod} 上，软 d + \textit{stød}）
      \item \textit{derimod} \ipa{ˈdɛɐ̯iˌmoðˀ}（重音在前面的 \textit{der-} 上，刚好与 \textit{hvorimod} 相反！）
      \item \textit{til gengæld} \ipa{tˢe ˈɡɛŋˌɡɛlˀ}（\textit{til} 弱读 \ipa{tˢe}；\textit{gæld} 中的 \textit{d} 不发音！）
      \item \textit{sammenlignet} \ipa{ˈsɑmənˌliːˀnəð}（\textit{g} 不发音，词尾 \textit{-et} 读轻弱的 \ipa{əð}）
    \end{itemize}
  \item \textbf{城乡生活高频词}：\textit{Udkantsdanmark} \ipa{ˈuðkʰanˀdsˌdanmɑɡ}（\textit{ud} 读软 d，\textit{kants} 带 \textit{stød}）；\textit{storbydrømme} \ipa{ˈsdoɐ̯byˌdʁœmə}；\textit{et lokalsamfund} \ipa{loˈkʰæːˀlˌsɑmfɔnˀ}（词尾 \textit{d} 不发音）；\textit{en husleje} \ipa{ˈhuːsˌlɑɪ̯ə}。
\end{itemize}
\end{tipbox}

\begin{transbox}
\small
\textbf{\color{NCDTeal}[Tekst B 中文对照]}\\
\textbf{聊天对话（搬离大城市吗？）}：
\textbf{乔纳斯 (0)}：嗨萨拉！你有没有抽空看我昨天发给你的锡尔克堡那栋小房子？\\
\textbf{萨拉 (0/Z)}：看了，我午休时刚看！带着大花园看起来真的好温馨！\\
\textbf{乔纳斯}：是吧？但我还是有点担心距离问题。我每天得往奥胡斯通勤差不多一小时呢。\\
\textbf{萨拉 (1)}：\textbf{[空 1]} \quad \textbf{乔纳斯}：你说得对。而且我老板今天早上还真说了，我每周可以居家办公两天。\\
\textbf{萨拉 (2)}：\textbf{[空 2]} \quad \textbf{乔纳斯}：那是当然！不过另一方面，我们也会有点想念市中心的朋友和咖啡馆。\\
\textbf{萨拉 (3)}：\textbf{[空 3]} \quad \textbf{乔纳斯}：好主意！我现在就给房产中介写信，预约周六去看房。
\end{transbox}

% ---------------------------------------------------------------------
\ncdsubsec{三、Gloser（核心生词表）}

\begin{tipbox}{词表重点}
注意 \textit{at tiltrække}（吸引）的不规则过去时是 \dabold{tiltrak}（不是 \textit{tiltræk}），过去分词是 \dabold{har tiltrukket}。另外，\textit{at trives}（生活舒心/茁壮适应）是一个\textbf{异态动词（Deponent verbum）}——外形永远带 \textit{-s}，但意思完全是主动的（如 \textit{Vi trives godt på landet}）。
\end{tipbox}

\begin{glosertable}
名词 & et landdistrikt & \textit{(landdistriktet, landdistrikter, -erne)} & 乡村地区，农村区域 \\
名词 & en storby & \textit{(storbyen, storbyer, storbyerne)} & 大城市，都市 \\
名词 & et lokalsamfund & \textit{(lokalsamfundet, lokalsamfund, -ene)} & 地方社区，本地社群 \\
名词 & en husleje & \textit{(huslejen, huslejer, huslejerne)} & 房租 \\
名词 & en boligpris & \textit{(boligprisen, boligpriser, -erne)} & 房价 \\
名词 & en fordel & \textit{(fordelen, fordele, fordelene)} & 优点，优势 \\
名词 & en ulempe & \textit{(ulempen, ulemper, ulemperne)} & 缺点，劣势 \\
名词 & et kulturliv & \textit{(kulturlivet, ingen flertal)} & 文化生活 \\
名词 & en pendlertid & \textit{(pendlertiden, pendlertider, -erne)} & 通勤时间 \\
名词 & en livskvalitet & \textit{(livskvaliteten, ingen flertal)} & 生活质量 \\
\midrule
动词 [har] & at tiltrække & \textit{(tiltrækker, tiltrak, har tiltrukket)} & 吸引 \\
动词 \dabold{[er]} & at flytte & \textit{(flytter, flyttede, \dabold{er flyttet})} & 搬家，迁居 \\
动词 [har] & at sammenligne (med) & \textit{(sammenligner, sammenlignede, har sammenlignet)} & 把……与……作比较 \\
动词 [har] & at overveje & \textit{(overvejer, overvejede, har overvejet)} & 考虑，斟酌 \\
异态动词 & at trives & \textit{(trives, trivedes, har trivedes)} & 舒心生活，适应良好 \\
\midrule
形容词 & rolig & \textit{(roligt, rolige)} & 安静的，平缓的 \\
形容词 & travl & \textit{(travlt, travle)} & 忙碌的，繁忙的 \\
\end{glosertable}

% ---------------------------------------------------------------------
\ncdsubsec{四、Grammatik \& Faste vendinger（核心语法与地道表达精讲）}

\begin{grammatikbox}{语法点 1：城乡对比说明文的两大黄金段落架构}
在 DU3 写作（如邮件建议、社会议题短文）中对比 A 与 B 两个选项时，切忌想到哪写到哪，请根据字数选择下面两种标准架构之一：

\vspace{0.3em}
\renewcommand{\arraystretch}{1.24}
{\small
\begin{tabularx}{\linewidth}{@{} >{\bfseries\raggedright\arraybackslash}p{3.0cm} >{\raggedright\arraybackslash}p{5.5cm} >{\raggedright\arraybackslash}X @{}}
\toprule
\textbf{段落架构} & \textbf{段落排布模板} & \textbf{适用场景与段落路标} \\
\midrule
\textbf{1. 平行对比法}\newline （\textit{Blokmetoden}） &
第 1 段：引言提出选择\newline
第 2 段：\textbf{集中写完方案 \textit{A}（如大城市）}的优缺点\newline
第 3 段：\textbf{集中写完方案 \textit{B}（如乡村）}的优缺点\newline
第 4 段：权衡总结 &
\textbf{最适合 100--150 词短文与回信}。\newline
第 3 段开头强转场：\textit{På den anden side / Hvis man derimod vælger at bo på landet...} \\
\midrule
\textbf{2. 逐点对比法}\newline （\textit{Punkt-for-punkt}） &
第 1 段：引言提出对比背景\newline
第 2 段：\textbf{维度一（如房价/空间）}：A vs B\newline
第 3 段：\textbf{维度二（如工作/交通）}：A vs B\newline
第 4 段：\textbf{维度三（如社交/结论）}：A vs B &
\textbf{最适合 200--300 词长说明文}（见本课 Tekst A）。\newline
每段开头亮维度：\textit{Når det gælder... / Hvad angår... / Sammenlignet med...} \\
\bottomrule
\end{tabularx}
}
\end{grammatikbox}

\begin{grammatikbox}{语法点 2：五大高阶对比衔接词的“词性与语序矩阵”}
同样表示“相反 / 而 / 另一方面”，它们放进句子里触发的语序规则完全不同！务必分清\textbf{从属连词}和\textbf{副词/介词短语}：

\vspace{0.3em}
\renewcommand{\arraystretch}{1.24}
{\small
\begin{tabularx}{\linewidth}{@{} >{\bfseries\raggedright\arraybackslash}p{2.9cm} >{\raggedright\arraybackslash}p{3.5cm} >{\raggedright\arraybackslash}X @{}}
\toprule
\textbf{衔接词} & \textbf{词性与核心语序} & \textbf{标准句式示范（注意加粗词位置）} \\
\midrule
\da{hvorimod} & \textbf{从属连词}（“而，反之”）\newline 引导\textbf{从句：不倒装}，\textit{ikke} 放动词前 & \textit{I byen er huslejen høj, \da{hvorimod} man på landet \dabold{ikke betaler} så meget.} \\
\midrule
\da{mens ... derimod} & \textit{mens} 引导从句，\textit{derimod} 插在从句主语后（或句中） & \textit{Byen er travl, \da{mens} livet på landet \dabold{derimod er} meget roligere.} \\
\midrule
\da{Derimod} /\newline \da{Til gengæld} & \textbf{主句副词}\newline 放句首\textbf{必须 V2 倒装}（V + S） & \textit{\da{Til gengæld} \dabold{får man} mere plads.}\newline \textit{\da{Derimod} \dabold{er der} færre busser på landet.} \\
\midrule
\da{På den ene side..., på den anden side...} & \textbf{成对副词短语}\newline 放句首\textbf{两句都要 V2 倒装} & \textit{\da{På den ene side} \dabold{har byen} mange job, \da{på den anden side} \dabold{koster} boligerne en formue.} \\
\midrule
\da{Sammenlignet med + 名词} & \textbf{分词介词短语}\newline 放句首\textbf{主句必须 V2 倒装} & \textit{\da{Sammenlignet med København} \dabold{er} Silkeborg meget billigere.} \\
\bottomrule
\end{tabularx}
}
\end{grammatikbox}

\begin{grammatikbox}{英丹对比排雷 \& 地道搭配（Faste vendinger）}
\begin{enumerate}
  \item \textbf{别把 \textit{On the other hand} 的无倒装带进丹麦语！}：英语写 \textit{On the other hand, you get more space}；丹麦语 \textit{På den anden side} 占第 1 位，动词必须第 2 位：\ding{55}\ \textit{På den anden side man får mere plads} $\rightarrow$ \ding{51}\ \textit{På den anden side \dabold{får man} mere plads}！
  \item \textbf{严格区分 \textit{hvorimod} 与 \textit{derimod}}：\textit{hvorimod} 是连接两半句的\textbf{从属连词}（不单独领起独立句，后接主语 + 动词）；\textit{derimod} 是\textbf{副词}（可放独立句首触发 V2 倒装，或跟在 \textit{mens / og} 后用）。
  \item \textbf{\textit{til gengæld} 的地道用法}：意为“作为回报/相对应的是（有得也有失，或坏事里的弥补优点）”，极度高频！例如房贵但离朋友近，或通勤远但房子大，都要用 \textit{til gengæld}。
\end{enumerate}

\textbf{本课高频地道搭配（Faste vendinger）：}
\begin{itemize}
  \item \da{at få mere plads for pengene}（花同样的钱获得更大的空间）/ \da{at bo tæt på naturen}（住得亲近自然）
  \item \da{at have en høj livskvalitet}（拥有高质量生活）/ \da{at trives i lokalsamfundet}（在本地社区生活得融洽）
  \item \da{Når det gælder...} / \da{Hvad angår...}（谈到……方面——逐点对比法段落起手式！）
\end{itemize}
\end{grammatikbox}

% ---------------------------------------------------------------------
\ncdsubsec{五、Mønsterøvelser（句型操练）}

\textbf{练习 A：语序与对比衔接词填空（\textit{hvorimod, derimod, til gengæld, på den ene side, på den anden side, sammenlignet med}）}
\begin{enumerate}
  \item \underline{\hspace{3.2cm}} København er en lille by som Vejle meget roligere og billigere.
  \item I storbyen har man mange caféer og museer, \underline{\hspace{2.4cm}} man i landdistrikterne ofte har naturen lige uden for døren.
  \item \underline{\hspace{2.8cm}} er det nemt at finde arbejde i byen, \underline{\hspace{2.8cm}} kan huslejen være alt for høj.
  \item Der er langt til den nærmeste biograf fra vores nye hus. \underline{\hspace{2.4cm}} har vi nu en stor have og søde naboer!
  \item Mange unge flytter til København for at studere, mens børnefamilier \underline{\hspace{2.0cm}} ofte søger ud mod grønne forstæder.
  \item 找出并改正 1 处语序错误：\textit{På den anden side man skal køre i bil næsten hver dag på landet.} $\rightarrow$
  \item 找出并改正 1 处从句语序错误：\textit{Byen har mange koncerter, hvorimod på landet har man ikke så mange tilbud.} $\rightarrow$
\end{enumerate}

\textbf{练习 B：段落衔接“建筑手术”——按要求重组合并句子}
\begin{enumerate}
  \item 【逐点对比法｜住房维度】用 \textit{På den ene side..., på den anden side...} 合并以下两句：\newline
        \textit{En treværelses lejlighed i København koster ofte meget om måneden.} + \textit{Man kan leje et helt hus med have i Nordjylland for samme pris.} $\rightarrow$
  \item 【逐点对比法｜交通维度】用 \textit{hvorimod} 合并（注意从句中 \textit{ikke} 的位置！）：\newline
        \textit{I storbyen går metroen hvert femte minut.} + \textit{Man kan ikke altid stole på hyppige busser i små landsbyer.} $\rightarrow$
  \item 【平行对比法｜转场起句】用 \textit{Sammenlignet med} 触发 V2 倒装合并：\newline
        \textit{Vi ser på det travle byliv.} + \textit{Hverdagen i et lille lokalsamfund føles meget mere afslappet.} $\rightarrow$
  \item 【弥补转折】用 \textit{Til gengæld} 改写并连成通顺两句：\newline
        \textit{Det tager 40 minutter at pendle med tog. Vi bor kun fem minutter fra skoven og stranden.} $\rightarrow$
\end{enumerate}

\textbf{练习 C：中译丹（重点演练城乡对比连词与 V2 倒装）}
\begin{enumerate}
  \item 与大城市相比，丹麦乡村的房价通常低得多。
  \item 一方面大学和工作机会都在城市里，另一方面大城市的房租非常高。
  \item 许多年轻人喜欢热闹的城市生活，而家庭则反而更喜欢安静的环境。（用 \textit{mens ... derimod}）
  \item 通勤时间有点长；不过相应地，我们在小城里拥有很高的生活质量。（用 \textit{Til gengæld}）
  \item 在乡村邻里关系很紧密，而在大城市里人们有时根本不认识自己的邻居。（用 \textit{hvorimod} + \textit{ikke}）
\end{enumerate}

% ---------------------------------------------------------------------
\ncdsubsec{六、DU3 Modul 3 实战微任务与参考答案}

\begin{emailbox}{Skrivning Opgave 2 实战模拟｜给朋友的城乡选择建议信}
\textbf{Situation}: Din ven, Mikkel, er snart færdig med sin uddannelse. Han overvejer, om han skal blive boende i København eller flytte til en mindre by i Jylland for at få sit første job. Han har skrevet til dig for at høre din mening.\\
\textbf{Opgave}: Skriv et svar til Mikkel på \textbf{115--145 ord}. Sammenlign de to muligheder ved enten at bruge \textit{blokmetoden} (平行对比法) eller \textit{punkt-for-punkt-metoden} (逐点对比法), og giv ham et godt råd til sidst. Brug mindst tre kontrastudtryk fra denne lektion.
\end{emailbox}

\begin{answerbox}
\textbf{Før du læser 参考回答}
\begin{enumerate}
  \item \textit{Storbyen har universiteter, virksomheder og et bredt kulturliv, men høje boligpriser og et travlt tempo. Udkantsdanmark tilbyder mere plads for pengene, frisk luft, natur og et stærkt lokalsamfund, men kræver ofte længere pendlertid og færre busser.}
  \item \textit{De overvejer at flytte fra Aarhus til et lille hus i Silkeborg. Da Jonas må arbejde hjemme to dage om ugen, og huslejen er meget lavere, beslutter de at bestille en fremvisning af huset på lørdag.}
\end{enumerate}

\textbf{Tekst B 阅读答案（Opgave 2A）}
\begin{enumerate}
  \item[(1)] \textbf{C} —— Jonas 担心通勤一小时，Sara 用 \textit{Det er rigtigt, men til gengæld er huslejen jo næsten halvt så dyr...} 回应并提出补偿优势，自然引出 Jonas 下句 \textit{Det har du ret i.}
  \item[(2)] \textbf{E} —— Jonas 说老板允许每周居家办公两天，Sara 惊喜回应 \textit{Ej, hvor godt! Det gør jo en kæmpe forskel...}（B 项语义相反）。
  \item[(3)] \textbf{A} —— Jonas 担心想念市中心的朋友和咖啡馆，Sara 用 \textit{Til gengæld ligger mine forældres hus kun ti minutter derfra, og vores venner kan jo altid komme på besøg...} 宽慰并引出“Super idé!”。
\end{enumerate}

\textbf{练习 A 答案}
\begin{enumerate}
  \item \da{Sammenlignet med}（后接名词短语，注意句中 \textit{er en lille by...} 倒装）
  \item \da{hvorimod}（从属连词，从句中主语 \textit{man} 在动词 \textit{har} 前）
  \item \da{På den ene side}；\da{på den anden side}（两句均已 V2 倒装）
  \item \da{Til gengæld}（独立句首触发倒装 \textit{har vi}）
  \item \da{derimod}（位于 \textit{mens} 从句主语之后）
  \item 缺少 V2 倒装 $\rightarrow$ \textit{På den anden side \dabold{skal man} køre i bil næsten hver dag på landet.}
  \item 从句内部不能用主句倒装！$\rightarrow$ \textit{Byen har mange koncerter, \da{hvorimod} man på landet \dabold{ikke har} så mange tilbud.}
\end{enumerate}

\textbf{练习 B 答案}
\begin{enumerate}
  \item \textit{\da{På den ene side} koster en treværelses lejlighed i København ofte meget om måneden, \da{på den anden side} kan man leje et helt hus med have i Nordjylland for samme pris.}
  \item \textit{I storbyen går metroen hvert femte minut, \da{hvorimod} man i små landsbyer \dabold{ikke} altid kan stole på hyppige busser.}
  \item \textit{\da{Sammenlignet med} det travle byliv \dabold{føles} hverdagen i et lille lokalsamfund meget mere afslappet.}
  \item \textit{Det tager 40 minutter at pendle med tog. \da{Til gengæld} \dabold{bor vi} kun fem minutter fra skoven og stranden.}
\end{enumerate}

\textbf{练习 C 答案}
\begin{enumerate}
  \item \textit{Sammenlignet med storbyerne er boligpriserne i landdistrikterne som regel meget lavere.}
  \item \textit{På den ene side ligger universiteterne og mange arbejdspladser i byerne, på den anden side er huslejen i storbyen meget høj.}
  \item \textit{Mange unge kan lide det travle byliv, mens familier derimod foretrækker rolige omgivelser.}
  \item \textit{Pendlertiden er lidt lang. Til gengæld har vi en høj livskvalitet i den lille by.}
  \item \textit{På landet kender naboerne hinanden godt, hvorimod man i storbyen nogle gange slet ikke kender sine naboer.}
\end{enumerate}

\textbf{Modul 3 城乡建议高分回信（139 词 · 采用平行对比法）}\\
\textit{Til: Mikkel \quad|\quad Emne: Sv: Storby eller mindre by efter studiet?}

\textit{Hej Mikkel}

\textit{Stort tillykke med, at du snart er færdig med din uddannelse! Jeg kan godt forstå, at du overvejer begge muligheder.}

\textit{\da{På den ene side} er det nemmere at finde det første spændende job i København, og du har allerede et stort netværk og et rigt kulturliv lige uden for døren. \da{Sammenlignet med} mindre byer er huslejen dog meget høj, og det kan være svært at spare penge op.}

\textit{\da{På den anden side} får du i en mindre jysk by meget mere plads for pengene og et roligt tempo tæt på naturen. \da{Til gengæld} skal du måske regne med at bruge en bil, \da{hvorimod} du i København bare kan cykle overalt.}

\textit{Hvis du får et godt jobtilbud i Jylland, synes jeg, at du skal prøve det --- du kan jo altid flytte tilbage senere!}

\textit{Mange hilsner\\
\mbox{[Dit navn]}}

\small\textbf{考官点评}：全信结构严格按照“平行对比法”先集中论述哥本哈根、再用 \textit{På den anden side} 切换至日德兰小城，中间精准穿插了 \textit{Sammenlignet med}, \textit{Til gengæld}（V2 倒装）和 \textit{hvorimod}（从句语序），篇幅 139 词（远超 Modul 3 的 90 词底线），论证自然有力。
\end{answerbox}

\ncdlesson{20}{At finde sin plads i et nyt sprog}{在新语言中找到自己的位置 [综合总复习]}
{DU3 Modul 3（B1）结业总汇演 $\rightarrow$ 无缝衔接 NCD 4（Modul 4 / B1+）｜对标 NCE 2 结业向 NCE 3 过渡}
{\textbf{第三册“六大核心语法支柱”全景汇演}（\textit{-s} vs \textit{blive} 被动 · 介词 + 从句 · 分词定语 · 无 \textit{hvis} 条件倒装 · 复合词 \textit{fuge-s} · 复合介词与分裂排比）· \textbf{250--300 词社会话题说明文写作架构}}
{祝贺你走到《新概念丹麦语 3》的收官课！如果说 NCD 2 让你用正确的从句讲好“个人经历”，NCD 3 则真正交给了你“分析丹麦社会制度与文化现象”的 B1 工具箱。精读本课 280 词社会长文与 140 词反思邮件，并拿下 130+ 词综合写作，你已完全具备通关 DU3 Modul 3 并迈入 NCD 4（思辨与辩论）的实力！}

% ---------------------------------------------------------------------
\begin{foerduleaeser}
带着下面两个宏观问题阅读 Tekst A 和 Tekst B：
\begin{enumerate}
  \item \textit{Hvorfor er det danske sprog mere end bare et praktisk redskab til indkøb ifølge Tekst A, og hvordan hjælper det på arbejdspladsen?}（根据 Tekst A，为什么丹麦语不只是用来购物的实用工具？它如何在职场中起到帮助？）
  \item \textit{Hvordan har skribentens dansk udviklet sig i løbet af Modul 3 i Tekst B, og hvad glæder han sig til næste gang?}（在 Tekst B 中，作者的丹麦语在 Modul 3 期间有了怎样的变化？他对下一阶段有什么期待？）
\end{enumerate}
\end{foerduleaeser}

% ---------------------------------------------------------------------
\ncdsubsec{二、Tekst A \& Tekst B（社会深读与 Modul 3 真题场景）}

\subsection*{Tekst A：\textit{At finde sin plads i et nyt sprog}（社会深读 · 约 280 词）}

\begin{tipbox}{语篇说明}
本篇收官长文约 280 词，完整示范\textbf{“250--300 词社会话题说明文”}的起承转合，并将全册核心语法融于一炉：\textbf{\color{NCDNavy}深蓝色标注被动语态、分词定语、整句关系从句与复合词}；\textbf{\color{NCDRed}深红色标注无 \textit{hvis} 条件倒装、介词从句、复合介词与分裂排比}！
\end{tipbox}

Når man flytter til Danmark fra udlandet, kan man i begyndelsen klare mange praktiske ting på engelsk. Alligevel opdager de fleste hurtigt, at sproget spiller en helt afgørende rolle, hvis man virkelig ønsker at finde sin plads i det danske \da{samfund}. Et nyt sprog er nemlig meget mere end grammatik og gloser; \dabold{det er først og fremmest} nøglen til at forstå både de \da{skrevne} og \da{uskrevne} regler i hverdagen.

\dabold{I forbindelse med} arbejdsmarkedet er dansk ofte en stor fordel. Selvom der i mange internationale virksomheder \da{tales} engelsk til de officielle møder, foregår den uformelle snak ved kaffemaskinen næsten altid på dansk. Hvis man ikke forstår de små pauser og den afslappede humor, kan man let føle sig udenfor, \da{hvilket} påvirker ens \da{arbejdsglæde}. \dabold{Tør man} derimod øve sig med kollegerne fra første dag, \da{bliver} ens indsats næsten altid \da{værdsat}.

Også i kontakten med offentlige institutioner som \da{sundhedsvæsenet}, skoler og kommunen mærker man en tydelig forskel. \da{I forhold til} at være passiv modtager af oversatte beskeder giver et stærkere dansk en \da{voksende} selvtillid og kontrol over eget liv. Man begynder at \dabold{regne med, at} man selv kan løse små problemer, og man er ikke længere i tvivl om, \dabold{hvordan} man skal udtrykke sine behov.

Selvfølgelig kræver det tålmodighed at mestre dansk udtale, som i starten kan virke svær. \dabold{Ikke alene skal} man vænne sig til bløde d'er og mange vokaler, \dabold{men man skal også} lære at bruge sproget i et aktivt \da{lokalsamfund}. \da{På baggrund af} mine egne erfaringer vil jeg dog sige, at hvert eneste nye ord åbner en dør. Når man endelig tør grine med på en dansk joke i kantinen, ved man, at det hele er \da{lykkedes}.

\begin{transbox}
\small
\textbf{\color{NCDTeal}[Tekst A 中文精译]}\\
从国外搬到丹麦时，起初人们可以用英语应付许多生活琐事。然而，大多数人很快就会发现，如果真想在丹麦社会中找到属于自己的位置，语言起着绝对关键的作用。因为一门新语言远不止是语法和单词表；它首先正是理解日常生活中成文与不成文规则的钥匙。

在职场方面，会说丹麦语往往是一项巨大优势。尽管在许多跨国企业的正式会议上使用的是英语，但茶水间咖啡机旁的非正式闲聊几乎总是用丹麦语进行的。如果听不懂这些小憩时的交流和轻松的幽默，人就很容易感到被隔绝在外，而这会影响个人的工作幸福感。相反，如果一个人敢于从第一天起就和同事一起练习，他的这份努力几乎总是会得到认可与赞赏。

同样地，在与医疗卫生体系、学校和市政厅等公共机构打交道时，人也能感受到明显的区别。相比于被动接收翻译过来的通知，更扎实的丹麦语会赋予人不断增长的自信以及对自身生活的掌控感。人们开始相信自己能够独立解决小问题，也不再不确定该如何表达自己的需求。

当然，掌握起初可能显得棘手的丹麦语发音需要耐心。人不仅需要习惯软 \textit{d} 和众多元音，而且还必须学会在活跃的本地社区中运用这门语言。不过基于我自己的亲身经历，我想说，学会的每一个新词都会推开一扇门。当某天你终于敢在食堂里跟着大家一起为一句丹麦语笑话放声大笑时，你就知道这一切都成功了。
\end{transbox}

\subsection*{Tekst B：\textit{Tak for Modul 3 --- og mit nye liv på dansk}（对标 DU3 Modul 3 结业反思长邮件 · 约 140 词）}

\begin{emailbox}{DU3 Modul 3 结业高分长邮件（141 词）｜Til: Lærer Sofie \quad Emne: Tak for Modul 3 --- og mine næste skridt}
\textbf{Kære Sofie}

Nu har jeg lige afleveret min sidste opgave på Modul 3, og jeg vil bare sige tusind tak for et fantastisk halvt år i klassen!

Da jeg startede i januar, var jeg lidt i tvivl om, \dabold{hvorvidt} jeg kunne følge med i alle de nye tekster om det danske samfund. \dabold{I forhold til} begyndelsen føler jeg nu en helt anden tryghed, når jeg læser nyheder eller skriver formelle mails. \dabold{Det er især} vores mange samtaler om \da{arbejdspladskultur}, \da{sundhedsvæsenet} og grøn omstilling, \da{der} har givet mig et meget større ordforråd.

I sidste uge \da{blev jeg} faktisk \da{inviteret} til min første jobsamtale på dansk! \dabold{Ikke alene forstod} jeg alle chefens spørgsmål, \dabold{men jeg kunne også} fortælle om mine tidligere erfaringer uden at skifte til engelsk. Nu glæder jeg mig til at fortsætte på Modul 4 efter sommerferien, hvor der \da{diskuteres} endnu flere aktuelle emner.

\textbf{De bedste hilsner}\\
\textbf{Daniel}
\end{emailbox}

\begin{tipbox}{本课发音与流利度总览要点（Udtale \& Flydende dansk）}
\begin{itemize}
  \item \textbf{异态动词与 \textit{-s} 被动结尾轻读}：\textit{lykkes} \ipa{ˈløɡəs}（不是读成重重的 \ipa{-es}！）；\textit{tales} \ipa{ˈtˢæːləs}；\textit{diskuteres} \ipa{disɡuˈtˢeːˀɐs}（注意 \textit{-eres} 融合读 \ipa{eːˀɐs}！）；\textit{værdsat} \ipa{ˈvɛɐ̯ˌsɑd}（\textit{værd} 中的 d 不发音）。
  \item \textbf{社会抽象词中的不发音辅音（-d / -v / -g）}：
    \begin{itemize}
      \item \textit{et samfund} \ipa{ˈsɑmˌfɔnˀ}（词尾 \textit{d} 不发音，带 \textit{stød}）
      \item \textit{en selvtillid} \ipa{ˈsɛlˌtˢeliðˀ}（\textit{selv} 的 \textit{v} 不发音，结尾是软 d）
      \item \textit{en mulighed} \ipa{ˈmuːliˌheðˀ}（\textit{g} 不发音）
      \item \textit{uskreven} \ipa{ˈuˌsgʁɛːˀvən} / \textit{uskrevne} \ipa{ˈuˌsgʁɛʊ̯ˀnə}（复数中 \textit{v} 滑音化为 \ipa{ʊ̯}！）
    \end{itemize}
  \item \textbf{固定短语自然连读}：\textit{gøre sig umage} \ipa{ˈɡøːɐ sɑ ˈuˌmæːə}（\textit{sig} 弱读 \ipa{sɑ}，\textit{g} 不发音）；\textit{i højere grad} \ipa{i ˈhʌɪ̯ɐɐ ˈɡʁɑːˀð}。
\end{itemize}
\end{tipbox}

\begin{transbox}
\small
\textbf{\color{NCDTeal}[Tekst B 中文对照]}\\
\textbf{给索菲老师的结业致谢与反思邮件}：亲爱的索菲，我刚刚提交了 Modul 3 的最后一份作业，只想向你万分感谢这半年在班里的精彩学习！

一月份刚开课时，我还有点怀疑自己能否跟上所有那些关于丹麦社会的新课文。与最初相比，如今在阅读新闻或撰写正式邮件时，我感受到了完全不同的笃定感。尤其是我们在课堂上围绕职场文化、医疗体系与绿色转型进行的众多讨论，赋予了我庞大得多的词汇量。

上周我甚至收到了人生第一次全丹麦语求职面试的邀请！我不仅听懂了主管的所有问题，而且还能在不切换回英语的情况下介绍自己以往的工作经验。现在我满怀期待着暑假后继续升入 Modul 4，去在那里探讨更多热点时事话题。致以最美好的问候，丹尼尔。
\end{transbox}

% ---------------------------------------------------------------------
\ncdsubsec{三、Gloser（核心生词表）}

\begin{tipbox}{词表重点}
收官词表精选了社会参与、沟通自信与语言掌控的核心 B1 词汇。特别留意 \textit{at lykkes}（成功、顺利办成）——它既是带 \textit{-s} 的异态动词，构成完成时还必须使用 \dabold{er lykkedes}（例如：\textit{Det er lykkedes mig at bestå}）。
\end{tipbox}

\begin{glosertable}
名词 & et samfund & \textit{(samfundet, samfund, samfundene)} & 社会，共同体 \\
名词 & en selvtillid & \textit{(selvtilliden, ingen flertal)} & 自信，自信心 \\
名词 & en arbejdspladskultur & \textit{(arbejdspladskulturen, -kulturer, -erne)} & 职场文化 \\
名词 & en sammenhæng & \textit{(sammenhængen, sammenhænge, -ene)} & 语境，上下文；联系 \\
名词 & en adgang & \textit{(adgangen, adgange, adgangene)} & 通道，准入，机会 \\
名词 & et fællesskab & \textit{(fællesskabet, fællesskaber, -erne)} & 集体，社群归属感 \\
名词 & en indsats & \textit{(indsatsen, indsatser, indsatserne)} & 努力，付出，举措 \\
名词 & et redskab & \textit{(redskabet, redskaber, redskaberne)} & 工具，手段 \\
名词 & en nuance & \textit{(nuancen, nuancer, nuancerne)} & 细微差别，语气层次 \\
名词 & en arbejdsglæde & \textit{(arbejdsglæden, ingen flertal)} & 工作愉悦感 \\
\midrule
动词 [har] & at deltage (i) & \textit{(deltager, deltog, har deltaget)} & 参与，参加 \\
动词 [har] & at udtrykke (sig) & \textit{(udtrykker, udtrykte, har udtrykt)} & 表达（自己） \\
异态动词 \dabold{[er]} & at lykkes & \textit{(lykkes, lykkedes, \dabold{er lykkedes})} & 成功，顺利实现 \\
动词 [har] & at mestre & \textit{(mestrer, mestrede, har mestret)} & 精通，掌握 \\
动词 [har] & at værdsætte & \textit{(værdsætter, værdsatte, har værdsat)} & 重视，赏识 \\
\midrule
形容词 & uskreven & \textit{(uskrevet, uskrevne)} & 不成文的（\textit{uskrevne regler}） \\
分词形容词 & voksende & \textit{(不变)} & 不断增长的 \\
\end{glosertable}

% ---------------------------------------------------------------------
\ncdsubsec{四、Grammatik \& Skriveteknik（第三册语法总汇演与 250--300 词说明文写作）}

\begin{grammatikbox}{全景复盘：NCD 3（DU3 Modul 3）六大核心语法支柱总表}
在合上《新概念丹麦语 3》之前，请用下面这张总表完成全册 20 课核心语法的最终盘点：

\vspace{0.3em}
\renewcommand{\arraystretch}{1.24}
{\small
\begin{tabularx}{\linewidth}{@{} >{\bfseries\raggedright\arraybackslash}p{2.9cm} >{\raggedright\arraybackslash}p{5.5cm} >{\raggedright\arraybackslash}X @{}}
\toprule
\textbf{核心语法支柱} & \textbf{决定 B1 水准的黄金铁律} & \textbf{代表例句} \\
\midrule
\textbf{1. 双轨被动语态}\newline （L1） & \textbf{\textit{-s} 被动}表规则、习惯与通用现象；\textbf{\textit{blive} 被动}（\textit{bliver/blev/er blevet} + 分词）表具体一次事件。 & \textit{På kontoret \da{tales} der dansk.}\newline \textit{Jeg \dabold{blev inviteret} til samtale.} \\
\midrule
\textbf{2. 介词从句与整句指代}\newline （L2, L11, L13） & 动词+介词接从句时留意介词保留（必要时加形式宾语 \textit{det}，接 \textit{hv-} 可直连）；逗号 + \textit{hvilket} 指代前面整句话。 & \textit{Man er i tvivl om, \da{hvordan}...}\newline \textit{Det påvirker arbejdet, \da{hvilket} er trist.} \\
\midrule
\textbf{3. 分词定语与名词化}\newline （L3, L6, L16） & 现在分词（\textit{-ende}）不变形，表“正在/持续……”；过去分词（\textit{-et/-ede}）定指加 \textit{-e}；复合词最右端定词性，派生词常加 \textit{fuge-s}。 & \textit{en \da{voksende} selvtillid}\newline \textit{de \dabold{uskrevne} regler}\newline \textit{en \da{arbejdsplads}kultur} \\
\midrule
\textbf{4. 无 \textit{hvis} 条件与分裂排比}\newline （L4, L17） & 省掉 \textit{hvis} 时，从句动词放第 1 位（\textbf{V + S}），后面主句紧接着也要倒装！\textit{Ikke alene} 放句首必须 V2 倒装。 & \textit{\dabold{Tør man} øve sig, \da{bliver} man rost.}\newline \textit{\dabold{Ikke alene skal} man øve, men man får også selvtillid.} \\
\midrule
\textbf{5. 逻辑介词与对比架构}\newline （L7, L10, L18, L19） & 复合介词（\textit{i forhold til, i forbindelse med, med hensyn til, på baggrund af}）后不接 \textit{-ende}；\textit{hvorimod} 不倒装 vs \textit{Derimod/Til gengæld} 倒装。 & \textit{\da{På baggrund af} mine erfaringer...}\newline \textit{Her er dyrt, \da{hvorimod} landet er billigt.} \\
\midrule
\textbf{6. 推测、趋势与比例}\newline （L8, L12, L14, L15） & \textit{jo + 比较级..., desto + 比较级...}；趋势动词 \textit{stige/falde}；比例搭配单复数随后面中心名词变化。 & \textit{Jo mere man øver, desto bedre bliver man.} \\
\bottomrule
\end{tabularx}
}
\end{grammatikbox}

\begin{grammatikbox}{高阶写作：250--300 词社会话题说明文“四段黄金架构”}
无论是在 Modul 3 的扩展写作，还是在即将迈入的 NCD 4（Modul 4 / PD3）报刊说明文中，写 250--300 词的社会话题短文都可以套用 Tekst A 示范的四段式结构：

\vspace{0.3em}
\renewcommand{\arraystretch}{1.24}
{\small
\begin{tabularx}{\linewidth}{@{} >{\bfseries\raggedright\arraybackslash}p{2.9cm} >{\raggedright\arraybackslash}p{5.5cm} >{\raggedright\arraybackslash}X @{}}
\toprule
\textbf{段落结构} & \textbf{核心任务与目标词数} & \textbf{推荐调用的 NCD 3 武器} \\
\midrule
\textbf{1. Indledning}\newline （引言聚焦） & 从普遍现象切入，点出话题社会意义（约 45--55 词） & 客观引述 + 分裂句 \textit{Det er først og fremmest..., der...} \\
\midrule
\textbf{2. Hoveddel 1}\newline （制度或首要维度） & 分析现象在职场或制度中的具体影响（约 85--95 词） & \textit{I forbindelse med...} + \textit{-s} 被动 + \textit{hvilket} 整句定语从句 + 无 \textit{hvis} 条件句 \\
\midrule
\textbf{3. Hoveddel 2}\newline （对比或深层挑战） & 切入另一维度或平衡正反面（约 85--95 词） & \textit{I forhold til... / Med hensyn til...} + 分词定语 + 形式宾语/介词从句 \\
\midrule
\textbf{4. Afslutning}\newline （升华收束） & 基于前文总结立场与未来展望（约 45--55 词） & \textit{Ikke alene..., men også...} + \textit{På baggrund af...} \\
\bottomrule
\end{tabularx}
}
\end{grammatikbox}

\begin{grammatikbox}{英丹终极排雷（NCD 3 收官自检）\& NCD 4 衔接预告}
\begin{enumerate}
  \item \textbf{被动语态别全写成 \textit{er + 分词}}：英语 \textit{is spoken / was invited} 全部用 \textit{be}，但丹麦语表“规则习惯”用 \da{tales}，表“动作发生”必须用 \da{bliver talt / blev inviteret}；\textit{er + 分词} 只表示静态结果（如 \textit{Døren er lukket} 门是关着的）。
  \item \textbf{复合词与介词后的动词形式}：名词复合一律连写（\textit{arbejdspladskultur}，不拆开）；介词后接动作一律用名词或 \textit{at + 不定式}，绝不能写 \textit{-ende}。
  \item \textbf{牢记四大句首触发 V2 倒装的高频陷阱}：\ding{172} 省略 \textit{hvis} 的条件从句（\textit{Har man tid, \dabold{kan man}...}）；\ding{173} 句首 \textit{Ikke alene}（\textit{Ikke alene \dabold{er} det...}）；\ding{174} 句首复合介词或对比副词（\textit{På baggrund af det \dabold{besluttede jeg}... / Til gengæld \dabold{får man}...}）；\ding{175} 倒装同时别忘了\textbf{从句内部永远不倒装，\textit{ikke} 放动词前}！
  \item \textbf{NCD 4《思辨篇 · Holdning og Debat》抢先看}：进入第四册（Modul 4），我们将解锁丹麦语最迷人也最微妙的“灵魂情态小品词”（\textit{jo, da, nok, vist, lige, vel}）、非真实虚拟条件句以及报刊读者来信（\textit{Debatindlæg}）的礼貌反驳艺术！
\end{enumerate}
\end{grammatikbox}

% ---------------------------------------------------------------------
\ncdsubsec{五、Mønsterøvelser（全册综合句型操练）}

\textbf{练习 A：NCD 3 八关综合诊错与结构手术（每题找出 1 处典型语法错误并改正，或按要求改写）}
\begin{enumerate}
  \item 【被动语态】在办公室每天说丹麦语（改错）：\textit{På vores kontor er der talt dansk hver dag.} $\rightarrow$
  \item 【条件倒装】去掉 \textit{hvis} 改写（注意双倒装）：\textit{Hvis du øver dig lidt hver dag, forstår du snart danskerne bedre.} $\rightarrow$
  \item 【整句从句】用 \textit{hvilket} 合并：\textit{Jeg har fået et studiejob på en dansk café. Det hjælper mit sprog meget.} $\rightarrow$
  \item 【复合介词】改错：\textit{I forbindelse med lærende et nyt sprog skal man være tålmodig.} $\rightarrow$
  \item 【排比倒装】改错：\textit{Ikke alene sproget hjælper dig på arbejdet, men det også giver nye venner.} $\rightarrow$
  \item 【分词定语】用括号内词的复数定指形式填空：\textit{Man skal kende både de skrevne og de \underline{\hspace{2.0cm}} (uskreven) regler.}
  \item 【对比连词】改错：\textit{Her taler vi meget dansk, hvorimod i min gamle afdeling snakkede vi kun engelsk.} $\rightarrow$
  \item 【复合名词】合并并确定冠词（\textit{en/et}）：\textit{en uddannelse + et forløb}（培训流程）$\rightarrow$
\end{enumerate}

\textbf{练习 B：综合语篇完形（对标 DU3 Modul 3 \textit{Læsning Opgave 1} 进阶，从方框中选 8 个词填空，有 2 个多余词）}

\textit{备选词：deltage · hvilket · forhold · bliver · selvtillid · uskrevne · tales · baggrund · hvis · voksende}

\vspace{0.2em}
\begin{emailbox}{Læsning Opgave 1 进阶｜Sprogskolen som bro til samfundet}
For mange nye tilflyttere er sprogskolen den første vigtige bro til det danske samfund. Her (1) \underline{\hspace{2.0cm}} der ikke kun grammatik, men kursisterne lærer også om danske traditioner og (2) \underline{\hspace{2.0cm}} regler på arbejdspladsen. På (3) \underline{\hspace{2.0cm}} af diskussionerne i klassen får man en (4) \underline{\hspace{2.0cm}} forståelse af, hvordan velfærdsstaten fungerer. Det giver større (5) \underline{\hspace{2.0cm}} i hverdagen, (6) \underline{\hspace{2.0cm}} gør det lettere at (7) \underline{\hspace{2.0cm}} i både fritidsklubber og forældremøder. I (8) \underline{\hspace{2.0cm}} til i starten føler de fleste sig efter et år langt mere hjemme i Danmark.
\end{emailbox}

\textbf{练习 C：中译丹（综合调用全册六大语法支柱）}
\begin{enumerate}
  \item 只要大家互相倾听，很多误会就会被很快解决。（省略 \textit{hvis} + \textit{-s} 被动）
  \item 在申请新工作之际，学会用丹麦语介绍自己非常重要。（\textit{I forbindelse med} + 复合词）
  \item 我不确定自己是否能听懂丹麦语新闻，这让我有点紧张。（\textit{i tvivl om, hvorvidt...} + \textit{hvilket}）
  \item 与在课堂上相比，在真实职场中说丹麦语需要更大的勇气。（\textit{I forhold til}）
  \item 学一门新语言不仅需要时间，而且还能让你更好地理解一个国家的文化。（\textit{Ikke alene..., men også...}）
  \item 越是在日常生活中使用丹麦语，它就变得越自然。（\textit{Jo..., desto...}）
\end{enumerate}

% ---------------------------------------------------------------------
\ncdsubsec{六、DU3 Modul 3 毕业综合写作任务与参考答案}

\begin{emailbox}{B1 Capstone Writing Task（DU3 Modul 3 结业综合论述任务）}
\textbf{Opgave}: Skriv en sammenhængende tekst på \textbf{140--180 ord} (mindst 130 ord) om emnet:\newline
\textbf{\textit{”Sprogets betydning for at finde sin plads i Danmark”}}

Du skal komme ind på følgende fire punkter:
\begin{itemize}
  \item Hvorfor dansk er vigtigt i hverdagen ud over at handle ind på engelsk
  \item Hvad der ofte er sværest ved at lære sproget i starten
  \item Hvordan arbejdspladsen eller lokalsamfundet påvirkes, når man bruger dansk aktivt
  \item Et godt råd til nye udlændinge, der lige er startet på danskkursus.
\end{itemize}
\textit{（挑战目标：请在这篇 130+ 词的毕业作文里，主动调用至少 \textbf{5 项} NCD 3 核心语法武器！）}
\end{emailbox}

\begin{answerbox}
\textbf{Før du læser 参考回答}
\begin{enumerate}
  \item \textit{Dansk er nøglen til at forstå både skrevne og uskrevne regler i hverdagen. På arbejdspladsen foregår den uformelle snak ved kaffemaskinen næsten altid på dansk, og hvis man tør øve sig med kollegerne, bliver ens indsats næsten altid værdsat.}
  \item \textit{Skribenten føler nu en helt anden tryghed, når han læser nyheder eller skriver formelle mails, og har fået et større ordforråd. Han blev inviteret til sin første jobsamtale på dansk og glæder sig nu til at fortsætte på Modul 4 efter sommerferien.}
\end{enumerate}

\textbf{练习 A 答案（八关综合诊错）}
\begin{enumerate}
  \item 表达通用习惯须用 \textit{-s} 被动 $\rightarrow$ \textit{På vores kontor \dabold{tales} der dansk hver dag.}
  \item \textit{\dabold{Øver du dig} lidt hver dag, \dabold{forstår du} snart danskerne bedre.}（条件动词放第 1 位，主句也须倒装！）
  \item \textit{Jeg har fået et studiejob på en dansk café, \dabold{hvilket} hjælper mit sprog meget.}（用逗号 + \textit{hvilket} 取代 \textit{. Det}）
  \item 介词后严禁用 \textit{-ende} $\rightarrow$ \textit{I forbindelse med \dabold{at lære} et nyt sprog...}（或 \textit{I forbindelse med indlæringen af...}）
  \item 句首 \textit{Ikke alene} 缺倒装 + \textit{også} 错位 $\rightarrow$ \textit{Ikke alene \dabold{hjælper sproget} dig på arbejdet, men det \dabold{giver også} nye venner.}
  \item \da{uskrevne}（复数定指变化：\textit{uskreven $\rightarrow$ uskrevne}）
  \item \textit{hvorimod} 从句中禁止用主句倒装 $\rightarrow$ \textit{Her taler vi meget dansk, hvorimod vi i min gamle afdeling \dabold{kun snakkede} engelsk.}
  \item \da{et uddannelsesforløb}（\textit{-else} 结尾加 \textit{fuge-s}，右端词 \textit{et forløb} 决定中性 \textit{et}）
\end{enumerate}

\textbf{练习 B 答案}
(1) \da{tales} \quad (2) \da{uskrevne} \quad (3) \da{baggrund} \quad (4) \da{voksende} \quad (5) \da{selvtillid} \quad (6) \da{hvilket} \quad (7) \da{deltage} \quad (8) \da{forhold}

\textbf{练习 C 答案}
\begin{enumerate}
  \item \textit{Lytter man til hinanden, løses mange misforståelser hurtigt.}（或 \textit{...bliver mange misforståelser hurtigt løst.}）
  \item \textit{I forbindelse med at søge et nyt job er det meget vigtigt at kunne præsentere sig selv på dansk.}
  \item \textit{Jeg er i tvivl om, hvorvidt jeg kan forstå de danske nyheder, hvilket gør mig lidt nervøs.}
  \item \textit{I forhold til i klassen kræver det større mod at tale dansk på en rigtig arbejdsplads.}
  \item \textit{Ikke alene kræver det tid at lære et nyt sprog, men det giver dig også en bedre forståelse af landets kultur.}
  \item \textit{Jo mere man bruger dansk i hverdagen, desto mere naturligt bliver det.}
\end{enumerate}

\textbf{DU3 Modul 3 结业满分综合范文（171 词）}\\
\textit{\textbf{Sprogets betydning for at finde sin plads i Danmark}}

\textit{Når man lige er flyttet til Danmark, klarer man sig ofte fint på engelsk i supermarkedet. Men på længere sigt opdager man, at \da{det først og fremmest er} det danske sprog, \da{der} giver ægte adgang til samfundet. I hverdagen \da{tales} der nemlig dansk i alt fra børnehaver til nabomøder.}

\textit{I begyndelsen kan udtalen dog virke svær, og mange er i tvivl om, \dabold{hvorvidt} danskerne overhovedet forstår dem. \da{I forhold til} grammatikken tager det ofte længere tid at vænne sig til de bløde konsonanter og den hurtige rytme.}

\textit{\dabold{Tør man} alligevel bruge sit dansk på arbejdspladsen og i foreningslivet, sker der noget positivt. Man opdager hurtigt de \da{uskrevne} regler og fælles referencer, \da{hvilket} skaber en \da{voksende} følelse af tryghed og fællesskab. Samtidig \da{bliver} ens indsats næsten altid \da{værdsat} af både kolleger og naboer.}

\textit{\dabold{Ikke alene giver} et nyt sprog bedre jobmuligheder, \dabold{men det åbner også} døren til dybere venskaber. Mit bedste råd til nye kursister er derfor enkelt: vær ikke bange for at begå fejl --- øv jer lidt hver eneste dag!}

\small\textbf{结业总评}：这篇 171 词的结业佳作严格按四段结构展开，自然且不堆砌地调度了全册七大王牌武器：分裂句（\textit{Det er først og fremmest..., der...}）、\textit{-s} 被动（\textit{tales}）与 \textit{blive} 被动（\textit{bliver værdsat}）、介词接 \textit{hv-} 从句（\textit{i tvivl om, hvorvidt...}）、复合介词（\textit{I forhold til}）、无 \textit{hvis} 条件倒装（\textit{Tør man...}）、过去/现在分词定语（\textit{uskrevne, voksende}）、整句关系从句（\textit{hvilket}）以及句首排比倒装（\textit{Ikke alene giver..., men det åbner også...}）。祝贺你圆满通关《新概念丹麦语 3》，我们第四册《思辨篇》见！
\end{answerbox}


\end{document}
